% !TEX root = ../print.tex
% Chapter 3 == section 3 of the article (notes/PLAN-article-structure.md, "The outline"): the
% isotropic theory in dimension d. Two sections. The first, sec:isotropic-origin, is new material
% around the draft's section 2.1: the origin behaviour in dimension d and the dimension filtration.
% The second, sec:isotropic-classification, is the draft's section 2.1 proper: the classification
% and the bridge, by reduction to the line. The plan's outline names the classification first; here
% it comes second because its nodes read the first section's vocabulary: A_d
% (def:filtration-thresholds), the isotropic origin datum and its normalisation
% (def:isotropic-origin-datum), and the strictness of the filtration (thm:filtration-strict). Until
% Q-0278 (2026-10-04) the two sections were the parts 02-dimension.tex and 02-isotropic.tex; no
% label, statement or proof changed in the move (scripts/check-statements-unchanged.py).
%
% The origin and the filtration. Provenance: Q-0043 (hub queue), from P-0013. The source reading
% the item required first is records/affine/READING-sato1999levy-dimension-d.md, and every page
% anchor cited in that section is verified there, at the page, on 2026-09-25.
%
% The classification. The draft's Definition 2.1, Theorem 2.2 and Proposition 2.3, the port of
% [V]'s classification and of [B]'s bridge to R^d. Provenance: Q-0208 (hub queue), on the author's
% answer to Q-0205 (2026-09-28): option (c), done BY REDUCTION to [V, Thm. 7.3] through the
% one-dimensional marginal, the route of notes/PLAN-d-dimensional-port.md (Q-0210), its steps 1-3.
%
% Design, and where it departs from the draft. The draft plans a NATIVE port, redoing each step of
% [V] on R^d, and its port notes are claims. None of them is used here. A rotation-invariant law is
% determined by its one-dimensional marginal (lem:radial-marginal), the marginals of an isotropic
% family satisfy [V, Def. 3.1] as they stand (lem:isotropic-marginal-family), and [V, Thm. 7.3] is
% then applied to them unchanged, restated where it is used and cited by its v0.1 number; no \uses
% edge crosses into [V]'s blueprint. What is left in dimension d is the class A_d of
% def:filtration-thresholds, and membership in it is read off the cascade. Everything is stated on
% KERNELS, as def:affine-cascade-family is: the passage from the operator form of Definition 2.1
% to kernels ([V, Lem. 4.1] proved in R^d) is lem:operators-to-kernels(5) in 04-classical.tex
% (Q-0207), which no node of this chapter reads.
%
% The explicit form (2.1) is a separate [A] node, prop:isotropic-polar-reading, on ledger A2, A3
% and A4. It fixes the normalisation of the profile against def:isotropic-origin-datum: this
% module's k_d is |S^{d-1}| times the draft's k, its Gaussian term is (1/2) a |xi|^2 where the
% draft writes a r^2, and at d = 1 its k_1 is [V]'s folded profile.
%
% The scheme. def:isotropic-kernel-family is kept as it stands, with its free gauge and (ND) at
% every pair: prop:isotropic-corner in 04-classical.tex (Q-0280) proves it is the slice {sI} of the
% (G, W) scheme's isotropic corner up to the variance gauge, and no node of this chapter reads it.
%
% Stays prose, and why: the draft's port notes after Theorem 2.2 (claims about a native port that
% this route does not take; the plan's section 2 checks them against the source), and the
% sentences of Proposition 2.3 that name densities (the Bessel potentials of R^d, the elliptical
% Student-t): no node computes a density, and none is needed. Decision 3 of the plan makes the
% draft's other port claims one paragraph of the article, with no claims beyond it; no node here.

\chapter{The isotropic theory}
\label{ch:dimension}

This chapter is the article's \S 3: the isotropic theory in $\RR^d$, which is the draft's \S 2.1 and
new material around it. The behaviour of an isotropic aperture at its own centre, and the dimension
filtration it is read against, make up \S\ref{sec:isotropic-origin}. The isotropic classification and
the bridge, proved by reduction to the line's theorem, are \S\ref{sec:isotropic-classification}, which
comes second because it names the filtration's classes. \ref{thm:origin-boundedness-d}, with
\ref{prop:isotropic-polar-reading}(1) for the reading in self-decomposable laws, answers the
fourth open question of the released cone module [B, \S 8]: \emph{whether the boundedness
threshold in $\RR^d$ is the dimension}. It is, for every isotropic self-decomposable law with no
Gaussian part, in both directions, and the reason is not a property of the origin ---
the quantity read there does not move with the dimension at all
(\ref{cor:origin-profile-exponent}) --- but of the transform: the
threshold is the dimension of the Fourier integral that turns a decay rate into a bounded
density, and the laws being symmetric, their transforms are positive, so that the integral
decides boundedness in both directions. The five nodes after them answer the two questions [B]'s Remark~4.10 and its Figure~4
leave open about the chain $\mathcal{A}_1 \supseteq \mathcal{A}_2 \supseteq \dots$ itself: the
chain is strict at \emph{every} step, because the thresholds $w_d$ of [B, (A.2)] are strictly
decreasing, and no $\Cin$ ray of the line reaches the plane, so that of the extreme rays of the
line's cone only the Gaussian extends. The thresholds also tend to $0$, at the rate $d^{-1/2}$,
which is the rest of [B, Rem.~4.10]'s numerical sentence.

% CHANGED Q-0429 (2026-10-12): this paragraph rewritten for the split of cor:dictionary-filtration
% and prop:separating-family-mixtures (ten nodes where there were eight), and for the state of
% their Lean, which it had stated as "not yet formalized" for all eight since Q-0425 although
% three were proved in Q-0427. The second part of Q-0429 changes its last three sentences.
% CHANGED Q-0429, part 2: the last three sentences, now that the shadow lemma, the consequences
% on data and the theorem on the separating family have their Lean.
Ten further nodes describe these results in terms a reader can compute with; each has a proof of
record. \ref{prop:isotropic-exponent-profile}, \ref{lem:exponent-profile-constant} and
\ref{prop:origin-rate-profile-d} give the transform at infinity and the rate at the centre
through the polar profile, with their constants. \ref{lem:threshold-asymptotic} gives the limit
of the critical values, $w_d \sim \sqrt{\pi/d}$. \ref{lem:shadow-polar-profile} computes the
polar profile of the shadow of a datum one dimension down, and
\ref{prop:profile-one-dimension-down} draws the consequences on data: the profile one dimension
down is continuous, so a datum whose profile has a jump goes no further, the step datum
separates $\mathcal{A}_d$ from $\mathcal{A}_{d+1}$ at every $d$, and no $\Cin$ ray reaches the
plane. \ref{thm:separating-family-critical-values} is the statement on the separating family:
$F^{(w)}$ belongs to $\mathcal{A}_D$ exactly when $w \le w_D$, the critical values decrease
strictly, and $\sqrt D\,w_D \to \sqrt\pi$. \ref{prop:radial-dictionary} relates the classes
$\mathcal{A}_d$ to the characteristic functions of the rotation-invariant laws on $\RR^d$, which
explains why the classes shrink with the dimension; \ref{cor:dictionary-filtration} reads the
consequences on data in its terms, and \ref{prop:separating-family-mixtures} says what the
members of the separating family are as laws. Six of the ten are machine-checked: the two
propositions on the polar profile, the limit of the critical values, the shadow lemma, the
consequences on data and the theorem on the separating family. So every statement by which a
reader decides whether an exponent is admitted in a dimension is machine-checked. The closed
form of the constant, the dictionary, its readings and the family as laws are proved in print
and are not formalized.

% Q-0434 (2026-10-13): this paragraph, for the five nodes of the second review round.
Five nodes come from the second external review of the isotropic paper.
\ref{cor:isotropic-scale-zero-kernels} states the chapter's main result once: the
rotation-invariant self-decomposable laws on $\RR^d$ other than the point mass at the origin are
exactly the kernels from scale zero to a positive scale of the isotropic cascade families, on
kernels and on operators, and a family with a given kernel is determined by the gauge. It is
machine-checked. Three nodes describe the stages $\mu_{s,t}$ with $s > 0$, which are what a
cascade applies. \ref{lem:isotropic-stage-levy-measure} gives the \Levy{} measure of a stage and
its mass, $c\,\log(\chi(t)/\chi(s))$. \ref{prop:isotropic-stage-atom} says that a stage of a
family with no Gaussian part and a finite $c$ is a point mass of weight $(\chi(s)/\chi(t))^c$ at
the centre plus an absolutely continuous part, and is not self-decomposable.
\ref{prop:isotropic-stage-density} says that every stage is infinitely divisible, and has a
density when $c$ is infinite or there is a Gaussian part. \ref{prop:step-exponent-dimensions}
says how far the step exponent of dimension $d$ goes: it is the exponent of an infinitely
divisible law exactly up to dimension $d + 2$, and of a self-decomposable one exactly up to
$d$. These four are proved in print. The three on the stages are typed and not proved in Lean,
and whether they are formalized is not decided; the last is stated in terms of the dictionary
and has no Lean.

\section{The origin and the dimension filtration}
\label{sec:isotropic-origin}

% draft: new material (not in the draft; the draft's 2.1 is a port of [V] and says nothing about
% the origin). The objects are [B]'s, restated so that this chapter stands on its own: an
% isotropic origin datum is the polar form of [B, Prop. 4.9(2)]'s Levy measure with the delay
% datum forgotten, and its polar profile is [B, (4.6)] when the datum comes from a delay datum.
% CHANGED Q-0281: typed (the statement skeleton). The displayed claim eq:polar-decomposition ("so
% that ... for every nonnegative Borel g vanishing at the origin") is split to lem:polar-decomposition
% (ADR-0025; records/formalization/PRICING.md, finding 4); the text is unchanged and the \lean tag
% names its declaration too, with two notation lemmas identifying the Lean definitions of |S^{d-1}|
% and of k_d(0+) with the ones written here. "Its law" is read by specification (PRICING finding 1):
% any probability measure with transform e^{-F_d} (IsLawOf), its existence being A3(iii).
\begin{definition}[isotropic origin data in dimension $d$, and the polar profile]
\label{def:isotropic-origin-datum}
\lean{AffineHemigroup.OriginDatum, AffineHemigroup.OriginDatum.levyMeasure,
AffineHemigroup.OriginDatum.exponentL, AffineHemigroup.OriginDatum.exponent,
AffineHemigroup.OriginDatum.profileAtOrigin, AffineHemigroup.OriginDatum.IsSubordinated,
AffineHemigroup.IsLawOf, AffineHemigroup.sphereArea, AffineHemigroup.polarProfileOf,
AffineHemigroup.sphereArea_eq_toSphere, AffineHemigroup.OriginDatum.tendsto_profileAtOrigin,
AffineHemigroup.OriginDatum.polar_decomposition}\leanok
Fix an integer $d \ge 1$ and write $\sigma_{d-1}$ for the surface measure of the unit sphere
$S^{d-1} \subset \RR^d$ and $|S^{d-1}| = \sigma_{d-1}(S^{d-1})$ for its surface area. An
\emph{isotropic origin datum in dimension $d$} is a pair
$(a, k_d)$ with $a \ge 0$ and $k_d : (0,\infty) \to [0,\infty)$ nonincreasing, subject to
\[
  \int_0^1 k_d(r)\,r\,dr < \infty
  \qquad\text{and}\qquad
  \int_1^\infty k_d(r)\,\frac{dr}{r} < \infty .
\]
Its \Levy{} measure is the rotation-invariant measure $\nu_d$ on $\RR^d \setminus \{0\}$ with radial
density $f_d(r) = k_d(r)\,\bigl(|S^{d-1}|\,r^d\bigr)^{-1}$, so that
\begin{equation}\label{eq:polar-decomposition}
  \int_{\RR^d} g \,d\nu_d
   = \frac{1}{|S^{d-1}|}\int_{S^{d-1}}\!\!\sigma_{d-1}(d\eta)\int_0^\infty g(r\eta)\,k_d(r)\,\frac{dr}{r}
\end{equation}
for every nonnegative Borel $g$ vanishing at the origin. Its \emph{exponent} is
\begin{equation}\label{eq:isotropic-exponent}
  F_d(\xi) = \tfrac12 a|\xi|^2
   + \int_{\RR^d}\bigl(1 - \cos\langle \xi, x\rangle\bigr)\,\nu_d(dx), \qquad \xi \in \RR^d ,
\end{equation}
and its \emph{law} is the symmetric infinitely divisible law $\mu$ on $\RR^d$ with
$\widehat\mu = e^{-F_d}$. We call $k_d$ the datum's \emph{polar profile} and
$c := k_d(0{+}) \in [0,\infty]$ its \emph{profile at the origin}. The datum is
\emph{subordinated} when there is a delay datum $(a, k_{\mathrm I})$ in the sense of
\restated{B}{Def.~4.11}{sec:line-bridge}, with $k_{\mathrm I}$ nonincreasing, whose polar profile \prog{B}{(4.6)} is $k_d$:
\begin{equation}\label{eq:polar-profile-d}
  k_d(r) = \frac{2}{\Gamma(d/2)}\int_0^\infty v^{d/2-1}e^{-v}\,
    k_{\mathrm I}\Bigl(\frac{r^2}{2v}\Bigr)\,dv , \qquad r > 0 .
\end{equation}
\statusT\quad\emph{(Vocabulary, and no claim. Three readings are deliberately \textbf{not}
asserted here. First, that the laws of the isotropic origin data with $a = 0$ are exactly the
isotropic self-decomposable laws on $\RR^d$ with no Gaussian part: that is \sscite{sato1999levy}'s
polar criterion, Thm.~15.10, p.~95, read for a rotation-invariant \Levy{} measure, and it is
\ref{prop:isotropic-polar-reading}(1), which spends \ledger{A2} for it. Second, that a
subordinated datum's exponent is $F_d(\xi) = F_{\mathrm I}(|\xi|^2/2)$, $F_{\mathrm I}$ being the
delay datum's exponent, and third, that its law is the Gaussian variance mixture with the delay law
as mixing law; those are [B, Prop.~4.9(2)] and [B, Rem.~4.3], restated where they are used, in
\ref{thm:origin-boundedness-d}(3). Until Q-0286 the second stood in this definition, and the
words ``nonincreasing in the sense of [B, Def.~4.11]'' did too: [B, Def.~4.11] asks only for a
measurable profile, so monotonicity is this definition's own condition, as it now reads. The
definition itself reads only the displayed formulas.
\ref{eq:polar-decomposition} is the polar decomposition of a rotation-invariant measure and is
why $k_d$, and not $f_d$, is the object the origin behaviour is stated on: comparing
\ref{eq:polar-decomposition} with \sscite{sato1999levy}~(15.8), p.~95, identifies $k_d$ with
$|S^{d-1}|$ times the source's own polar profile $k_\eta$, that is with the
\textbf{sphere-total} $\int_{S^{d-1}}k_\eta(r)\,\sigma_{d-1}(d\eta)$ of it; at $d = 1$ the sphere is
$S^0 = \{\pm1\}$, $\sigma_0$ is counting measure, and the sphere-total is the folded profile
$k(r) + k(-r)$ of [B, (2.3)], whose value at the origin is the constant $c = k(0{+}) + k(0{-})$
of \sscite{sato1999levy}~Thm.~28.4, p.~191 and Thm.~53.8, pp.~410--411. That identification, and
nothing else, is what makes the theorem below the $d$-dimensional form of
[B, Cor.~3.9]; it is recorded at the page in
\texttt{records/affine/READING-sato1999levy-dimension-d.md} \S 2.)}
\end{definition}

% CHANGED Q-0281: split from def:isotropic-origin-datum (ADR-0025; records/formalization/PRICING.md,
% finding 4): the claim the definition displays as eq:polar-decomposition, stated as a lemma. The
% definition's text, and the equation's label, are unchanged.
\begin{lemma}[the polar decomposition of an isotropic origin datum's \Levy{} measure]
\label{lem:polar-decomposition}
\uses{def:isotropic-origin-datum}
\lean{AffineHemigroup.OriginDatum.polar_decomposition}\leanok
Let $d \ge 1$ and let $(a, k_d)$ be an isotropic origin datum in dimension $d$ with \Levy{} measure
$\nu_d$. Then \ref{eq:polar-decomposition} holds for every nonnegative Borel $g$ on $\RR^d$
vanishing at the origin, $\sigma_{d-1}$ being the surface measure of $S^{d-1}$ and $|S^{d-1}|$ its
surface area.
\statusT\quad\emph{(Cites nothing: polar coordinates in $\RR^d$, Lebesgue measure being the image
of $\sigma_{d-1} \otimes r^{d-1}dr$ under $(\eta, r) \mapsto r\eta$. Split from
\ref{def:isotropic-origin-datum} on 2026-10-04 (Q-0281, ADR-0025); the definition's text is
unchanged.)}
\end{lemma}

\begin{proof}
\leanok
By definition $\nu_d$ has the density $f_d(|x|) = k_d(|x|)(|S^{d-1}||x|^d)^{-1}$ on
$\RR^d \setminus \{0\}$. In polar coordinates $x = r\eta$, $dx = r^{d-1}dr\,\sigma_{d-1}(d\eta)$, so
for $g \ge 0$ Borel and vanishing at the origin
$\int g\,d\nu_d = \int_{S^{d-1}}\sigma_{d-1}(d\eta)\int_0^\infty g(r\eta)\,
k_d(r)(|S^{d-1}|r^d)^{-1}r^{d-1}\,dr$, which is \ref{eq:polar-decomposition}.
\end{proof}

% Withdrawn 2026-10-03 (the author, after Q-0151): def:radial-origin-asymptotic, hypothesis (H),
% the pointwise radial lower bound p(x) >= c_1|x|^{c-d} (log(1/|x|) at c = d) with a radially
% nonincreasing representative. It was introduced (Q-0043, kept a definition by Q-0148) as the
% hypothesis of thm:origin-boundedness-d(4), because Sato's Prop. 28.1 is one-way for a general
% law on R^d; Q-0151 proved clause (4) for every datum without it (lem:positive-transform-origin),
% and no node used it after that. Its lower bound p(x) >= c_1|x|^{c-d} is proved for every d >= 2 and
% every datum by prop:origin-lower-bound-d (Q-0272); the log(1/|x|) at c = d and the slowly varying
% factor stay open for d >= 2 (the isotropic paper's section 4). The monotone representative is
% prop:isotropic-radial-monotone since Q-0316 (Wolfe's theorem, ledger A8; records/affine/GATE-paper-iso.md
% item 18).

% Q-0151. Route (ii) of the queue entry: "a sign condition on the transform along rays that lets
% Prop. 28.1's criterion be read as an equivalence". The sign condition is nonnegativity, and the
% equivalence is Gaussian smoothing; nothing about rays is needed. Ledger-free.
\begin{lemma}[a nonnegative transform: bounded near the origin only if integrable]
\label{lem:positive-transform-origin}
\lean{AffineHemigroup.positive_transform_heat, AffineHemigroup.positive_transform_bound,
AffineHemigroup.density_bound_ball, AffineHemigroup.heatKernel}\leanok
Let $d \ge 1$ and let $\mu$ be a probability measure on $\RR^d$ whose characteristic function
$\widehat\mu(\xi) = \int e^{i\langle\xi,x\rangle}\mu(dx)$ is real and nonnegative at every
$\xi \in \RR^d$. For $t > 0$ let $g_t(x) = (2\pi t)^{-d/2}e^{-|x|^2/2t}$ be the heat kernel.
\begin{enumerate}
\item For every $t > 0$,
\[
  \int_{\RR^d} g_t\,d\mu = (2\pi)^{-d}\int_{\RR^d}\widehat\mu(\xi)\,e^{-t|\xi|^2/2}\,d\xi ,
\]
and as $t \downarrow 0$ this increases to $(2\pi)^{-d}\int_{\RR^d}\widehat\mu(\xi)\,d\xi \in (0,\infty]$.
\item If there are $\delta > 0$ and $M < \infty$ with $\mu(B) \le M|B|$ for every Borel set
$B \subseteq \{|x| < \delta\}$, then $\int_{\RR^d}\widehat\mu(\xi)\,d\xi \le (2\pi)^d M$. The
hypothesis holds in particular when $\mu$ has a density with a representative bounded by $M$ on
$\{0 < |x| < \delta\}$.
\end{enumerate}
\statusT\quad\emph{(Cites nothing: the Gaussian Fourier integral, Fubini's theorem and monotone
convergence; ledger entry A1 is named below for comparison and not spent. \textbf{What it is
worth.} With entry A1 at $n = 0$, which \ref{thm:origin-boundedness-d}(2) spends, it makes the
Fourier-integrability criterion an equivalence for a nonnegative transform: an integrable transform,
a bounded continuous density, and $\mu \le M\cdot$Lebesgue near the origin are then the same
condition. The law of every isotropic origin datum has the positive transform $e^{-F_d}$, so this is
the step that makes \ref{thm:origin-boundedness-d}(4) unconditional; on the subordinated slice
clause (3) proves the same conclusion directly, with an explicit monotone density. It is route (ii) of Q-0151,
and the sign condition that route asks for along rays holds everywhere for these laws. \textbf{Not
claimed}: anything for a transform that changes sign, where entry A1 remains one-way; nor
anything about the value of a representative at a point: clause (2) is a statement about the
measure.)}
\end{lemma}

\begin{proof}
\leanok
\emph{Clause (1).} The one-dimensional Gaussian integral
$\int_{\RR}e^{-isy}e^{-ts^2/2}\,ds = (2\pi/t)^{1/2}e^{-y^2/2t}$, multiplied over the $d$
coordinates, gives $g_t(x) = (2\pi)^{-d}\int_{\RR^d}e^{-i\langle\xi,x\rangle}e^{-t|\xi|^2/2}\,d\xi$.
The integrand has modulus $e^{-t|\xi|^2/2}$, which is integrable against $d\xi \otimes \mu(dx)$
because $\mu$ is finite, so Fubini's theorem gives
\[
  \int g_t\,d\mu = (2\pi)^{-d}\int e^{-t|\xi|^2/2}\int e^{-i\langle\xi,x\rangle}\mu(dx)\,d\xi
   = (2\pi)^{-d}\int e^{-t|\xi|^2/2}\,\widehat\mu(-\xi)\,d\xi ,
\]
and the substitution $\xi \mapsto -\xi$ gives the display. Since $\widehat\mu \ge 0$, the integrand
increases as $t$ decreases, and monotone convergence gives the limit. It is positive because
$\widehat\mu$ is continuous with $\widehat\mu(0) = 1$.

\emph{Clause (2).} Put $U = \{|x| < \delta\}$. The hypothesis gives
$\int_U h\,d\mu \le M\int_U h\,dx$ for every nonnegative Borel $h$, for indicator functions of
Borel subsets of $U$ directly and in general by simple functions and monotone convergence. With
$h = g_t$, $\int_U g_t\,d\mu \le M\int_{\RR^d}g_t\,dx = M$. Off $U$,
$g_t \le (2\pi t)^{-d/2}e^{-\delta^2/2t}$, so
$\int_{\RR^d\setminus U}g_t\,d\mu \le (2\pi t)^{-d/2}e^{-\delta^2/2t}$, which tends to $0$ as
$t \downarrow 0$. By clause (1), $(2\pi)^{-d}\int\widehat\mu = \lim_{t\downarrow0}\int g_t\,d\mu \le
M$. For the last sentence: if $q$ is a density of $\mu$ with $q \le M$ on $U \setminus \{0\}$, then
$\mu(B) = \int_{B\setminus\{0\}}q \le M|B|$ for every Borel $B \subseteq U$, since the origin is a
Lebesgue-null set.
\end{proof}

% Q-0151, and the link the brief asked to check: the upper half of the transform's behaviour at
% infinity, which is all thm:origin-boundedness-d(4) needs, and its regular variation, which is
% the Abelian half of the "Tauberian link" prop:log-orbit-rays' annotation left open. Both read
% only k_d <= c, k_d -> c at the origin, and the integrability at infinity.
% CHANGED Q-0281: \uses re-pointed to the split child it reads (eq:polar-decomposition).
% CHANGED Q-0425 (2026-10-12): annotation only. "Not claimed: an explicit form of L; nor that L is
% bounded above" now points to prop:isotropic-exponent-profile, which gives both for c > 0. The
% statement, its proof and the \lean tag are unchanged, and so are the papers' shared copies.
\begin{lemma}[a profile bounded at the origin gives a logarithmic exponent]
\label{lem:isotropic-exponent-log}
\uses{def:isotropic-origin-datum,lem:polar-decomposition}
\lean{AffineHemigroup.isotropic_exponent_log_bound, AffineHemigroup.isotropic_exponent_log_limit,
AffineHemigroup.OriginDatum.tailConst}\leanok
Let $d \ge 1$, let $(0,k_d)$ be an isotropic origin datum in dimension $d$ with Gaussian
coefficient $0$ and $c = k_d(0{+}) < \infty$, let $F_d$ be its exponent, and put
$C_k := \int_1^\infty k_d(r)\,dr/r$.
\begin{enumerate}
\item \emph{(The exponent grows at most like $c\log|\xi|$.)} For every $\xi$ with $|\xi| \ge 1$,
\[
  F_d(\xi) \le c\log|\xi| + \tfrac{13}{4}\,c + 2C_k ,
  \qquad\text{so}\qquad
  e^{-F_d(\xi)} \ge e^{-13c/4 - 2C_k}\,|\xi|^{-c} .
\]
\item \emph{(The exponent is logarithmic.)} For every $a > 0$,
$F_d(a\xi) - F_d(\xi) \to c\log a$ as $|\xi| \to \infty$. Equivalently, for a unit vector $e$ the
function $g(u) := F_d(\sqrt u\,e)$, $u \ge 0$, satisfies $g(au) - g(u) \to \tfrac c2\log a$ as
$u \to \infty$, for every $a > 0$.
\end{enumerate}
\statusT\quad\emph{(Cites nothing: the same integration by parts as in
\ref{thm:origin-boundedness-d}(1)'s proof, redone in the proof below, Frullani's integral in its monotone form, the Riemann--Lebesgue lemma on $(0,\infty)$ and
dominated convergence. The statement is about the exponent alone; where the datum's law $\mu$
is named, as in \ref{thm:origin-boundedness-d}, $\widehat\mu = e^{-F_d}$, and the law's existence is
entry A3(iii), which this lemma does not spend. \textbf{What it is worth.} Clause (1) is the half of
the transform's behaviour at infinity that \ref{thm:origin-boundedness-d}(4) needs. Clause (2) says
$e^{-F_d(\xi)} = |\xi|^{-c}L(|\xi|)$ with $L$
slowly varying at infinity and, by clause (1), bounded below. On the line the first step of
\sscite{sato1999levy}'s proof of Thm.~53.8, (53.33) on p.~412, has more: the transform is
asymptotic to an explicit constant times $|z|^{-c}K(|z|^{-1})$, $K \ge 1$ being the factor of
(53.25); clause (1) is the $d$-dimensional trace of $K \ge 1$, and Lemma~53.9, p.~411, is the
opposite bound, which \ref{thm:origin-boundedness-d}(1) replaces by a bound with a loss in the
exponent. For $c > 0$, clause (2) is the hypothesis of \ref{prop:log-orbit-rays} with
$\kappa = c/2$ for the orbit of the datum's law; that is the link between a profile bounded at the
origin and a logarithmic exponent which that node's annotation names, in the direction it needs, and
the link is Abelian, not Tauberian. It is spent in \ref{prop:bounded-profile-rays}. \textbf{Not
claimed}: the converse, that a logarithmic exponent comes from a profile bounded at the origin.
Until Q-0425 this annotation also listed as not claimed an explicit form of $L$, and whether $L$
is bounded above, that is whether $e^{-F_d} \asymp |\xi|^{-c}$. For $c > 0$ both are now
\ref{prop:isotropic-exponent-profile}: $L(\rho)$ is asymptotic to a constant times $K(1/\rho)$,
with $K(x) = \exp\int_x^1(c - k_d(u))\,du/u$ the factor of Sato's (53.25) written in the polar
profile, and $L$ is bounded above exactly when $\int_0^1(c - k_d(u))\,du/u < \infty$.)}
\end{lemma}

\begin{proof}
\leanok
Since the Gaussian coefficient is $0$, \ref{eq:isotropic-exponent} and \ref{eq:polar-decomposition}
give, with $\sigma_{d-1}$ the surface measure of $S^{d-1}$,
\[
  F_d(\xi) = \frac{1}{|S^{d-1}|}\int_{S^{d-1}}\varphi\bigl(\langle\xi,\eta\rangle\bigr)\,\sigma_{d-1}(d\eta),
  \qquad
  \varphi(s) := \int_0^\infty\bigl(1 - \cos(rs)\bigr)\,k_d(r)\,\frac{dr}{r} \quad (s \in \RR).
\]
$\varphi$ is even and finite, because $1 - \cos(rs) \le \min\{2, r^2s^2/2\}$ and the two
integrability conditions of \ref{def:isotropic-origin-datum} hold. As $k_d$ is nonincreasing with
right limit $c$ at the origin, $k_d \le c$ on $(0,\infty)$. And $k_d(r) \to 0$ as $r \to \infty$,
since a nonincreasing $k_d \ge \ell > 0$ would make $\int_1^\infty k_d(r)\,dr/r$ diverge.

\emph{The logarithmic bounds on $\Cin$.} Write $\Cin(X) = \int_0^X(1-\cos s)\,s^{-1}ds$ for
$X \ge 0$, a nondecreasing function. For $X \ge 1$,
$\Cin(X) = \Cin(1) + \log X - \int_1^X s^{-1}\cos s\,ds$. Here $0 \le \Cin(1) \le \int_0^1 \tfrac s2\,ds
= \tfrac14$, because $1 - \cos s \le s^2/2$. Integrating by parts,
$\bigl|\int_1^X s^{-1}\cos s\,ds\bigr| = \bigl|[s^{-1}\sin s]_1^X + \int_1^X s^{-2}\sin s\,ds\bigr|
\le 1 + \sin 1 + 1 \le 3$. For $0 < X < 1$ the lower bound below holds because its right side is
negative. Hence
\begin{equation}\label{eq:cin-log}
  \log X - 3 \le \Cin(X) \quad (X > 0),
  \qquad
  \Cin(X) \le \log X + \tfrac{13}{4} \quad (X \ge 1).
\end{equation}

\emph{Clause (1).} Write $\rho = |\xi| \ge 1$. For $|s| \le \rho$, split $\varphi(s)$ at $r = 1$.
On $(0,1)$ use $k_d \le c$, and on $[1,\infty)$ use $1 - \cos \le 2$:
\[
  \varphi(s) \le c\,\Cin(|s|) + 2C_k \le c\,\Cin(\rho) + 2C_k .
\]
On the sphere $|\langle\xi,\eta\rangle| \le \rho$, and the normalized surface measure has mass $1$.
So $F_d(\xi) \le c\,\Cin(\rho) + 2C_k$, and the upper bound of \ref{eq:cin-log} gives the bound on
$F_d$. The bound on $e^{-F_d}$ follows.

\emph{Clause (2), one variable.} Let $a \ge 1$ and $s > 0$. Substituting $u = ar$ in
$\varphi(as)$,
\[
  \varphi(as) - \varphi(s) = \int_0^\infty \bigl(1 - \cos(us)\bigr)\,h_a(u)\,du ,
  \qquad h_a(u) := \frac{k_d(u/a) - k_d(u)}{u} \ge 0 ,
\]
where both terms on the left are finite. For $0 < \varepsilon < N$,
\[
  \int_\varepsilon^N h_a(u)\,du = \int_{\varepsilon/a}^{N/a}k_d(v)\,\frac{dv}{v}
     - \int_\varepsilon^N k_d(v)\,\frac{dv}{v}
   = \int_{\varepsilon/a}^{\varepsilon}k_d(v)\,\frac{dv}{v} - \int_{N/a}^{N}k_d(v)\,\frac{dv}{v} ,
\]
and every integral here is finite because $k_d \le c$. The first term lies between
$k_d(\varepsilon)\log a$ and $c\log a$, and tends to $c\log a$ as $\varepsilon \downarrow 0$. The
second lies between $0$ and $k_d(N/a)\log a$, and tends to $0$ as $N \to \infty$. By monotone
convergence $\int_0^\infty h_a = c\log a$, so $h_a$ is integrable on $(0,\infty)$, and the
Riemann--Lebesgue lemma gives $\int_0^\infty\cos(us)\,h_a(u)\,du \to 0$ as $s \to \infty$. Hence
\[
  \varphi(as) - \varphi(s) \to c\log a \quad (s \to \infty),
  \qquad 0 \le \varphi(as) - \varphi(s) \le 2c\log a \quad (s > 0).
\]
For $0 < a < 1$ apply this with $1/a$ at the point $as$, which tends to infinity with $s$:
$\varphi(as) - \varphi(s) = -\bigl(\varphi(a^{-1}\cdot as) - \varphi(as)\bigr) \to c\log a$, and
$|\varphi(as) - \varphi(s)| \le 2c|\log a|$.

\emph{Clause (2), the sphere average.} Write $\xi = \rho\widehat\xi$ with $|\widehat\xi| = 1$.
Since $\varphi$ is even,
\[
  F_d(a\xi) - F_d(\xi) = \frac{1}{|S^{d-1}|}\int_{S^{d-1}}
    \Bigl[\varphi\bigl(a\rho|\langle\widehat\xi,\eta\rangle|\bigr)
      - \varphi\bigl(\rho|\langle\widehat\xi,\eta\rangle|\bigr)\Bigr]
    \sigma_{d-1}(d\eta) .
\]
For every $\eta$ with $\langle\widehat\xi,\eta\rangle \ne 0$ the bracket tends to $c\log a$ as
$\rho \to \infty$. The remaining $\eta$ form a $\sigma_{d-1}$-null set: a great subsphere when
$d \ge 2$, and no point at all when $d = 1$. The bracket is bounded by $2c|\log a|$. By dominated
convergence $F_d(a\rho\widehat\xi) - F_d(\rho\widehat\xi) \to c\log a$ as $\rho \to \infty$. The
left side does not depend on $\widehat\xi$, because $F_d$ is rotation invariant, so the limit holds
as $|\xi| \to \infty$. For the equivalent form, $g(au) - g(u) = F_d(\sqrt a\,\sqrt u\,e) -
F_d(\sqrt u\,e) \to c\log\sqrt a = \tfrac c2\log a$.
\end{proof}

% draft: new material. Provenance: P-0013, the conjecture accepted 2026-09-22 as Q-0043. The
% statement is the one P-0013 asked for -- "an isotropic admissible kernel with no Gaussian part
% is bounded at the origin exactly when its polar profile at the origin exceeds d" -- with the
% two halves separated, because the source reading showed they were not equally available: the
% sufficiency half is proved in every dimension with no extra hypothesis, the necessity half was
% proved on the subordinated slice and assumed (H) outside it. Since Q-0151 (2026-10-03) the
% necessity half is clause (4) for every datum, by lem:positive-transform-origin; clause (3) stays
% for its explicit monotone density and the exact limit at the origin.
%
% Since Q-0286 no clause needs A2: the old clause (5), the polar reading of the data's laws,
% duplicated prop:isotropic-polar-reading(1) and was dropped after its citers were re-pointed
% there; clause (5) is now the dichotomy, which closed clause (4) until then.
% CHANGED Q-0281: typed (the statement skeleton). Clause (1) is split to lem:origin-decay-exponent and
% clause (3) to prop:origin-boundedness-subordinated (ADR-0025; records/formalization/PRICING.md: the
% clauses differ in cost and in trust, (1) and (3) citing nothing, (2) and (5) spending A1); the
% \lean tag names the children's declarations too. The law is any probability measure with transform
% e^{-F_d} (PRICING finding 1), so A3(iii) is not spent by the Lean statements. "c in (0,infty]" is
% read as the hypothesis c > 0 (clause (3) is false for the zero datum, whose law delta_0 has no
% density). Clause (3) is NARROWED: the delay law is a hypothesis of the clause (PRICING finding 2),
% its existence, [B, Rem. 4.3], resting on Bernstein's theorem, which this module's ledger does not
% carry; the clause's text and its proof are changed accordingly, nothing else.
% CHANGED Q-0281: \uses re-pointed to the split child it reads (eq:polar-decomposition).
% CHANGED Q-0281: blind review R2-2: the tag names measurable_gaussVar, the measurability of the
% Gaussian kernel that clause (3)'s mixture (a Measure.bind in Lean) presupposes.
% CHANGED Q-0281 wave 1: blind check (the survey's cleanup, filed with KS-1): in the annotation's
% reading for self-decomposable laws, "hence with no zeros" now goes through prop:gw-ray-families(1),
% which spends nothing; the reading still spends A2, A3 and A4 inside prop:isotropic-polar-reading(1).
% CHANGED Q-0281 wave 2: blind review W2A-5: the Assignment's sentence on clause (3) names the nodes
% its rewritten proof reads (lem:isotropic-exponent-log(1), lem:positive-transform-origin(2),
% lem:delay-polar-profile), where it said "Tonelli and one Gamma integral" with [B, Prop. 4.9(2)].
% CHANGED Q-0281 wave 2: proved (AffineHemigroup.OriginBoundedness, with the children; guarded), every
% declaration of the tag now proved. Clause (2) (origin_density_smooth, origin_density_bounded_smooth)
% was the last: it spends A1, which was restated at its page's letter on 2026-10-05, the density of a
% probability measure being nonnegative (the author's decision; blueprint/AXIOMS.md A1). The text of
% the statement is unchanged; the proof of clause (2) is completed on the Lean route.
% CHANGED Q-0381 (2026-10-10; the fidelity review's paper-facing points 3 and 6, R17): three
% sentences of the statement reworded so that each has one reading, each to what the Lean states
% and the proof of record already proves; no proof changes. (i) The preamble's "put c := k_d(0+)
% in (0,infty]" becomes "put c := k_d(0+) in [0,infty] and suppose c > 0": c > 0 is a hypothesis
% (hc : 0 < D.profileAtOrigin in every declaration of the tag), which a blind restater took for a
% consequence. (ii) Clause (3) states "W > 0 almost surely", which eq:mixture-density needs, the
% proof derives ("So W > 0 almost surely") and the Lean proves (delayLaw_zero : rho {0} = 0 under
% the clause's hypotheses, with rho (Iio 0) = 0); the \lean tag gains that declaration, which
% prints Lean core. (iii) Clause (5) says what "a density bounded near the origin" is: some density
% q of mu with q <= M on a ball {|x| < delta} (origin_dichotomy's own form), not "every density
% essentially bounded there" (equivalent, by changing q on a null set). The children
% lem:origin-decay-exponent and prop:origin-boundedness-subordinated take the same preamble and
% clause. Owed to the Lean session: a guard line for delayLaw_zero in CIAxiomGuard.lean.
\begin{theorem}[the boundedness threshold in dimension $d$ is the dimension]
\label{thm:origin-boundedness-d}
\uses{def:isotropic-origin-datum,lem:positive-transform-origin,lem:isotropic-exponent-log,lem:origin-decay-exponent,prop:origin-boundedness-subordinated,lem:polar-decomposition,lem:delay-polar-profile,thm:cited-A1,thm:cited-A3}
\lean{AffineHemigroup.origin_decay_exponent, AffineHemigroup.origin_density_smooth,
AffineHemigroup.origin_density_bounded_smooth, AffineHemigroup.subordinated_profileAtOrigin,
AffineHemigroup.subordinated_law, AffineHemigroup.delayLaw_zero,
AffineHemigroup.subordinated_density_limit,
AffineHemigroup.subordinated_unbounded, AffineHemigroup.origin_unbounded,
AffineHemigroup.origin_unbounded_rate, AffineHemigroup.origin_dichotomy,
AffineHemigroup.measurable_gaussVar}\leanok
Let $d \ge 1$, let $(0, k_d)$ be an isotropic origin datum in dimension $d$ with Gaussian
coefficient $0$, let $\mu$ be its law, put $c := k_d(0{+}) \in [0,\infty]$ and suppose $c > 0$.
\begin{enumerate}
\item \emph{(The profile at the origin is a decay exponent.)} For every finite $\beta < c$ there
is $C_\beta < \infty$ with
\[
  \bigl|\widehat\mu(\xi)\bigr| \le C_\beta\,(1+|\xi|)^{-\beta}, \qquad \xi \in \RR^d .
\]
\item \emph{(Above the threshold: bounded, and a ladder of derivatives.)} Let $N \ge 0$ be an
integer. If $c > d + N$ then $\mu$ has a density of class $C^N$ on the whole of $\RR^d$, all of
whose partial derivatives of orders $0,\dots,N$ tend to $0$ at infinity. In particular $c > d$
makes the density continuous and bounded, and $c = \infty$ makes it $C^\infty$.
\item \emph{(At and below the threshold, on the subordinated slice: unbounded.)} Suppose the
datum is subordinated, with delay datum $(0,k_{\mathrm I})$ of exponent $F_{\mathrm I}$, that $W$
is a random delay whose law, a probability measure on $[0,\infty)$, has Laplace transform
$e^{-F_{\mathrm I}}$, and that $Z$ is standard normal on $\RR^d$, independent of $W$. Then
$c = 2\,k_{\mathrm I}(0{+})$, $W > 0$ almost surely, $\mu$ is the law of $\sqrt{W}Z$, and $\mu$
has the radial density
\begin{equation}\label{eq:mixture-density}
  p(x) = \int_0^\infty (2\pi u)^{-d/2}e^{-|x|^2/2u}\,\Pr(W \in du) .
\end{equation}
This $p$ is nonincreasing in $|x|$, and
\[
  \lim_{x \to 0} p(x) = (2\pi)^{-d}\int_{\RR^d}\widehat\mu(\xi)\,d\xi
   = \frac{(2\pi)^{-d/2}}{\Gamma(d/2)}\int_0^\infty \sigma^{d/2-1}e^{-F_{\mathrm I}(\sigma)}\,d\sigma ,
\]
finite exactly when $c > d$. If $c \le d$, then for every representative $q$ of $\mu$'s
density, every $\delta > 0$ and every $M$ there is an $x$ with $0 < |x| < \delta$ and
$q(x) > M$: no representative is bounded near the origin.
\item \emph{(At and below the threshold, for every datum: unbounded.)} Suppose $c \le d$. Then
$\int_{\RR^d}\widehat\mu(\xi)\,d\xi = \infty$, and there are no $\delta > 0$ and $M < \infty$ with
$\mu(B) \le M|B|$ for every Borel set $B \subseteq \{|x| < \delta\}$. In particular, for every
representative $q$ of a density of $\mu$, every $\delta > 0$ and every $M$ there is an $x$ with
$0 < |x| < \delta$ and $q(x) > M$. Quantitatively, with $g_t$ the heat kernel of
\ref{lem:positive-transform-origin}, there is $C > 0$ with
\[
  \int_{\RR^d} g_t\,d\mu \ge C\,t^{-(d-c)/2} \quad (c < d), \qquad
  \int_{\RR^d} g_t\,d\mu \ge C\,\log\frac1t \quad (c = d), \qquad 0 < t \le 1 ,
\]
and for every $\beta$ with $0 < \beta < c$ there is $C'_\beta$ with
$\int g_t\,d\mu \le C'_\beta\,t^{-(d-\beta)/2}$ for all $t > 0$. And
$\log\int g_t\,d\mu / \log(1/t) \to (d-c)/2$ as $t \downarrow 0$.
\item \emph{(The dichotomy.)} The law $\mu$ has a density bounded near the origin, that is, a
density $q$ with $q(x) \le M$ whenever $|x| < \delta$, for some $\delta > 0$ and $M < \infty$,
exactly when $c > d$.
\end{enumerate}
\statusA\quad\ledger{A1}\quad\ledger{A3}\quad\emph{(\textbf{Assignment.} Every clause
speaks of the datum's law $\mu$, whose existence is \ledger{A3}(iii), the converse half of the
\Levy--Khintchine representation. Clause (2) is the only clause that cites analysis: it is
\ledger{A1}, \sscite{sato1999levy}~Prop.~28.1, p.~190, the $\RR^d$ Fourier-integrability
criterion, applied to clause (1)'s bound, and it carries exactly what that entry carries and
nothing about the origin. Clauses (1), (3) and (4) cite nothing beyond the law's existence:
clause (1) is the sphere average of an elementary $\Cin$ inequality, clause (3) is Tonelli and one
Gamma integral, with clause (1), the lower bound of \ref{lem:isotropic-exponent-log}(1),
\ref{lem:positive-transform-origin}(2) and \ref{lem:delay-polar-profile} (which is
[B, Prop.~4.9(2)] restated), and the delay law a hypothesis (its
existence, [B, Rem.~4.3], rests on Bernstein's theorem, which no entry of this module's ledger
carries; Q-0281 made it a hypothesis), and clause (4) is
\ref{lem:positive-transform-origin} applied to the positive transform $e^{-F_d}$ with the lower
bound of \ref{lem:isotropic-exponent-log}(1). Clause (5) joins clauses (2) and (4), so it spends
\ledger{A1} through clause (2). \textbf{The reading for self-decomposable laws} is not a clause of
this theorem: that the laws of the data with $a = 0$ are exactly the isotropic self-decomposable
laws on $\RR^d$ with no Gaussian part is \ref{prop:isotropic-polar-reading}(1), on \ledger{A2},
\ledger{A3} and \ledger{A4}, and with it clauses (2)--(5) say that such a law is bounded near the
origin exactly when its polar profile exceeds $d$ there. The law is
infinitely divisible by \ledger{A4}, so by \ledger{A3}(i) its transform is an exponential, real
because the law is symmetric; it is therefore the law of a datum by that clause,
and the uniqueness of \ledger{A3}(ii) makes the datum's \Levy{} measure its own and identifies
``no Gaussian part'' with $a = 0$. The profile at the origin is not changed by modifying the polar
profile on a null set, since $k_d(0{+})$ is the right limit at $0$ of a nonincreasing function and two
such functions equal almost everywhere have the same right limits. \textbf{History.} Until Q-0286 the
reading was this theorem's clause (5), and the dichotomy was the last sentence of clause (4); the
old clause (5) duplicated \ref{prop:isotropic-polar-reading}(1), with its proof, and was dropped
after its citers (\ref{thm:filtration-strict}, \ref{thm:no-cin-ray-in-the-plane} and
\ref{def:isotropic-origin-datum}'s annotation) were re-pointed to that clause, so the theorem no
longer spends \ledger{A2} or \ledger{A4}. The last sentence of clause (3), the dichotomy on the
subordinated slice, was dropped with it, being clause (5) restricted to that slice. Until Q-0151
(2026-10-03) clause (4) assumed a hypothesis (H), a pointwise radial lower bound on the density
near the origin, because \texttt{records/affine/READING-sato1999levy-dimension-d.md} \S 4 found
that \ledger{A1} is one-way when the transform changes sign; the transform of a datum's law never
does, and the clause holds for every datum. (H) was then withdrawn from the blueprint (the author,
2026-10-03). \textbf{Not claimed}: any pointwise asymptotic. Clause (2) gives no rate of
convergence at infinity, and clauses (3) and (4) give no pointwise rate of divergence at the
origin: clause (4)'s rate is that of the density at the origin of $\mu$ blurred by the heat kernel,
$\int g_t\,d\mu$, and the $d$-dimensional analogue of \sscite{sato1999levy}~(53.28) with its slowly
varying factor $K$, which is what (H) asks for pointwise, is neither used nor asserted. Nor is
$\mu$ claimed to have a density when $c \le d$: clause (4) is a statement about the measure, and
clause (5) reads ``has a density bounded near the origin'' in that sense. Nor anything about the
$\Cin$ rays for $d \ge 2$ beyond \ref{thm:no-cin-ray-in-the-plane}: their unit members have $c = 1$
at $d = 1$ and sit at that dimension's threshold [B, Cor.~3.9(2)], they are not subordinated
[B, Prop.~4.9(3)], and they are not in $\mathcal{A}_d$ for any $d \ge 2$.
\textbf{What the clauses are worth.} At $d = 1$ clause (2) is
[B, Cor.~3.9(1)]'s smooth regime with its $C^{N-1}$ read as $C^N$ at $N < c-1$, and clauses (4)
and (5) are that corollary's singular and threshold regimes for every datum, as a dichotomy and
without a pointwise rate (the proof of that corollary reads one off
\sscite{sato1999levy}~Thm.~53.8, through [B]'s behaviour at the origin); clause (3) is their
restriction to the subordinated slice, with an explicit monotone density. In $\RR^d$ the
dichotomy needs no hypothesis because the laws are symmetric and infinitely divisible, so
$\widehat\mu = e^{-F_d} > 0$: the sign that Prop.~28.1 lacks for a general law is supplied by the
structure of the datum.
The \Matern{} member of index $\gamma$ has $k_{\mathrm I}(0{+}) = \gamma$, hence $c = 2\gamma$
\textbf{in every dimension}, and clause (3) puts its threshold at $\gamma > d/2$, which is the
smoothness parameter range $\gamma - d/2 > 0$ of \sscite{porcu2024matern} and
\sscite{guttorp2006studies} (who write $\nu$ for it): the condition under which the \Matern{}
covariance in $\RR^d$ is valid and the boundedness of the aperture at its own centre are the same
inequality (\texttt{records/affine/READING-porcu2024matern.md},
\texttt{READING-guttorp2006studies.md}). \textbf{Why the threshold moves with the dimension while
the quantity does not.} Clause (1) says $c$ is a decay exponent for $|\widehat\mu|$; clause (2)
says a decay exponent buys a bounded density once it exceeds the dimension of the transform
integral. The origin quantity is dimension-free (clause (3) gives $k_d(0{+}) = 2k_{\mathrm I}(0{+})$
for every $d$), and the dimension enters through $d\xi$ alone. The sentence [B, \S 8] left open,
that in the plane an admissible aperture may be unbounded at its own centre, is clause (3) at
$d = 2$ with $c \le 2$: the \Matern{} member of index $\gamma \le 1$ is such an aperture.)}
\end{theorem}

\begin{proof}
\leanok
% CHANGED Q-0281 wave 2: proof of clause (1) rewritten to the Lean route (lem:origin-decay-exponent,
% AffineHemigroup.origin_decay_exponent). The previous route bounded the sphere average J_d of
% log|<xi-hat, eta>| through the density (1-u^2)^{(d-3)/2} of the image of the surface measure; the
% checked route reads only caps {|<e_1, eta>| >= eps}, whose normalized measure tends to 1 because a
% great subsphere is null, at the price of a slightly smaller exponent beta < beta_1 < c.
\emph{Clause (1).} Write $\rho = |\xi|$, and let $\Cin$ be as in the proof of
\ref{lem:isotropic-exponent-log}. Since $a = 0$, \ref{eq:isotropic-exponent} and
\ref{eq:polar-decomposition} give
\[
  -\log\bigl|\widehat\mu(\xi)\bigr| = F_d(\xi)
   = \frac{1}{|S^{d-1}|}\int_{S^{d-1}}\varphi\bigl(\langle\xi,\eta\rangle\bigr)\,\sigma_{d-1}(d\eta),
  \qquad
  \varphi(s) := \int_0^\infty\bigl(1 - \cos(rs)\bigr)\,k_d(r)\,\frac{dr}{r} .
\]
$F_d$ is radial: if $|\xi| = |\xi'|$, the reflection $Q$ in $(\xi - \xi')^\perp$ carries $\xi$ to
$\xi'$, preserves Lebesgue measure and the radial density of $\nu_d$, so substituting $x = Qy$ in
\ref{eq:isotropic-exponent} gives $F_d(\xi') = F_d(\xi)$. It is therefore enough to bound
$F_d(\rho e_1)$. Fix a finite $\beta < c$. If $\beta \le 0$ the bound holds with $C_\beta = 1$,
since $|\widehat\mu| \le 1 \le (1+|\xi|)^{-\beta}$. Let $\beta > 0$ and choose $\beta_1$ with
$\beta < \beta_1 < c$. As $k_d$ is nonincreasing with right limit $c$ at the origin, there is
$r_\beta \in (0,1]$ with $k_d \ge \beta_1$ on $(0,r_\beta)$. The integrand of $\varphi$ is
nonnegative, so discarding $(r_\beta,\infty)$ and substituting gives
$\varphi(s) \ge \beta_1\,\Cin(r_\beta|s|)$.

For $\varepsilon > 0$ let $K_\varepsilon := \{\eta \in S^{d-1} : |\langle e_1,\eta\rangle| \ge
\varepsilon\}$. As $\varepsilon \downarrow 0$ these caps increase to the complement of the great
subsphere $\{\langle e_1,\eta\rangle = 0\}$, which is $\sigma_{d-1}$-null: $\sigma_{d-1}(E)$ is $d$
times the Lebesgue measure of the cone $(0,1)\cdot E$, and the cone over the great subsphere lies in
the hyperplane $e_1^\perp$. So there is $\varepsilon \in (0,1]$ with
$p := \sigma_{d-1}(K_\varepsilon)/|S^{d-1}| \ge \beta/\beta_1$. Put
$L := 3 - \log(\varepsilon r_\beta) \ge 0$. For $\eta \in K_\varepsilon$,
$r_\beta|\langle\rho e_1,\eta\rangle| \ge \rho\varepsilon r_\beta$, so the lower bound of
\ref{eq:cin-log} gives $\varphi(\langle\rho e_1,\eta\rangle) \ge \beta_1(\log\rho - L)$. Discarding
the rest of the sphere, $F_d(\rho e_1) \ge p\,\beta_1(\log\rho - L) \ge \beta(\log\rho - L)$ when
$\log\rho \ge L$, and the outer inequality holds as well when $\log\rho < L$, because $F_d \ge 0$.
Hence $|\widehat\mu(\xi)| \le e^{\beta L}\rho^{-\beta}$ for every $\xi \ne 0$. For $\rho \ge 1$,
$1 + \rho \le 2\rho$ gives $\rho^{-\beta} \le 2^\beta(1+\rho)^{-\beta}$; for $\rho < 1$,
$|\widehat\mu(\xi)| \le 1 \le 2^\beta(1+\rho)^{-\beta}$. This is the stated bound with
$C_\beta = 2^\beta(1 + e^{\beta L})$.

\emph{Clause (2).} Let $c > d+N$ and choose a finite $\beta$ with $d + N < \beta < c$. By clause
(1),
\[
  \int_{\RR^d}\bigl|\widehat\mu(\xi)\bigr|\,|\xi|^N\,d\xi
   \le C_\beta\int_{\RR^d}(1+|\xi|)^{-\beta}|\xi|^N\,d\xi < \infty .
\]
The integral converges at infinity because $\beta - N > d$, and at the origin because $N \ge 0$.
That is the hypothesis of \ledger{A1}, whose conclusion is the clause: a density $g \ge 0$ of
class $C^N$ whose partial derivatives of orders $0,\dots,N$ tend to $0$ at infinity.
% CHANGED Q-0281 wave 2: the "in particular" argued on the Lean route (origin_density_bounded_smooth).
At $N = 0$ the density is continuous and tends to $0$ at infinity, so it is bounded: it is at most
$1$ outside a compact set, and continuous on that set. For $c = \infty$ every $N$ is admissible,
and the densities obtained for different $N$ coincide: two continuous densities of $\mu$ agree
almost everywhere, hence everywhere, since a nonempty open set has positive Lebesgue measure. So the
density at $N = 0$ is of class $C^N$ for every $N$, that is, of class $C^\infty$.

\emph{Clause (3).} First, $c = 2\,k_{\mathrm I}(0{+})$. As $r \downarrow 0$ the integrand of
\ref{eq:polar-profile-d} increases to $v^{d/2-1}e^{-v}k_{\mathrm I}(0{+})$, because $k_{\mathrm I}$
is nonincreasing. So monotone convergence gives
$k_d(0{+}) = 2\,\Gamma(d/2)^{-1}\Gamma(d/2)\,k_{\mathrm I}(0{+})$. Next, \prog{B}{Prop.~4.9(2)},
restated, which is \ref{lem:delay-polar-profile} for the delay datum $(0,k_{\mathrm I})$: the datum's exponent is $F_d(\xi) = F_{\mathrm I}(|\xi|^2/2)$. The law of $W$ has Laplace
transform $e^{-F_{\mathrm I}}$ by hypothesis. So
$\widehat\mu(\xi) = e^{-F_{\mathrm I}(|\xi|^2/2)} = \int_0^\infty e^{-u|\xi|^2/2}\Pr(W \in du)$,
which is the transform of the law of $\sqrt{W}Z$, and $\mu$ is that law. Since
$k_{\mathrm I}(0{+}) = c/2 > 0$, clause (1) read along a ray, at $|\xi| = \sqrt{2\sigma}$, gives
$e^{-F_{\mathrm I}(\sigma)} \to 0$, so
$\Pr(W = 0) = \lim_{\sigma\to\infty}e^{-F_{\mathrm I}(\sigma)} = 0$. So $W > 0$ almost surely,
and conditioning on $W$ gives \ref{eq:mixture-density}. For $x \ne 0$ the integrand is bounded in
$u$, so $p(x) < \infty$. Each integrand is radial and nonincreasing in $|x|$, and so is $p$.
Monotone convergence as $x \to 0$ gives
\[
  p_0 := \lim_{x\to0}p(x) = \int_0^\infty (2\pi u)^{-d/2}\Pr(W \in du) \in (0,\infty] .
\]
The two other forms of $p_0$ in the clause follow from Tonelli. Since $\widehat\mu \ge 0$, its
integral is defined in $[0,\infty]$, and $\int_{\RR^d}e^{-u|\xi|^2/2}d\xi = (2\pi/u)^{d/2}$ gives
\[
  (2\pi)^{-d}\int_{\RR^d}\widehat\mu(\xi)\,d\xi
   = (2\pi)^{-d}\int_0^\infty\Bigl(\frac{2\pi}{u}\Bigr)^{d/2}\Pr(W \in du) = p_0 .
\]
And $\int_0^\infty\sigma^{d/2-1}e^{-\sigma u}d\sigma = \Gamma(d/2)u^{-d/2}$ gives
\[
  \int_0^\infty\sigma^{d/2-1}e^{-F_{\mathrm I}(\sigma)}\,d\sigma
   = \Gamma(d/2)\int_0^\infty u^{-d/2}\Pr(W \in du) = \Gamma(d/2)\,(2\pi)^{d/2}\,p_0 .
\]

% CHANGED Q-0281 wave 2: the last two paragraphs of clause (3) rewritten to the Lean route
% (prop:origin-boundedness-subordinated, AffineHemigroup.Subordinated). The previous route decided
% the finiteness of p_0 through bounds on F_I and compared measures with the monotone representative
% p; the checked route reads the whole transform, as clause (4) does for every datum.
It remains to decide when $p_0$ is finite, and the argument reads the transform $\widehat\mu$
itself. If $c > d$, pick a finite $\beta \in (d,c)$. By clause (1),
$\widehat\mu(\xi) \le C_\beta(1+|\xi|)^{-\beta}$, which is integrable over $\RR^d$ because
$\beta > d$; so $p_0 < \infty$. If $c \le d$, then $c$ is finite, and
\ref{lem:isotropic-exponent-log}(1) gives $\widehat\mu(\xi) \ge C_0|\xi|^{-c}$ for $|\xi| \ge 1$,
with $C_0 = e^{-13c/4 - 2C_k}$. In polar coordinates
$\int_{\RR^d}\widehat\mu \ge C_0|S^{d-1}|\int_1^\infty\rho^{d-1-c}\,d\rho = \infty$, because
$d - 1 - c \ge -1$. So $p_0 = \infty$, the threshold $c = d$ included.

That makes \emph{every} representative unbounded near the origin, by comparing measures rather than
functions, as in \restated{B}{Cor.~3.9}{cor:line-threshold}. If a representative $q$ of $\mu$'s density had $q \le M$ on
$\{0 < |x| < \delta\}$, then $\mu(B) = \int_B q \le M|B|$ for every Borel
$B \subseteq \{|x| < \delta\}$, the origin being a Lebesgue-null set, and
\ref{lem:positive-transform-origin}(2) would give $\int_{\RR^d}\widehat\mu \le (2\pi)^dM < \infty$,
a contradiction.

\emph{Clause (4).} By \ref{def:isotropic-origin-datum}, $\widehat\mu = e^{-F_d}$ with $F_d$ real
and finite, so $\widehat\mu > 0$ everywhere and \ref{lem:positive-transform-origin} applies. By
\ref{lem:isotropic-exponent-log}(1), $\widehat\mu(\xi) \ge C_0|\xi|^{-c}$ for $|\xi| \ge 1$, with
$C_0 = e^{-13c/4 - 2C_k} > 0$ and $C_k$ as in that lemma. In polar coordinates,
\[
  \int_{\RR^d}\widehat\mu(\xi)\,d\xi \ge C_0\,|S^{d-1}|\int_1^\infty \rho^{d-1-c}\,d\rho = \infty ,
\]
because $d - 1 - c \ge -1$. If there were $\delta$ and $M$ with $\mu(B) \le M|B|$ for every Borel
$B \subseteq \{|x| < \delta\}$, \ref{lem:positive-transform-origin}(2) would make this integral at
most $(2\pi)^dM$. So there are none. A representative $q$ of a density with $q \le M$ on
$\{0 < |x| < \delta\}$ would give such $\delta$ and $M$, by the last sentence of
\ref{lem:positive-transform-origin}(2). So there is none.

For the rate, \ref{lem:positive-transform-origin}(1) and the same lower bound give, for
$0 < t \le 1$,
\[
  \int g_t\,d\mu \ge (2\pi)^{-d}C_0\,|S^{d-1}|\int_1^\infty \rho^{d-1-c}e^{-t\rho^2/2}\,d\rho .
\]
If $c < d$, the substitution $\rho = \sigma/\sqrt t$ turns the last integral into
$t^{-(d-c)/2}\int_{\sqrt t}^\infty\sigma^{d-1-c}e^{-\sigma^2/2}\,d\sigma \ge
t^{-(d-c)/2}\int_1^\infty\sigma^{d-1-c}e^{-\sigma^2/2}\,d\sigma$, a positive multiple of
$t^{-(d-c)/2}$. If $c = d$, it is at least
$e^{-1/2}\int_1^{t^{-1/2}}\rho^{-1}\,d\rho = \tfrac12 e^{-1/2}\log(1/t)$, since
$t\rho^2 \le 1$ on that range. For the upper bound let $0 < \beta < c$. Clause (1) and
$(1+|\xi|)^{-\beta} \le |\xi|^{-\beta}$ give, with $\xi = y/\sqrt t$,
\[
  \int g_t\,d\mu \le (2\pi)^{-d}C_\beta\int_{\RR^d}|\xi|^{-\beta}e^{-t|\xi|^2/2}\,d\xi
   = (2\pi)^{-d}C_\beta\,t^{-(d-\beta)/2}\int_{\RR^d}|y|^{-\beta}e^{-|y|^2/2}\,dy ,
\]
and the last integral is finite because $\beta < c \le d$. Now take logarithms and divide by
$\log(1/t)$. By the lower bound, the quotient $\log\int g_t\,d\mu/\log(1/t)$ is at least
$(d-c)/2 - \varepsilon$ for every $\varepsilon > 0$ once $t$ is small enough. At $c = d$ this
uses $\log\log(1/t)/\log(1/t) \to 0$. By the upper bound, the quotient is at most
$(d-\beta)/2 + \varepsilon$ for every $\beta < c$ and $\varepsilon > 0$ once $t$ is small enough. So it
tends to $(d-c)/2$.

\emph{Clause (5).} If $c > d$, clause (2) at $N = 0$ gives $\mu$ a continuous density that tends
to $0$ at infinity, hence a bounded one. If $c \le d$, clause (4) excludes every bound
$\mu(B) \le M|B|$ on the Borel subsets $B$ of a ball $\{|x| < \delta\}$, and with it every
representative of a density bounded near the origin.
\end{proof}

% CHANGED Q-0281: split from thm:origin-boundedness-d (ADR-0025; records/formalization/PRICING.md: the
% clauses differ in cost and in trust): the theorem's preamble and its clause (1), copied verbatim,
% with the parent's clause number. It cites nothing, where the parent spends A1 in (2) and (5); the
% parent \uses this node, and its text and numbering are unchanged.
% CHANGED Q-0381 (2026-10-10): the preamble follows the parent's ("suppose c > 0", the hypothesis
% hc of the declaration, said outright; the fidelity review's paper-facing point 3).
% CHANGED Q-0409 (2026-10-11): the \lean tag gains OriginDatum.exponent_decay, the same bound stated
% for e^{-F_d}, with no law and without c > 0. It was the body of origin_decay_exponent's proof,
% which reads the law only through |mu-hat| = e^{-F_d}; split off so that
% cor:origin-profile-exponent reads the bound on the exponent. The text is unchanged.
\begin{lemma}[the profile at the origin is a decay exponent]
\label{lem:origin-decay-exponent}
\uses{def:isotropic-origin-datum,lem:polar-decomposition,lem:isotropic-exponent-log}
\lean{AffineHemigroup.origin_decay_exponent, AffineHemigroup.OriginDatum.exponent_decay}\leanok
Let $d \ge 1$, let $(0, k_d)$ be an isotropic origin datum in dimension $d$ with Gaussian
coefficient $0$, let $\mu$ be its law, put $c := k_d(0{+}) \in [0,\infty]$ and suppose $c > 0$.
\begin{enumerate}
\item[(1)] For every finite $\beta < c$ there is $C_\beta < \infty$ with
\[
  \bigl|\widehat\mu(\xi)\bigr| \le C_\beta\,(1+|\xi|)^{-\beta}, \qquad \xi \in \RR^d .
\]
\end{enumerate}
\statusT\quad\emph{(Cites nothing: the sphere average of an elementary $\Cin$ inequality. The law is
any probability measure with transform $e^{-F_d}$, so its existence, ledger entry A3(iii), is not
read.
Clause (1) of \ref{thm:origin-boundedness-d}, split off on 2026-10-04 (Q-0281, ADR-0025), with the
parent's clause number; the parent's clauses (2)--(5) read it.)}
\end{lemma}

\begin{proof}
\leanok
This is the part of the proof of \ref{thm:origin-boundedness-d} headed clause (1).
\end{proof}

% CHANGED Q-0281: split from thm:origin-boundedness-d (ADR-0025; records/formalization/PRICING.md: the
% clauses differ in cost and in trust): the theorem's preamble and its clause (3), copied from the
% parent with the parent's clause number. NARROWED with the parent (PRICING finding 2): the delay law
% is a hypothesis, "let W be a random delay with the delay law [B, Rem. 4.3]" becoming "that W is a
% random delay whose law ... has Laplace transform e^{-F_I}"; the existence of such a law is
% Bernstein's theorem, which this module's ledger does not carry (Paper V admits it as its A11,
% SpatialLine.bernstein_completely_monotone, and derives no delay law from it). The parent \uses this
% node.
% CHANGED Q-0281: blind review R2-2: the tag names measurable_gaussVar, which the mixture "the law of
% sqrt(W) Z" (a Measure.bind in Lean) presupposes.
% CHANGED Q-0281 wave 2: blind review W2A-5: the status annotation names the nodes the rewritten
% proof reads (lem:isotropic-exponent-log(1), lem:positive-transform-origin(2),
% lem:delay-polar-profile), where it said "Tonelli and one Gamma integral, with clause (1)".
% CHANGED Q-0381 (2026-10-10): the preamble and the clause follow the parent's ("suppose c > 0";
% "W > 0 almost surely", which is delayLaw_zero, added to the tag; the fidelity review's
% paper-facing point 3, R17).
\begin{proposition}[the subordinated slice: the density at the origin]
\label{prop:origin-boundedness-subordinated}
\uses{def:isotropic-origin-datum,lem:origin-decay-exponent,lem:delay-polar-profile,lem:isotropic-exponent-log,lem:positive-transform-origin,lem:polar-decomposition}
\lean{AffineHemigroup.subordinated_profileAtOrigin, AffineHemigroup.subordinated_law,
AffineHemigroup.delayLaw_zero,
AffineHemigroup.subordinated_density_limit, AffineHemigroup.subordinated_unbounded,
AffineHemigroup.gaussVar, AffineHemigroup.mixtureDensity,
AffineHemigroup.measurable_gaussVar}\leanok
Let $d \ge 1$, let $(0, k_d)$ be an isotropic origin datum in dimension $d$ with Gaussian
coefficient $0$, let $\mu$ be its law, put $c := k_d(0{+}) \in [0,\infty]$ and suppose $c > 0$.
\begin{enumerate}
\item[(3)] Suppose the datum is subordinated, with delay datum $(0,k_{\mathrm I})$ of exponent
$F_{\mathrm I}$, that $W$ is a random delay whose law, a probability measure on $[0,\infty)$, has
Laplace transform $e^{-F_{\mathrm I}}$, and that $Z$ is standard normal on $\RR^d$, independent of
$W$. Then $c = 2\,k_{\mathrm I}(0{+})$, $W > 0$ almost surely, $\mu$ is the law of $\sqrt{W}Z$, and
$\mu$ has the radial density $p$ of \ref{eq:mixture-density}. This $p$ is nonincreasing in $|x|$, and
\[
  \lim_{x \to 0} p(x) = (2\pi)^{-d}\int_{\RR^d}\widehat\mu(\xi)\,d\xi
   = \frac{(2\pi)^{-d/2}}{\Gamma(d/2)}\int_0^\infty \sigma^{d/2-1}e^{-F_{\mathrm I}(\sigma)}\,d\sigma ,
\]
finite exactly when $c > d$. If $c \le d$, then for every representative $q$ of $\mu$'s
density, every $\delta > 0$ and every $M$ there is an $x$ with $0 < |x| < \delta$ and
$q(x) > M$: no representative is bounded near the origin.
\end{enumerate}
\statusT\quad\emph{(Cites [B, Prop.~4.9(2)], restated as \ref{lem:delay-polar-profile}, and no
ledger entry: Tonelli and one Gamma integral, with clause (1) of the parent,
\ref{lem:origin-decay-exponent}, the lower bound of \ref{lem:isotropic-exponent-log}(1) and
\ref{lem:positive-transform-origin}(2). The law $\mu$ is
any probability measure with transform $e^{-F_d}$, and the delay law is a \textbf{hypothesis}. Its
existence for a causally admissible $F_{\mathrm I}$ is [B, Rem.~4.3], which rests on Bernstein's
theorem; this module's ledger does not carry it, so the clause is conditional on it (Q-0281, the
narrowing of PRICING finding 2; the author may instead spend Paper V's A11 by reference). Clause (3)
of \ref{thm:origin-boundedness-d}, split off on 2026-10-04 (Q-0281, ADR-0025), with the parent's
clause number.)}
\end{proposition}

\begin{proof}
\leanok
This is the part of the proof of \ref{thm:origin-boundedness-d} headed clause (3), with its two
following paragraphs.
\end{proof}

% not in the draft: queue item Q-0409 (hub), the author's decision of 2026-10-10 on pull request
% #103 ("file the corollary please"), from Q-0408's report: the isotropic paper could say that the
% profile at the origin is the same number in every dimension only of the subordinated data
% (thm:origin-boundedness-d(3), c = 2 k_I(0+)), and the general statement follows from two bounds
% the chapter already has. This node is that statement. No statement of another node changes;
% lem:origin-decay-exponent's \lean tag gains the exponent form of its bound, which this node reads.
% The environment is `theorem` because the blueprint has no `corollary` environment, as for
% cor:isotropic-classification-operators; the isotropic paper prints it as a corollary.
%
% The form, decided on paper before anything was typed.
% (i) The statement is about EXPONENTS, and names no law except in the last sentence of clause (2).
% The two bounds it reads are bounds on F_d: lem:isotropic-exponent-log(1) is stated so, and
% lem:origin-decay-exponent is stated for |mu-hat| = e^{-F_d} while its proof bounds F_d. So the node
% needs neither the existence of the law (A3(iii)) nor c > 0, which the theorem supposes.
% (ii) The Gaussian coefficient is clause (1), kept. Without it clause (3) would have to SUPPOSE
% a = a' = 0: for a finite c the bound F <= c log r + C excludes a Gaussian part in the other
% dimension, and for c = infinity nothing does. With it the two numbers are both read off F, clause
% (2) holds for every datum (the jump part F_d - a|xi|^2/2 is the exponent of (0, k_d)), and clause
% (3) has no hypothesis on a. Its proof is the quadratic bound behind the finiteness of F_d, split
% at a radius tau instead of at 1; priced at about 120 Lean lines when the type was written.
% (iii) The limit of clause (2) is in [0, infinity]: c = infinity is allowed (the quotient tends to
% +infinity, by the lower bound alone) and so is c = 0 (the upper bound alone; with a = 0 that is the
% trivial datum, F_d = 0). The Lean states the limit in the extended reals, of the real quotient,
% so no truncation is hidden in the type.
% (iv) Clause (3) is stated on two data with one radial function F, the form of
% prop:isotropic-polar-reading(1) and of polar_reading_unique. It does not need the nesting
% A_{d+1} in A_d: it compares any two dimensions in which F has a datum. The reading for the
% classes A_d, and the application of thm:origin-boundedness-d(5) dimension by dimension, are in
% the annotation and not clauses: a clause on F in A_D would spend A2-A4 through
% prop:isotropic-polar-reading, and one on the density A1, where this node spends nothing. Both are
% machine-checked as probes (probe_origin_profile_classA and probe_origin_profile_threshold,
% Probes.lean), under the trust base.
% Axioms, printed before the annotation was written: the four declarations of the tag and
% OriginDatum.exponent_decay print propext, Classical.choice, Quot.sound, and nothing else. The four
% are pinned in CIAxiomGuard.lean and are headline declarations of linkage.toml [boundary].
\begin{theorem}[the exponent determines the Gaussian coefficient and the profile at the origin]
\label{cor:origin-profile-exponent}
\uses{def:isotropic-origin-datum,lem:polar-decomposition,lem:isotropic-exponent-log,lem:origin-decay-exponent}
\lean{AffineHemigroup.exponent_gaussian_limit, AffineHemigroup.exponent_profile_limit,
AffineHemigroup.transform_profile_limit, AffineHemigroup.origin_profile_dimension_free}\leanok
Let $d \ge 1$, let $(a, k_d)$ be an isotropic origin datum in dimension $d$, let $F_d$ be its
exponent, and put $c := k_d(0{+}) \in [0,\infty]$.
\begin{enumerate}
\item \emph{(The Gaussian coefficient.)} $F_d(\xi)/|\xi|^2 \to \tfrac12 a$ as $|\xi| \to \infty$.
\item \emph{(The profile at the origin.)} As $|\xi| \to \infty$,
\[
  \frac{F_d(\xi) - \tfrac12 a|\xi|^2}{\log|\xi|} \to c ,
\]
where the limit is $+\infty$ when $c = \infty$. In particular, if $a = 0$ and $\mu$ is the datum's
law, then $-\log|\widehat\mu(\xi)|/\log|\xi| \to c$.
\item \emph{(The same numbers in every dimension.)} Let $d' \ge 1$, let $(a', k'_{d'})$ be an
isotropic origin datum in dimension $d'$ with exponent $F'_{d'}$, and suppose that the two exponents
are one function of the radius: there is a function $F$ on $[0,\infty)$ with $F_d(\xi) = F(|\xi|)$
for every $\xi \in \RR^d$ and $F'_{d'}(\xi') = F(|\xi'|)$ for every $\xi' \in \RR^{d'}$. Then
$a' = a$ and $k'_{d'}(0{+}) = k_d(0{+})$.
\end{enumerate}
\statusT\quad\emph{(Cites nothing: the quadratic bound behind the finiteness of $F_d$, split at a
radius $\tau$, for clause (1); the two logarithmic bounds of \ref{lem:origin-decay-exponent} and
\ref{lem:isotropic-exponent-log}(1) for clause (2); the uniqueness of a limit for clause (3). The
statement is about exponents. No law is named except in the last sentence of clause (2), where
$\mu$ is any probability measure with transform $e^{-F_d}$, so the law's existence, ledger entry
A3(iii), is not read, and $c > 0$ is not supposed. \textbf{What it is worth.} Clause (2) makes
the profile at the origin a number read off the exponent: for a datum with no Gaussian part, the
power at which $e^{-F_d}$ falls off. Before this node the chapter had that reading as two
one-sided bounds only, \ref{lem:isotropic-exponent-log}(1) for a finite $c$ and
\ref{thm:origin-boundedness-d}(1) for every finite $\beta < c$. Clause (3) removes the dimension
from both numbers. On the subordinated slice \ref{thm:origin-boundedness-d}(3) gave
$k_d(0{+}) = 2k_{\mathrm I}(0{+})$ for every $d$, and this is the statement for every datum.
\textbf{The reading for the classes} is not a clause. By \ref{prop:isotropic-polar-reading}(1)
and~(2) a function $F$ is in $\mathcal{A}_d$ exactly when $F(|\cdot|)$ is the exponent of a datum in
dimension $d$, and that datum is unique up to a null set of its profile, which changes neither $a$
nor $k_d(0{+})$. Since $\mathcal{A}_D \subseteq \mathcal{A}_d$ for $d \le D$ ([B, Prop.~4.9]),
clause (3) says that an $F \in \mathcal{A}_D$ has one Gaussian coefficient and one profile at the
origin $c$, the same in every dimension $d \le D$. If $a = 0$ and $c > 0$,
\ref{thm:origin-boundedness-d}(5) then applies in each of these dimensions to that one number: the
law with exponent $F(|\cdot|)$ on $\RR^d$ has a density bounded near the origin exactly when
$c > d$. This reading spends what those two nodes spend, the entries A2, A3 and A4 for the data and
A1 for the density, where the node itself spends nothing. It is machine-checked as two probes,
\texttt{probe\_origin\_profile\_classA} and \texttt{probe\_origin\_profile\_threshold}.
\textbf{Edge cases.} $c = \infty$ is allowed in clauses (2) and (3): the quotient tends to
$+\infty$, by the bound from below alone. So is $c = 0$, by the bound from above alone; with
$a = 0$ that is the trivial datum, $k_d \equiv 0$ and $F_d \equiv 0$. The quotient of clause (2) is
nonnegative for $|\xi| > 1$. \textbf{Readings.} In the Lean, ``as $|\xi| \to \infty$'' is the
cocompact filter of $\RR^d$, the limit of clause (2) is that of the real quotient in the extended
reals, with $c \in [0,\infty]$ embedded, and the $F$ of clause (3) is a real function read at
$|\xi| \ge 0$ only. \textbf{Not claimed}: a rate in clause (2). For a finite $c$ the difference
form $F_d(a\xi) - F_d(\xi) \to c\log a$ of \ref{lem:isotropic-exponent-log}(2) says more, and it
is what \ref{prop:log-orbit-rays} reads. Nor the converse, that a function with the two limits is
the exponent of a datum; nor any relation between the two profiles of clause (3) beyond their
limits at radius zero.)}
\end{theorem}

\begin{proof}
\leanok
Let $\nu_d$ be the datum's \Levy{} measure and put
$G(\xi) := \int_{\RR^d}(1 - \cos\langle\xi,x\rangle)\,\nu_d(dx)$, so that
$F_d(\xi) = \tfrac12 a|\xi|^2 + G(\xi)$ by \ref{eq:isotropic-exponent}. Then $G \ge 0$, and $G$ is the
exponent of the datum $(0,k_d)$, which has the same polar profile and the same profile at the origin
$c$. As in the proof of \ref{lem:isotropic-exponent-log}, \ref{eq:polar-decomposition} gives
\[
  G(\xi) = \frac{1}{|S^{d-1}|}\int_{S^{d-1}}\varphi\bigl(\langle\xi,\eta\rangle\bigr)\,\sigma_{d-1}(d\eta),
  \qquad
  \varphi(s) := \int_0^\infty\bigl(1 - \cos(rs)\bigr)\,k_d(r)\,\frac{dr}{r} .
\]

\emph{Clause (1).} Fix $\tau \in (0,1]$ and put $I(\tau) := \int_0^\tau k_d(r)\,r\,dr$ and
$T(\tau) := \int_\tau^\infty k_d(r)\,dr/r$. Both are finite: $I(\tau)$ is at most the first integral
of \ref{def:isotropic-origin-datum}, and
$T(\tau) \le k_d(\tau)\log(1/\tau) + \int_1^\infty k_d(r)\,dr/r$ because $k_d$ is nonincreasing.
Since $1 - \cos(rs) \le r^2s^2/2$ on $(0,\tau)$ and $1 - \cos \le 2$ on $[\tau,\infty)$,
$\varphi(s) \le \tfrac12 s^2\,I(\tau) + 2\,T(\tau)$. On the sphere
$|\langle\xi,\eta\rangle| \le |\xi|$, and the normalized surface measure has mass $1$. So
\[
  0 \le G(\xi) \le \tfrac12|\xi|^2\,I(\tau) + 2\,T(\tau), \qquad \xi \in \RR^d .
\]
Let $\varepsilon > 0$. The integrand of $I(1)$ is integrable on $(0,1)$, so $I(\tau) \to 0$ as
$\tau \downarrow 0$, and there is a $\tau$ with $I(\tau) \le \varepsilon$. For that $\tau$,
$G(\xi)/|\xi|^2 \le \tfrac12\varepsilon + 2T(\tau)/|\xi|^2 \le \varepsilon$ once
$|\xi|^2 \ge 4T(\tau)/\varepsilon$. Hence $G(\xi)/|\xi|^2 \to 0$, and
$F_d(\xi)/|\xi|^2 = \tfrac12 a + G(\xi)/|\xi|^2 \to \tfrac12 a$.

\emph{Clause (2).} The quotient is $G(\xi)/\log|\xi|$, and it is nonnegative for $|\xi| > 1$.
\emph{From below.} Let $\beta$ be finite with $0 < \beta < c$. By \ref{lem:origin-decay-exponent}
for the datum $(0,k_d)$ there is $C_\beta$ with $e^{-G(\xi)} \le C_\beta(1+|\xi|)^{-\beta}$ for
every $\xi$, and $C_\beta \ge 1$ by the bound at $\xi = 0$. That lemma states the bound for
$|\widehat\mu|$, $\mu$ being a law of the datum $(0,k_d)$. Its proof bounds $e^{-G}$, which is
$|\widehat\mu|$ by the definition of the law, so the bound holds whether or not a law is given.
Taking logarithms, $G(\xi) \ge \beta\log|\xi| - \log C_\beta$ for $|\xi| \ge 1$. \emph{From above.}
Let $c < \infty$. By \ref{lem:isotropic-exponent-log}(1) for the datum $(0,k_d)$,
$G(\xi) \le c\log|\xi| + \tfrac{13}{4}c + 2C_k$ for $|\xi| \ge 1$, with
$C_k = \int_1^\infty k_d(r)\,dr/r$. \emph{The limit.} Divide by $\log|\xi|$, which tends to
infinity. For every $\varepsilon > 0$ the quotient is at least $\beta - \varepsilon$ once $|\xi|$ is
large enough, for every finite $\beta$ with $0 < \beta < c$, and at most $c + \varepsilon$ once
$|\xi|$ is large enough when $c < \infty$. If $c = \infty$ the first bound makes the quotient tend to
$+\infty$. If $0 < c < \infty$ the two bounds give the limit $c$. If $c = 0$ the quotient lies
between $0$ and $\varepsilon$, and tends to $0$. For the last sentence of the clause, $a = 0$ gives
$F_d = G$, and $|\widehat\mu| = e^{-F_d}$ by the definition of the law.

\emph{Clause (3).} Since $d, d' \ge 1$ there are unit vectors $e \in \RR^d$ and $e' \in \RR^{d'}$,
and $F_d(re) = F(r) = F'_{d'}(re')$ for every $r > 0$. Clause (1) for the two data, read along
$\xi = re$ and $\xi' = re'$, gives $F(r)/r^2 \to \tfrac12 a$ and $F(r)/r^2 \to \tfrac12 a'$ as
$r \to \infty$. So $a' = a$. Clause (2), read along the same rays, then gives
$\bigl(F(r) - \tfrac12 a r^2\bigr)/\log r \to k_d(0{+})$ and
$\bigl(F(r) - \tfrac12 a r^2\bigr)/\log r \to k'_{d'}(0{+})$ as $r \to \infty$. A limit in
$[0,\infty]$ is unique, so $k'_{d'}(0{+}) = k_d(0{+})$.

Every step is machine-checked, and no cited fact is spent.
\end{proof}

% Q-0272. The pointwise half that thm:origin-boundedness-d(4) leaves out: a lower bound on the
% density itself, not on its heat blur. The brief proposed a radial Abelian theorem for the
% transform (about (d-1)/2 radial derivatives of e^{-F_d} against Bessel kernels, "risky at d = 2");
% the proof below takes the other side of the Fourier transform instead and reads only the Levy
% measure, one scale at a time: the jumps of size >= eps are an atom e^{-m} >= const eps^c at the
% origin, the jumps of size in [tau eps, eps) a twofold convolution that fills a ball of radius
% 3eps/2, and the jumps below tau eps move the result by less than eps with probability >= 1/2. It
% works for every d >= 2 alike (d >= 2 is used once, for a vector orthogonal to z), so no case is
% left open; the rate's slowly varying factor and a monotone representative are not obtained.
% CHANGED Q-0281 wave 2: typed (the statement skeleton of the three nodes added since). The bound is
% typed with its constants: for every eps_0 in (0,1] with k_d(eps_0) >= c/2, delta = eps_0/2 and
% c_1 = lowerBoundConst; "the density q of the absolutely continuous part" is the Radon-Nikodym
% derivative of mu against Lebesgue measure, "every density" every real q with mu = q dx. The law is
% any probability measure with transform e^{-F_d} (IsLawOf). The text is unchanged.
% CHANGED Q-0281 wave 2, group B: proved, in AffineHemigroup.OriginLowerBound, names and types
% unchanged; every declaration of the tag prints A3(iii) (levy_khintchine_converse) and Lean core
% only, and is guarded. The text of the statement is unchanged. Two sentences of the statusA follow
% the rewritten proof: "the two compound Poisson laws of the proof are written down as series, and the
% rest is Fubini's theorem, the uniqueness of the Fourier transform and polar coordinates" now reads
% "the compound Poisson law ... a series ... the uniqueness of the Fourier transform, a Gaussian
% average and polar coordinates"; "This proof uses d >= 2 at one step, the choice of a vector
% orthogonal to z ..." now names the second step, the bound 3/(16d) <= 3/32 on the small jumps.
% CHANGED Q-0415 (2026-10-11): annotation only. "Not claimed: an upper bound q(x) <= C|x|^{c-d}"
% now points to prop:origin-upper-bound-d for the bound with every exponent beta < c. The
% statement, its proof and the \lean tag are unchanged.
\begin{proposition}[a pointwise lower bound at the origin]
\label{prop:origin-lower-bound-d}
\uses{def:isotropic-origin-datum,lem:polar-decomposition,thm:cited-A3}
\lean{AffineHemigroup.origin_lower_bound, AffineHemigroup.origin_lower_bound_density,
AffineHemigroup.origin_lower_bound_ball, AffineHemigroup.lowerBoundConst,
AffineHemigroup.lowerBoundTau, AffineHemigroup.lowerBoundKappa}\leanok
Let $d \ge 2$, let $(0, k_d)$ be an isotropic origin datum in dimension $d$ with Gaussian
coefficient $0$, let $\mu$ be its law, and suppose that $c := k_d(0{+})$ is finite, $0 < c < \infty$.
Then there are $\delta > 0$ and $c_1 > 0$ with
\begin{equation}\label{eq:origin-lower-bound}
  \mu(E) \ge c_1\int_E |x|^{c-d}\,dx \qquad\text{for every Borel set } E \subseteq \{|x| < \delta\} .
\end{equation}
In particular:
\begin{enumerate}
\item the density $q$ of the absolutely continuous part of $\mu$ satisfies
$q(x) \ge c_1|x|^{c-d}$ for Lebesgue-almost every $x$ with $|x| < \delta$; so does every density of
$\mu$ when $\mu$ has one, and every representative of $q$ is at least
$c_1|x|^{c-d}$ almost everywhere near the origin, which tends to $\infty$ as $x \to 0$ when $c < d$;
\item $\mu(\{|x| < r\}) \ge c_1|S^{d-1}|\,c^{-1}\,r^{c}$ for $0 < r \le \delta$.
\end{enumerate}
The constants can be taken as follows. Let $\varepsilon_0 \in (0,1]$ satisfy
$k_d(\varepsilon_0) \ge c/2$, put $C_k := \int_1^\infty k_d(r)\,dr/r$, let $V_d$ be the volume of the
unit ball, and put
\[
  \tau := \min\Bigl\{\tfrac12,\ \bigl(\tfrac{3}{4d^2c}\bigr)^{1/2}\Bigr\},\quad
  \kappa_0 := \Bigl(\frac{c}{2|S^{d-1}|}\Bigr)^{2}\frac{V_d}{8^d},\quad
  \delta := \frac{\varepsilon_0}{2},\quad
  c_1 := \frac{\kappa_0}{4}\,\tau^{c}e^{-C_k}\min\{2^{c-d},4^{c-d}\} .
\]
\statusA\quad\ledger{A3}\quad\emph{(\textbf{Assignment.} The law $\mu$ and the law $\mu_1$ of the
truncated datum in the proof exist by \ledger{A3}(iii), the converse half of the \Levy--Khintchine
representation, which is what \ref{thm:origin-boundedness-d} spends for $\mu$; nothing else is
cited: the compound Poisson law of the proof is written down as a series, and the rest is
Fubini's theorem, the uniqueness of the Fourier transform, a Gaussian average and polar
coordinates (\ref{lem:polar-decomposition}). No new ledger entry, and \ledger{A2} is not read: the proof uses
only that $k_d$ is nonincreasing with $k_d(0{+}) = c$, not that $\mu$ is self-decomposable.
\textbf{What it is worth.} It is the pointwise half that \ref{thm:origin-boundedness-d}(4) does not
give: that clause bounds the density of $\mu$ blurred by the heat kernel, $\int g_t\,d\mu \ge
C\,t^{-(d-c)/2}$, and this node bounds the density itself, from below, at the same exponent; for
$c < d$ the clause's rate follows again from \ref{eq:origin-lower-bound} by integrating against
$g_t$; and for $c \le d$ the clause's upper bound $\int g_t\,d\mu \le C'_\beta t^{-(d-\beta)/2}$,
$\beta < c$, shows that no bound $c_1|x|^{\beta-d}$ with $\beta < c$ holds, so the exponent is sharp
there. It is the lower half
of the withdrawn hypothesis (H) (see the comment before \ref{lem:positive-transform-origin}),
now proved for every datum and every $d \ge 2$. For the line see [B, Cor.~3.9]. This proof uses
$d \ge 2$ at two steps: the choice of a vector orthogonal to $z$ where the twofold convolution fills
a ball, and the bound $3/(16d) \le 3/32$ on the small jumps. For $c > d$ the bound is true
but weak, $\mu$ then having a bounded continuous density by \ref{thm:origin-boundedness-d}(2).
\textbf{Not claimed}: that $\mu$ has a density when $c \le d$ (clause (1) is about its absolutely
continuous part); an upper bound $q(x) \le C|x|^{c-d}$ (on the line Sato's rate carries the factor
$K$, which is unbounded exactly when $\int_0^1(c - k(y) - k(-y))\,dy/y = \infty$; the upper bound
with every exponent $\beta < c$ in place of $c$, for $c \le d$ and the nonincreasing density, is
\ref{prop:origin-upper-bound-d}); the factor
$\log(1/|x|)$ at $c = d$, which (H) asked for and the bound $|x|^0 = 1$ does not give; the
$d$-dimensional slowly varying factor itself; and a radially nonincreasing representative, which
is \ref{prop:isotropic-radial-monotone}'s, not this node's. Nor
anything at $c = \infty$, where \ref{thm:origin-boundedness-d}(2) makes the density smooth.)}
\end{proposition}

\begin{proof}
\leanok
% CHANGED Q-0281 wave 2, group B: proof rewritten to the Lean route. One compound Poisson law of all
% the jumps of size at least tau eps replaces the two of the earlier text (for [tau eps, eps) and for
% [eps, infty)); the twofold convolution is bounded through Tonelli's theorem rather than a pointwise
% density; the small jumps are controlled by a Gaussian average of the transform in place of the
% coordinate truncation inequality; the dyadic shells are closed at the top. The constants are
% unchanged.
Write $|E|$ for the Lebesgue measure of $E$, $B(z,r)$ and $\bar B(z,r)$ for the open and the closed
ball, and $\nu_d$ for the datum's \Levy{} measure. As $k_d$ is nonincreasing with right limit $c$ at
the origin, $k_d \le c$ on $(0,\infty)$ and $k_d(r) \uparrow c$ as $r \downarrow 0$, so
$\varepsilon_0$ exists. Every integral of $\nu_d$ below is computed by \ref{eq:polar-decomposition},
applied to a nonnegative Borel function vanishing at the origin.

\emph{The decomposition at scale $\varepsilon$.} Fix $\varepsilon \in (0,\varepsilon_0]$ and let
$\nu_1$ and $\nu_{\ge}$ be the restrictions of $\nu_d$ to $\{0 < |x| < \tau\varepsilon\}$ and to
$\{|x| \ge \tau\varepsilon\}$. Each has a radial density, so each is symmetric, and
$\psi_1(\xi) := \int(1 - \cos\langle\xi,x\rangle)\,\nu_1(dx)$ and
$\psi_{\ge}(\xi) := \int(1 - \cos\langle\xi,x\rangle)\,\nu_{\ge}(dx)$ are finite with
$F_d = \psi_1 + \psi_{\ge}$, $F_d$ being \ref{eq:isotropic-exponent} at $a = 0$. The mass of
$\nu_{\ge}$ is finite: as $\tau\varepsilon \le 1$ and $k_d \le c$,
\[
  m := \nu_{\ge}(\RR^d) = \int_{\tau\varepsilon}^{\infty}k_d(r)\,\frac{dr}{r}
   \le c\log\frac{1}{\tau\varepsilon} + C_k ,
  \qquad\text{so}\qquad e^{-m} \ge \tau^c\varepsilon^c e^{-C_k} .
\]
For a finite symmetric measure $\lambda$ of mass $m_\lambda$ put
$\pi_\lambda := e^{-m_\lambda}\sum_{n\ge0}\lambda^{*n}/n!$, with $\lambda^{*0} = \delta_0$; it is a
probability measure, and, summing the series of transforms by dominated convergence
($|\widehat\lambda^n| \le m_\lambda^n$), its transform is
$\exp(\widehat\lambda(\xi) - m_\lambda) = \exp\bigl(-\int(1 - \cos\langle\xi,x\rangle)\,\lambda(dx)\bigr)$,
the sine part of $\widehat\lambda$ vanishing by symmetry. Put $\pi := \pi_{\nu_{\ge}}$, of transform
$e^{-\psi_{\ge}}$. The profile $k_d\mathbf 1_{(0,\tau\varepsilon)}$ is nonincreasing and satisfies the
two integrability conditions of \ref{def:isotropic-origin-datum}, so
$(0, k_d\mathbf 1_{(0,\tau\varepsilon)})$ is an isotropic origin datum, with \Levy{} measure $\nu_1$
and exponent $\psi_1$; let $\mu_1$ be its law (\ledger{A3}(iii)), so $\widehat\mu_1 = e^{-\psi_1}$.
Then $\mu_1 * \pi$ has transform $e^{-F_d} = \widehat\mu$, and by the uniqueness of the Fourier
transform
\[
  \mu = \mu_1 * \pi .
\]
Keeping the term $n = 2$ of the series, $\pi \ge \tfrac12 e^{-m}\,\nu_{\ge} * \nu_{\ge}$, as measures.

\emph{The twofold convolution fills a ball.} Let $S := \{\varepsilon/2 \le |x| < \varepsilon\}$, which
lies in $\{|x| \ge \tau\varepsilon\}$ because $\tau \le 1/2$. On $S$ the density
$k_d(|x|)\,(|S^{d-1}||x|^d)^{-1}$ of $\nu_{\ge}$ is at least
$a := c\,(2|S^{d-1}|\varepsilon^d)^{-1}$, because $k_d(|x|) \ge k_d(\varepsilon_0) \ge c/2$ and
$|x|^d < \varepsilon^d$; so $\nu_{\ge}(A) \ge a\,|A \cap S|$ for every Borel $A$. Let
$|z| \le 3\varepsilon/2$. As $d \ge 2$ there is a unit vector $w$ orthogonal to $z$; put
$y_0 := z/2 + sw$ with $s := (9\varepsilon^2/16 - |z|^2/4)^{1/2}$, so that
$|y_0| = |z - y_0| = 3\varepsilon/4$. Every $y \in B(y_0,\varepsilon/8)$ has $|y|$ and $|z - y|$ in
$(5\varepsilon/8, 7\varepsilon/8)$, inside $S$, so
$|\{y \in S : z - y \in S\}| \ge |B(0,\varepsilon/8)|$. For a Borel $G \subseteq \bar B(0,3\varepsilon/2)$,
Tonelli's theorem and the substitution $x = z - y$ give
\[
\begin{aligned}
  (\nu_{\ge} * \nu_{\ge})(G) &\ge a^2\iint\mathbf 1_S(x)\,\mathbf 1_S(y)\,\mathbf 1_G(x + y)\,dx\,dy \\
   &= a^2\int_G\bigl|\{y \in S : z - y \in S\}\bigr|\,dz
   \ge a^2\,\bigl|B(0,\varepsilon/8)\bigr|\,|G| = \kappa_0\,\varepsilon^{-d}|G| .
\end{aligned}
\]
Hence $\pi(G) \ge \tfrac12 e^{-m}\kappa_0\varepsilon^{-d}|G|$ for every Borel
$G \subseteq \bar B(0, 3\varepsilon/2)$.

\emph{The small jumps move little.} By $1 - \cos u \le u^2/2$ and \ref{eq:polar-decomposition},
\[
  \psi_1(\xi) \le \frac{|\xi|^2}{2}\int|x|^2\,\nu_1(dx)
   = \frac{|\xi|^2}{2}\int_0^{\tau\varepsilon}k_d(r)\,r\,dr
   \le \frac{c\,\tau^2\varepsilon^2}{4}\,|\xi|^2 .
\]
Let $\gamma$ be the centred Gaussian law on $\RR^d$ with covariance $\varepsilon^{-2}I$, whose
transform at $w$ is $e^{-|w|^2/2\varepsilon^2}$ and whose second moment is
$\int|\xi|^2\,\gamma(d\xi) = d\,\varepsilon^{-2}$. By Fubini's theorem, $\widehat\mu_1$ being real,
and by $1 - e^{-\psi} \le \psi$,
\[
  \int\bigl(1 - e^{-|w|^2/2\varepsilon^2}\bigr)\,\mu_1(dw)
   = \int\bigl(1 - \widehat\mu_1(\xi)\bigr)\,\gamma(d\xi)
   \le \int\psi_1\,d\gamma
   \le \frac{c\,\tau^2\varepsilon^2}{4}\cdot\frac{d}{\varepsilon^2} = \frac{c\,\tau^2 d}{4} .
\]
The integrand on the left is at least $1 - e^{-1/2} \ge 1/3$ where $|w| \ge \varepsilon$, and
$c\tau^2 d/4 \le 3/(16d) \le 3/32$ by the choice of $\tau$ and $d \ge 2$, so
\[
  \mu_1\bigl(\{|w| \ge \varepsilon\}\bigr) \le \tfrac{9}{32} \le \tfrac12 .
\]

\emph{One scale.} Let $E \subseteq \bar B(0,\varepsilon/2)$ be Borel. For $|w| < \varepsilon$ the
translate $E - w$ lies in $\bar B(0,3\varepsilon/2)$ and has measure $|E|$. So, the integrand being
nonnegative,
\[
  \mu(E) = \int\pi(E - w)\,\mu_1(dw)
   \ge \int_{\{|w|<\varepsilon\}}\pi(E - w)\,\mu_1(dw)
   \ge \tfrac12 e^{-m}\kappa_0\varepsilon^{-d}|E|\cdot\tfrac12 .
\]
With $e^{-m} \ge \tau^c\varepsilon^c e^{-C_k}$ this is
\[
  \mu(E) \ge \kappa_1\,\varepsilon^{c-d}\,|E| ,\qquad
  \kappa_1 := \tfrac14\kappa_0\tau^ce^{-C_k},\qquad
  E \subseteq \bar B(0,\varepsilon/2),\ 0 < \varepsilon \le \varepsilon_0 .
\]

\emph{All scales.} For $n \ge 0$ put $\varepsilon_n := 2^{-n}\varepsilon_0$ and
$A_n := \{\varepsilon_n/4 < |x| \le \varepsilon_n/2\}$; these sets are disjoint and cover
$B(0,\delta)\setminus\{0\}$. For $x \in A_n$, $2|x| \le \varepsilon_n < 4|x|$, so
$\varepsilon_n^{c-d} \ge \theta|x|^{c-d}$ with $\theta := \min\{2^{c-d},4^{c-d}\}$ (the first bound
when $c \ge d$, the second when $c < d$). For a Borel $E \subseteq B(0,\delta)$, the origin being a
Lebesgue-null set, the previous step at $\varepsilon = \varepsilon_n$ gives
\[
  \mu(E) \ge \sum_{n\ge0}\mu(E \cap A_n)
   \ge \kappa_1\sum_{n\ge0}\varepsilon_n^{c-d}|E \cap A_n|
   \ge \kappa_1\theta\sum_{n\ge0}\int_{E\cap A_n}|x|^{c-d}\,dx
   = c_1\int_E|x|^{c-d}\,dx ,
\]
which is \ref{eq:origin-lower-bound}.

% CHANGED Q-0281 wave 2: blind review W2B-6 (the reason for the comparison of set integrals).
\emph{The consequences.} Write $\mu = \mu_{\mathrm{ac}} + \mu_{\mathrm s}$ for the Lebesgue
decomposition and let $N$ be a Borel set with $\mu_{\mathrm s}(N) = 0$ and
$|\RR^d \setminus N| = 0$. For a Borel $E \subseteq B(0,\delta)$,
$\int_E q\,dx \ge \int_{E\cap N} q\,dx = \mu_{\mathrm{ac}}(E \cap N) = \mu(E \cap N)
\ge c_1\int_{E\cap N}|x|^{c-d}\,dx = c_1\int_E|x|^{c-d}\,dx$. A function whose integral over every
Borel subset of $B(0,\delta)$ is at least that of the measurable $c_1|x|^{c-d}$ is at least
$c_1|x|^{c-d}$ almost everywhere on $B(0,\delta)$, both functions being nonnegative and measurable
and Lebesgue measure being $\sigma$-finite: on a set $E \subseteq B(0,\delta)$ of finite measure
where $q + 1/n \le c_1|x|^{c-d} \le n$, the inequality gives
$\int_E q\,dx + |E|/n \le \int_E q\,dx$, the integral being finite as $q \le n$ on $E$, so $|E| = 0$, and such sets exhaust the set
where $q < c_1|x|^{c-d}$. That is clause (1); for a density $q$ of $\mu$ the
same argument runs with $\int_E q\,dx = \mu(E)$, and $c_1|x|^{c-d} > 0$ off the origin, which is
null. As $x \to 0$, $|x| \to 0$ through positive values, and $|x|^{c-d} \to \infty$ when $c < d$.
Clause (2) is \ref{eq:origin-lower-bound} at $E = \{|x| < r\}$, in polar coordinates:
$\int_{\{|x|<r\}}|x|^{c-d}dx = |S^{d-1}|\int_0^r\rho^{c-1}d\rho = |S^{d-1}|\,r^c/c$.

\emph{What is machine-checked.} Everything above, with \ledger{A3}(iii), for the law $\mu_1$, as the
one interface it cites.
\end{proof}

% Q-0316 (brief item 1; records/affine/GATE-paper-iso.md item 18, the author's call between prose
% and a node; the blind register review's section 3 flag 1 and section 4 flags 1, 3, 4, 15): Wolfe's
% unimodality theorem, ledger A8, made a statement. The paper printed the consequence in its open
% problems as a two-sentence argument that did not handle the uniform laws on lower-dimensional
% sets, which have no density, and did not say where absolute continuity comes from. Here both are
% done: the radial part of a k-dimensional set's uniform law is r^{k-d} times a nonincreasing
% fraction, and the atom at the origin is excluded by the decay of the transform, so absolute
% continuity is proved, not cited. No \lean tag (the brief): the statement skeleton session types
% it, with A8's interface. Cost, priced at the statement: the measure theory of the mixture (an
% integral over a space of measures, Kanter's Borel structure) is the expensive part in Lean; the
% concavity argument is elementary; the decay reads lem:origin-decay-exponent at d = 1.
% CHANGED Q-0281 wave 2: typed (the statement skeleton of the three nodes added since). The statements
% need no vocabulary of A8, which has no Lean interface yet (AXIOMS.md: its Lean line waits for the
% author's approval of the entry; no axiom is added here), so they are typed in full and the proof
% waits for the interface. "A density p with p(x) = q(|x|) for x != 0" is mu = q(|x|) dx with q >= 0
% nonincreasing on (0,infty); "not trivial" is a > 0 or k_d not a.e. 0 on (0,infty). The text is
% unchanged.
% CHANGED Q-0281 wave 2, group B: proved, after the author admitted ledger A8 as the Lean interface
% AffineHemigroup.wolfe_symmetric_unimodal (2026-10-05, on the librarian's audit). The text of the
% statement is unchanged. radial_monotone_law prints A8, A2 (sato_polar_criterion), A3(iii)
% (levy_khintchine_converse) and Lean core, and is guarded. radial_monotone_family prints A8 and,
% through AffineHemigroup.isotropic_classification, Paper V's SpatialLine.fourier_toolbox_levy_unique
% (this module's A2-A4 are not spent by clause (1)); it is guarded, V's interface admitted through
% the include line of blueprint/trust-boundary.txt (unguarded until article-kit read an include,
% Q-0298, by the author's decision of 2026-10-05, blind review W1E-3).
% CHANGED Q-0365 (blind review S65-3): the two sentences above are superseded. Since Q-0365 the
% classification's necessity direction takes Paper V's uniqueness clause as the hypothesis
% ScaleSpace.SymLevyUnique, discharged by AffineHemigroup.line_levy_unique from A3(ii)-(iii); so
% radial_monotone_family prints A8 and A3(ii)-(iii) (through isotropic_classification), and no
% interface of Paper V. The text of the statement is unchanged.
% CHANGED Q-0415 (2026-10-11): annotation only. "Not claimed: a pointwise bound on q (the open
% problem of the isotropic paper's section 4)" was true when written; since
% prop:origin-upper-bound-d the bounds c_1 r^{c-d} <= q(r) <= C_beta r^{beta-d} are proved for a
% datum with no Gaussian part and 0 < c <= d (the lower one for d >= 2), and what stays open is the
% sharp rate. The statement, its proof and the \lean tag are unchanged.
% CHANGED Q-0416 (2026-10-11): annotation only. The sharp rate is no longer open: it is
% prop:origin-rate-d, and the "Not claimed" sentence points there.
\begin{proposition}[the kernels from scale zero have a radially nonincreasing density]
\label{prop:isotropic-radial-monotone}
\uses{def:isotropic-kernel-family,thm:isotropic-classification,def:filtration-thresholds,def:isotropic-origin-datum,lem:polar-decomposition,lem:radial-marginal,lem:origin-decay-exponent,lem:positive-transform-origin,prop:isotropic-polar-reading,thm:cited-A2,thm:cited-A3,thm:cited-A8}
\lean{AffineHemigroup.radial_monotone_family, AffineHemigroup.radial_monotone_law}\leanok
Let $d \ge 1$.
\begin{enumerate}
\item Let $(\mu_{s,t})$ be an isotropic cascade family in dimension $d$ and let $t > 0$. Then
$\mu_{0,t}$ has a density $p$ with $p(x) = q(|x|)$ for $x \ne 0$, where
$q : (0,\infty) \to [0,\infty)$ is nonincreasing.
\item The same holds for the law of every isotropic origin datum $(a,k_d)$ in dimension $d$ that is
not trivial ($a > 0$, or $k_d$ not almost everywhere $0$).
\end{enumerate}
\statusA\quad\ledger{A8}\quad\ledger{A2}\quad\ledger{A3}\quad\emph{(\textbf{Assignment.} Clause (1)
spends \ledger{A8} alone, Wolfe's theorem with Kanter's definition: a symmetric self-decomposable
law is a mixture of uniform laws on compact convex symmetric sets. Everything else is proved; the
Lean of clause (1) reaches self-decomposability through \ref{thm:isotropic-classification} and,
since Q-0365, prints \ledger{A3}(ii)--(iii) there, where it printed Paper V's interface. The
kernel is self-decomposable by the cascade, as in the necessity half of
\ref{thm:isotropic-classification}; it has no atom at the origin because its transform tends to $0$,
by \ref{thm:origin-boundedness-d}(1) read on the line (\ref{lem:origin-decay-exponent}, which cites
nothing) and \ref{lem:positive-transform-origin}(1); and the density comes from the mixture by
rotation invariance. Clause (2) spends in addition \ledger{A2} and \ledger{A3}, through
\ref{prop:isotropic-polar-reading}(1) in its converse direction, which makes the datum's law
self-decomposable. \textbf{What \ledger{A8} carries and what it does not.} It carries the mixture,
and nothing about densities: Kanter's class contains the uniform law on a segment of $\RR^2$ and
the point mass. The proof below treats the sets of every dimension and the point mass itself.
\textbf{Not claimed}: a pointwise bound on $q$ (for a datum with no Gaussian part and
$0 < c \le d$ the bounds are \ref{prop:origin-upper-bound-d}'s, and the rate is
\ref{prop:origin-rate-d}'s);
continuity of $q$; uniqueness of $q$ beyond a Lebesgue-null set (the proof produces the version
read off a right derivative); anything about laws that are not rotation invariant. \textbf{What it
is worth.} On the subordinated slice \ref{thm:origin-boundedness-d}(3) writes the density down as a
Gaussian mixture, and it is visibly nonincreasing; off the slice this node is the only source of a
monotone representative. On the line it is the symmetric case of the unimodality of
self-decomposable laws (Wintner; \sscite{sato1999levy}, pp.~403--404). Provenance: Q-0316.)}
\end{proposition}

\begin{proof}
\leanok
% CHANGED Q-0281 wave 2, group B: proof rewritten to the Lean route. Clause (1)'s four properties are
% read from the self-decomposable law that the definition of A_d supplies, with the residual
% (sI)_* rho in place of the cascade's mu_{s,t}, and its decay from the datum (2a_V, k) in dimension 1
% whose law is that law's image on a line; the radial part of one uniform law is computed ray by ray
% in polar coordinates of V, in place of the sphere fraction phi_K and h_K; the density is read
% through the images of mu and of Lebesgue measure under the volume coordinate, in place of the
% reversed substitution, and the uniqueness from the radial image is the one the polar reading
% proves, in place of the Haar average. Clause (2) and the paragraphs on the atom and the mixture
% follow the earlier text; clause (2) now cites lem:origin-decay-exponent, which is
% thm:origin-boundedness-d(1) and is what the node uses.
The argument reads four properties of a law $\mu$ on $\RR^d$: it is rotation invariant, it is
self-decomposable, its transform is real, and $\widehat\mu(\xi) \to 0$ as $|\xi| \to \infty$. The
first two paragraphs check them for the two clauses. The rest derives the density from them.

\emph{Clause (1): the four properties.} Put $\mu := \mu_{0,t}$. By
\ref{thm:isotropic-classification}, $\widehat\mu(\xi) = e^{-F(\chi(t)|\xi|)}$ for a gauge $\chi$
and an $F \in \mathcal{A}_d$ that does not vanish identically on $[0,\infty)$, and
$s := \chi(t) > 0$. By the definition of $\mathcal{A}_d$ there is a self-decomposable law $\mu'$
on $\RR^d$ with $\widehat{\mu'}(\xi) = e^{-F(|\xi|)}$, so $\widehat\mu(\xi) = \widehat{\mu'}(s\xi)$.
So $\widehat\mu$ is real and radial: for $Q \in \Orth(d)$,
$\widehat{Q_*\mu}(\xi) = \widehat\mu(Q^{\tp}\xi) = \widehat\mu(\xi)$, and a law is determined by its
transform, so $\mu$ is rotation invariant. For $\lambda \in (0,1)$ let $\rho_\lambda$ be a residual
of $\mu'$, $\widehat{\mu'}(\xi) = \widehat{\mu'}(\lambda\xi)\,\widehat\rho_\lambda(\xi)$. Then
\[
  \widehat\mu(\xi) = \widehat{\mu'}(s\xi) = \widehat{\mu'}(\lambda s\xi)\,\widehat\rho_\lambda(s\xi)
   = \widehat\mu(\lambda\xi)\,\widehat{(sI)_*\rho_\lambda}(\xi) ,
\]
which is \ref{eq:self-decomposable} with residual $(sI)_*\rho_\lambda$. It remains to see that
$e^{-F(r)} \to 0$ as $r \to \infty$. $F$ is admissible, of the form \ref{eq:admissible-line} with
Gaussian coefficient $a_{\mathrm V}$ and folded profile $k$. The pair $(2a_{\mathrm V}, k)$ is an
isotropic origin datum in dimension $1$: its two integrability conditions are those of
\ref{eq:admissible-line}, and by \ref{eq:polar-decomposition} at $d = 1$, where $|S^0| = 2$, its
exponent at $\omega$ is
$a_{\mathrm V}\omega^2 + \int_{\RR}(1 - \cos\omega x)\,k(|x|)\,(2|x|)^{-1}dx = F(|\omega|)$. Its
law is the image of $\mu'$ under $x \mapsto \langle e_1, x\rangle$, whose transform is
$\widehat{\mu'}(\omega e_1) = e^{-F(|\omega|)}$, as in \ref{lem:radial-marginal}(1). The datum is
not trivial: a trivial datum has exponent $0$, and $F$ does not vanish identically on $[0,\infty)$.
So the last two sentences of the next paragraph, at $d = 1$, give $e^{-F(r)} \to 0$, and
$\widehat\mu(\xi) \to 0$ as $|\xi| \to \infty$.

\emph{Clause (2): the four properties.} Let $\mu$ be the datum's law, so that
$\widehat\mu = e^{-F_d}$ is real, positive and radial by \ref{eq:isotropic-exponent}. For
$Q \in \Orth(d)$, $\widehat{Q_*\mu}(\xi) = \widehat\mu(Q^{\tp}\xi) = \widehat\mu(\xi)$, and a law is
determined by its transform, so $\mu$ is rotation invariant. Put $F(r) := F_d(re_1)$. Then
$\xi \mapsto F(|\xi|)$ is the exponent $F_d$ of the datum, so by
\ref{prop:isotropic-polar-reading}(1), in its converse direction, it is the exponent of a
self-decomposable law, which is $\mu$. Finally $F_d(\xi) \ge \tfrac12 a|\xi|^2$, which tends to
infinity if $a > 0$. If $a = 0$, then $k_d$ is not almost everywhere $0$, so
$c = k_d(0{+}) > 0$, $k_d$ being at most $c$ on $(0,\infty)$, and
\ref{lem:origin-decay-exponent} with a finite $\beta \in (0,c)$ gives
$\widehat\mu(\xi) \le C_\beta(1+|\xi|)^{-\beta}$.

\emph{No atom at the origin.} Let $\mu$ have the four properties. For $\varepsilon > 0$, multiply
the identity of \ref{lem:positive-transform-origin}(1) at $t = \varepsilon$ by
$(2\pi\varepsilon)^{d/2}$ and substitute $\xi = y/\sqrt\varepsilon$:
\[
  \int_{\RR^d} e^{-|x|^2/2\varepsilon}\,\mu(dx)
   = (2\pi)^{-d/2}\int_{\RR^d}\widehat\mu\bigl(y/\sqrt\varepsilon\bigr)\,e^{-|y|^2/2}\,dy .
\]
Let $\varepsilon \downarrow 0$. On the left the integrand tends to the indicator of the origin, and
on the right $\widehat\mu(y/\sqrt\varepsilon) \to 0$ for every $y \ne 0$. The integrand on the left is bounded
by $1$, which is integrable against $\mu$, and the one on the right by $e^{-|y|^2/2}$, which is
integrable against $dy$, as $|\widehat\mu| \le 1$; so by dominated convergence $\mu(\{0\}) = 0$.

\emph{The mixture.} The law $\mu$ is symmetric, by rotation invariance at $Q = -I$, and
self-decomposable. By \ledger{A8} it is symmetric unimodal: there is a probability measure
$\Lambda$ on the uniform laws $u_K$ of the compact convex symmetric sets $K \subset \RR^d$ with
\[
  \mu(A) = \int u_K(A)\,\Lambda(du_K) \qquad \text{for every Borel } A \subseteq \RR^d .
\]
Here $u_K$ is uniform for the Lebesgue measure $\mathrm{Vol}_V$ of the linear span $V$ of $K$, whose
dimension $k$ may be anything from $0$ to $d$, and $K$ is nonempty. $\Lambda$ is a probability
measure on the measurable space of the probability measures on $\RR^d$, whose $\sigma$-algebra is
generated by the evaluations $\nu \mapsto \nu(A)$, and almost every measure under it is a $u_K$;
nothing below asks the set $W_d$ of the $u_K$ to be measurable. The source puts its mixing measure
on $W_d$ itself, with the Borel sets of weak convergence, which on the probability measures of
$\RR^d$ are the sets of that $\sigma$-algebra. Reading it as almost every measure being a $u_K$
uses that $W_d$ is such a set. It is one, being a countable union of weakly compact sets, as the
paragraph on $W_d$ at the end of the proof shows. The formal statement puts
$\Lambda$ on all measures on $\RR^d$, which loses nothing: the probability measures form a
measurable set, and almost every measure under $\Lambda$, being a $u_K$, lies in it. A set with
$k < d$ carries no density in $\RR^d$, and the next two paragraphs deal with that.
% CHANGED Q-0281 wave 2: blind review W2B-1 (the reading step: W_d universally measurable; Lambda
% on all measures, carried by the probability measures).
% CHANGED Q-0424 (2026-10-12; the external review of the isotropic paper, round 1, required change
% on the proof of Proposition 3.14; records/affine/reviews/VERIFY-round1-mathematics.md, item 6(a)):
% the reading step is proved. Until now the proof asserted that K -> u_K is a Borel map and that
% W_d, being analytic, is universally measurable, with no argument and no reference for either.
% The paragraph "The set W_d is a Borel set" below shows more by less: W_d is a countable union of
% weakly compact sets. It uses Blaschke's selection theorem and no descriptive set theory, so the
% open question of a reference for "an analytic set is universally measurable" is gone.

\emph{The radial part of one uniform law.} Let $\omega_d$ be the volume of the unit ball and
$\varphi(x) := \omega_d|x|^d$ the volume of the ball of radius $|x|$. For $v \ge 0$ put
$r_v := (v/\omega_d)^{1/d}$ and $B(v) := \varphi^{-1}((0,v]) = \{0 < |x| \le r_v\}$, the
punctured closed ball of volume $v$, and put $G_K(v) := u_K(B(v))$. We show that $G_K$ is concave
on $[0,\infty)$. If $k = 0$, then $K = \{0\}$, $u_K = \delta_0$ and $G_K \equiv 0$, as $B(v)$
misses the origin. Let $k \ge 1$, let $S_V$ be the unit sphere of $V$ and $\sigma_V$ its surface
measure, so that $\mathrm{Vol}_V(A) = \int_{S_V}\sigma_V(d\theta)\int_0^\infty
\mathbf{1}_A(s\theta)\,s^{k-1}\,ds$ for every Borel $A \subseteq V \setminus \{0\}$ (polar
coordinates in $V$). Fix $\theta \in S_V$. The set $J_\theta := \{s > 0 : s\theta \in K\}$ contains
$(0, s]$ with each of its points $s$: $K$ contains $0 = \tfrac12 y + \tfrac12(-y)$ for any
$y \in K$, so $s\theta \in K$ and $0 < s' < s$ give
$s'\theta = \tfrac{s'}{s}\,s\theta + \bigl(1 - \tfrac{s'}{s}\bigr)0 \in K$ by convexity. It is
bounded, $K$ being compact. With $\rho_\theta := \sup J_\theta$ ($\rho_\theta := 0$ if
$J_\theta$ is empty), $(0,\rho_\theta) \subseteq J_\theta \subseteq (0,\rho_\theta]$, and for
$s > 0$, $s\theta \in B(v)$ exactly when $s \le r_v$. So the ray in the direction $\theta$
contributes
\[
  \int_0^\infty\mathbf{1}_{K \cap B(v)}(s\theta)\,s^{k-1}\,ds
   = \int_0^{\min\{\rho_\theta,\,r_v\}}s^{k-1}\,ds
   = \frac1k\min\Bigl\{\rho_\theta^k,\ \Bigl(\frac{v}{\omega_d}\Bigr)^{k/d}\Bigr\} .
\]
As $k \le d$, $v \mapsto (v/\omega_d)^{k/d}$ is concave on $[0,\infty)$, and so is its minimum with
a constant. Integrating over $\theta$ against $\sigma_V$ and dividing by
$\mathrm{Vol}_V(K) \in (0,\infty)$, $G_K(v) = \mathrm{Vol}_V(K \cap B(v))/\mathrm{Vol}_V(K)$ is
concave in $v$. The volume is positive because $K$ spans $V$, so it contains a basis
$v_1,\dots,v_k$ of $V$ and, being convex and symmetric, the convex hull of the $\pm v_i$, which has
positive volume in $V$; it is finite because $K$ is compact. (Concavity survives multiplication by
any nonnegative constant, so the formal proof does not need this; that $u_K$ is a probability
measure is checked separately.) In words: each ray of $K$ contributes the minimum of a constant and
a concave power of the volume, the power being $k/d \le 1$ because $K$ has dimension at most $d$.
% CHANGED Q-0281 wave 2: blind review W2B-4 (the reason for Vol_V(K) in (0, infty); the check
% lemma AffineHemigroup.isProbabilityMeasure_kanterUniform, guarded, Lean core).

\emph{The density.} Put $G(v) := \mu(B(v)) = \int G_K(v)\,\Lambda(du_K)$. A mixture of concave
functions is concave, the inequality $aG_K(v) + bG_K(w) \le G_K(av + bw)$ integrating against
$\Lambda$, so $G$ is concave on $[0,\infty)$. Let $m$ be the image of $\mu$ under $\varphi$, a
finite measure on $\RR$ with $m((0,v]) = G(v)$ and $m((-\infty,0]) = \mu(\{0\}) = 0$. $G$ is
nondecreasing; let $D(v)$ be its right derivative at $v > 0$. It is finite, nonnegative and
nonincreasing, and $G$ is continuous on $(0,\infty)$, so by the fundamental theorem of calculus
for right derivatives $m((a,b]) = G(b) - G(a) = \int_a^b D(w)\,dw$ for $0 < a \le b$. Letting
$a \downarrow 0$ by continuity from below, and as $m$ does not charge $(-\infty,0]$, the measures
$m$ and $D(w)\,dw$ on $(0,\infty)$ agree on every interval $(a,b]$, so $m = D(w)\,dw$ on
$(0,\infty)$. Put $q(r) := D(\omega_d r^d)$, a nonincreasing function on $(0,\infty)$ with values
in $[0,\infty)$, so that $q(|x|) = D(\varphi(x))$. The image of Lebesgue measure on $\RR^d$ under
$\varphi$ is Lebesgue measure on $(0,\infty)$, because $\varphi^{-1}((-\infty,v])$ is the closed
ball of volume $v$. So the image of $q(|x|)\,dx$ under $\varphi$ is $D(w)\,dw = m$, the image of
$\mu$, and as $|x| = (\varphi(x)/\omega_d)^{1/d}$ the two measures have the same image under
$x \mapsto |x|$. Both are finite and rotation invariant ($q(|x|)\,dx$ because Lebesgue measure is and
the density is radial), and both vanish at the origin. A measure with these properties is
determined by its image under $x \mapsto |x|$. For let $m$ be one of them and
$\varepsilon > 0$, and let $m_\varepsilon$ be its restriction to
$U_\varepsilon := \{|x| > \varepsilon\}$, a finite measure invariant under every $Q \in \Orth(d)$,
as $U_\varepsilon$ is. Then $\widehat{m_\varepsilon}(\xi) = \widehat{m_\varepsilon}(|\xi|\theta)$
for every unit vector $\theta$, a reflection carrying $\xi$ to $|\xi|\theta$. Averaging over
$\theta$ against $\sigma_{d-1}/|S^{d-1}|$ and exchanging the integrals gives
\[
  \widehat{m_\varepsilon}(\xi)
   = \frac{1}{|S^{d-1}|}\int_{U_\varepsilon}\Phi(|\xi|,|x|)\,m(dx) ,
  \qquad
  \Phi(s,r) := \int_{S^{d-1}} e^{i s r\langle e_1,\theta\rangle}\,\sigma_{d-1}(d\theta) ,
\]
since $\int e^{i s\langle x,\theta\rangle}\sigma_{d-1}(d\theta)$ depends on $x$ only through $|x|$,
by the invariance of $\sigma_{d-1}$. The right side depends on $m$ only through its image under
$x \mapsto |x|$. So the two measures have restrictions to $U_\varepsilon$ with the same transform,
hence equal; letting $\varepsilon \downarrow 0$ along $1/(n+1)$, they agree off the origin, and as
neither charges it, they are equal. Hence $\mu = q(|x|)\,dx$, and $p(x) := q(|x|)$ is a density of
$\mu$.
% CHANGED Q-0281 wave 2: blind review W2B-3 (the determination by the radial image restated, in
% place of a step borrowed from inside the proof of prop:isotropic-polar-reading; the Lean is
% AffineHemigroup.eq_of_rotationInvariant).

% ADDED Q-0424 (2026-10-12): the argument of records/affine/reviews/VERIFY-round1-mathematics.md,
% item 6(a), written out and checked again step by step. Two differences from the record: the
% subspaces converge through their orthogonal projections, so that the Grassmannian is not named,
% and the continuity is proved on each compact family, which is all that the last step uses.
% CHANGED Q-0430 (2026-10-12): citations only, no step of the argument changes. Schneider, Convex
% Bodies: The Brunn-Minkowski Theory, 2nd ed., is held and was read at the page images
% (records/affine/READING-schneider2014convex.md; CITATIONS-primary-sources-B.md, item 1):
% Thm. 1.8.7, p. 63 (Blaschke's selection theorem, for convex bodies, which there are the
% nonempty compact convex sets, p. 8); Lemma 1.8.14, p. 66 (the Hausdorff distance is the
% supremum of the difference of the support functions); p. 45 (h_K <= h_L iff K in L, and
% h(lambda K) = lambda h(K)). Not cited: (1.60), p. 61, the Hausdorff distance in the form of
% the first inclusion, read in the text layer only.
% The identification of the Borel sets of weak convergence with the sets of the evaluation
% sigma-algebra, used in the paragraph on the mixture, is in the annotation of thm:cited-A8 as
% Kanter's convention and has no reference (it is not in Schneider).
\emph{The set $W_d$ is a Borel set.} This paragraph proves the reading of the source made in the
paragraph on the mixture, and no other step uses it. Let $\bar B$ be the closed unit ball of
$\RR^d$, and give the nonempty compact convex sets $K = -K$ of $\RR^d$ the Hausdorff metric. For
integers $k \ge 1$ and $n \ge 1$ let $\mathcal{K}_{k,n}$ be the set of those $K$ whose span $V$
has dimension $k$ and which satisfy
\[
  \tfrac1n\bar B \cap V \subseteq K \subseteq n\bar B .
\]
Every $K$ of dimension $k \ge 1$ is in $\mathcal{K}_{k,n}$ for some $n$. It is bounded, and it
contains the convex hull of the points $\pm v_i$ of the paragraph on one uniform law, which
contains a ball of $V$ round the origin.

Each $\mathcal{K}_{k,n}$ is compact. Let $K_j \in \mathcal{K}_{k,n}$, with spans $V_j$. By
Blaschke's selection theorem (\sscite{schneider2014convex}, Thm.~1.8.7, p.~63, where a convex
body is any nonempty compact convex set, p.~8), the nonempty compact convex subsets of a ball
form a compact space in the Hausdorff metric. So a subsequence of $(K_j)$ converges to a nonempty compact convex set
$K$, and $K = -K$. Let $P_j$ be the orthogonal projection onto $V_j$. Along a further subsequence
the matrices $P_j$ converge, and their limit $P$ is the orthogonal projection onto a subspace $V'$
of dimension $k$, since $P_j^2 = P_j = P_j^{\tp}$ and $\tr P_j = k$ pass to the limit. A point
$x \in K$ is a limit of points $x_j \in K_j$, and $x_j = P_jx_j \to Px$, so $x \in V'$. Hence
$K \subseteq V' \cap n\bar B$. For $y \in V'$ with $|y| \le 1/n$ the points $P_jy$ lie in
$\tfrac1n\bar B \cap V_j \subseteq K_j$ and converge to $y$, so $y \in K$. Hence the span of $K$
is $V'$, and $K \in \mathcal{K}_{k,n}$.

The map $K \mapsto u_K$ is continuous on $\mathcal{K}_{k,n}$, for the weak convergence of measures.
Let $K_j \to K$ in $\mathcal{K}_{k,n}$, with spans $V_j$ and $V$. Choose a basis $y_1,\dots,y_k$ of
$V$ in $K$ and points $y_{i,j} \in K_j$ with $y_{i,j} \to y_i$ as $j \to \infty$. For large $j$
the points $y_{1,j},\dots,y_{k,j}$ are linearly independent, their Gram determinant converging to
that of the $y_i$, which is positive. So they are a basis of $V_j$. Let $T_j$ be the linear map
from $V$ onto $V_j$ with $T_jy_i = y_{i,j}$, and let $\varepsilon_j$ be the supremum of $|T_jz - z|$
over the $z \in V$ with $|z| \le 1$. Then $\varepsilon_j \to 0$. A linear bijection carries the
uniform law of a set to the uniform law of its image, so $u_{K_j}$ is the image of $u_{K'_j}$ under
$T_j$, where $K'_j := T_j^{-1}K_j \subseteq V$. For $w \in K'_j$, $|T_jw - w| \le \varepsilon_j|w|$
and $|T_jw| \le n$, so $|w| \le n/(1-\varepsilon_j)$, and the Hausdorff distance of $K'_j$ and $K_j$
is at most $\varepsilon_jn/(1-\varepsilon_j)$. So $K'_j \to K$. Let $\eta_j$ be the Hausdorff
distance of $K'_j$ and $K$, and put $r := 1/n$ and $B_V := \bar B \cap V$, so that
$rB_V \subseteq K$. Then
\[
  K'_j \subseteq K + \eta_jB_V \subseteq \bigl(1 + \eta_j/r\bigr)K ,
\]
$K$ being convex. In the other direction, write $h_A(e) := \sup_{x \in A}\langle e,x\rangle$ for the
support function of a compact convex set $A \subseteq V$ at a unit vector $e$ of $V$. Then
$h_K \le h_{K'_j} + \eta_j$, the Hausdorff distance of two compact convex sets being the supremum
over the unit vectors of the absolute difference of their support functions
(\sscite{schneider2014convex}, Lemma~1.8.14, p.~66), and $h_K \ge r$, so
$h_{K'_j} \ge (1 - \eta_j/r)\,h_K$. A compact
convex set contains every compact convex set with a smaller support function, and
$h_{\lambda A} = \lambda h_A$ for $\lambda \ge 0$ (\sscite{schneider2014convex}, p.~45), so
$(1 - \eta_j/r)K \subseteq K'_j$ once $\eta_j < r$. The two inclusions hold for $K$ in place of
$K'_j$ as well. So $K'_j$ and $K$ differ by a set of volume at most
$\bigl((1+\eta_j/r)^k - (1-\eta_j/r)^k\bigr)\mathrm{Vol}_V(K)$, which tends to $0$, and
$u_{K'_j}(A) \to u_K(A)$ uniformly in the Borel set $A$. For a bounded continuous function
$\varphi$ on $\RR^d$,
\[
  \int\varphi\,du_{K_j} = \int\varphi(T_jw)\,u_{K'_j}(dw) ,
\]
and $\varphi(T_jw) \to \varphi(w)$ uniformly for $w$ in the bounded set $2K$, which contains $K'_j$
for large $j$. So $\int\varphi\,du_{K_j} \to \int\varphi\,du_K$.

So each set $\{u_K : K \in \mathcal{K}_{k,n}\}$ is the continuous image of a compact set. It is
weakly compact, hence weakly closed. With $\{\delta_0\}$, the case $k = 0$, these countably many
sets make up $W_d$. So $W_d$ is a Borel set for weak convergence.

\emph{What is machine-checked.} Everything above but the paragraph on $W_d$, with \ledger{A8} as
the interface it cites, and,
for clause (2), \ledger{A2} and \ledger{A3}(iii) through \ref{prop:isotropic-polar-reading}(1). The
paragraph on $W_d$ is about the passage from the source's form of \ledger{A8} to the form of
\ref{thm:cited-A8}, which the formal proof takes as stated. It is proved here in prose, with
Blaschke's selection theorem as a classical theorem, and is not formalized.
% CHANGED Q-0377: the last sentence was stale since Q-0365, which moved the uniqueness of the Levy
% pair from Paper V's interface to this module's A3 (prop:line-levy-unique); the guard prints
% A3(ii)-(iii) and A8 on radial_monotone_family. Annotation only; no step of the proof changes.
Clause (1) reads \ref{thm:isotropic-classification}, whose formal proof spends \ledger{A3}(ii)--(iii)
at $d = 1$ through \ref{prop:line-levy-unique}.
\end{proof}

% Q-0415 (the author's decision of 2026-10-11, "it should be at least part of the blueprint and
% preferably Lean-checked"). The upper half that prop:origin-lower-bound-d leaves out, from two
% statements the chapter already has: the upper bound of thm:origin-boundedness-d(4) on the heat
% average, and the radially nonincreasing density of prop:isotropic-radial-monotone(2). The
% derivation is the scope audit's of the same day (finding 1 of
% records/affine/SCOPE-AUDIT-paper-iso-2026-10-11.md, branch queue/q-0414), checked on paper before
% anything was typed. The node sits here, after prop:isotropic-radial-monotone, because it reads
% that node and prop:origin-lower-bound-d; the isotropic paper prints it after the lower bound.
%
% The form, decided on paper before anything was typed.
% (i) A new node, not a clause of an existing one. The upper bound reads the monotone density
% (A8), which thm:origin-boundedness-d and prop:origin-lower-bound-d do not spend; a clause there
% would have regraded a shared statement of the paper.
% (ii) The proof goes through the mass of a ball, and that step is clause (1): the heat kernel
% g_{r^2} is at least (2 pi)^{-d/2} e^{-1/2} r^{-d} on {|x| < r}, so mu{|x| < r} is at most a
% constant times r^d times the heat average at t = r^2. This is a statement about the measure. It
% needs no density and does not read A8, and it is the counterpart of
% prop:origin-lower-bound-d(2). The audit's Step 1 integrated the heat kernel over the ball
% (P(|Z| < 1)); bounding it below on the ball gives a larger constant and needs no change of
% variables.
% (iii) Clause (2) is stated for EVERY nonincreasing q with mu = q(|x|) dx, and at every r > 0.
% The existence of such a q is prop:isotropic-radial-monotone(2); the bound for a given q spends
% nothing. Two such q agree off a countable set of radii, and the bound holds at every radius for
% each.
% (iv) The hypothesis is 0 < c <= d, that of thm:origin-boundedness-d(4), whose upper bound the
% proof reads. For c > d the same argument gives the bounds for beta < d only, which is weaker
% than the boundedness of the density, and the limit of clause (3) is false there (the density is
% bounded above, and below by a positive constant, near the origin; the quotient tends to 0).
% c = infinity is excluded with it.
% (v) Clause (3) needs d >= 2 and nothing else new: it is where prop:origin-lower-bound-d enters.
% Its lower bound is an almost-everywhere statement for a general density; for a nonincreasing q it
% holds at every radius, which the clause says and the limit needs. The limit for the mass of a
% ball is stated beside the one for the density because it names no representative and does not
% read A8.
% Priced when the types were written: about 400 Lean lines, none of them a new measure-theoretic
% bridge (the density enters through withDensity_apply on a ball, which asks no measurability of q).
% Axioms, printed before the annotation was written (lake env lean CIAxiomGuard.lean, exit 0):
% origin_upper_bound and origin_upper_bound_of_heat print propext, Classical.choice, Quot.sound;
% origin_density_lower, origin_density_log_order and origin_ball_log_order add
% levy_khintchine_converse (A3(iii)); origin_density_upper and origin_density_order add to that
% sato_polar_criterion (A2) and wolfe_symmetric_unimodal (A8).
% CHANGED Q-0416 (2026-10-11): annotation only, two sentences. "Not claimed ... these are the open
% problem of the isotropic paper's section 4" and "A sharper form, not stated ... needs the theory
% of regular variation, which neither the chapter nor the Lean library has" were true when written;
% both now point to prop:origin-rate-d. The statement, its proof and the \lean tag are unchanged.
% CHANGED Q-0425 (2026-10-12): annotation only, one sentence added to "Not claimed": the bound at
% beta = c is decided for c < d by prop:origin-rate-profile-d(2). The statement, its proof and the
% \lean tag are unchanged.
\begin{proposition}[an upper bound at the origin, and the order of the density there]
\label{prop:origin-upper-bound-d}
\uses{def:isotropic-origin-datum,thm:origin-boundedness-d,lem:positive-transform-origin,prop:isotropic-radial-monotone,prop:origin-lower-bound-d,thm:cited-A2,thm:cited-A3,thm:cited-A8}
\lean{AffineHemigroup.origin_upper_bound, AffineHemigroup.origin_upper_bound_of_heat,
AffineHemigroup.origin_density_upper, AffineHemigroup.origin_density_lower,
AffineHemigroup.origin_density_log_order, AffineHemigroup.origin_ball_log_order,
AffineHemigroup.origin_density_order}\leanok
Let $d \ge 1$, let $(0, k_d)$ be an isotropic origin datum in dimension $d$ with Gaussian
coefficient $0$, let $\mu$ be its law, and suppose that $c := k_d(0{+})$ satisfies $0 < c \le d$.
Let $V_d$ be the volume of the unit ball of $\RR^d$.
\begin{enumerate}
\item \emph{(The mass of a ball.)} For every $\beta$ with $0 < \beta < c$ there is
$A_\beta < \infty$ with
\[
  \mu(\{|x| < r\}) \le A_\beta\,r^{\beta} \qquad\text{for every } r > 0 .
\]
One may take $A_\beta = (2\pi)^{d/2}e^{1/2}\,C'_\beta$, for every constant $C'_\beta$ with which
the upper bound of \ref{thm:origin-boundedness-d}(4) holds.
\item \emph{(The density.)} The law $\mu$ has a density $p$ with $p(x) = q(|x|)$ for $x \ne 0$,
where $q : (0,\infty) \to [0,\infty)$ is nonincreasing. For every such $q$, every $0 < \beta < c$, and
$A_\beta = (2\pi)^{d/2}e^{1/2}\,C'_\beta$ with $C'_\beta$ as in clause (1),
\[
  q(r) \le \frac{A_\beta}{V_d}\,r^{\beta-d} \qquad\text{for every } r > 0 .
\]
\item \emph{(The order at the origin, for $d \ge 2$.)} Let $d \ge 2$, let $\delta$ and $c_1$ be
the constants displayed at the end of \ref{prop:origin-lower-bound-d}, and let $q$ be as in
clause (2). Then
$q(r) \ge c_1\,r^{c-d}$ for every $r$ with $0 < r < \delta$, with no exceptional set. With
$p(x) = q(|x|)$,
\[
  \frac{\log p(x)}{\log(1/|x|)} \to d - c \quad (x \to 0), \qquad
  \frac{\log \mu(\{|x| < r\})}{\log r} \to c \quad (r \downarrow 0) .
\]
\end{enumerate}
\statusA\quad\ledger{A8}\quad\ledger{A2}\quad\ledger{A3}\quad\emph{(\textbf{Assignment.} The node
cites nothing of its own and inherits what the two nodes it joins spend. The density of clause (2)
exists by \ref{prop:isotropic-radial-monotone}(2), which spends \ledger{A8} (Wolfe's theorem),
\ledger{A2} and \ledger{A3}(iii). The lower bound of clause (3) is
\ref{prop:origin-lower-bound-d}, which spends \ledger{A3}(iii). Clause (1), and the bound of
clause (2) for a given $q$, spend nothing: the law is any probability measure with transform
$e^{-F_d}$ and $q$ is any nonincreasing function with $\mu = q(|x|)\,dx$, so the existence of
neither is read. The formal statements are graded the same way. \texttt{origin\_upper\_bound}
(clause (1), and clause (2) for a given $q$, with the constants) prints Lean core;
\texttt{origin\_density\_lower}, \texttt{origin\_density\_log\_order} and
\texttt{origin\_ball\_log\_order} (clause (3), for a given $q$) print \ledger{A3}(iii);
\texttt{origin\_density\_upper} and \texttt{origin\_density\_order}, which assert that $q$ exists,
print \ledger{A8}, \ledger{A2} and \ledger{A3}(iii). \textbf{What the citations carry and what they
do not.} \ledger{A8} carries the mixture of uniform laws from which
\ref{prop:isotropic-radial-monotone} derives a nonincreasing $q$, and nothing about the size of
$q$. \ledger{A3}(iii) carries the existence of the law of a truncated datum inside
\ref{prop:origin-lower-bound-d}. The bounds on $q$ are proved, the lower one in that node and the
upper one here. \textbf{What it is
worth.} With \ref{prop:origin-lower-bound-d} the spike of the density is bounded from both sides
near the origin, $c_1|x|^{c-d} \le p(x) \le (A_\beta/V_d)\,|x|^{\beta-d}$ for every $\beta < c$,
and its logarithmic order is exactly $d - c$. Before this node the chapter had that order for the
density blurred by the heat kernel only, the last sentence of \ref{thm:origin-boundedness-d}(4),
and the annotation of \ref{prop:origin-lower-bound-d} used the same upper bound to show that the
exponent of the lower bound is sharp. The second limit of clause (3) says the same of the measure
and names no density: the mass of a small ball has the logarithmic order $c$. \textbf{Where it
comes from.} The scope audit of the isotropic paper of 2026-10-11 found that the paper's first
open problem asked, as its weaker alternative, for an upper bound that
\ref{thm:origin-boundedness-d}(4) and \ref{prop:isotropic-radial-monotone} already give
(\texttt{records/affine/SCOPE-AUDIT-paper-iso-2026-10-11.md}, finding 1). The proof below is that
derivation, taken through the mass of a ball. \textbf{Edge cases.} At $c = d$ the bounds read
$c_1 \le p(x) \le (A_\beta/V_d)\,|x|^{\beta-d}$ for every $\beta < d$ and the first limit is $0$.
The density is unbounded there by \ref{thm:origin-boundedness-d}(4), and the limit does not decide
whether it grows like $\log(1/|x|)$. The case $c > d$ is excluded, and $c = \infty$ with it. There
the argument of clause (1) gives the two bounds for $\beta < d$ only, which is weaker than the
boundedness of \ref{thm:origin-boundedness-d}(2), and the first limit of clause (3) is false: for
a finite $c > d$ and $d \ge 2$ the density is bounded near the origin, and bounded below there by
a positive constant, $q$ being nonincreasing and positive at small radii, so the quotient tends to
$0$ and not to $d - c < 0$. At $d = 1$ clauses (1) and (2) hold. Clause (3) is for $d \ge 2$
because \ref{prop:origin-lower-bound-d} is, and on the line [B, Cor.~3.9] has more. The bounds of
clauses (1) and (2) hold at every $r > 0$. For large $r$ they say less than
$\mu(\{|x| < r\}) \le 1$ and $q(r) \le (V_d r^d)^{-1}$, which is the first inequality of the proof
of clause (2). \textbf{Not claimed}: the bound at $\beta = c$, that is
$p(x) \le C|x|^{c-d}$. The constant $A_\beta$ is not bounded as $\beta \uparrow c$, and on the line
Sato's rate carries the factor $K$, which is unbounded exactly when
$\int_0^1(c - k(y) - k(-y))\,dy/y = \infty$. Since Q-0425,
\ref{prop:origin-rate-profile-d}(2) decides the bound at $\beta = c$ for $c < d$ in every
dimension: it holds exactly when $\int_0^1(c - k_d(u))\,du/u < \infty$. Nor an asymptotic equivalence
$p(x) \sim C|x|^{c-d}L(|x|)$ with $L$ slowly varying, nor an explicit slowly varying factor in the
upper bound, nor the factor $\log(1/|x|)$ at $c = d$: since Q-0416 the asymptotic equivalence
and the logarithmic growth at $c = d$ are \ref{prop:origin-rate-d}, on one cited fact more
(\ledger{A11}) and without explicit constants or radii. Nor anything for a datum with a Gaussian
part, whose law has a bounded smooth density. \textbf{Prior art.} The method is in print, and so is an upper bound; the statement at
one law is not, as far as was read. The reading is Q-0414's
(\texttt{records/affine/PRIOR-ART-origin-threshold-reading.md} \S 5 and
\texttt{READING-grzywny2019asymptotic.md}, on that item's branch when this was written), and the
two pages cited here were read again at the image for this node. The source is Grzywny, Ryznar
and Trojan, \emph{Asymptotic behaviour and estimates of slowly varying convolution semigroups},
International Mathematics Research Notices 2019, held as the advance-access article paginated 1--66 and cited by its
own numbers. It treats the transition densities $p(t,\cdot)$ of the isotropic unimodal \Levy{}
processes, and the law of a datum of this node is such a process at time $1$. (i) The step from
a radially nonincreasing density to the mass nearby is its ``elementary observation'' (5.3),
p.~43, $r^d p(t,r) \le C\,\Pr(r/2 \le |X_t| < r)$ for all $t, r > 0$; the first inequality of the
proof of clause (2) is the same remark with a ball in place of the shell. (ii) Its (5.4), p.~44,
bounds $r^d p(t,r)$ at every $t$ and $r$ by a Gaussian average of a difference of the transform,
and with \ref{thm:origin-boundedness-d}(1) the bound of clause (2) follows from (5.4) in two
lines. So the comparison of a monotone density with a Gaussian average of the transform is that
paper's method. (iii) Its Theorem~5.4, p.~43, is $p(t,x) \le C\,t\,|x|^{-d}K_d(|x|)$ at all times,
which for a datum with a finite $c$ is an upper bound with the exponent $-d$ (the memo computes
$K_d \le c/2$). So neither this node nor the paper may say that no upper bound on the density
near the centre was known. (iv) The bound with the decay of the transform, hence the exponent
$c - d$ with a slowly varying factor, is in its Proposition~5.6 and Theorem~5.7 for small times
only, $t < t_0$ with $t_0 \le d/(2c)$ by the memo's reading of the proofs; on the line the exact
rate is Sato and Yamazato's (\sscite{sato1999levy}~Thm.~53.8). What this node adds is a statement
at one law, for every $0 < c \le d$ and every exponent $\beta < c$, with an explicit constant,
proved from this chapter's nodes and machine-checked. \textbf{A sharper form, not stated.} The
memo records that for $0 < c < d$ the same first step, with \ref{lem:isotropic-exponent-log}(2)
and Potter's bound for regularly varying functions, gives
$q(r) \le C\,r^{-d}\,\widehat\mu(e/r)$ for small $r$ and a unit vector $e$, that is
$|x|^{c-d}$ times a slowly varying factor. It has no explicit constant, so it is not a clause
here. Since Q-0416 it is contained in \ref{prop:origin-rate-d}(1), where the quotient
$q(r)\,r^{d}/\widehat\mu(e/r)$ has the limit $R(d,c)$; that node proves the bound on the ratios of
the transform that stands for Potter's. \textbf{Readings.} In the Lean the ball is the open ball, the
measure of clause (1) is bounded in the extended nonnegative reals, ``$x \to 0$'' is the punctured
neighbourhood filter of the origin, ``$r \downarrow 0$'' the filter of positive reals at $0$, and
the two quotients are quotients of real logarithms. The numbers under the logarithms are positive
for small arguments, by the lower bounds. A $q$ of clause (2) is a real function, nonincreasing on
$(0,\infty)$, with $\mu$ the measure of density $x \mapsto \max\{q(|x|),0\}$; its nonnegativity is
not asked, which only widens the class.)}
\end{proposition}

\begin{proof}
\leanok
Write $B_r := \{|x| < r\}$ for the open ball and $|E|$ for the Lebesgue measure of $E$, so that
$|B_r| = V_d\,r^d$.

\emph{Clause (1).} Fix $\beta$ with $0 < \beta < c$, and let $C'_\beta$ be a constant with
$\int g_t\,d\mu \le C'_\beta\,t^{-(d-\beta)/2}$ for all $t > 0$, $g_t$ being the heat kernel of
\ref{lem:positive-transform-origin}. There is one by \ref{thm:origin-boundedness-d}(4), since
$c \le d$. Let $r > 0$. For $|x| < r$,
\[
  g_{r^2}(x) = (2\pi r^2)^{-d/2}e^{-|x|^2/2r^2} \ge (2\pi)^{-d/2}e^{-1/2}\,r^{-d} .
\]
As $g_{r^2} \ge 0$,
\[
  \int_{\RR^d} g_{r^2}\,d\mu \ge \int_{B_r} g_{r^2}\,d\mu
   \ge (2\pi)^{-d/2}e^{-1/2}\,r^{-d}\,\mu(B_r) .
\]
The left side is at most $C'_\beta\,(r^2)^{-(d-\beta)/2} = C'_\beta\,r^{\beta-d}$. Hence
$\mu(B_r) \le (2\pi)^{d/2}e^{1/2}\,C'_\beta\,r^{\beta}$, which is the bound with the stated
$A_\beta$.

\emph{Clause (2).} As $c = \sup_{r>0}k_d(r) > 0$, there is $r_0 > 0$ with $k_d(r_0) > 0$, and
$k_d \ge k_d(r_0)$ on $(0,r_0]$ because $k_d$ is nonincreasing. So $k_d$ is not almost everywhere
$0$, the datum is not trivial, and \ref{prop:isotropic-radial-monotone}(2) gives a density
$p(x) = q(|x|)$ with $q$ nonincreasing. Now let $q : (0,\infty) \to [0,\infty)$ be any
nonincreasing function such that $x \mapsto q(|x|)$ is a density of $\mu$, and let $r > 0$. For
$0 < |x| < r$, $q(|x|) \ge q(r)$. The origin is a Lebesgue-null set, so
\[
  \mu(B_r) = \int_{B_r} q(|x|)\,dx \ge q(r)\,|B_r| = q(r)\,V_d\,r^d .
\]
With clause (1), $q(r) \le A_\beta V_d^{-1}\,r^{\beta-d}$.

\emph{Clause (3), the lower bound at every radius.} Let $d \ge 2$. Since $c \le d$, $c$ is finite,
and \ref{prop:origin-lower-bound-d} applies. By its clause (1) the density $x \mapsto q(|x|)$ of
$\mu$ is at least $c_1|x|^{c-d}$ at almost every $x \in B_\delta$. Suppose that
$q(r_0) < c_1 r_0^{c-d}$ for some $r_0$ with $0 < r_0 < \delta$. The function
$r \mapsto c_1 r^{c-d}$ is continuous at $r_0$, so there is $r_1$ with $r_0 < r_1 \le \delta$ and
$q(r_0) < c_1 r^{c-d}$ for $r_0 < r < r_1$. For such $r$, $q(r) \le q(r_0) < c_1 r^{c-d}$,
because $q$ is nonincreasing. So $q(|x|) < c_1|x|^{c-d}$ on the shell $\{r_0 < |x| < r_1\}$. The
shell is a nonempty open subset of $B_\delta$, hence of positive Lebesgue measure, and this
contradicts the bound at almost every point. So $q(r) \ge c_1 r^{c-d}$ for every $0 < r < \delta$.

\emph{Clause (3), the limits.} Both follow from one remark. Let $L$ be real and let $f$ be a
function on $(0,\infty)$ such that, for small $r > 0$, $f(r) \ge C r^{-L}$ with a $C > 0$ and, for
every $\varepsilon > 0$, $f(r) \le C_\varepsilon r^{-(L+\varepsilon)}$ with a
$C_\varepsilon > 0$. Then $\log f(r)/\log(1/r) \to L$ as $r \downarrow 0$. Indeed, for
$0 < r < 1$ small, $f(r) > 0$ and $\log(1/r) > 0$, and taking logarithms gives
\[
  L + \frac{\log C}{\log(1/r)} \le \frac{\log f(r)}{\log(1/r)}
   \le L + \varepsilon + \frac{\log C_\varepsilon}{\log(1/r)} .
\]
As $\log(1/r) \to \infty$, the quotient is at least $L - \varepsilon$ and at most
$L + 2\varepsilon$ once $r$ is small enough, for every $\varepsilon > 0$.

For the density take $f = q$ and $L = d - c$. The lower bound is the one just proved,
$q(r) \ge c_1 r^{-(d-c)}$ for $0 < r < \delta$. For the upper bound let $\varepsilon > 0$ and
choose $\beta$ with $0 < \beta < c$ and $\beta \ge c - \varepsilon$. By clause (2), for
$0 < r < 1$,
$q(r) \le A_\beta V_d^{-1}\,r^{\beta-d} \le A_\beta V_d^{-1}\,r^{-(d-c+\varepsilon)}$, as
$\beta - d \ge -(d - c + \varepsilon)$ and $r < 1$. So $\log q(r)/\log(1/r) \to d - c$, and with
$r = |x|$ this is the first limit, since $|x| \downarrow 0$ through positive values as $x \to 0$.

For the mass of a ball take $f(r) = \mu(B_r)$ and $L = -c$. By
\ref{prop:origin-lower-bound-d}(2), $\mu(B_r) \ge c_1|S^{d-1}|\,c^{-1}\,r^{c}$ for
$0 < r \le \delta$. By clause (1), with $\beta$ as before,
$\mu(B_r) \le A_\beta\,r^{\beta} \le A_\beta\,r^{c-\varepsilon}$ for $0 < r < 1$. So
$\log\mu(B_r)/\log(1/r) \to -c$, and $\log(1/r) = -\log r$ gives the second limit.

\emph{What is machine-checked.} Everything above. The cited facts enter through the two nodes the
proof reads: \ledger{A8}, \ledger{A2} and \ledger{A3}(iii) for the existence of $q$ in clause (2),
through \ref{prop:isotropic-radial-monotone}(2), and \ledger{A3}(iii) for the lower bounds of
clause (3), through \ref{prop:origin-lower-bound-d}.
\end{proof}

% Q-0416 (the author's decision of 2026-10-11, "adopt the rate as a proposition"). New node: the
% exact rate of the mass of a small ball and of the radially nonincreasing density at the origin,
% for 0 < c <= d. The derivation is the reading memo's (records/affine/
% PRIOR-ART-origin-threshold-reading.md section 6, on Q-0414's branch), refereed blind before this
% item (records/affine/REFEREE-rate-at-the-centre.md: correct as stated, nine presentational
% repairs R1-R9). The proof of record below is the derivation as the referee leaves it, on the
% route the Lean takes, which differs from the memo's in four places, each in the direction of
% citing less.
% (i) Potter's bound (Bingham-Goldie-Teugels Thm. 1.5.6) is not cited. The referee's substitute
% (report section 6, item 3) is used: the transform is nonincreasing along rays, so the limit of
% lem:isotropic-exponent-log(2) at the one ratio a = 2 gives, by dyadic iteration, the ratio bound
% of Step 1. The referee's claim that F_d is nondecreasing along rays was checked: it is the
% inequality 0 <= phi(as) - phi(s), a >= 1, of that lemma's proof, averaged over the sphere. The
% lower bound of lem:isotropic-exponent-log(1) is not needed for the limits either (the referee
% uses it on the small ball |eta| < X/sqrt(lambda)): the ratio bound of Step 1 covers that range.
% (ii) Proposition 1.5.9a of the same book (the integral of a slowly varying function against
% dt/t is slowly varying) is not cited: Step 3' proves what clause (2) needs of it, for the one
% function at hand, from the monotonicity of the transform along rays. The integration by parts
% of the memo's section 6.2 and of the referee's R6 is replaced by a two-sided comparison of the
% Gaussian average with the truncated integral.
% (iii) The monotone density theorem (the same book, Thm. 1.7.2b) is not cited: Step 5 is its
% proof, (1.7.5) on p. 39 of the book, carried out on spherical shells in the radial variable. So
% the volume variable of the memo's Step C, and the referee's R1 and R4 on it, do not arise.
% (iv) One cited fact remains: Karamata's Tauberian theorem at the origin (Thm. 1.7.1' of the
% book, p. 38), the new ledger entry A11, used once in each clause.
% The form, decided when the types were written.
% (a) The normaliser is the transform itself, mu-hat(e/s), not s^c times a slowly varying factor:
% the two are the same by lem:isotropic-exponent-log(2), and the statement with the transform
% names no auxiliary function.
% (b) Each clause states the rate of the mass of a ball first. It names no density and does not
% read A8, as clause (1) of prop:origin-upper-bound-d does not; it is Step 4.
% (c) The density clause is for EVERY nonincreasing q with mu = q(|x|) dx, at every radius, as in
% prop:origin-upper-bound-d(2).
% (d) Clause (2) is normalised by the integral of the transform over the ball of radius 1/s, which
% is what the density of an integrable transform would be at the origin, cut off at the frequency
% 1/s. Its logarithmic lower bound is lem:isotropic-exponent-log(1) in polar coordinates.
% Priced when the types were written: about 1400 Lean lines in two modules. It came to 1640:
% AffineHemigroup.RegularVariation (530 lines, Mathlib types only: Steps 1, 3 and 3') and
% AffineHemigroup.OriginRate (1110 lines: the transform along a ray, Step 2, Step 4 on A11,
% Step 5, the constants and the eight declarations of the tag). What writing the types found:
% (e) the lower bound of lem:isotropic-exponent-log(1) is not needed for the limits (it is used
% for the logarithmic bound of clause (2) only), because Step 1's ratio bound also covers the
% range Ru < X; (f) the monotone density step needs no change to the volume variable when it is
% written on shells; (g) a q of the statement need not be nonnegative, and Step 5 then works with
% max(q,0), which equals q near the origin once the limit is known to be positive; (h) the ray
% profile is taken on all of R (positive and at most 1 everywhere, nonincreasing on the positive
% half-line), which is what the transform gives and what makes the truncated moment a continuous
% function with no case distinction at 0.
% Axioms, printed before the annotation was written (lake env lean CIAxiomGuard.lean, exit 0):
% origin_ball_rate, origin_density_rate, origin_ball_rate_critical, origin_density_rate_critical
% and origin_density_log_lower print propext, Classical.choice, Quot.sound and
% karamata_tauberian_zero (A11); origin_transform_ball_log prints the first three only;
% origin_rate and origin_rate_critical add to A11 levy_khintchine_converse (A3(iii)),
% sato_polar_criterion (A2) and wolfe_symmetric_unimodal (A8). tendsto_radial_density_of_ball
% (Step 5), antitoneOn_dyadic_bound (Step 1) and IsRayProfile.tendsto_truncMoment_dilate
% (Step 3'(ii)) print Lean core, and are pinned.
% CHANGED Q-0425 (2026-10-12): annotation only. "Not claimed: an explicit form of the slowly varying
% factor L" now points to prop:isotropic-exponent-profile and prop:origin-rate-profile-d. The
% statement, its proof and the \lean tag are unchanged.
% CHANGED Q-0430 (2026-10-12): annotation only, the passage Prior art. Clause (1) is cited to its
% primary statement (Leonenko's book, Thms. 1.4.3 and 1.4.4, with Leonenko and Olenko 1991), and
% the Hankel route to Soni and Soni 1975 with Bingham and Inoue, each anchor read at the page
% image (records/affine/CITATIONS-primary-sources-A.md). The sentence on what was not found in
% what was read is gone. The statement, its proof and the lean tag are unchanged.
\begin{proposition}[the rate at the origin]
\label{prop:origin-rate-d}
\uses{def:isotropic-origin-datum,lem:positive-transform-origin,lem:isotropic-exponent-log,prop:isotropic-radial-monotone,thm:cited-A11,thm:cited-A2,thm:cited-A3,thm:cited-A8}
\lean{AffineHemigroup.origin_ball_rate, AffineHemigroup.origin_density_rate,
AffineHemigroup.origin_ball_rate_critical, AffineHemigroup.origin_density_rate_critical,
AffineHemigroup.origin_transform_ball_log, AffineHemigroup.origin_density_log_lower,
AffineHemigroup.origin_rate, AffineHemigroup.origin_rate_critical}\leanok
Let $d \ge 1$, let $(0, k_d)$ be an isotropic origin datum in dimension $d$ with Gaussian
coefficient $0$, let $\mu$ be its law, and suppose that $c := k_d(0{+})$ satisfies $0 < c \le d$.
Let $e$ be a unit vector of $\RR^d$, let $|S^{d-1}|$ be the surface measure of the unit sphere and
$V_d$ the volume of the unit ball. The law $\mu$ has a density $p$ with $p(x) = q(|x|)$ for
$x \ne 0$, where $q : (0,\infty) \to [0,\infty)$ is nonincreasing. Let $q$ be any such function.
All limits are as $s \downarrow 0$.
\begin{enumerate}
\item \emph{(Below the threshold.)} Let $c < d$ and put
\[
  R(d,c) := \frac{\Gamma\bigl(\tfrac{d-c}{2}\bigr)}{2^{c}\,\pi^{d/2}\,\Gamma\bigl(\tfrac c2\bigr)} .
\]
Then
\[
  \frac{\mu(\{|x| \le s\})}{R(d,c)\,|S^{d-1}|\,c^{-1}\,\widehat\mu(e/s)} \to 1
  \qquad\text{and}\qquad
  \frac{q(s)}{R(d,c)\,s^{-d}\,\widehat\mu(e/s)} \to 1 .
\]
\item \emph{(At the threshold.)} Let $c = d$ and put
$I(s) := \int_{\{|\xi| \le 1/s\}}\widehat\mu(\xi)\,d\xi$. Then
\[
  \frac{\mu(\{|x| \le s\})}{(2\pi)^{-d}\,V_d\,s^{d}\,I(s)} \to 1
  \qquad\text{and}\qquad
  \frac{q(s)}{(2\pi)^{-d}\,I(s)} \to 1 .
\]
Moreover $I(s) \ge |S^{d-1}|\,e^{-13d/4 - 2C_k}\log(1/s)$ for $0 < s \le 1$, with
$C_k = \int_1^\infty k_d(r)\,dr/r$. So for every
$C < (2\pi)^{-d}\,|S^{d-1}|\,e^{-13d/4 - 2C_k}$ there is $s_0 > 0$ with
$q(s) \ge C\log(1/s)$ for $0 < s < s_0$.
\end{enumerate}
\statusA\quad\ledger{A11}\quad\ledger{A8}\quad\ledger{A2}\quad\ledger{A3}\quad\emph{(\textbf{Assignment.}
\ledger{A11} is Karamata's Tauberian theorem at the origin (\ref{thm:cited-A11}). It carries one
step of the proof, Step 4: from the behaviour of $\int e^{-\lambda|x|^2}\mu(dx)$ as
$\lambda \to \infty$ to the mass of small balls. It is spent once in each clause. The density of
the statement exists by \ref{prop:isotropic-radial-monotone}(2), which spends \ledger{A8},
\ledger{A2} and \ledger{A3}(iii). The limits for a given $q$ do not read its existence. The formal
statements are graded accordingly. \texttt{origin\_ball\_rate}, \texttt{origin\_density\_rate},
\texttt{origin\_ball\_rate\_critical}, \texttt{origin\_density\_rate\_critical} and
\texttt{origin\_density\_log\_lower}, in which the law and $q$ are given, print \ledger{A11} and
no other entry. \texttt{origin\_transform\_ball\_log}, the bound on $I(s)$, prints Lean core.
\texttt{origin\_rate} and \texttt{origin\_rate\_critical}, which assert that $q$ exists, print
\ledger{A11}, \ledger{A8}, \ledger{A2} and \ledger{A3}(iii). \textbf{What the citation carries and
what it does not.} \ledger{A11} carries the Tauberian implication for a finite measure on
$[0,\infty)$, with a positive index and a positive constant. It does not carry the behaviour of
the Gaussian average, which is Steps 1 to 3 below, the regular variation of the transform, which
is \ref{lem:isotropic-exponent-log}(2), or the passage from the mass of balls to the density,
which is Step 5. Three results of the same book that the derivation could cite are proved
instead, each in the form the proof needs: Potter's bound (Thm.~1.5.6) is replaced by Step 1,
Proposition~1.5.9a by Step 3$'$, and the monotone density theorem (Thm.~1.7.2b) by Step 5, which
is its proof on spherical shells. \textbf{What it is worth.} By
\ref{lem:isotropic-exponent-log}, $\widehat\mu(\rho e) = \rho^{-c}L(\rho)$ with $L$ slowly varying
at infinity and $L \ge e^{-13c/4 - 2C_k}$ on $[1,\infty)$. So clause (1) says
$p(x) \sim R(d,c)\,|x|^{c-d}L(1/|x|)$ as $x \to 0$: the density is asymptotic to $|x|^{c-d}$ times
a slowly varying factor, and the factor is the transform's. $R(d,c)$ is the constant of the Riesz
kernel, $(2\pi)^{-d}\int|\xi|^{-c}e^{-i\langle\xi,x\rangle}d\xi = R(d,c)|x|^{c-d}$ as a tempered
distribution: Fourier inversion holds asymptotically as if the transform were the pure power.
This is the asymptotic equivalence that \ref{prop:origin-upper-bound-d} does not claim, and it
contains the two bounds of that node and of \ref{prop:origin-lower-bound-d} for $c < d$ near the
origin, in every dimension $d \ge 1$, without their explicit constants and radii. Clause (2)
settles the boundary: the density grows like the truncated integral of the transform, which is
at least a constant times $\log(1/|x|)$ and is slowly varying in $1/|x|$. \textbf{Checks.} For the
\Matern{} law, $\widehat\mu(\xi) = (1 + |\xi|^2/2)^{-\gamma}$ and $c = 2\gamma < d$, clause (1)
gives $p(x) \sim \Gamma(d/2 - \gamma)\,2^{-\gamma}\pi^{-d/2}\Gamma(\gamma)^{-1}\,|x|^{2\gamma - d}$,
which is the leading term of the closed form with the Bessel function $K_{d/2-\gamma}$. At $d = 1$,
$R(1,c) = 1/(2\Gamma(c)\cos(c\pi/2))$ by the duplication and reflection formulas, and the two
clauses are the symmetric case of \sscite{sato1999levy}~Thm.~53.8(i) and (ii). The referee's
report has both, and numerical checks in dimensions $1$, $2$, $3$ and $5$, one with an unbounded
slowly varying factor. In the Lean the constant is read at the \Matern{} law of the library,
whose transform is $(1 + |\xi|^2)^{-\gamma}$: \texttt{probe\_origin\_density\_rate\_matern} derives
$q(s)\,s^{d-2\gamma} \to R(d,2\gamma)$ from clause (1), and $R(d,2\gamma)$ is the leading
coefficient of that law's closed form (compared by hand; the closed form is not in the library).
\textbf{Edge cases.} $c > d$ is excluded: the density is then bounded and
continuous (\ref{thm:origin-boundedness-d}(2)), and $\Gamma((d-c)/2)$ is negative or undefined.
$c = \infty$ is excluded with it, and so are the rotation-invariant stable laws, whose profile is
unbounded at the origin. $c = 0$ is the trivial datum, whose law is the point mass. A Gaussian
part is excluded: the law then has a bounded smooth density. $d = 1$ is included; nothing in the
proof uses $d \ge 2$. The limits are for every nonincreasing $q$ with $\mu = q(|x|)\,dx$ and at
every point: two such functions differ at most at countably many radii, and Step 5 applies to
each. \textbf{Not claimed}: a rate of convergence in either clause; anything for
$c > d$, where the question is the order at which the density fails to be smooth; anything for a
datum with a Gaussian part. Until Q-0425 this annotation also listed as not claimed an explicit
form of the slowly varying factor $L$, such as the $d$-dimensional analogue of the factor $K$ of
\sscite{sato1999levy}~(53.25). That form is now \ref{prop:isotropic-exponent-profile}, which
gives $L(\rho)$ as asymptotic to a constant times $K(1/\rho)$ with $K$ computed from the polar
profile, and \ref{prop:origin-rate-profile-d} composes it with the two clauses of this node.
\textbf{Prior art.} The argument is Sato and Yamazato's for one-sided
laws on the line: the proof of their Theorem~1.6 passes from the Laplace transform to the
distribution function by Karamata's theorem and from there to the density by monotonicity
(\sscite{sato1978distribution}, p.~298).
Here the Gaussian average of an isotropic law takes the place of the Laplace transform of a
one-sided one. On the line the statement is their Theorem~1.7, p.~277
(\sscite{sato1999levy}~Thm.~53.8). In $\RR^d$, Grzywny, Ryznar and Trojan obtain by the same
method the asymptotics of the transition density $p(t,x)$ of an isotropic unimodal \Levy{}
process as $t \to 0$ with $t\psi(1/|x|) \to \infty$, $\psi$ being the exponent
(\sscite{grzywny2019asymptotic}, Thm.~3.9, p.~26 of the advance-access article), with the
constant $\Gamma(d/2)/(2\pi^{d/2})$, which is the limit of $R(d,tc)/(tc)$ as $tc \to 0$. The two
pages were read again at the image for this node. Clause (1) is in print for a
rotation-invariant probability law on $\RR^d$ whose transform varies regularly with index $-c$,
$0 < c < d$, as a theorem on the spectral measure of a homogeneous isotropic random field and
its covariance function. The mass of a ball is \sscite{leonenko1999limit}, Thm.~1.4.3, p.~65:
the implication from the covariance function to the spectral function, stated there for every
$d \ge 1$ and every $0 < c < d$ with no condition of monotonicity, with the constant
$c_4(d,c) = c/(R(d,c)\,|S^{d-1}|)$. The density, when it is a nonincreasing function of the
radius, is Thm.~1.4.4, pp.~66--67, of the same book, with the constant $c_5(d,c) = R(d,c)$. The
proof printed there (p.~104) starts from the identity that the Laplace--Stieltjes transform of
the mass of balls is the average of the covariance function against the Poisson kernel
(Lemma~1.4.11, p.~72): it is the route of the proof below, with the Poisson kernel in place of
the Gaussian one. For a transform asymptotic to a pure power the implication for the ball is
\sscite{leonenko1991tauberian}, Corollary, p.~1542 of the translation, with its constant given
as an integral. The book's comments (p.~351) give the theorems for a pure power to Leonenko and
Olenko and their proofs with a slowly varying factor to Olenko, whose thesis was not read. Both
theorems are collected in the survey \sscite{leonenko2013tauberian}, as Thms.~3 and~4 in the
numbering of its arXiv version, which is the one cited. The two hypotheses of these theorems are statements of this chapter for a law of
this node: the regular variation of the transform is \ref{lem:isotropic-exponent-log}(2), and
the density is nonincreasing in the radius by \ref{prop:isotropic-radial-monotone}(2). For
$c > (d-3)/2$ the mass of a ball also follows from the Abel--Tauber theorem for Hankel
transforms, in the direction from a function to its transform, applied to
$t^{(d-3)/2}\,\widehat\mu(te)$, whose Hankel transform of order $d/2$ is the mass of a ball up
to an explicit factor. For a bounded function, which here means $d \ge 3$, that direction is
\sscite{soni1975slowlya}, Thm.~1, p.~169, with the remark of p.~171 that the function need only
decrease ultimately, and \sscite{soni1975slowlyc}, Thm.~17(a), p.~613, an equivalence for a
function decreasing on the half-line, with pp.~617--618 for the Hankel kernel. With a condition
of local integrability in place of boundedness, and so for every $d$, it is
\sscite{bingham1997abel}, Thm.~1.2 of the preprint. The boundary case of the same theorem gives
the first limit of clause (2): \sscite{soni1975slowlya}, Thm.~2, p.~169, and
\sscite{soni1975slowlyc}, Thm.~18(a), pp.~613--614, for a bounded function, and
\sscite{bingham1997abel}, Thm.~1.1 with (1.4). What each source states, and under which
hypotheses, is the record \texttt{records/affine/CITATIONS-primary-sources-A.md} with its
reading notes. None of these is used by the proof below.
\textbf{Readings.} In the Lean the ball is
the closed ball, its mass is taken as a real number, ``$s \downarrow 0$'' is the filter of
positive reals at $0$, and each limit is the limit of the displayed quotient of real numbers.
$\widehat\mu(e/s)$ is the real part of the characteristic function at $s^{-1}e$, the
characteristic function being real, and $I(s)$ is the integral of that real part over the closed
ball. A $q$ of the statement is a real function, nonincreasing on $(0,\infty)$, with $\mu$ the
measure of density $x \mapsto \max\{q(|x|),0\}$; its nonnegativity is not asked, which only widens
the class.)}
\end{proposition}

\begin{proof}
\leanok
As $c = \sup_{r>0}k_d(r) > 0$ and $k_d$ is nonincreasing, $k_d$ is positive on an interval
$(0,r_0]$, so the datum is not trivial and \ref{prop:isotropic-radial-monotone}(2) gives a density
$p(x) = q(|x|)$ of $\mu$ with $q$ nonincreasing. From here on $q$ is any nonincreasing function
on $(0,\infty)$ such that $x \mapsto q(|x|)$ is a density of $\mu$.

Write $\chi(\rho) := \widehat\mu(\rho e) = e^{-F_d(\rho e)}$ for $\rho > 0$. As $F_d$ is rotation
invariant, $\widehat\mu(\xi) = \chi(|\xi|)$ for $\xi \ne 0$. Three properties of $\chi$ are used.

(a) $0 < \chi \le 1$, and $\chi$ is continuous, because $F_d \ge 0$ and $\widehat\mu$ is a
characteristic function.

(b) $\chi$ is nonincreasing. In the notation of the proof of
\ref{lem:isotropic-exponent-log}, $F_d(\xi)$ is the average over the sphere of
$\varphi(\langle\xi,\eta\rangle)$, and the substitution $u = ar$ made there gives
$\varphi(as) = \int_0^\infty(1 - \cos(us))\,k_d(u/a)\,du/u \ge \varphi(s)$ for $a \ge 1$ and every
real $s$, because $k_d(u/a) \ge k_d(u)$. So $F_d(a\xi) \ge F_d(\xi)$ for $a \ge 1$.

(c) $\chi(a\rho)/\chi(\rho) \to a^{-c}$ as $\rho \to \infty$, for every $a > 0$. This is
\ref{lem:isotropic-exponent-log}(2), which applies because $c \le d$ is finite.

\emph{Step 1: a bound on the ratios of $\chi$.} Let $\gamma > c$. By (c) at $a = 2$,
$\chi(2y)/\chi(y) \to 2^{-c} > 2^{-\gamma}$, so there is $X \ge 1$ with
$\chi(y) \le 2^{\gamma}\chi(2y)$ for every $y \ge X$. Let $X \le y \le x$, and let $n \ge 0$ be the
least integer with $2^n y \ge x$, so that $n < \log_2(x/y) + 1$. Iterating the inequality and
using (b),
\[
  \chi(y) \le 2^{\gamma n}\,\chi(2^n y) \le 2^{\gamma n}\,\chi(x)
   \le 2^{\gamma}\,(x/y)^{\gamma}\,\chi(x) .
\]
It follows that, with $C := 2^{\gamma}/\chi(X)$,
\begin{equation}\label{eq:rate-ratio-bound}
  \frac{\chi(Ru)}{\chi(R)} \le 1 + C\,u^{-\gamma} \qquad (R \ge X,\ u > 0) .
\end{equation}
Indeed, for $u \ge 1$ the quotient is at most $1$ by (b). For $u < 1$ and $Ru \ge X$ the display
with $y = Ru$ and $x = R$ gives the bound $2^{\gamma}u^{-\gamma}$, and $2^{\gamma} \le C$ because
$\chi(X) \le 1$. For $u < 1$ and $Ru < X$ the numerator is at most $1$, and the display with
$y = X$ and $x = R$ gives $1/\chi(R) \le 2^{\gamma}(R/X)^{\gamma}/\chi(X) \le C\,u^{-\gamma}$.

\emph{Step 2: the Gaussian average of the law.} For $\lambda > 0$ put
$M(\lambda) := \int_{\RR^d}e^{-\lambda|x|^2}\mu(dx)$. At $t = 1/(2\lambda)$ the heat kernel is
$g_t(x) = (\lambda/\pi)^{d/2}e^{-\lambda|x|^2}$, and \ref{lem:positive-transform-origin}(1) gives
$M(\lambda) = (\pi/\lambda)^{d/2}(2\pi)^{-d}\int_{\RR^d}\widehat\mu(\xi)\,e^{-|\xi|^2/4\lambda}\,d\xi$.
In polar coordinates, at $\lambda = R^2$,
\begin{equation}\label{eq:rate-gaussian-average}
  M(R^2) = (4\pi)^{-d/2}\,|S^{d-1}|\,R^{-d}\,T(R),
  \qquad
  T(R) := \int_0^\infty \chi(\rho)\,\rho^{d-1}\,e^{-\rho^2/4R^2}\,d\rho ,
\end{equation}
and the substitution $\rho = Ru$ gives
$T(R) = R^{d}\int_0^\infty \chi(Ru)\,u^{d-1}\,e^{-u^2/4}\,du$.

\emph{Step 3: the Gaussian average for $c < d$.} Let $c < d$ and take $\gamma := (c+d)/2$ in
Step 1, so that $c < \gamma < d$. By the second form of $T(R)$ in Step 2,
\[
  \frac{T(R)}{R^{d}\chi(R)} = \int_0^\infty \frac{\chi(Ru)}{\chi(R)}\,u^{d-1}\,e^{-u^2/4}\,du .
\]
For each $u > 0$ the quotient under the integral tends to $u^{-c}$ as $R \to \infty$, by (c). For
$R \ge X$ the integrand is at most $(1 + C\,u^{-\gamma})\,u^{d-1}e^{-u^2/4}$ by
\ref{eq:rate-ratio-bound}, which is integrable on $(0,\infty)$ because $\gamma < d$. By dominated
convergence
\[
  \frac{T(R)}{R^{d}\chi(R)} \to \int_0^\infty u^{d-1-c}\,e^{-u^2/4}\,du
   = 2^{d-c-1}\,\Gamma\bigl(\tfrac{d-c}{2}\bigr) \qquad (R \to \infty),
\]
and with $|S^{d-1}| = 2\pi^{d/2}/\Gamma(d/2)$,
\begin{equation}\label{eq:rate-abelian}
  \frac{M(R^2)}{\chi(R)} \to I_c := (4\pi)^{-d/2}\,|S^{d-1}|\,2^{d-c-1}\,\Gamma\bigl(\tfrac{d-c}{2}\bigr)
   = \frac{2^{-c}\,\Gamma\bigl(\tfrac{d-c}{2}\bigr)}{\Gamma\bigl(\tfrac d2\bigr)} .
\end{equation}

\emph{Step 3$'$: the Gaussian average for $c = d$.} Let $c = d$, and put
$P(y) := \int_0^y\chi(\rho)\,\rho^{d-1}\,d\rho$ for $y > 0$ and $L(\rho) := \rho^{d}\chi(\rho)$.
By (b), for $0 < x \le y$,
\begin{equation}\label{eq:rate-P-increments}
  \chi(y)\,\frac{y^{d} - x^{d}}{d} \le P(y) - P(x) \le \chi(x)\,\frac{y^{d} - x^{d}}{d} .
\end{equation}
\emph{(i) $L(R)/P(R) \to 0$ as $R \to \infty$.} The left inequality of
\ref{eq:rate-P-increments} at $x = y/2$ reads $P(y) - P(y/2) \ge \kappa\,L(y)$ with
$\kappa := (1 - 2^{-d})/d$. For each integer $j \ge 0$,
$L(R/2^j)/L(R) = 2^{-jd}\chi(2^{-j}R)/\chi(R) \to 1$ as $R \to \infty$, by (c) with $c = d$. Let
$n \ge 1$. For all large $R$, $L(R/2^j) \ge L(R)/2$ for every $j < n$, and then
\[
  P(R) \ge \sum_{j=0}^{n-1}\bigl(P(R/2^j) - P(R/2^{j+1})\bigr) \ge \kappa\sum_{j=0}^{n-1}L(R/2^j)
   \ge \tfrac{n\kappa}{2}\,L(R) .
\]
As $n$ is arbitrary, $L(R)/P(R) \to 0$.
\emph{(ii) $P(aR)/P(R) \to 1$ as $R \to \infty$, for every $a > 0$.} For $a \ge 1$ the right
inequality of \ref{eq:rate-P-increments} gives
$0 \le P(aR) - P(R) \le (a^{d} - 1)\,L(R)/d$, and (i) gives the limit. For $a < 1$ the quotient is
the reciprocal of $P(a^{-1}R')/P(R')$ at $R' = aR$, which tends to infinity with $R$.
\emph{(iii) $T(R)/P(R) \to 1$ as $R \to \infty$.} Let $\varepsilon > 0$. On $(0, \varepsilon R]$
the Gaussian factor of $T(R)$ is at least $e^{-\varepsilon^2/4}$, so
$T(R) \ge e^{-\varepsilon^2/4}P(\varepsilon R)$. On $(0, R]$ it is at most $1$, and on $(R,\infty)$,
$\chi(\rho) \le \chi(R)$ by (b), so
\[
  T(R) \le P(R) + \chi(R)\int_0^\infty\rho^{d-1}e^{-\rho^2/4R^2}\,d\rho
   = P(R) + 2^{d-1}\Gamma\bigl(\tfrac d2\bigr)\,L(R) .
\]
Hence
$e^{-\varepsilon^2/4}\,P(\varepsilon R)/P(R) \le T(R)/P(R) \le 1 + 2^{d-1}\Gamma(d/2)\,L(R)/P(R)$.
By (ii) the left side tends to $e^{-\varepsilon^2/4}$ and by (i) the right side tends to $1$. As
$\varepsilon$ is arbitrary, $T(R)/P(R) \to 1$. With \ref{eq:rate-gaussian-average},
\begin{equation}\label{eq:rate-abelian-critical}
  \frac{M(R^2)}{(4\pi)^{-d/2}\,|S^{d-1}|\,R^{-d}\,P(R)} \to 1 \qquad (R \to \infty) .
\end{equation}

\emph{Step 4: the mass of a ball, by Karamata's theorem.} Let $\nu$ be the image of $\mu$ under
$x \mapsto |x|^2$. It is a probability measure on $[0,\infty)$ with
$\int_{[0,\infty)}e^{-\lambda r}\,\nu(dr) = M(\lambda)$ and $\nu([0,r]) = \mu(\{|x| \le \sqrt r\})$.

Let $c < d$ and put $\ell(\lambda) := \lambda^{c/2}\chi(\sqrt\lambda)$ for $\lambda > 0$. It is
positive and continuous, and
$\ell(a\lambda)/\ell(\lambda) = a^{c/2}\chi(\sqrt a\sqrt\lambda)/\chi(\sqrt\lambda) \to 1$ as
$\lambda \to \infty$ for every $a > 0$, by (c). By \ref{eq:rate-abelian} at $R = \sqrt\lambda$,
$M(\lambda)/(I_c\,\lambda^{-c/2}\,\ell(\lambda)) \to 1$. Karamata's Tauberian theorem
(\ref{thm:cited-A11}, with the index $c/2 > 0$ and the constant $I_c > 0$) gives
$\nu([0,r]) \sim I_c\,r^{c/2}\,\ell(1/r)/\Gamma(1 + c/2)$ as $r \downarrow 0$, and
$r^{c/2}\ell(1/r) = \chi(r^{-1/2})$. At $r = s^2$,
\begin{equation}\label{eq:rate-ball}
  \frac{\mu(\{|x| \le s\})}{B_c\,\chi(1/s)} \to 1 \quad (s \downarrow 0),
  \qquad B_c := \frac{I_c}{\Gamma(1 + c/2)} = \frac{R(d,c)\,|S^{d-1}|}{c} ,
\end{equation}
the last equality by $\Gamma(1 + c/2) = (c/2)\Gamma(c/2)$ and $|S^{d-1}| = 2\pi^{d/2}/\Gamma(d/2)$.
This is the first limit of clause (1).

Let $c = d$ and put $\ell(\lambda) := P(\sqrt\lambda)$. It is positive and continuous, and
$\ell(a\lambda)/\ell(\lambda) \to 1$ by Step 3$'$(ii). By \ref{eq:rate-abelian-critical},
$M(\lambda)/\bigl((4\pi)^{-d/2}|S^{d-1}|\,\lambda^{-d/2}\,\ell(\lambda)\bigr) \to 1$. Karamata's
theorem with the index $d/2$ gives, at $r = s^2$,
$\mu(\{|x| \le s\}) \sim (4\pi)^{-d/2}\,|S^{d-1}|\,s^{d}\,P(1/s)/\Gamma(1 + d/2)$. In polar
coordinates $|S^{d-1}|\,P(1/s) = I(s)$, and
$(4\pi)^{-d/2}/\Gamma(1 + d/2) = (2\pi)^{-d}\,\pi^{d/2}/\Gamma(1 + d/2) = (2\pi)^{-d}\,V_d$. This is
the first limit of clause (2).

\emph{Step 5: from the mass of balls to the density.} Write $m(s) := \mu(\{|x| \le s\})$. Let
$W$ be a positive function on $(0,\infty)$, $\varkappa > 0$ and $B > 0$, and suppose that
$W(as)/W(s) \to a^{\varkappa}$ for every $a > 0$ and $m(s)/W(s) \to B$, as $s \downarrow 0$. Then
\begin{equation}\label{eq:rate-monotone-density}
  \frac{q(s)\,s^{d}}{W(s)} \to \frac{B\,\varkappa}{|S^{d-1}|} \qquad (s \downarrow 0) .
\end{equation}
Indeed, let $0 < a < b$. On the shell $\{as < |x| \le bs\}$, whose volume is
$V_d\,s^{d}(b^{d} - a^{d})$, the density $q(|x|)$ lies between $q(bs)$ and $q(as)$, because $q$ is
nonincreasing. So
\[
  q(bs)\,V_d\,s^{d}\,(b^{d} - a^{d}) \le m(bs) - m(as) \le q(as)\,V_d\,s^{d}\,(b^{d} - a^{d}) ,
\]
and $(m(bs) - m(as))/W(s) \to B\,(b^{\varkappa} - a^{\varkappa})$. With $b = 1$ and $a < 1$ the left
inequality shows that, for every $\varepsilon > 0$, $q(s)s^{d}/W(s) \le
B(1 - a^{\varkappa})/(V_d(1 - a^{d})) + \varepsilon$ for all small $s$. With $a = 1$ and $b > 1$ the
right one shows $q(s)s^{d}/W(s) \ge B(b^{\varkappa} - 1)/(V_d(b^{d} - 1)) - \varepsilon$ for all small
$s$. Now $(t^{\varkappa} - 1)/(t^{d} - 1) \to \varkappa/d$ as $t \to 1$, the quotient of the
derivatives of $t^{\varkappa}$ and $t^{d}$ at $1$. So the first bound tends to
$B\varkappa/(V_d\,d) + \varepsilon$ as $a \uparrow 1$ and the second to
$B\varkappa/(V_d\,d) - \varepsilon$ as $b \downarrow 1$. As $\varepsilon$ is arbitrary and
$V_d\,d = |S^{d-1}|$, this proves \ref{eq:rate-monotone-density}. It is the argument of the
monotone density theorem (\sscite{bingham1987regular}, Thm.~1.7.2b, pp.~39--40), on spherical
shells.

For $c < d$ take $W(s) := \chi(1/s)$, $\varkappa := c$ and $B := B_c$: the first hypothesis is (c)
at the ratio $1/a$, the second is \ref{eq:rate-ball}. So
$q(s)\,s^{d}/\chi(1/s) \to B_c\,c/|S^{d-1}| = R(d,c)$, the second limit of clause (1). For $c = d$
take $W(s) := s^{d}I(s)$, $\varkappa := d$ and $B := (2\pi)^{-d}V_d$: the first hypothesis is
Step 3$'$(ii), since $I(s) = |S^{d-1}|P(1/s)$, and the second is the first limit of clause (2).
So $q(s)/I(s) \to (2\pi)^{-d}\,V_d\,d/|S^{d-1}| = (2\pi)^{-d}$, the second limit of clause (2).

\emph{The logarithmic bound of clause (2).} Let $c = d$ and $0 < s \le 1$. By
\ref{lem:isotropic-exponent-log}(1), $\chi(\rho) \ge e^{-13d/4 - 2C_k}\rho^{-d}$ for $\rho \ge 1$.
In polar coordinates
\[
  I(s) \ge |S^{d-1}|\int_1^{1/s}\chi(\rho)\,\rho^{d-1}\,d\rho
   \ge |S^{d-1}|\,e^{-13d/4 - 2C_k}\int_1^{1/s}\frac{d\rho}{\rho}
   = |S^{d-1}|\,e^{-13d/4 - 2C_k}\log(1/s) .
\]
Let $C < C_* := (2\pi)^{-d}|S^{d-1}|e^{-13d/4-2C_k}$. If $C \le 0$ there is nothing to prove, as
$q(s) > 0$ for small $s$ by the second limit. Otherwise $q(s) \ge (C/C_*)(2\pi)^{-d}I(s)$ for all
small $s$, by the same limit, and the bound on $I(s)$ gives $q(s) \ge C\log(1/s)$.

\emph{What is machine-checked.} Everything above except Karamata's theorem, which is the cited
fact \ledger{A11} and enters in Step 4 only. The existence of $q$ is
\ref{prop:isotropic-radial-monotone}(2), on \ledger{A8}, \ledger{A2} and \ledger{A3}(iii).
\end{proof}

% not in the draft: queue item Q-0425 (hub), batch R3 of the response to the external review of
% the isotropic paper, round 1 (records/affine/PLAN-review-response-paper-iso.md, D3 option (2)).
% The slowly varying factor of prop:origin-rate-d through the polar profile, in three nodes, so
% that each is true on its own and the Lean of each can cover all of it (ADR-0025):
%   prop:isotropic-exponent-profile   the asymptotic of the exponent, its constant written as a
%                                     limit and a sphere integral; no cited fact; holds for every
%                                     finite c > 0;
%   lem:exponent-profile-constant     the two closed forms of the constant (Euler's constant, the
%                                     digamma mean); under D3 option (2) computations of the prose;
%   prop:origin-rate-profile-d        the composition with prop:origin-rate-d, 0 < c <= d, which
%                                     inherits that node's cited facts.
% The derivation is the verification record's (records/affine/reviews/
% VERIFY-round1-open-problems.md, part B), checked line by line; where the text below departs from
% it is recorded in records/affine/NODES-q0425.md. No \lean tag on any of the three: nothing is
% typed yet, and the statements are refereed blind first (the item's brief).
% CHANGED Q-0425 (the blind check; records/formalization/blind-check-q0425-report-A.md, verdict
% "correct as stated" for the three nodes, repairs R1 to R6): in this node, the statement of
% clause (1) says where K(ax)/K(x) is defined (R1) and gives the limit of K at 0 as two cases, with
% K(0+) defined in each (R4); clause (3) states the two-sided bound for every rho > 0, with the
% tail integral of the profile, which the proof gave already (R6); the proof no longer says that
% the layer-cake identity is the one place where monotonicity is used, which was false (R2), says
% why the inner integral of the second term is a Borel function of eta (R3), and closes with
% "of the ledger" (R5); the annotation says that c > 0 is not used by the asymptotic. In
% lem:exponent-profile-constant: the closing sentence of the proof and the first sentence of the
% annotation name the two facts used as known (R5). In prop:origin-rate-profile-d: the limits of
% clauses (2) and (3) with values in (0,infty] are two cases each (R4), in statement and proof.
% No clause of the three changed in content, and none is printed by a paper yet.
% CHANGED Q-0427: typed and proved (Formalization/AffineHemigroup/SlowFactor.lean,
% ProfileLayerCake.lean, SphereLogMean.lean, ExponentProfile.lean). The statement is unchanged.
% Added: the \lean and \leanok tags, the passage "Formalization" of the annotation, and the closing
% sentence of the proof. lem:exponent-profile-constant keeps no \lean tag (D3 option (2)).
\begin{proposition}[the exponent at infinity through the polar profile]
\label{prop:isotropic-exponent-profile}
\uses{def:isotropic-origin-datum,lem:polar-decomposition,lem:isotropic-exponent-log}
\lean{AffineHemigroup.slowFactor_basic, AffineHemigroup.slowFactor_ratio,
AffineHemigroup.slowFactor_tendsto_of_integrable, AffineHemigroup.slowFactor_tendsto_atTop,
AffineHemigroup.cin_sub_log_limit, AffineHemigroup.sphere_log_mean,
AffineHemigroup.sphereLogMean_one, AffineHemigroup.exponent_profile_asymptotic,
AffineHemigroup.exponent_profile_bounds, AffineHemigroup.transform_power_comparison,
AffineHemigroup.transform_power_bounded_iff, AffineHemigroup.transform_power_limit,
AffineHemigroup.OriginDatum.slowFactor, AffineHemigroup.OriginDatum.slowFactorZero,
AffineHemigroup.sphereLogMean, AffineHemigroup.exponentConst}\leanok
Let $d \ge 1$, let $(0,k_d)$ be an isotropic origin datum in dimension $d$ with Gaussian
coefficient $0$ and $0 < c := k_d(0{+}) < \infty$, let $F_d$ be its exponent, put
$C_k := \int_1^\infty k_d(r)\,dr/r$, and let $e$ be a unit vector of $\RR^d$. Write
$\Cin(X) = \int_0^X (1 - \cos u)\,du/u$ and put
\[
  K(x) := \exp\Bigl[\int_x^1 \bigl(c - k_d(u)\bigr)\,\frac{du}{u}\Bigr] \qquad (0 < x \le 1) .
\]
\begin{enumerate}
\item \emph{(The factor.)} $K$ is continuous and nonincreasing on $(0,1]$, with
$1 = K(1) \le K(x) \le x^{-c}$. For every $a > 0$, $K(ax)/K(x) \to 1$ as $x \downarrow 0$; the
quotient is defined for $0 < x \le \min\{1, 1/a\}$. If $\int_0^1 (c - k_d(u))\,du/u < \infty$,
then $K(x)$ increases, as $x \downarrow 0$, to the finite limit
\[
  K(0{+}) := \exp\int_0^1 \bigl(c - k_d(u)\bigr)\,\frac{du}{u} \;\ge\; 1 .
\]
If the integral diverges, then $K(x) \to \infty$ as $x \downarrow 0$, and $K(0{+}) := \infty$.
\item \emph{(Two constants.)} The limit
\[
  g_\infty := \lim_{X\to\infty}\bigl(\Cin(X) - \log X\bigr)
\]
exists, and $-3 \le g_\infty \le \tfrac{13}{4}$. The function
$\eta \mapsto \log|\langle e,\eta\rangle|$ is integrable on $S^{d-1}$ for $\sigma_{d-1}$, and
\[
  \lambda_d := \frac{1}{|S^{d-1}|}\int_{S^{d-1}}\log|\langle e,\eta\rangle|\,\sigma_{d-1}(d\eta)
  \in (-\infty, 0]
\]
does not depend on $e$, with $\lambda_1 = 0$. Put $\varkappa_d := g_\infty + \lambda_d$.
\item \emph{(The exponent.)} As $\rho \to \infty$,
\[
  F_d(\rho e) - \bigl(c\log\rho - \log K(1/\rho) + C_k\bigr) \to c\,\varkappa_d ,
  \qquad\text{that is}\qquad
  \frac{e^{-F_d(\rho e)}}{e^{-C_k - c\varkappa_d}\,\rho^{-c}\,K(1/\rho)} \to 1 .
\]
Moreover, for every $\rho \ge 1$,
\[
  c\,(\lambda_d - 3) \;\le\; F_d(\rho e) - \bigl(c\log\rho - \log K(1/\rho) + C_k\bigr)
  \;\le\; \tfrac{13}{4}\,c .
\]
More generally, for every $\rho > 0$,
\[
  c\,(\lambda_d - 3) \;\le\; F_d(\rho e) - \int_{1/\rho}^{\infty}k_d(r)\,\frac{dr}{r}
  \;\le\; \tfrac{13}{4}\,c ,
\]
and for $\rho \ge 1$ the integral equals $c\log\rho - \log K(1/\rho) + C_k$.
\item \emph{(When the transform is comparable to a power.)} For every $\xi$ with $|\xi| \ge 1$,
\[
  e^{-13c/4 - C_k}\,K(1/|\xi|) \;\le\; |\xi|^{c}\,e^{-F_d(\xi)} \;\le\;
  e^{c(3-\lambda_d) - C_k}\,K(1/|\xi|) .
\]
So $|\xi|^{c}e^{-F_d(\xi)}$ is bounded on $\{|\xi| \ge 1\}$ if and only if
$\int_0^1 (c - k_d(u))\,du/u < \infty$, and then it tends to
$e^{-C_k - c\varkappa_d}\,K(0{+})$ as $|\xi| \to \infty$.
\end{enumerate}
\statusT\quad\emph{(Cites nothing: the bounds on $\Cin$ in the proof of
\ref{lem:isotropic-exponent-log}, Tonelli's theorem in the layer-cake form, dominated convergence
and polar coordinates in $\RR^d$. The statement is about the exponent alone, and no law is named.
\textbf{What it is worth.} \ref{lem:isotropic-exponent-log} writes
$e^{-F_d(\rho e)} = \rho^{-c}L(\rho)$ with $L$ slowly varying and bounded below. Clause (3) gives
$L$: it is asymptotic to the constant $e^{-C_k - c\varkappa_d}$ times $K(1/\rho)$, and $K$ is
computed from the polar profile near the origin by one integral. Clause (4) decides when $L$ is
bounded, that is when $e^{-F_d(\xi)}$ is comparable to $|\xi|^{-c}$: exactly when $c - k_d(u)$ is
integrable against $du/u$ at $0$. A profile equal to $c$ near $0$, or approaching $c$ like a power
of $u$, gives a bounded $K$. The profile $c\,(1 - 1/\log(e/u))$ on $(0,1)$ gives
$K(x) = (1 + \log(1/x))^{c}$, which is unbounded. \textbf{Hypotheses.} $0 < c < \infty$ and nothing
else. No relation between $c$ and $d$ is needed, so the proposition holds above the threshold
$c > d$ of \ref{thm:origin-boundedness-d} too. The positivity of $c$ is not used by the
asymptotic of the exponent either: at $c = 0$ the profile vanishes, $F_d = 0$, $K = 1$,
$C_k = 0$, and every clause holds trivially. The hypothesis $c > 0$ is kept, because the node is
read for a positive profile at the origin and \ref{prop:origin-rate-profile-d} needs it. The
monotonicity of $k_d$, which the definition of a datum supplies, enters the exponent through the
layer-cake identity of the proof. It is also used for $k_d \le c$ and $k_d(r) \to 0$, in clause
(1), and in the first term of the sphere average. No rate of approach of $k_d$ to $c$ and no
regularity are assumed. The two-sided bound of clause (3) holds for every $\rho > 0$ when it is
written with the tail integral of the profile; $\rho \ge 1$ is needed only to write that integral
through $K$. \textbf{What it contains.}
\ref{lem:isotropic-exponent-log}(2) follows from clauses (1) and (3), and the upper bound of
clause (3) is smaller than that of \ref{lem:isotropic-exponent-log}(1) by $C_k + \log K(1/\rho)$.
Neither earlier statement is changed; \ref{prop:origin-rate-d} and
\ref{thm:origin-boundedness-d} read them as they stand. \textbf{The constants.} $g_\infty$ is
defined as a limit and $\lambda_d$ as an integral, so that the statement names no special
function. Their values, $g_\infty = \gamma_{\mathrm E}$ and
$\lambda_d = \tfrac12[\psi(\tfrac12) - \psi(\tfrac d2)]$, are
\ref{lem:exponent-profile-constant}. \textbf{Where it comes from, and attribution.} The form in
$\RR^d$, with the mean of $\log|\langle e,\eta\rangle|$ over the sphere in the constant, was
stated by the external referee of the isotropic paper (round 1) and derived in full in
\texttt{records/affine/reviews/VERIFY-round1-open-problems.md}, part~B, from which the proof below
is taken. On the line, $d = 1$, the statement is Sato and Yamazato's, in the symmetric case: the
profile $k_1$ is the folded profile, $K$ is the factor of \sscite{sato1978distribution}~(1.9),
p.~277, and of \sscite{sato1999levy}~(53.25), p.~411, and clause (3) is
\sscite{sato1978distribution}, Lemma~5.1, (5.9), p.~298, restated as (53.33) on p.~412 of
\sscite{sato1999levy}, in the proof of its Thm.~53.8. The anchors were read at the page image for
that record (its B.5) and for the attribution item Q-0426 (its record
\texttt{records/affine/PRIOR-ART-review-round1-additions.md}, section~3). In $\RR^d$ the
factor is theirs, read on the polar profile; the statement is proved here and nothing is
attributed for it. \textbf{Not claimed}: a rate of
convergence in clause (3), which depends on how fast $k_d$ approaches $c$ (for the logarithmic
profile above the record finds, numerically, an error of order $1/\log\rho$); anything for
$c = \infty$, where $K$ is not defined, or for a datum with a Gaussian part; that the constants
$-3$ and $\tfrac{13}{4}$ are sharp, they are those of \ref{eq:cin-log} (numerically the function
$G$ of the proof lies in $[0, 0.776]$, so they are not:
\texttt{records/formalization/blind-check-q0425-A/t1\_constants.py}); the converse, that an
exponent with this asymptotic comes from a datum. \textbf{Formalization} (Q-0427). $K$ is the
total function \texttt{OriginDatum.slowFactor}, the exponential of an oriented interval integral.
It is the function of the statement for $0 < x \le 1$, and no clause reads it elsewhere: the
limit of $K(ax)/K(x)$ is along $x \downarrow 0$, where both arguments are eventually in $(0,1]$.
The condition $\int_0^1 (c - k_d(u))\,du/u < \infty$ is the integrability of the integrand on
$(0,1]$, the integrand being nonnegative and measurable. $K(0{+})$ in the convergent case is
\texttt{OriginDatum.slowFactorZero}, and the divergent case is a limit to $+\infty$. $g_\infty$
is the constant \texttt{ScaleSpace.cinConst} of the shared library, with the limit as a theorem,
which determines it. $\lambda_d$ is \texttt{sphereLogMean}, written at the first basis vector,
and its equality with the mean for every unit vector is a theorem. The surface measure is
Mathlib's, of total mass $|S^{d-1}|$, which is $2$ at $d = 1$. The limits in $\rho$ are along the
real numbers, and the limit of clause (4) is along $|\xi| \to \infty$. $d$ is a natural number,
$d \ge 1$. The hypothesis $c > 0$ is not assumed by the Lean statements, which need only
$c < \infty$.)}
\end{proposition}

\begin{proof}
\leanok
Write $k := k_d$ and $\bar\sigma := \sigma_{d-1}/|S^{d-1}|$ for the normalized surface measure.
Three things are taken from the proof of \ref{lem:isotropic-exponent-log}. First,
\[
  F_d(\xi) = \int_{S^{d-1}}\varphi\bigl(\langle\xi,\eta\rangle\bigr)\,\bar\sigma(d\eta),
  \qquad
  \varphi(s) := \int_0^\infty\bigl(1 - \cos(rs)\bigr)\,k(r)\,\frac{dr}{r} \quad (s \in \RR),
\]
with $\varphi$ even and finite. Second, $k \le c$ on $(0,\infty)$ and $k(r) \to 0$ as
$r \to \infty$. Third, the bounds \ref{eq:cin-log}, and $0 \le \Cin(X) \le \tfrac14$ for
$0 \le X \le 1$, because $\Cin$ is nondecreasing and $\Cin(1) \le \tfrac14$.

\emph{Clause (1).} For $0 < u \le 1$ the integrand $(c - k(u))/u$ lies in $[0, c/u]$. So
$\int_x^1 (c - k(u))\,du/u$ lies in $[0, c\log(1/x)]$, it is continuous and nonincreasing in $x$,
and it vanishes at $x = 1$. Hence $K$ is continuous and nonincreasing with
$1 = K(1) \le K(x) \le x^{-c}$. Let $a > 0$, and let $0 < x \le \min\{1, 1/a\}$, so that $x$ and
$ax$ lie in $(0,1]$. Between $x$ and $ax$ the integrand is at most $(c - k(\max\{a,1\}x))/u$,
because $k$ is nonincreasing. So
\[
  \bigl|\log K(ax) - \log K(x)\bigr| = \Bigl|\int_{ax}^{x}\bigl(c - k(u)\bigr)\frac{du}{u}\Bigr|
   \le \bigl(c - k(\max\{a,1\}\,x)\bigr)\,|\log a| ,
\]
which tends to $0$ as $x \downarrow 0$ because $k(0{+}) = c$. By monotone convergence
$\int_x^1 (c - k(u))\,du/u$ increases to $\int_0^1 (c - k(u))\,du/u$ as $x \downarrow 0$. If this
integral is finite, $K(x)$ increases to its exponential, which is $K(0{+})$. If it is infinite,
$K(x) \to \infty$.

\emph{Clause (2), the limit $g_\infty$.} Put $\log^+ X := \max\{\log X, 0\}$ for $X > 0$,
$\log^+ 0 := 0$, and
\[
  G(X) := \Cin(X) - \log^+ X \qquad (X \ge 0) .
\]
Then $-3 \le G \le \tfrac{13}{4}$ on $[0,\infty)$: for $X \ge 1$ this is \ref{eq:cin-log}, and for
$0 \le X < 1$, $G(X) = \Cin(X) \in [0,\tfrac14]$. For $X \ge 1$,
$\Cin(X) = \Cin(1) + \log X - \int_1^X u^{-1}\cos u\,du$, so
\[
  G(X) = \Cin(1) - \int_1^X \frac{\cos u}{u}\,du
       = \Cin(1) - \Bigl[\frac{\sin u}{u}\Bigr]_1^X - \int_1^X \frac{\sin u}{u^2}\,du .
\]
The bracket tends to $-\sin 1$ as $X \to \infty$, and $u^{-2}\sin u$ is absolutely integrable on
$[1,\infty)$. So $G(X)$ has a limit as $X \to \infty$. As $G(X) = \Cin(X) - \log X$ for $X \ge 1$,
this limit is $g_\infty$, and $-3 \le g_\infty \le \tfrac{13}{4}$.

\emph{Clause (2), the mean of the logarithm over the sphere.} Let $B$ be the closed unit ball of
$\RR^d$ and $e_1$ the first basis vector. Lebesgue measure is the image of
$\sigma_{d-1} \otimes r^{d-1}dr$ under $(\eta, r) \mapsto r\eta$, so for every Borel
$h : S^{d-1} \to [0,\infty]$,
\[
  \int_B h\bigl(x/|x|\bigr)\,dx = \int_0^1 r^{d-1}\,dr\int_{S^{d-1}}h\,d\sigma_{d-1}
   = \frac1d\int_{S^{d-1}}h\,d\sigma_{d-1} .
\]
Take $h(\eta) := -\log|\langle e_1,\eta\rangle| \in [0,\infty]$. For $x \in B$ with $x_1 \ne 0$,
$h(x/|x|) = -\log|x_1| + \log|x| \le -\log|x_1|$, because $|x| \le 1$. As $B$ lies in the cube
$[-1,1]^d$, on which $-\log|x_1| \ge 0$,
\[
  \frac1d\int_{S^{d-1}}h\,d\sigma_{d-1} \le \int_{[-1,1]^d}\bigl(-\log|x_1|\bigr)\,dx
   = 2^{d-1}\int_{-1}^{1}\bigl(-\log|t|\bigr)\,dt = 2^{d} .
\]
So $\log|\langle e_1,\eta\rangle|$ is $\sigma_{d-1}$-integrable, and its mean is at most $0$
because $|\langle e_1,\eta\rangle| \le 1$. For a unit vector $e$ take $Q \in \Orth(d)$ with
$Qe_1 = e$. Then $\langle e,\eta\rangle = \langle e_1, Q^{\tp}\eta\rangle$, and $\sigma_{d-1}$ is
invariant under $Q^{\tp}$, so the integral is the same for $e$ as for $e_1$. At $d = 1$ the
sphere is $\{\pm1\}$, $|\langle e,\eta\rangle| = 1$, and $\lambda_1 = 0$.

\emph{Clause (3), the tail integral.} For $s > 0$ put
\[
  T(s) := \int_{1/s}^{\infty}k(r)\,\frac{dr}{r} .
\]
It is finite: it is at most $C_k$ for $s \le 1$ and at most $c\log s + C_k$ for $s \ge 1$. For
$s \ge 1$,
\begin{equation}\label{eq:profile-tail}
  T(s) = \int_{1/s}^{1}k(r)\,\frac{dr}{r} + C_k
       = c\log s - \int_{1/s}^{1}\bigl(c - k(r)\bigr)\frac{dr}{r} + C_k
       = c\log s - \log K(1/s) + C_k .
\end{equation}

\emph{Clause (3), the layer-cake identity.} For $0 < y < c$ put
$\tau(y) := \sup\{r > 0 : k(r) > y\}$. The set is not empty, because $\sup k = c > y$. It is
bounded, because $k(r) \to 0$ as $r \to \infty$ and $y > 0$. So $0 < \tau(y) < \infty$. As $k$ is
nonincreasing, the set contains every $r' < r$ with each of its points $r$, so it is
$(0,\tau(y))$ or $(0,\tau(y)]$; and $\tau$ is nonincreasing in $y$, hence Borel. Since
$0 \le k \le c$,
\[
  k(r) = \int_0^c \mathbf 1\{y < k(r)\}\,dy \qquad (r > 0),
\]
and the indicator is a Borel function of $(r,y)$, $k$ being monotone. For a Borel
$w : (0,\infty) \to [0,\infty)$, Tonelli's theorem gives
\[
  \int_0^\infty w(r)\,k(r)\,dr = \int_0^c\Bigl(\int_0^{\tau(y)}w(r)\,dr\Bigr)dy .
\]
Let $s > 0$. With $w(r) = (1 - \cos rs)/r$ the inner integral is $\Cin(s\,\tau(y))$, by the
substitution $u = rs$. With $w(r) = \mathbf 1\{r > 1/s\}/r$ it is $\log^+(s\,\tau(y))$. So
\[
  \varphi(s) = \int_0^c \Cin\bigl(s\,\tau(y)\bigr)\,dy, \qquad
  T(s) = \int_0^c \log^+\bigl(s\,\tau(y)\bigr)\,dy .
\]
Both are finite, and their integrands differ by $G(s\,\tau(y))$, which is bounded. Hence
\begin{equation}\label{eq:profile-layer-cake}
  \varphi(s) = T(s) + \int_0^c G\bigl(s\,\tau(y)\bigr)\,dy \qquad (s > 0) .
\end{equation}
This is an identity, with no remainder. It is where the monotonicity of $k$ enters the exponent.
Monotonicity is also used for $k \le c$ and $k(r) \to 0$, in clause (1), and in the first term
below.

\emph{Clause (3), the sphere average.} Let $\rho > 0$ and write
$\eta_1 := \langle e,\eta\rangle$ for $\eta \in S^{d-1}$. As $\varphi$ is even,
$F_d(\rho e) = \int\varphi(\rho|\eta_1|)\,\bar\sigma(d\eta)$. The set $\{\eta_1 = 0\}$ is
$\bar\sigma$-null: it is a great subsphere when $d \ge 2$ and empty when $d = 1$. By
\ref{eq:profile-layer-cake} at $s = \rho|\eta_1|$,
\begin{equation}\label{eq:profile-sphere}
  F_d(\rho e) - T(\rho) = \int_{S^{d-1}}\bigl[T(\rho|\eta_1|) - T(\rho)\bigr]\,\bar\sigma(d\eta)
   + \int_{S^{d-1}}\int_0^c G\bigl(\rho|\eta_1|\,\tau(y)\bigr)\,dy\,\bar\sigma(d\eta) ,
\end{equation}
where both integrals converge absolutely by the bounds that follow.

\emph{The first term.} Let $0 < t \le 1$. Then
$T(\rho) - T(\rho t) = \int_{1/\rho}^{1/(\rho t)}k(r)\,dr/r$ lies between
$k(1/(\rho t))\log(1/t)$ and $c\log(1/t)$. So $-c\log(1/t) \le T(\rho t) - T(\rho) \le 0$ for
every $\rho > 0$, and $T(\rho t) - T(\rho) \to c\log t$ as $\rho \to \infty$, because
$k(0{+}) = c$. At $t = |\eta_1|$ the bound $c\,|\log|\eta_1||$ is $\bar\sigma$-integrable by
clause (2). By dominated convergence the first term of \ref{eq:profile-sphere} tends to
$c\lambda_d$ as $\rho \to \infty$, and it lies in $[c\lambda_d, 0]$ for every $\rho > 0$.

\emph{The second term.} The integrand lies in $[-3, \tfrac{13}{4}]$. For every $y \in (0,c)$ and
every $\eta$ with $\eta_1 \ne 0$, $\rho|\eta_1|\tau(y) \to \infty$ as $\rho \to \infty$, because
$\tau(y) > 0$, so the integrand tends to $g_\infty$. The function
$(\eta,y) \mapsto G(\rho|\eta_1|\tau(y))$ is Borel on $S^{d-1} \times (0,c)$ and bounded: $\tau$
is Borel, being monotone, and $G$ is continuous on $[0,\infty)$. So the inner integral
$\int_0^c G(\rho|\eta_1|\tau(y))\,dy$ is a Borel function of $\eta$, and it lies in
$[-3c, \tfrac{13}{4}c]$ for every $\rho > 0$ and every $\eta$. For $\eta$ with $\eta_1 \ne 0$ it
tends to $c\,g_\infty$ as $\rho \to \infty$, by bounded convergence on $(0,c)$. By bounded
convergence for $\bar\sigma$, the second term tends to $c\,g_\infty$, and it lies in
$[-3c, \tfrac{13}{4}c]$ for every $\rho > 0$.

Adding the two, $F_d(\rho e) - T(\rho) \to c(\lambda_d + g_\infty) = c\varkappa_d$ as
$\rho \to \infty$, and $c(\lambda_d - 3) \le F_d(\rho e) - T(\rho) \le \tfrac{13}{4}c$ for every
$\rho > 0$. As $T(\rho) = \int_{1/\rho}^\infty k(r)\,dr/r$, this is the bound of clause (3) for
every $\rho > 0$. For $\rho \ge 1$, \ref{eq:profile-tail} replaces $T(\rho)$ by
$c\log\rho - \log K(1/\rho) + C_k$. This is the limit and the two-sided bound of clause (3) for
$\rho \ge 1$. The
quotient form follows by taking exponentials: the quotient is $\exp$ of minus the difference
between $F_d(\rho e) - (c\log\rho - \log K(1/\rho) + C_k)$ and its limit.

\emph{Clause (4).} $F_d$ is rotation invariant, since $\sigma_{d-1}$ is, so
$F_d(\xi) = F_d(|\xi|e)$. Let $\rho = |\xi| \ge 1$ and write
$F_d(\rho e) = c\log\rho - \log K(1/\rho) + C_k + \Delta(\rho)$. Then
$\rho^{c}e^{-F_d(\rho e)} = e^{-C_k - \Delta(\rho)}K(1/\rho)$, and the two-sided bound of clause
(3) on $\Delta(\rho)$ gives the display. By the display, $\rho^{c}e^{-F_d(\rho e)}$ is bounded on
$[1,\infty)$ exactly when $K$ is bounded on $(0,1]$, which by clause (1) is the condition on the
integral. In that case $K(1/\rho) \to K(0{+}) < \infty$ and $\Delta(\rho) \to c\varkappa_d$, which
gives the limit.

No step cites a fact of the ledger, and every step is machine-checked. The Lean follows this
proof with three changes of bookkeeping. The limit $g_\infty$ is read from the shared library,
which has it with the rate $|\Cin(X) - \log X - g_\infty| \le 2/X$ for $X \ge 1$. The layer-cake
identity is Mathlib's, applied to the measure $w(r)\,dr$, so that $\tau$ is used as a monotone
function only. And the two terms of \ref{eq:profile-sphere} are kept in one integrand,
$\Cin(\rho|\eta_1|\tau(y)) - \log^+(\rho\,\tau(y))$, which lies in
$[\log|\eta_1| - 3, \tfrac{13}{4}]$ and tends to $g_\infty + \log|\eta_1|$ when $\eta_1 \ne 0$
and $0 < y < c$; one dominated convergence on $(0,c)$ and one on the sphere give the limit, and
the same two bounds give the interval.
\end{proof}

% not in the draft: Q-0425, the second of the three nodes (see the comment before
% prop:isotropic-exponent-profile). The record's Steps 2 and 3. Under D3 option (2) these two
% evaluations are computations of the prose: the Lean of the proposition states the constant as
% g_infty + lambda_d. Mathlib at the pin has Real.eulerMascheroniConstant, digamma and
% digamma_one_half (the record's B.8), so a Lean proof is possible later; it is priced in
% records/affine/NODES-q0425.md. No \lean tag.
% CHANGED Q-0430 (2026-10-12): annotation only. The pointer to Chen and Weth is read at the page
% (records/affine/READING-chen2019dirichlet.md). The paper does not print it.
\begin{lemma}[the constant of the exponent in closed form]
\label{lem:exponent-profile-constant}
\uses{prop:isotropic-exponent-profile}
Let $\psi := \Gamma'/\Gamma$ be the digamma function on $(0,\infty)$ and
$\gamma_{\mathrm E} := -\psi(1) = -\Gamma'(1)$ Euler's constant. For an integer $d \ge 1$ let
$g_\infty$, $\lambda_d$ and $\varkappa_d = g_\infty + \lambda_d$ be the constants of
\ref{prop:isotropic-exponent-profile}(2).
\begin{enumerate}
\item $g_\infty = \gamma_{\mathrm E}$.
\item $\lambda_d = \tfrac12\bigl[\psi(\tfrac12) - \psi(\tfrac d2)\bigr]$.
\item Hence
\[
  \varkappa_d = \gamma_{\mathrm E} + \tfrac12\bigl[\psi(\tfrac12) - \psi(\tfrac d2)\bigr]
   = \tfrac12\bigl(\gamma_{\mathrm E} - \psi(\tfrac d2)\bigr) - \log 2 ,
\]
so that $\varkappa_1 = \gamma_{\mathrm E}$, $\varkappa_2 = \gamma_{\mathrm E} - \log 2$,
$\varkappa_3 = \gamma_{\mathrm E} - 1$, and $\varkappa_{d+2} = \varkappa_d - 1/d$ for every
$d \ge 1$.
\item \emph{(The constants in closed form.)} For real numbers $c > 0$ and $C$,
\[
  e^{-C - c\varkappa_d}
   = 2^{c}\exp\Bigl(-C - \tfrac c2\bigl(\gamma_{\mathrm E} - \psi(\tfrac d2)\bigr)\Bigr) .
\]
Consequently, with $R(d,c)$ the constant of \ref{prop:origin-rate-d}(1) for $0 < c < d$,
\[
  R(d,c)\,e^{-C - c\varkappa_d}
   = \frac{\Gamma\bigl(\tfrac{d-c}{2}\bigr)}{\pi^{d/2}\,\Gamma\bigl(\tfrac c2\bigr)}\,
     \exp\Bigl(-C - \tfrac c2\bigl(\gamma_{\mathrm E} - \psi(\tfrac d2)\bigr)\Bigr),
\]
and
\[
  (2\pi)^{-d}\,|S^{d-1}|\,e^{-C - d\varkappa_d}
   = \frac{2}{\pi^{d/2}\,\Gamma\bigl(\tfrac d2\bigr)}\,
     \exp\Bigl(-C - \tfrac d2\bigl(\gamma_{\mathrm E} - \psi(\tfrac d2)\bigr)\Bigr) .
\]
\end{enumerate}
\statusT\quad\emph{(Cites nothing of the ledger: the integral for $\Gamma'$, the functional
equation and the duplication formula of $\Gamma$, Fubini's theorem, and the standard Gaussian
vector in polar coordinates. Two of these are used as known and not proved in the text:
Legendre's duplication formula and the surface measure
$|S^{d-1}| = 2\pi^{d/2}/\Gamma(d/2)$, which the chapter uses elsewhere. Both are in Mathlib at
the pin (\texttt{Real.Gamma\_mul\_Gamma\_add\_half}, \texttt{EuclideanSpace.volume\_ball}), so
neither would be an interface of a formal proof. \textbf{What it is worth.} The constant of
\ref{prop:isotropic-exponent-profile}(3)
depends on the dimension through $\psi(d/2)$ alone, and falls by $1/d$ from dimension $d$ to
$d + 2$. Clause (4) gives the three constants as they are used: $A = e^{-C_k - c\varkappa_d}$ of
\ref{prop:isotropic-exponent-profile}(3), and $Q$ and $B$ of
\ref{prop:origin-rate-profile-d}, each at $C = C_k$. \textbf{Why a separate node.} The proposition
states its constant as a limit and an integral, and holds as it stands. This lemma evaluates
them. Under option (2) of decision D3 of the response plan
(\texttt{records/affine/PLAN-review-response-paper-iso.md}) the two evaluations are computations
of the prose, as the \Matern{} readings of \ref{prop:origin-rate-d} are, and the Lean of the
proposition does not need them. \textbf{Checks}
(\texttt{records/affine/reviews/verify-round1-scripts/b\_slowly\_varying\_factor.py}, rerun for
this node). The mean of $\log|\langle e,\eta\rangle|$ by quadrature against clause (2), for
$d = 2,\dots,7$, to 29 digits. For the \Matern{} datum of index $\gamma$, whose transform
$(1 + \rho^2/2)^{-\gamma}$ is known, $c = 2\gamma$ and the limit of
\ref{prop:isotropic-exponent-profile}(4) must be $2^{\gamma}$: the product
$e^{-C_k - c\varkappa_d}K(0{+})$ equals $2^{\gamma}$ to the 15 digits printed for six pairs $(d,\gamma)$ with
$d$ from $1$ to $5$, one of them with $c > d$. This checks the dependence on $d$, with
$\psi(d/2)$. \textbf{Attribution.} At $d = 1$,
$e^{-C_k - c\varkappa_1} = e^{-c\gamma_{\mathrm E} - C_k}$ is the constant $c$ of
\sscite{sato1978distribution}~(1.8), p.~277, and $\kappa$ of \sscite{sato1999levy}~(53.27),
p.~411, in the symmetric case, with $C_k$ formed from the folded profile; both were compared at
the page image (the record's B.5). The values $\psi(1) = -\gamma_{\mathrm E}$ and
$\psi(\tfrac12) = -\gamma_{\mathrm E} - 2\log 2$ and the recurrence are
\sscite{abramowitz1964handbook}~6.3.2--6.3.5, p.~258, read at the page for the attribution item
Q-0426; the proof derives them and cites nothing. \textbf{A pointer, read at the page.}
$\varkappa_d = -\tfrac12\rho_d$, where
$\rho_N = 2\log 2 + \psi(N/2) - \gamma_{\mathrm E}$ is the constant of the logarithmic Laplacian
in \sscite{chen2019dirichlet}, Thm.~1.1 with (1.4) and (1.5), cited by the page of its arXiv
version (p.~2 of 1710.03416v6). There the operator of symbol $2\log|\xi|$
(part~(ii) of the theorem) is written with the kernel $c_N|z|^{-N}$, $c_N = 2/|S^{N-1}|$, the
term $u(x)$ being kept inside the unit ball only, plus $\rho_N$ times the identity. Inside the
unit ball that kernel is the \Levy{} density of the step datum $(0, 2\cdot\mathbf 1_{(0,1)})$,
whose exponent is $2\log\rho + 2\varkappa_N + o(1)$ by
\ref{prop:isotropic-exponent-profile}(3). The identity $\varkappa_d = -\tfrac12\rho_d$ is a computation of the
reading note \texttt{records/affine/READING-chen2019dirichlet.md}, not a statement of Chen and
Weth, and the paper does not print it. \textbf{Not claimed}: numerical values of the
constants; any
other characterisation of Euler's constant than $-\Gamma'(1)$, which is the one the proof
uses.)}
\end{lemma}

\begin{proof}
\emph{Three facts about $\Gamma$.} For $a > 0$, $\Gamma(a) = \int_0^\infty t^{a-1}e^{-t}\,dt$. On
a compact interval $[a_0,a_1] \subset (0,\infty)$ the derivative of the integrand in $a$,
$t^{a-1}e^{-t}\log t$, is bounded in absolute value by
$(t^{a_0-1} + t^{a_1-1})\,e^{-t}\,|\log t|$, which is integrable on $(0,\infty)$. So one may
differentiate under the integral sign:
\[
  \Gamma'(a) = \int_0^\infty t^{a-1}e^{-t}\log t\,dt ,
\]
the integral converging absolutely. Consequently, if a random variable $T$ has the gamma law of
shape $a$, with density $t^{a-1}e^{-t}/\Gamma(a)$ on $(0,\infty)$, then $\log T$ is integrable and
$\mathbb{E}\log T = \psi(a)$. At $a = 1$ this reads
$\gamma_{\mathrm E} = -\int_0^\infty e^{-v}\log v\,dv$. Next, $\Gamma(a+1) = a\,\Gamma(a)$ gives
$\psi(a+1) = \psi(a) + 1/a$. Last, the duplication formula
$\Gamma(a)\,\Gamma(a+\tfrac12) = 2^{1-2a}\sqrt\pi\,\Gamma(2a)$ gives, on taking logarithms and
differentiating, $\psi(a) + \psi(a+\tfrac12) = 2\psi(2a) - 2\log 2$. At $a = \tfrac12$,
\[
  \psi(\tfrac12) = \psi(1) - 2\log 2 = -\gamma_{\mathrm E} - 2\log 2 .
\]

\emph{Clause (1).} Let $G(X) = \Cin(X) - \log^+ X$ be as in the proof of
\ref{prop:isotropic-exponent-profile}, so that $-3 \le G \le \tfrac{13}{4}$ and
$G(X) \to g_\infty$ as $X \to \infty$. For $0 < \varepsilon < 1$ put
\[
  J(\varepsilon) := \int_0^\infty \bigl(1 - \cos u\bigr)\,e^{-\varepsilon u}\,\frac{du}{u} .
\]
\emph{(a)} $(1 - \cos u)/u = \int_0^1 \sin(tu)\,dt$, and
$\int_0^\infty \sin(tu)\,e^{-\varepsilon u}\,du = t/(\varepsilon^2 + t^2)$. The function
$|\sin(tu)|\,e^{-\varepsilon u}$ is integrable on $(0,1) \times (0,\infty)$, so Fubini's theorem
gives
\[
  J(\varepsilon) = \int_0^1 \frac{t\,dt}{\varepsilon^2 + t^2}
   = \tfrac12\log\bigl(1 + \varepsilon^{-2}\bigr)
   = \log\frac1\varepsilon + \tfrac12\log\bigl(1 + \varepsilon^{2}\bigr) .
\]
\emph{(b)} $e^{-\varepsilon u} = \int_u^\infty \varepsilon e^{-\varepsilon x}\,dx$. The integrand
$(1 - \cos u)\,u^{-1}\,\varepsilon e^{-\varepsilon x}\,\mathbf 1\{u < x\}$ is nonnegative, so
Tonelli's theorem gives
\[
  J(\varepsilon) = \int_0^\infty \varepsilon e^{-\varepsilon x}\,\Cin(x)\,dx
   = \int_0^\infty e^{-v}\,\Cin(v/\varepsilon)\,dv .
\]
\emph{(c)} For $v > 0$,
\[
  \Cin(v/\varepsilon) - \log\frac1\varepsilon
   = G(v/\varepsilon) + \Bigl[\log^+\frac v\varepsilon - \log\frac1\varepsilon\Bigr] .
\]
The bracket equals $\log v$ when $v \ge \varepsilon$. When $v < \varepsilon$ it equals
$\log\varepsilon$, which lies between $\log v$ and $0$. So the bracket is bounded by $|\log v|$ in
absolute value and tends to $\log v$ as $\varepsilon \downarrow 0$. As $|G| \le \tfrac{13}{4}$ and
$G(v/\varepsilon) \to g_\infty$, dominated convergence with the integrable function
$e^{-v}(\tfrac{13}{4} + |\log v|)$ gives, by (b) and $\int_0^\infty e^{-v}dv = 1$,
\[
  J(\varepsilon) - \log\frac1\varepsilon \to g_\infty + \int_0^\infty e^{-v}\log v\,dv
   = g_\infty - \gamma_{\mathrm E} \qquad (\varepsilon \downarrow 0) .
\]
\emph{(d)} By (a) the left side is $\tfrac12\log(1 + \varepsilon^2)$, which tends to $0$. Hence
$g_\infty = \gamma_{\mathrm E}$.

\emph{Clause (2).} Let $Z$ be a standard Gaussian vector of $\RR^d$, with density
$(2\pi)^{-d/2}e^{-|x|^2/2}$, and $e$ a unit vector. In polar coordinates, for a Borel
$h \ge 0$ on $S^{d-1}$ and a Borel $f \ge 0$ on $(0,\infty)$,
\[
  \mathbb{E}\bigl[h(Z/|Z|)\,f(|Z|)\bigr]
   = \int_{S^{d-1}}h\,d\sigma_{d-1}\int_0^\infty f(r)\,(2\pi)^{-d/2}e^{-r^2/2}\,r^{d-1}\,dr .
\]
With $h = 1$ and the substitution $t = r^2/2$, using $|S^{d-1}| = 2\pi^{d/2}/\Gamma(d/2)$, this
shows that $|Z|^2/2$ has the gamma law of shape $d/2$. With $f = 1$ it shows that $Z/|Z|$ has the
law $\sigma_{d-1}/|S^{d-1}|$. The coordinate $\langle e,Z\rangle$ is a standard Gaussian variable,
so $\langle e,Z\rangle^2/2$ has the gamma law of shape $\tfrac12$, by the case $d = 1$. By the
first fact, $\log|Z|$ and $\log|\langle e,Z\rangle|$ are integrable, with
\[
  \mathbb{E}\log|Z| = \tfrac12\bigl[\log 2 + \psi(\tfrac d2)\bigr], \qquad
  \mathbb{E}\log|\langle e,Z\rangle| = \tfrac12\bigl[\log 2 + \psi(\tfrac12)\bigr] .
\]
Since $\langle e,Z\rangle = |Z|\,\langle e, Z/|Z|\rangle$,
\[
  \lambda_d = \mathbb{E}\log\bigl|\langle e, Z/|Z|\rangle\bigr|
   = \mathbb{E}\log|\langle e,Z\rangle| - \mathbb{E}\log|Z|
   = \tfrac12\bigl[\psi(\tfrac12) - \psi(\tfrac d2)\bigr] .
\]

\emph{Clause (3).} The first expression is clauses (1) and (2). The second follows from
$\psi(\tfrac12) = -\gamma_{\mathrm E} - 2\log 2$. At $d = 1$ the first expression is
$\gamma_{\mathrm E}$. At $d = 2$ the second is
$\tfrac12(\gamma_{\mathrm E} - \psi(1)) - \log 2 = \gamma_{\mathrm E} - \log 2$. By
$\psi(a+1) = \psi(a) + 1/a$ at $a = d/2$,
$\varkappa_{d+2} = \varkappa_d - \tfrac12\cdot\tfrac2d = \varkappa_d - 1/d$, and so
$\varkappa_3 = \varkappa_1 - 1$.

\emph{Clause (4).} By the second expression of clause (3),
$-c\varkappa_d = c\log 2 - \tfrac c2(\gamma_{\mathrm E} - \psi(\tfrac d2))$, which is the first
display. The second follows from
$R(d,c) = \Gamma(\tfrac{d-c}{2})/(2^{c}\pi^{d/2}\Gamma(\tfrac c2))$. For the third take $c = d$
in the first display and use $|S^{d-1}| = 2\pi^{d/2}/\Gamma(\tfrac d2)$:
$(2\pi)^{-d}\cdot 2\pi^{d/2}\,\Gamma(\tfrac d2)^{-1}\cdot 2^{d} = 2/(\pi^{d/2}\Gamma(\tfrac d2))$.

No step cites a fact of the ledger. The proof uses Legendre's duplication formula and
$|S^{d-1}| = 2\pi^{d/2}/\Gamma(d/2)$ as known. Nothing here is machine-checked yet.
\end{proof}

% not in the draft: Q-0425, the third of the three nodes: prop:isotropic-exponent-profile(3)
% composed with the two clauses of prop:origin-rate-d. The form, decided when the statement was
% written.
% (i) A new node, not clauses of prop:origin-rate-d: that statement is shared with the isotropic
% paper (pinned), its Lean is complete, and the profile form holds for it as it stands.
% (ii) The constants are written in varkappa_d, the constant of prop:isotropic-exponent-profile,
% so that the node is true and checkable without the closed forms; lem:exponent-profile-constant(4)
% gives Q and B with Euler's constant and the digamma function.
% (iii) Clause (2), when a pure power holds, is the consequence for prop:origin-upper-bound-d that
% the brief asks for, stated as a clause because it is an equivalence with a limit, not a remark.
% (iv) Clause (3) has the limit of q(s)/log(1/s), finite or not, beside the equivalence with
% B L(s): it is what a reader compares with the Matern kernel of index d/2.
% No \lean tag.
% CHANGED Q-0427: typed and proved (Formalization/AffineHemigroup/OriginRateProfile.lean). The
% statement is unchanged. Added: the \lean and \leanok tags, the passage "Formalization" of the
% annotation, and the closing sentence of the proof.
\begin{proposition}[the rate at the origin through the polar profile]
\label{prop:origin-rate-profile-d}
\uses{def:isotropic-origin-datum,prop:isotropic-exponent-profile,prop:origin-rate-d,prop:isotropic-radial-monotone,thm:cited-A11,thm:cited-A2,thm:cited-A3,thm:cited-A8}
\lean{AffineHemigroup.origin_rate_profile, AffineHemigroup.origin_rate_profile_critical,
AffineHemigroup.origin_density_rate_profile, AffineHemigroup.origin_ball_rate_profile,
AffineHemigroup.origin_density_power_of_integrable,
AffineHemigroup.origin_density_power_of_not_integrable,
AffineHemigroup.origin_density_rate_profile_critical,
AffineHemigroup.origin_ball_rate_profile_critical, AffineHemigroup.slowLog_bounds,
AffineHemigroup.slowLog_div_log, AffineHemigroup.origin_density_log_of_integrable,
AffineHemigroup.origin_density_log_of_not_integrable, AffineHemigroup.OriginDatum.profileConst,
AffineHemigroup.OriginDatum.slowLog}\leanok
Let $d \ge 1$, let $(0, k_d)$ be an isotropic origin datum in dimension $d$ with Gaussian
coefficient $0$, let $\mu$ be its law, and suppose that $c := k_d(0{+})$ satisfies $0 < c \le d$.
Let $C_k$, $K$ and $\varkappa_d$ be as in \ref{prop:isotropic-exponent-profile}, let $|S^{d-1}|$
be the surface measure of the unit sphere and $V_d$ the volume of the unit ball, and put
\[
  A := e^{-C_k - c\,\varkappa_d} .
\]
The law $\mu$ has a density $p$ with $p(x) = q(|x|)$ for $x \ne 0$, where
$q : (0,\infty) \to [0,\infty)$ is nonincreasing. Let $q$ be any such function. All limits are as
$s \downarrow 0$.
\begin{enumerate}
\item \emph{(Below the threshold.)} Let $c < d$, let $R(d,c)$ be the constant of
\ref{prop:origin-rate-d}(1), and put $Q := R(d,c)\,A$. Then
\[
  \frac{q(s)}{Q\,s^{c-d}\,K(s)} \to 1
  \qquad\text{and}\qquad
  \frac{\mu(\{|x| \le s\})}{Q\,|S^{d-1}|\,c^{-1}\,s^{c}\,K(s)} \to 1 .
\]
\item \emph{(When a pure power holds.)} Let $c < d$. If
$\int_0^1 (c - k_d(u))\,du/u < \infty$, then $q(s)\,s^{d-c} \to Q\,K(0{+})$, a finite positive
number, and there is $C < \infty$ with $q(s) \le C\,s^{c-d}$ for every $s > 0$. If the integral
diverges, then $q(s)\,s^{d-c} \to \infty$, and there is no such $C$.
\item \emph{(At the threshold.)} Let $c = d$, and put
\[
  \mathcal{L}(s) := \int_s^1 K(y)\,\frac{dy}{y} \quad (0 < s \le 1), \qquad
  B := (2\pi)^{-d}\,|S^{d-1}|\,A .
\]
Then
\[
  \frac{q(s)}{B\,\mathcal{L}(s)} \to 1
  \qquad\text{and}\qquad
  \frac{\mu(\{|x| \le s\})}{B\,V_d\,s^{d}\,\mathcal{L}(s)} \to 1 .
\]
Moreover $\log(1/s) \le \mathcal{L}(s) \le K(s)\log(1/s)$ for $0 < s \le 1$. If
$\int_0^1 (d - k_d(u))\,du/u < \infty$, then
\[
  \frac{\mathcal{L}(s)}{\log(1/s)} \to K(0{+})
  \qquad\text{and}\qquad
  \frac{q(s)}{\log(1/s)} \to B\,K(0{+}) ,
\]
both limits finite and positive. If the integral diverges, then
$\mathcal{L}(s)/\log(1/s) \to \infty$ and $q(s)/\log(1/s) \to \infty$.
\end{enumerate}
\statusA\quad\ledger{A11}\quad\ledger{A8}\quad\ledger{A2}\quad\ledger{A3}\quad\emph{(\textbf{Assignment.}
The node cites nothing of its own. It composes \ref{prop:isotropic-exponent-profile}(3), which
cites nothing, with the two clauses of \ref{prop:origin-rate-d}, and inherits what that node
spends: \ledger{A11}, Karamata's Tauberian theorem, for the limits, and \ledger{A8}, \ledger{A2}
and \ledger{A3}(iii) for the existence of $q$, through
\ref{prop:isotropic-radial-monotone}(2). The limits for a given $q$ do not read its existence.
\textbf{What the citation carries and what it does not.} \ledger{A11} carries the Tauberian step
inside \ref{prop:origin-rate-d}, and nothing about the factor $K$, which is proved in
\ref{prop:isotropic-exponent-profile}. \textbf{What it is worth.} The kernel at its centre is
computed from the datum. The power is $c - d$. The slowly varying factor is $K(s)$, one integral
of the polar profile over $(s,1)$. The constant is explicit in $c$, $d$ and $C_k$, the integral
of the profile beyond radius $1$; with \ref{lem:exponent-profile-constant}(4),
\[
  Q = \frac{\Gamma\bigl(\tfrac{d-c}{2}\bigr)}{\pi^{d/2}\,\Gamma\bigl(\tfrac c2\bigr)}\,
     \exp\Bigl(-C_k - \tfrac c2\bigl(\gamma_{\mathrm E} - \psi(\tfrac d2)\bigr)\Bigr), \qquad
  B = \frac{2}{\pi^{d/2}\,\Gamma\bigl(\tfrac d2\bigr)}\,
     \exp\Bigl(-C_k - \tfrac d2\bigl(\gamma_{\mathrm E} - \psi(\tfrac d2)\bigr)\Bigr) .
\]
\ref{prop:origin-rate-d} has the same two rates normalised by the transform; this node replaces
the transform by the profile. Clause (2) settles what \ref{prop:origin-upper-bound-d} leaves
open at $\beta = c$: for $c < d$ the density is bounded by a constant times the pure power
$|x|^{c-d}$ exactly when $c - k_d(u)$ is integrable against $du/u$ at $0$, and then
$p(x)\,|x|^{d-c}$ has the limit $Q\,K(0{+})$. For the other profiles the loss in the exponent of
\ref{prop:origin-upper-bound-d}(2) cannot be removed. As $K \ge 1$, the limit of clause (2) is at
least $Q > 0$, in every dimension $d \ge 1$. Clause (3) does the same at the
boundary: the density grows like $\log(1/|x|)$ exactly when that integral converges, with the
coefficient $B\,K(0{+})$, and strictly faster otherwise. \textbf{Checks}
(\texttt{records/affine/reviews/verify-round1-scripts/b\_slowly\_varying\_factor.py} and
\texttt{c\_critical\_case.py}, rerun for this node). For the \Matern{} law of index $\gamma$,
$c = 2\gamma < d$, the closed-form density over $Q\,s^{c-d}K(s)$ tends to $1$ in dimensions
$1$, $2$, $3$ and $5$, and $Q\,K(0{+})$ is the coefficient
$\Gamma(d/2 - \gamma)\,2^{-\gamma}\pi^{-d/2}\Gamma(\gamma)^{-1}$ of
\ref{prop:origin-rate-d}'s annotation. At $\gamma = d/2$, $B\,K(0{+})$ equals
$2\,(2\pi)^{-d/2}/\Gamma(d/2)$, the coefficient of $\log(1/|x|)$ in the expansion of the density
$2K_0(\sqrt2|x|)/((2\pi)^{d/2}\Gamma(d/2))$, for $d = 1, 2, 3$. \textbf{Attribution.} At $d = 1$
the statement is Sato and Yamazato's, in the symmetric case: clause (1) is
\sscite{sato1978distribution}~Thm.~1.7(i), with (1.12), p.~277, and clause (3) is its
Thm.~1.7(ii), with (1.14); both are restated as \sscite{sato1999levy}~Thm.~53.8, with (53.28) and
(53.30). The factors $K$ and $\mathcal{L}$ and the constants agree (the verification record's
B.5, at the page image; the record of the attribution item Q-0426, section~3). In $\RR^d$ the
passage from the transform to the density is \ref{prop:origin-rate-d}, with the attribution given
there, and the profile form is proved here; nothing is attributed for it. \textbf{Edge cases.} As in
\ref{prop:origin-rate-d}: $c > d$, $c = \infty$, $c = 0$ and a Gaussian part are excluded, and
$d = 1$ is included. \ref{prop:isotropic-exponent-profile} itself holds for every finite
$c > 0$; only the passage to the density needs $c \le d$. \textbf{Not claimed}: a rate of
convergence; an upper bound with an explicit constant at every radius, which is
\ref{prop:origin-upper-bound-d} with its loss in the exponent; anything for $c > d$.
\textbf{Formalization} (Q-0427). The limits for a given law and a given $q$ are declarations of
their own, and print \ledger{A11} alone; the two that assert that $q$ exists add \ledger{A8},
\ledger{A2} and \ledger{A3}(iii). $q$ is any real function nonincreasing on $(0,\infty)$ such
that $\mu$ has the density $q(|x|)$. It is not asked to be nonnegative; the existence statement
gives a nonnegative one. The ball is closed, the transform is not named, and every limit is
along $s \downarrow 0$ in the real numbers. $\mathcal{L}$ is the total function
\texttt{OriginDatum.slowLog}, read on $(0,1]$, and $A$ is \texttt{OriginDatum.profileConst}.
Each limit with a value in $(0,\infty]$ is two statements: a real limit when
$(c - k_d(u))/u$ is integrable on $(0,1]$, and a limit to $+\infty$ when it is not. The closed
forms of $Q$ and $B$ displayed above are not formalized; they are
\ref{lem:exponent-profile-constant}(4).)}
\end{proposition}

\begin{proof}
\leanok
A density $p(x) = q(|x|)$ with $q$ nonincreasing exists by
\ref{prop:isotropic-radial-monotone}(2), as in the proof of \ref{prop:origin-rate-d}. From here on
$q$ is any nonincreasing function on $(0,\infty)$ such that $x \mapsto q(|x|)$ is a density of
$\mu$. Let $e$ be a unit vector and write $\chi(\rho) := \widehat\mu(\rho e) = e^{-F_d(\rho e)}$
for $\rho > 0$. Since $0 < c \le d < \infty$, \ref{prop:isotropic-exponent-profile}(3) applies and
gives $\chi(\rho)/(A\,\rho^{-c}K(1/\rho)) \to 1$ as $\rho \to \infty$. At $\rho = 1/s$,
\begin{equation}\label{eq:profile-transform-at-s}
  \frac{\widehat\mu(e/s)}{A\,s^{c}\,K(s)} \to 1 \qquad (s \downarrow 0) .
\end{equation}

\emph{Clause (1).} Let $c < d$. By \ref{prop:origin-rate-d}(1),
$q(s)/(R(d,c)\,s^{-d}\,\widehat\mu(e/s)) \to 1$ and
$\mu(\{|x| \le s\})/(R(d,c)\,|S^{d-1}|\,c^{-1}\,\widehat\mu(e/s)) \to 1$. Multiplying each by
\ref{eq:profile-transform-at-s} gives the two limits, with $Q = R(d,c)\,A$.

\emph{Clause (2).} By clause (1), $q(s)\,s^{d-c}/(Q\,K(s)) \to 1$, and $Q > 0$. By
\ref{prop:isotropic-exponent-profile}(1), as $s \downarrow 0$ the factor $K(s)$ increases to the
finite number $K(0{+}) \ge 1$ if $\int_0^1(c - k_d(u))\,du/u < \infty$, and $K(s) \to \infty$ if
the integral diverges. In the second case $q(s)\,s^{d-c} \to \infty$, and no bound
$q(s) \le C s^{c-d}$ holds near $0$. In the first case $q(s)\,s^{d-c} \to Q\,K(0{+})$, which is
finite and positive, so there is $s_0 > 0$ with
$q(s)\,s^{d-c} \le 2Q\,K(0{+})$ for $0 < s \le s_0$. For $s \ge s_0$, as $q$ is nonincreasing,
$1 \ge \mu(\{|x| < s\}) \ge q(s)\,V_d\,s^{d}$, so
$q(s)\,s^{d-c} \le s^{-c}/V_d \le s_0^{-c}/V_d$. So
$C := \max\{2Q\,K(0{+}),\,s_0^{-c}/V_d\}$ serves for every $s > 0$.

\emph{Clause (3), the truncated integral of the transform.} Let $c = d$. In polar coordinates the
$I(s)$ of \ref{prop:origin-rate-d}(2) is
\[
  I(s) = \int_{\{|\xi| \le 1/s\}}\widehat\mu(\xi)\,d\xi = |S^{d-1}|\,P(1/s), \qquad
  P(R) := \int_0^R \chi(\rho)\,\rho^{d-1}\,d\rho .
\]
Put $f(\rho) := \chi(\rho)\,\rho^{d-1}$ and $h(\rho) := A\,K(1/\rho)/\rho$ for $\rho \ge 1$. Both
are positive and continuous, and $f(\rho)/h(\rho) = \chi(\rho)\,\rho^{d}/(A\,K(1/\rho)) \to 1$
as $\rho \to \infty$, by \ref{prop:isotropic-exponent-profile}(3) with $c = d$. The substitution
$y = 1/\rho$ gives, for $R \ge 1$,
\[
  H(R) := \int_1^R h(\rho)\,d\rho = A\int_{1/R}^{1}K(y)\,\frac{dy}{y} = A\,\mathcal{L}(1/R) ,
\]
and $H(R) \ge A\log R \to \infty$, because $K \ge 1$. Let $0 < \varepsilon < 1$ and choose
$R_0 \ge 1$ with $|f - h| \le \varepsilon h$ on $[R_0,\infty)$. For $R \ge R_0$,
\[
  \bigl|P(R) - H(R)\bigr| \le P(1) + \int_1^{R_0}|f - h|\,d\rho + \varepsilon\,H(R) ,
\]
and $P(1) \le 1/d$ because $\chi \le 1$. Dividing by $H(R)$, which tends to infinity,
$|P(R)/H(R) - 1| \le 2\varepsilon$ for all large $R$. So $P(R)/(A\,\mathcal{L}(1/R)) \to 1$ as
$R \to \infty$, that is
\[
  \frac{I(s)}{|S^{d-1}|\,A\,\mathcal{L}(s)} \to 1 \qquad (s \downarrow 0) .
\]

\emph{Clause (3), the two limits.} By \ref{prop:origin-rate-d}(2),
$q(s)/((2\pi)^{-d}I(s)) \to 1$ and $\mu(\{|x| \le s\})/((2\pi)^{-d}V_d\,s^{d}\,I(s)) \to 1$.
Multiplying each by the last display gives the two limits, with
$B = (2\pi)^{-d}|S^{d-1}|A$.

\emph{Clause (3), the logarithm.} Let $0 < s \le 1$. On $(s,1)$, $1 \le K(y) \le K(s)$, so
$\log(1/s) \le \mathcal{L}(s) \le K(s)\log(1/s)$. Let $0 < y_0 < 1$. For $0 < s < y_0$,
$\mathcal{L}(s) \ge \int_s^{y_0}K(y)\,dy/y \ge K(y_0)\log(y_0/s)$, and
$\log(y_0/s)/\log(1/s) \to 1$ as $s \downarrow 0$. So for every number $\kappa < K(y_0)$,
$\mathcal{L}(s)/\log(1/s) \ge \kappa$ for all small $s$. Now use
\ref{prop:isotropic-exponent-profile}(1) at $c = d$. If $\int_0^1(d - k_d(u))\,du/u < \infty$,
then $K(y_0)$ increases to the finite number $K(0{+})$ as $y_0 \downarrow 0$, so the lower bound
holds for every $\kappa < K(0{+})$; and $\mathcal{L}(s) \le K(s)\log(1/s) \le K(0{+})\log(1/s)$.
Hence $\mathcal{L}(s)/\log(1/s) \to K(0{+})$, and with the first limit of the clause
$q(s)/\log(1/s) \to B\,K(0{+})$, which is finite and positive. If the integral diverges, then
$K(y_0) \to \infty$ as $y_0 \downarrow 0$, so the lower bound holds for every real $\kappa$:
$\mathcal{L}(s)/\log(1/s) \to \infty$, and with the first limit $q(s)/\log(1/s) \to \infty$,
because $B > 0$.

The cited fact \ledger{A11} enters through \ref{prop:origin-rate-d} only, and the existence of
$q$ is \ref{prop:isotropic-radial-monotone}(2), on \ledger{A8}, \ledger{A2} and
\ledger{A3}(iii). Everything else is machine-checked: the composition, the comparison of
the two divergent integrals with the substitution $y = 1/\rho$, the bound at every radius of
clause (2), and the bounds and the limit of $\mathcal{L}$.
\end{proof}

% draft: new material (the draft does not mention the filtration; the chain is the released cone
% module's, and this chapter is where Paper VII opens on it). Provenance: Q-0044, from P-0014.
% Nothing below \uses a node of [B]: the released facts these nodes need -- the filtration, the
% witness family, the closed form of its polar profiles and the criterion (A.2) -- are restated
% here with their citations, as this repository's CLAUDE.md requires, and the restatement is the
% general form of a computation [B, App. A] carries out for all c and d while stating its
% conclusion at three parameter points only. That gap is exactly what P-0014 reports.
% CHANGED Q-0281: typed (the statement skeleton). Since Q-0286 the definition displays no claim; its
% one presupposition, that D^{(w)} is a delay datum ([B, Lem. 4.12]'s integrability), is the
% skeleton lemma witness_isDelayDatum, which the \lean tag names. The text is unchanged.
% CHANGED Q-0381 (2026-10-10; the fidelity review's paper-facing point 6, card T0.3): the members
% of S and of A_d are said to be functions on [0,infty), in three places ("the functions omega |->
% F_I(omega^2/2) on [0,infty)" for "the exponents F_I(.^2/2)", "functions on [0,infty)" and
% "for omega >= 0"). A blind restater found that the domain of F decides whether A_infty = S is even
% well posed; [0,infty) is the papers' reading (an admissible F and a causally admissible F_I are
% functions on [0,infty) where they are restated) and the Lean's (classA and classS read F on
% [0,infty) only). No claim changes.
% CHANGED Q-0417 (the freeze of the isotropic paper): the titles of this node and of
% lem:threshold-monotone, lem:threshold-criterion and lem:threshold-limit use the isotropic
% paper's names, "the separating family" for the witness family (D^{(w)}) and "critical values"
% for the thresholds w_d, so that the blueprint prints the titles the paper prints; "threshold"
% alone is the dimension d of thm:origin-boundedness-d. Titles only: the labels, the Lean names
% (witnessProfile, threshold, ...) and the annotations keep the earlier words, and no statement
% changes.
\begin{definition}[the dimension filtration, the separating family, and its critical values]
\label{def:filtration-thresholds}
\uses{def:isotropic-origin-datum,def:isotropic-kernel-family}
\lean{AffineHemigroup.classS, AffineHemigroup.classA, AffineHemigroup.classAInf,
AffineHemigroup.DelayDatum, AffineHemigroup.delayExponentL, AffineHemigroup.witnessProfile,
AffineHemigroup.witnessExponent, AffineHemigroup.witnessF, AffineHemigroup.polarProfileOf,
AffineHemigroup.thresholdFun, AffineHemigroup.thresholdSup, AffineHemigroup.threshold,
AffineHemigroup.witness_isDelayDatum}\leanok
Write $\mathcal{S}$ for the subordinated slice of \restated{B}{Prop.~4.9}{prop:line-filtration}, the set of the functions
$\omega \mapsto F_{\mathrm I}(\omega^2/2)$ on $[0,\infty)$ with $F_{\mathrm I}$ causally admissible
\restatedp{B}{Def.~4.1}{def:line-causal}. For an integer $d \ge 1$ write $\mathcal{A}_d$ for the set of the admissible
$F$ of \restated{V}{(2.3)}{eq:admissible-line}, functions on $[0,\infty)$, for which
$\xi \mapsto F(|\xi|)$, $\xi \in \RR^d$, is the exponent of a self-decomposable law on $\RR^d$, and
let $\mathcal{A}_\infty$ be the intersection of all the $\mathcal{A}_d$. For $w > 0$ let
$D^{(w)}$ be the delay datum \restatedp{B}{Def.~4.11}{sec:line-bridge} with drift $0$ and profile
$k^{(w)}_{\mathrm I}(u) = u^{-1/2} + w\,\mathbf{1}_{[1,2]}(u)$, let $F^{(w)}_{\mathrm I}$ be its
exponent, put $F^{(w)}(\omega) := F^{(w)}_{\mathrm I}(\omega^2/2)$ for $\omega \ge 0$, and write
$k^{(w)}_d$ for its polar profile \ref{eq:polar-profile-d} in dimension $d$. Finally put
\begin{equation}\label{eq:threshold-constant}
  M_d := \sup_{z>0}\, z^m\bigl(e^{-z} - 2^{-d/2}e^{-z/2}\bigr),
  \qquad w_d := \frac{\Gamma(m)}{2M_d}, \qquad m := \frac{d+1}{2} .
\end{equation}
\statusT\quad\emph{(Vocabulary, and no claim. The words \emph{admissible} and
\emph{self-decomposable} are restated in \ref{def:isotropic-kernel-family}'s annotation, as
\ref{eq:admissible-line} and \ref{eq:self-decomposable}. The family $(D^{(w)})$ is the one of
[B, Lem.~4.12], whose parameter [B] writes $c$; this module writes $w$, and $w_d$ for the
thresholds [B, Rem.~4.10] writes $c_d$, because $c$ is the profile at the origin of
\ref{def:isotropic-origin-datum}. What is new here is \ref{eq:threshold-constant}, which names
the threshold that the criterion [B, (A.2)] defines, and the two nodes that follow. The released
facts about these objects are restated where they are used: the filtration
$\mathcal{S} = \mathcal{A}_\infty \subseteq \dots \subseteq \mathcal{A}_2 \subseteq \mathcal{A}_1$,
with $\mathcal{A}_1$ the whole admissible cone, is [B, Prop.~4.9], of which the proof of
\ref{thm:filtration-strict}(3) cites (1) and (4); the closed form [B, (A.1)] of the polar profiles and the
criterion [B, (A.2)], both derived in the paragraph \emph{the polar profiles in closed form} of
[B, App.~A] for every parameter and every integer $d \ge 1$, are \ref{eq:witness-polar} and
\ref{eq:threshold-criterion} in the proof of \ref{lem:threshold-monotone}; and the admissibility
of $F^{(w)}$ for $w \le 6 + 4\sqrt2$ is [B, Lem.~4.12(1)], in the proof of
\ref{thm:filtration-strict}(1). [B]'s lemma states its conclusions only at the parameter points
$(4,1), (4,2), (8,1), (8,2)$ and $(4,3)$. $M_d \in (0,\infty)$, and the equivalence ``$k^{(w)}_d$
is nonincreasing exactly when $w \le w_d$'', are \ref{lem:threshold-monotone}'s first paragraph.
Until Q-0286 this definition also carried those released facts and that equivalence, as claims
inside a definition; they moved to the nodes that use them, and no statement of the chapter gained
or lost a claim in the move. \textbf{The numbers are not used.} [B, Rem.~4.10] reports
$w_1 \approx 12.61$, $w_2 \approx 5.24$, $w_3 \approx 3.21$ and $w_d \to 0$ (written $c_d$ there) as
computed numerically, and P-0014 continues the list to $d = 12$; no node below reads any of those values,
and the facts about them that are stated (that $(w_d)$ decreases, that $w_2$ is below
$6 + 4\sqrt2$, and that $w_d \to 0$) are proved. The values are drawn, from the closed form, by
\texttt{records/affine/fig-thresholds.py}.)}
\end{definition}

% draft: new material. Provenance: Q-0044, from P-0014, whose disagreement is that Remark 4.10 of
% the released module carries a numerical assertion in a load-bearing place. This is the statement
% that replaces the numerics. The proof is one variable and elementary: the first-order condition
% at the maximizer of one dimension's function, tested in the next dimension, against Wendel's
% Gamma inequality; the whole of it comes down to the identity
% (1 - q/sqrt2)^2 - (1-q)(1-q/2) = (3/2 - sqrt 2) q, and 3/2 - sqrt 2 is the reciprocal of
% 6 + 4 sqrt 2, the constant [B, App. A] already meets at d = 1.
% CHANGED Q-0281: typed (the statement skeleton). The clause "when d is an integer, k^{(w)}_d is
% nonincreasing on (0,infty) exactly when w <= w_d", with the closed form [B, (A.1)] and the criterion
% [B, (A.2)] that the proof restates for it, is split to lem:threshold-criterion (ADR-0025;
% records/formalization/PRICING.md, finding 4 and its first proving wave: the one clause that reads
% [B]'s closed form, against the one-variable rest); the text is unchanged and the \lean tag names
% the child's declarations too.
% CHANGED Q-0429 (2026-10-12): annotation only. "Nothing in this chapter needs it [w_d -> 0]" was
% no longer true since Q-0425; it now says which sentence reads the limit. The statement, its
% proof and the \lean tag are unchanged.
\begin{lemma}[the critical values are strictly decreasing]
\label{lem:threshold-monotone}
\uses{def:filtration-thresholds}
\lean{AffineHemigroup.thresholdSup_spec, AffineHemigroup.threshold_pos,
AffineHemigroup.thresholdSup_succ, AffineHemigroup.threshold_strictAnti,
AffineHemigroup.witness_polar_closed_form, AffineHemigroup.witness_threshold_criterion,
AffineHemigroup.witness_antitone_iff}\leanok
For a real parameter $d \ge 1$ put $m = (d+1)/2$, $\lambda_d = 2^{-d/2}$ and
$h_d(z) = z^m(e^{-z} - \lambda_d e^{-z/2})$, and let $M_d = \sup_{z>0}h_d(z)$ as in
\ref{eq:threshold-constant}. Then $M_d \in (0,\infty)$ and the supremum is attained at a point of
$(0,d\log 2)$. Consequently $w_d = \Gamma(m)/(2M_d)$ is a well-defined positive number, and when
$d$ is an integer, $k^{(w)}_d$ is nonincreasing on $(0,\infty)$ exactly when $w \le w_d$. Moreover
\[
  M_{d+1} > \sqrt{m}\;M_d \qquad\text{and}\qquad w_{d+1} < w_d
\]
for every $d \ge 1$. In particular $w_1 > w_2 > w_3 > \dots > 0$, so
$(w_{d+1}, w_d\,]$ is a nonempty interval for every integer $d \ge 1$.
\statusT\quad\emph{(Cites [B, App.~A] for the closed form [B, (A.1)] of the witness profiles and
the criterion [B, (A.2)], restated at the start of the proof, where the equivalence reads them;
nothing else. The two ingredients of the monotonicity are the first-order condition at a
maximizer and the log-convexity of $\Gamma$, and both are proved in the proof. Until Q-0286 the
two restated facts stood in \ref{def:filtration-thresholds}, and the equivalence was stated for a
real parameter $d$, for which $k^{(w)}_d$ is not defined; it is now stated for integers, where it
was always read. \textbf{Not
claimed here}: that $w_d \to 0$, which is the other half of [B, Rem.~4.10]'s numerical sentence and
is not a monotonicity --- it is the statement that $\Gamma(m)/M_d \to 0$, whose content is
Stirling's order $m^{-1/2}$ for $\Gamma(m)/(m^me^{-m})$ and not the bump term, the argument below
giving no rate at all. It is \ref{lem:threshold-limit}, with the rate. The strictness of the
chain does not need it: strictness at every step is \ref{thm:filtration-strict} and reads only
the strict decrease. One sentence of \ref{thm:separating-family-critical-values}(4) does, that
every $w \le w_1$ lies in one of the intervals $(w_{d+1}, w_d\,]$.)}
\end{lemma}

\begin{proof}
\leanok
\emph{The closed form and the criterion.} \prog{B}{App.~A}, restated (the paragraph \emph{the polar
profiles in closed form}, carried out there for every parameter and every integer $d \ge 1$). Let
$d \ge 1$ be an integer and $w > 0$, and write
$P(s,x) = \Gamma(s)^{-1}\int_0^x v^{s-1}e^{-v}dv$ for the regularized lower incomplete gamma
function. Split \ref{eq:polar-profile-d} at the two summands of $k^{(w)}_{\mathrm I}$. For the
power, $k_{\mathrm I}(r^2/2v) = (r^2/2v)^{-1/2}$ turns the integral into a Gamma integral and
gives $2^{3/2}r^{-1}\Gamma(m)/\Gamma(d/2)$. For the bump,
$\mathbf{1}_{[1,2]}(r^2/2v) = \mathbf{1}_{[r^2/4,\,r^2/2]}(v)$. Hence \prog{B}{(A.1)}:
\begin{equation}\label{eq:witness-polar}
  k^{(w)}_d(r) = \frac{2^{3/2}\,\Gamma(m)}{\Gamma(d/2)}\,\frac1r
   + 2w\Bigl[P\Bigl(\frac d2,\frac{r^2}2\Bigr) - P\Bigl(\frac d2,\frac{r^2}4\Bigr)\Bigr],
   \qquad r > 0 ,
\end{equation}
a smooth function on $(0,\infty)$. Since $\partial_x P(s,x) = x^{s-1}e^{-x}/\Gamma(s)$, the
derivative of $P(d/2, r^2/2b)$ in $r$ is $2\,(r\,\Gamma(d/2))^{-1}(r^2/2b)^{d/2}e^{-r^2/2b}$ for
$b > 0$. With $b = 1$ and $b = 2$, and $z = r^2/2$,
\[
  r^2\,\Gamma(d/2)\,\frac{d}{dr}k^{(w)}_d(r)
   = -2^{3/2}\Gamma(m) + 4wr\Bigl(z^{d/2}e^{-z} - \bigl(\tfrac z2\bigr)^{d/2}e^{-z/2}\Bigr) .
\]
As $r = \sqrt{2z}$, the last term is $4\sqrt2\,w\,z^m\bigl(e^{-z} - 2^{-d/2}e^{-z/2}\bigr)$. So the
derivative is nonpositive at $r$ if and only if \prog{B}{(A.2)} holds:
\begin{equation}\label{eq:threshold-criterion}
  w\,z^m\bigl(e^{-z} - 2^{-d/2}e^{-z/2}\bigr) \le \tfrac12\Gamma(m),
  \qquad z = \tfrac12 r^2 .
\end{equation}

\emph{The supremum, and the equivalence.} $h_d$ is smooth on $(0,\infty)$ with $h_d(0{+}) = 0$,
and $e^{-z} - \lambda_d e^{-z/2} > 0$ exactly when $e^{-z/2} > 2^{-d/2}$, that is exactly on
$(0, d\log 2)$, an interval that is nonempty because $d \ge 1$. So $h_d > 0$ somewhere,
$h_d \le 0$ on $[d\log 2,\infty)$, and $h_d(z) \le z^me^{-z} \to 0$ as $z \to \infty$. A positive
continuous function vanishing at both ends of $(0,\infty)$ attains its supremum, so
$M_d = h_d(z^*) \in (0,\infty)$ for some $z^* \in (0, d\log 2)$. For an integer $d$,
\ref{eq:threshold-criterion} reads $w\,h_d(z) \le \tfrac12\Gamma(m)$ at $z = r^2/2$. So the
criterion holds at every $r > 0$ exactly when $w M_d \le \tfrac12\Gamma(m)$, that is exactly when
$w \le w_d$. A differentiable function on $(0,\infty)$ is nonincreasing exactly when its
derivative is nonpositive throughout, which settles the equivalence for $k^{(w)}_d$.

\emph{The first-order condition.} Let $z^* \in (0,d\log 2)$ be a maximizer and put
$q := \lambda_d e^{z^*/2}$, so that $\lambda_d e^{-z^*/2} = q\,e^{-z^*}$. Since $z^* < d\log 2$,
$q \in (0,1)$. Then $M_d = (z^*)^m e^{-z^*}(1-q) > 0$. The maximizer is interior, so
$h_d'(z^*) = 0$, where
\[
  h_d'(z) = m z^{m-1}\bigl(e^{-z} - \lambda_d e^{-z/2}\bigr)
    + z^m\bigl(-e^{-z} + \tfrac12\lambda_d e^{-z/2}\bigr) .
\]
Dividing $h_d'(z^*) = 0$ by $(z^*)^{m-1}e^{-z^*} > 0$ turns it into $m(1-q) = z^*(1 - q/2)$. As
$1 - q/2 > 0$,
\[
  z^* = m\,\frac{1-q}{1-q/2} \qquad\text{and}\qquad
  \sqrt{z^*} = \sqrt{m}\,\sqrt{\frac{1-q}{1-q/2}} .
\]

\emph{Testing the next dimension at that point.} Since
$\lambda_{d+1} = 2^{-1/2}\lambda_d$ and $m_{d+1} = m + \tfrac12$,
\[
  h_{d+1}(z^*) = (z^*)^{m+1/2}\bigl(e^{-z^*} - 2^{-1/2}q\,e^{-z^*}\bigr)
   = \sqrt{z^*}\,\bigl(1 - q/\sqrt2\bigr)\,(z^*)^m e^{-z^*}
   = \sqrt{z^*}\,\frac{1 - q/\sqrt2}{1-q}\,M_d ,
\]
so that, substituting for $\sqrt{z^*}$,
\[
  M_{d+1} \;\ge\; h_{d+1}(z^*)
   \;=\; \sqrt{m}\,M_d\,\frac{1 - q/\sqrt2}{\sqrt{(1-q)(1-q/2)}} .
\]
The last fraction exceeds $1$, because its numerator and denominator are positive and
\[
  \bigl(1 - q/\sqrt2\bigr)^2 - (1-q)\bigl(1-q/2\bigr)
   = \Bigl(1 - \sqrt2\,q + \frac{q^2}{2}\Bigr)
     - \Bigl(1 - \frac32 q + \frac{q^2}{2}\Bigr)
   = \Bigl(\frac32 - \sqrt2\Bigr)q > 0
\]
for $q > 0$, because the constant $\tfrac32 - \sqrt2 = (1 - 2^{-1/2})^2$ is positive. (It is
the reciprocal of $6 + 4\sqrt2$.) Hence $M_{d+1} > \sqrt{m}\,M_d$.

\emph{The Gamma factor.} $\Gamma$ is log-convex on $(0,\infty)$: for $x,y > 0$ and
$\theta \in (0,1)$, H\"older's inequality applied to
$t^{\theta x + (1-\theta)y - 1}e^{-t} = (t^{x-1}e^{-t})^\theta (t^{y-1}e^{-t})^{1-\theta}$ with
exponents $1/\theta$ and $1/(1-\theta)$ gives
$\Gamma(\theta x + (1-\theta)y) \le \Gamma(x)^\theta\Gamma(y)^{1-\theta}$. At
$x = m$, $y = m+1$, $\theta = \tfrac12$ this is
$\Gamma(m + \tfrac12) \le \Gamma(m)^{1/2}\Gamma(m+1)^{1/2} = \sqrt{m}\,\Gamma(m)$, by
$\Gamma(m+1) = m\Gamma(m)$. Therefore
\[
  \frac{1}{2w_{d+1}} = \frac{M_{d+1}}{\Gamma(m + \frac12)}
   > \frac{\sqrt m\,M_d}{\Gamma(m+\frac12)}
   \ge \frac{\sqrt m\,M_d}{\sqrt m\,\Gamma(m)} = \frac{1}{2w_d} ,
\]
which is $w_{d+1} < w_d$. Every $w_d$ is positive, so the chain $w_1 > w_2 > \dots > 0$ and the
nonemptiness of $(w_{d+1},w_d\,]$ follow.
\end{proof}

% CHANGED Q-0281: split from lem:threshold-monotone (ADR-0025; records/formalization/PRICING.md,
% finding 4): the parent's clause "when d is an integer, k^{(w)}_d is nonincreasing on (0,infty)
% exactly when w <= w_d", stated with the two facts of [B, App. A] that the parent's proof restates for
% it, the closed form [B, (A.1)] and the criterion [B, (A.2)] (until Q-0286 claims inside
% def:filtration-thresholds). The parent's text and numbering are unchanged.
\begin{lemma}[the polar profiles of the separating family in closed form, and the monotonicity criterion]
\label{lem:threshold-criterion}
\uses{def:filtration-thresholds,lem:threshold-monotone}
\lean{AffineHemigroup.witness_polar_closed_form, AffineHemigroup.witness_threshold_criterion,
AffineHemigroup.witness_antitone_iff, AffineHemigroup.regIncGamma}\leanok
Let $d \ge 1$ be an integer and $w > 0$, put $m = (d+1)/2$, and write
$P(s,x) = \Gamma(s)^{-1}\int_0^x v^{s-1}e^{-v}dv$ for the regularized lower incomplete gamma
function. Then for $r > 0$
\[
  k^{(w)}_d(r) = \frac{2^{3/2}\,\Gamma(m)}{\Gamma(d/2)}\,\frac1r
   + 2w\Bigl[P\Bigl(\frac d2,\frac{r^2}2\Bigr) - P\Bigl(\frac d2,\frac{r^2}4\Bigr)\Bigr],
\]
a smooth function on $(0,\infty)$; its derivative is nonpositive at $r$ if and only if
$w\,z^m\bigl(e^{-z} - 2^{-d/2}e^{-z/2}\bigr) \le \tfrac12\Gamma(m)$ with $z = \tfrac12 r^2$; and
$k^{(w)}_d$ is nonincreasing on $(0,\infty)$ exactly when $w \le w_d$.
\statusT\quad\emph{(Cites [B, App.~A] for the closed form [B, (A.1)] and the criterion [B, (A.2)],
restated, and no ledger entry. The last sentence is the clause of \ref{lem:threshold-monotone} it
was split from on 2026-10-04 (Q-0281, ADR-0025); it reads that lemma's first paragraph, that $M_d$ is
attained and finite. The parent's text is unchanged.)}
\end{lemma}

\begin{proof}
\leanok
This is the part of the proof of \ref{lem:threshold-monotone} headed \emph{The closed form and the
criterion}, which derives \ref{eq:witness-polar} and \ref{eq:threshold-criterion}, and the second
half of the part headed \emph{The supremum, and the equivalence}.
\end{proof}

% draft: new material. Provenance: Q-0150, from the open item Q-0044 left: the limit in
% [B, Rem. 4.10]'s numerical sentence. The content is Stirling's order m^{-1/2}, proved here in
% the elementary two-sided form it needs rather than cited; the bump term only has to be
% evaluated at z = m, where it is a geometric factor (e/4)^{d/4}.
% CHANGED Q-0425 (2026-10-12): annotation only. "Not claimed: sharp constants ... nothing here says
% that limit exists" now points to lem:threshold-asymptotic, which has the limit. The statement,
% its proof and the \lean tag are unchanged, and so is the isotropic paper's verbatim copy.
% CHANGED Q-0429 (2026-10-12): annotation only. Its last sentence said that the exclusion of a
% fixed w from A_d for every large d "is not stated as a node here"; it is now a consequence of
% thm:separating-family-critical-values(1). The statement, its proof and the \lean tag are
% unchanged.
\begin{lemma}[the critical values tend to $0$ at the rate $d^{-1/2}$]
\label{lem:threshold-limit}
\uses{def:filtration-thresholds,lem:threshold-monotone}
\lean{AffineHemigroup.threshold_limit_bounds, AffineHemigroup.threshold_limit,
AffineHemigroup.thresholdEps, AffineHemigroup.stirlingK}\leanok
For a real parameter $d \ge 1$ let $m$, $h_d$, $M_d$ and $w_d = \Gamma(m)/(2M_d)$ be as in
\ref{lem:threshold-monotone}, and put
\[
  \varepsilon_d := 2^{-d/2}e^{m/2} = e^{1/4}\Bigl(\frac e4\Bigr)^{d/4},
  \qquad K := \sqrt{\pi/2} + \sqrt\pi + 4 .
\]
Then for every $d \ge 1$
\[
  w_d \;\ge\; \frac{e^{-1/2}}{2\sqrt m},
\]
and for every $d \ge 1$ with $\varepsilon_d < 1$
\[
  w_d \;\le\; \frac{K}{2(1-\varepsilon_d)\sqrt m} .
\]
Since $e < 4$, $\varepsilon_d$ decreases strictly to $0$ as $d \to \infty$. Hence
$w_d \to 0$, and $w_d$ is of exact order $d^{-1/2}$: $\sqrt m\,w_d \ge e^{-1/2}/2$ for every
$d \ge 1$, and $\sqrt m\,w_d \le K$ for every $d$ with $\varepsilon_d \le \tfrac12$, which is every
large $d$. In particular, for every fixed $w > 0$ there is a $D$ with $w > w_d$ for all
$d \ge D$.
\statusT\quad\emph{(Cites nothing. The Stirling bounds used, $e^{-1/2} \le
\Gamma(m)\,m^{1/2-m}e^m \le K$ for $m \ge 1$, are proved in the proof from the substitution
$t = m(1+u)$ and two elementary inequalities for $u - \log(1+u)$; the only constants read are
$e < 4$ and the Gaussian integral. No value of any $w_d$ is used, and neither is the monotonicity
of \ref{lem:threshold-monotone}, which the node \emph{uses} only for $M_d \in (0,\infty)$.
\textbf{Not claimed}: the limit of $\sqrt m\,w_d$. Neither $e^{-1/2}/2$ nor $K/2$ is that limit.
Since Q-0425 it is \ref{lem:threshold-asymptotic}: $\sqrt m\,w_d \to \sqrt{\pi/2}$, that is
$\sqrt d\,w_d \to \sqrt\pi$, along the real parameter $d$. The two bounds stated here are kept
beside it because they hold at every $d$, the lower one, and at every $d$ with
$\varepsilon_d < 1$, the upper one, which a limit does not give. Nor any statement about
$\mathcal{A}_d$. The last sentence is about
the numbers alone; that a fixed $w$ puts $F^{(w)}$ outside $\mathcal{A}_d$ for every $d \ge D$ is
[B, Rem.~4.10]'s other sentence, which the exclusion paragraph of \ref{thm:filtration-strict}'s
proof would give, since that paragraph reads $w > w_{d+1}$ and nothing else. Since Q-0429 it
follows from \ref{thm:separating-family-critical-values}(1), $F^{(w)} \in \mathcal{A}_D$ exactly
when $w \le w_D$, with the last sentence of this lemma.)}
\end{lemma}

\begin{proof}
\leanok
\emph{Stirling's bounds.} Let $m \ge 1$ and $\varphi(u) := u - \log(1+u)$ for $u > -1$, so that
$\varphi(0) = 0$ and $\varphi'(u) = u/(1+u)$. The substitution $t = m(1+u)$ in
$\Gamma(m+1) = \int_0^\infty t^m e^{-t}\,dt$ gives
\[
  \Gamma(m+1) = m^{m+1}e^{-m}\int_{-1}^\infty e^{-m\varphi(u)}\,du ,
  \qquad\text{so}\qquad
  \Gamma(m) = m^m e^{-m}\,I_m ,\quad I_m := \int_{-1}^\infty e^{-m\varphi(u)}\,du ,
\]
by $\Gamma(m+1) = m\Gamma(m)$. Four comparisons for $\varphi$ follow by comparing derivatives and
integrating from $0$. Put $\psi(u) := u^2/(2(1+u))$ for $u \ge 0$.
\begin{itemize}
\item On $(-1,0]$, $\varphi(u) \ge u^2/2$: there $\varphi'(s) = s/(1+s) \le s$, because
$s \le 0 < 1+s \le 1$. Integrate over $[u,0]$.
\item On $[0,\infty)$, $\varphi(u) \ge \psi(u)$: $\psi(0) = 0$ and
$\psi'(s) = (s^2+2s)/(2(1+s)^2) \le (2s^2+2s)/(2(1+s)^2) = \varphi'(s)$.
\item $\psi(u) \ge u^2/4$ on $[0,1]$, since $1+u \le 2$ there, and $\psi(u) \ge u/4$ on
$[1,\infty)$, since $u/(1+u) \ge \tfrac12$ there.
\item On $[0,\infty)$, $\varphi(u) \le u^2/2$, because $\varphi'(s) \le s$ there.
\end{itemize}
With $\int_0^\infty e^{-as^2}ds = \tfrac12\sqrt{\pi/a}$ for $a > 0$, the lower bounds for
$\varphi$, on $(-1,0]$, on $[0,1]$ and on $[1,\infty)$, give
\[
  I_m \le \int_{-\infty}^0 e^{-mu^2/2}\,du + \int_0^\infty e^{-mu^2/4}\,du
   + \int_1^\infty e^{-mu/4}\,du
  = \sqrt{\frac{\pi}{2m}} + \sqrt{\frac{\pi}{m}} + \frac4m e^{-m/4}
  \le \frac{K}{\sqrt m} ,
\]
using $4/m \le 4/\sqrt m$ for $m \ge 1$. The integrand is positive and $1/\sqrt m \le 1$, so the
upper bound $\varphi(u) \le u^2/2$ gives
\[
  I_m \ge \int_0^{1/\sqrt m} e^{-mu^2/2}\,du \ge \frac{1}{\sqrt m}\,e^{-1/2} .
\]
So $e^{-1/2}\,m^{m-1/2}e^{-m} \le \Gamma(m) \le K\,m^{m-1/2}e^{-m}$ for every $m \ge 1$, and
$m = (d+1)/2 \ge 1$ because $d \ge 1$.

\emph{$M_d$ from above.} Since $2^{-d/2}e^{-z/2} > 0$, $h_d(z) < z^m e^{-z}$, and
$z \mapsto z^m e^{-z}$ has derivative $z^{m-1}e^{-z}(m - z)$, so it is largest at $z = m$. Hence
$M_d \le m^m e^{-m}$, and by \ref{lem:threshold-monotone} $M_d > 0$, so
\[
  w_d = \frac{\Gamma(m)}{2M_d} \ge \frac{e^{-1/2}\,m^{m-1/2}e^{-m}}{2\,m^m e^{-m}}
   = \frac{e^{-1/2}}{2\sqrt m} .
\]

\emph{$M_d$ from below.} Evaluating at $z = m$,
\[
  M_d \ge h_d(m) = m^m\bigl(e^{-m} - 2^{-d/2}e^{-m/2}\bigr) = m^m e^{-m}\,(1 - \varepsilon_d) ,
\]
and $\varepsilon_d = 2^{-d/2}e^{(d+1)/4} = e^{1/4}(e/4)^{d/4}$. When $\varepsilon_d < 1$ this is
positive, and
\[
  w_d = \frac{\Gamma(m)}{2M_d} \le \frac{K\,m^{m-1/2}e^{-m}}{2\,m^m e^{-m}(1-\varepsilon_d)}
   = \frac{K}{2(1-\varepsilon_d)\sqrt m} .
\]

\emph{The limit.} Since $e < 4$, $e/4 \in (0,1)$ and $\varepsilon_d$ decreases strictly to $0$.
In particular $\varepsilon_d \le \tfrac12$ for all large $d$, and there the upper bound reads
$\sqrt m\,w_d \le K$. Since $m = (d+1)/2 \to \infty$, $w_d \to 0$ follows. For fixed $w > 0$ the
upper bound falls below $w$ from some $D$ on, and so does $w_d$.
\end{proof}

% not in the draft: queue item Q-0425 (hub), decision D4 option (1) of the response to the external
% review of the isotropic paper, round 1 (records/affine/PLAN-review-response-paper-iso.md): the
% limit that lem:threshold-limit's annotation said was not claimed. The derivation is the
% verification record's (records/affine/reviews/VERIFY-round1-open-problems.md, part A).
% A separate lemma, not a clause of lem:threshold-limit, for two reasons checked before anything
% was written. (i) The isotropic paper prints lem:threshold-limit verbatim (the marker
% `% shared with blueprint lem:threshold-limit` at paper-iso/appendix-isotropic.tex, unpinned), so
% a clause added to it would make a shared statement drift, which this item may not do. (ii) That
% node is \leanok on threshold_limit_bounds and threshold_limit, neither of which states a limit of
% sqrt(m) w_d; a clause without Lean under that tag would break ADR-0025. The paper is free to
% print this lemma as a further clause of its Lemma A.5 when it is mirrored.
% Clause (2) is stated because it is what the limit rests on and it holds at every d: the squeeze
% of w_d between two multiples of Gamma(m) e^m m^{-m}. No \lean tag.
% CHANGED Q-0425 (the blind check, report A, verdict "correct as stated"; repair R7 and one
% finding): the annotation says how sqrt(d) w_d fails to be monotone, with the numbers rerun
% (a_critical_values.py and the referee's t6_threshold.py reproduce their saved outputs), and that
% the upper bound of clause (2) says nothing at d = 1, 2. The referee's wording had the expansion
% sqrt(pi) - sqrt(d) w_d ~ sqrt(pi)/(3d) as a fact; it stays under "Not claimed", because it needs
% the second term of Stirling's series and clause (1) proves the first. Statement and proof
% unchanged, but for "of the ledger" in the proof's closing sentence.
% CHANGED Q-0427: typed and proved (Formalization/AffineHemigroup/ThresholdAsymptotic.lean). The
% statement is unchanged. Added: the \lean and \leanok tags, the passage "Formalization" of the
% annotation with the choices the Lean statement makes, and the closing sentence of the proof,
% which said that nothing was machine-checked.
\begin{lemma}[the limit of the critical values]
\label{lem:threshold-asymptotic}
\uses{def:filtration-thresholds,lem:threshold-monotone,lem:threshold-limit}
\lean{AffineHemigroup.tendsto_stirlingRatio, AffineHemigroup.threshold_asymptotic_bounds,
AffineHemigroup.threshold_asymptotic, AffineHemigroup.stirlingRatio}\leanok
For a real parameter $d \ge 1$ let $m = (d+1)/2$, $M_d$, $w_d = \Gamma(m)/(2M_d)$ and
$\varepsilon_d = 2^{-d/2}e^{m/2}$ be as in \ref{lem:threshold-limit}, and put
\[
  S_m := \Gamma(m)\,m^{1/2-m}\,e^{m} .
\]
\begin{enumerate}
\item \emph{(Stirling's formula.)} $S_m \to \sqrt{2\pi}$ as $m \to \infty$ through the real
numbers.
\item \emph{(The critical value against the Gamma function.)} For every $d \ge 1$,
\[
  \sqrt m\,w_d \;\ge\; \tfrac12\,S_m ,
\]
and for every $d \ge 1$ with $\varepsilon_d < 1$,
\[
  \sqrt m\,w_d \;\le\; \frac{S_m}{2\,(1-\varepsilon_d)} .
\]
\item \emph{(The limit.)} As $d \to \infty$ through the real numbers,
\[
  \sqrt m\,w_d \to \sqrt{\pi/2} \qquad\text{and}\qquad \sqrt d\,w_d \to \sqrt\pi .
\]
\end{enumerate}
\statusT\quad\emph{(Cites nothing. The proof reads the proof of \ref{lem:threshold-limit}: its
substitution, the three lower comparisons for $u - \log(1+u)$, and its two bounds on $M_d$. The
one new ingredient is the pointwise limit of the integrand. \textbf{What it is worth.}
\ref{lem:threshold-limit} shows that $w_d$ is of exact order $d^{-1/2}$, between the constants
$e^{-1/2}/2$ and $K/2$ for $\sqrt m\,w_d$, neither of which is the limit. Clause (3) gives the
limit: $w_d$ is asymptotic to $\sqrt{\pi/d}$. Clause (2) says more than the limit at a given $d$:
$w_d$ equals $\tfrac12\Gamma(m)\,e^{m}m^{-m}$ up to a factor in $[1, (1-\varepsilon_d)^{-1}]$, and
$\varepsilon_d = e^{1/4}(e/4)^{d/4}$ decays geometrically. The upper bound of clause (2) says
nothing at $d = 1$ and $d = 2$: $\varepsilon_d < 1$ exactly when $d > 1/\log(4/e) = 2.5887$, and
$\varepsilon_1 = 1.166$, $\varepsilon_2 = 1.058$, $\varepsilon_3 = 0.961$. At $d = 3$ it reads
$\sqrt m\,w_d \le 33.5$ against the value $4.53$. The bounds of
\ref{lem:threshold-limit} are kept: they hold at every $d$, which a limit does not give.
\textbf{Why a separate lemma.} The isotropic paper prints \ref{lem:threshold-limit} verbatim, and
that node's Lean covers all of it. A further clause there would have changed a shared statement
and put a clause with no Lean under a complete tag. \textbf{Attribution.} Clause (1) is
Stirling's formula for the Gamma function, \sscite{abramowitz1964handbook}~6.1.37, p.~257 (6.1.38
there is the two-sided form for a real argument), read at the page for the attribution item
Q-0426. It is proved here in the form the lemma needs and is not cited. \textbf{Numbers}
(\texttt{records/affine/reviews/verify-round1-scripts/a\_critical\_values.py} and the blind
referee's \texttt{records/formalization/blind-check-q0425-A/t6\_threshold.py}, both rerun for
this node; none is used in the proof). The convergence is slow at small $d$, because of
$\varepsilon_d$: $\sqrt d\,w_d$ is $12.61$ at $d = 1$, $2.056$ at $d = 20$, which is $16\,\%$
above $\sqrt\pi = 1.77245$, and $1.7784$ at $d = 50$. And $\sqrt d\,w_d$ is not monotone. Over
the integers it decreases strictly from $d = 1$ to $d = 81$ and increases strictly from $d = 81$
on, as far as computed, which is $d = 300$. It is above $\sqrt\pi$ for $d \le 55$ and below it
from $d = 56$ on ($1.77197$ at $d = 56$). Its least value over the integers is $1.76610$, at
$d = 81$; over the real parameter the minimum is near $d = 80.75$. It then returns to $\sqrt\pi$
from below: $\sqrt\pi - \sqrt d\,w_d$ is $5.906 \cdot 10^{-4}$ at $d = 10^3$ and
$5.908 \cdot 10^{-9}$ at $d = 10^8$, against $\sqrt\pi/(3d) = 0.5908/d$. By contrast $w_d$ and
$\sqrt m\,w_d$ decrease strictly over $d = 1,\dots,300$. So a table to $d = 12$ or $d = 20$ does
not show the limit. \textbf{Readings.} The limits are along the real parameter $d$, hence along
the integers. \textbf{Not claimed}: a rate of convergence, or the expansion
$\sqrt d\,w_d = \sqrt\pi\,(1 - 1/(3d) + \dots)$ that the numbers show. It would follow from the
second term of Stirling's series, $S_m = \sqrt{2\pi}\,(1 + 1/(12m) + \dots)$, with
$\sqrt{d/m} = \sqrt2\,(1 - 1/(2d) + \dots)$ and the geometric decay of $\varepsilon_d$; clause (1)
proves the first term only. Nor that $\sqrt d\,w_d$ is eventually increasing, or below $\sqrt\pi$
for every $d \ge 56$, which are numerical beyond what was computed; nor that $\sqrt m\,w_d$ is
monotone; nor any statement about $\mathcal{A}_d$. \textbf{Formalization} (Q-0427). $S_m$ is the
total function \texttt{stirlingRatio}, $\Gamma(m)\,m^{1/2-m}e^{m}$ with the real Gamma function
and a real power, for every real $m$. Clause (1) is a limit along the filter of the real numbers
at $+\infty$, with no hypothesis on $m$; the proof reads $m > 0$ for the integral and $m \ge 1$
for the domination, both of which hold eventually. Clause (2) is stated for every real $d \ge 1$,
and its upper bound carries the hypothesis $\varepsilon_d < 1$. The two limits of clause (3) are
along the real parameter $d$. $M_d$, $w_d$ and $\varepsilon_d$ are the functions of
\ref{lem:threshold-limit}'s Lean. Mathlib at the pin has Stirling's formula for the factorial
only; a probe derives that formula back from clause (1) along the integers, where $S_n$ is
$\sqrt2$ times Mathlib's sequence $n!/(\sqrt{2n}\,(n/e)^n)$.)}
\end{lemma}

\begin{proof}
\leanok
From the proof of \ref{lem:threshold-limit}: with $\varphi(u) := u - \log(1+u)$ for $u > -1$,
\[
  \Gamma(m) = m^m e^{-m}\,I_m , \qquad I_m := \int_{-1}^\infty e^{-m\varphi(u)}\,du ,
\]
so that $S_m = \sqrt m\,I_m$. Also $\varphi(u) \ge u^2/2$ on $(-1,0]$, $\varphi(u) \ge u^2/4$ on
$[0,1]$ and $\varphi(u) \ge u/4$ on $[1,\infty)$. And $0 < M_d \le m^m e^{-m}$ for every
$d \ge 1$, $M_d \ge m^m e^{-m}(1 - \varepsilon_d)$, and $\varepsilon_d \to 0$ as $d \to \infty$.

\emph{Clause (1).} The substitution $u = v/\sqrt m$ gives
\[
  S_m = \sqrt m\,I_m = \int_{\RR} f_m(v)\,dv ,
\]
where $f_m(v) := e^{-m\varphi(v/\sqrt m)}$ for $v > -\sqrt m$ and $f_m(v) := 0$ for
$v \le -\sqrt m$.
\emph{The pointwise limit.} Since $\varphi'(s) = s/(1+s)$ and $\varphi(0) = 0$,
$\varphi(u) - \tfrac12 u^2 = -\int_0^u s^2/(1+s)\,ds$. For $|u| < 1$ the denominator is at least
$1 - |u|$ on the interval of integration, so
\[
  \bigl|\varphi(u) - \tfrac12 u^2\bigr| \le \frac{|u|^3}{3\,(1 - |u|)} .
\]
Fix a real $v$. For $m > v^2$, the point $u = v/\sqrt m$ has $|u| < 1$, so $v > -\sqrt m$ and
\[
  \bigl|m\,\varphi(v/\sqrt m) - \tfrac12 v^2\bigr|
   \le \frac{|v|^3}{3\sqrt m\,\bigl(1 - |v|/\sqrt m\bigr)} ,
\]
which tends to $0$ as $m \to \infty$. Hence $f_m(v) \to e^{-v^2/2}$ for every real $v$.
\emph{The domination.} Let $m \ge 1$ and $v > -\sqrt m$, and apply the three comparisons at
$u = v/\sqrt m$. For $-\sqrt m < v \le 0$, $m\varphi(u) \ge v^2/2$. For $0 \le v \le \sqrt m$,
$m\varphi(u) \ge v^2/4$. For $v \ge \sqrt m$, $m\varphi(u) \ge \sqrt m\,v/4 \ge v/4$. So
\[
  0 \le f_m(v) \le g(v) := e^{-v^2/4} + e^{-|v|/4} \qquad (v \in \RR,\ m \ge 1),
\]
and $g$ is integrable on $\RR$. \emph{The limit.} By dominated convergence, along every sequence
$m_n \to \infty$, $S_{m_n} \to \int_{\RR}e^{-v^2/2}\,dv = \sqrt{2\pi}$. So $S_m \to \sqrt{2\pi}$ as
$m \to \infty$.

\emph{Clause (2).} $w_d = \Gamma(m)/(2M_d) = m^m e^{-m}I_m/(2M_d)$. As $0 < M_d \le m^m e^{-m}$,
$w_d \ge I_m/2$, that is $\sqrt m\,w_d \ge S_m/2$. If $\varepsilon_d < 1$, then
$M_d \ge m^m e^{-m}(1 - \varepsilon_d) > 0$ gives $w_d \le I_m/(2(1 - \varepsilon_d))$.

\emph{Clause (3).} As $d \to \infty$, $m = (d+1)/2 \to \infty$ and $\varepsilon_d \to 0$, so
$\varepsilon_d < 1$ for all large $d$, and both sides of clause (2) tend to $\sqrt{2\pi}/2 =
\sqrt{\pi/2}$ by clause (1). Hence $\sqrt m\,w_d \to \sqrt{\pi/2}$. Finally
$\sqrt d\,w_d = \sqrt{d/m}\,\sqrt m\,w_d$ and $d/m = 2d/(d+1) \to 2$, so
$\sqrt d\,w_d \to \sqrt2\,\sqrt{\pi/2} = \sqrt\pi$.

No step cites a fact of the ledger, and every step is machine-checked. The Lean takes the
substitution in the one step $t = m + \sqrt m\,v$ and proves the three comparisons at
$u = v/\sqrt m$ from those of \ref{lem:threshold-limit}'s Lean. Its pointwise limit reads the
bound $|\log(1+u) - u + \tfrac12 u^2| \le |u|^3/(1 - |u|)$ for $|u| < 1$ from Mathlib's estimate
of the logarithm series, without the factor $\tfrac13$, which the limit does not need. Dominated
convergence is applied along the filter of the real numbers at $+\infty$, which is countably
generated, so no sequence is named.
\end{proof}

% Q-0316 (the blind register review of the isotropic paper, records/affine/REVIEW-register-paper-iso.md,
% Appendix A flag 1, severity 3): the membership of the witness F^(w) in A_d rested on an unstated
% feature of [B]'s proof, that the polar-profile computation of [B, Prop. 4.9(2)] reads no
% monotonicity of the delay profile, while [B, Prop. 4.9(2)] as restated assumes a nonincreasing one.
% This lemma states that computation for every delay datum and proves it here; the exclusion step of
% thm:filtration-strict reads it too. No \lean tag (the brief of Q-0316): the statement skeleton
% session types it. Cost, priced at the statement: Mathlib has the Gaussian characteristic function,
% Tonelli (lintegral_lintegral_swap) and polar coordinates; the substitution u = r^2/2v and the
% Gamma integral are the only computations, so a few hundred lines, no new interface.
% CHANGED Q-0281 wave 2: typed (the statement skeleton of the three nodes added since). Readings:
% "k_d is finite" on (0,infty) is the integrability of the integrand of eq:polar-profile-d at every
% r > 0; "measurable on (0,infty)" is a.e.-measurability for Lebesgue measure on (0,infty), k_I being
% a.e.-measurable; the exponent is the [0,infty]-valued one, so the identity needs no finiteness
% proviso. The text is unchanged.
% CHANGED Q-0281 wave 2: blind review N3-1: "measurable on (0,infty)" is now read at its letter, the
% restriction of k_d to (0,infty) being measurable (it agrees there with the measurable
% |S^{d-1}| r^d f(r)); the a.e.-measurability typed first was weaker than the node, and the reason
% given for it (k_I is a.e.-measurable only) did not bind, k_d being an integral over k_I. The Lean
% statement is made stronger; the text is unchanged.
\begin{lemma}[the polar profile of a delay datum, without monotonicity]
\label{lem:delay-polar-profile}
\uses{def:isotropic-origin-datum,lem:polar-decomposition,lem:positive-transform-origin}
\lean{AffineHemigroup.delay_polar_profile, AffineHemigroup.delay_polar_profile_datum}\leanok
Let $d \ge 1$ and let $(b_0, k_{\mathrm I})$ be a delay datum in the sense of \restated{B}{Def.~4.11}{sec:line-bridge},
restated: $b_0 \ge 0$, and $k_{\mathrm I} : (0,\infty) \to [0,\infty)$ is measurable, not
necessarily monotone, with $\int_0^1 k_{\mathrm I}(u)\,du < \infty$ and
$\int_1^\infty k_{\mathrm I}(u)\,u^{-1}du < \infty$. Its exponent is
$F_{\mathrm I}(\sigma) = b_0\sigma + \int_0^\infty(1 - e^{-\sigma u})\,k_{\mathrm I}(u)\,du/u$.
Let $k_d$ be its polar profile \ref{eq:polar-profile-d}. Then $k_d$ is finite and measurable on
$(0,\infty)$, with
\[
  \int_0^1 k_d(r)\,r\,dr < \infty
  \qquad\text{and}\qquad
  \int_1^\infty k_d(r)\,\frac{dr}{r} < \infty ,
\]
and
\[
  F_{\mathrm I}\bigl(\tfrac12|\xi|^2\bigr) = \tfrac12 b_0|\xi|^2
   + \int_{\RR^d}\bigl(1 - \cos\langle\xi,x\rangle\bigr)\,\frac{k_d(|x|)}{|S^{d-1}|\,|x|^d}\,dx ,
   \qquad \xi \in \RR^d .
\]
In particular, if $k_d$ is nonincreasing, then $(b_0,k_d)$ is an isotropic origin datum in
dimension $d$ whose exponent is $\xi \mapsto F_{\mathrm I}(\tfrac12|\xi|^2)$.
\statusT\quad\emph{(Cites nothing: the Gaussian characteristic function, Tonelli's theorem, one
substitution and polar coordinates. \textbf{What it is.} The computation of [B, Prop.~4.9(2)], whose
polar profile is [B, (4.6)], carried out for a delay datum whose profile need not be monotone.
[B] states that proposition for a causally admissible exponent, with a nonincreasing delay profile,
and its App.~A reads the computation for the witness family of \ref{def:filtration-thresholds},
whose delay profile is not monotone, without stating that case. Here it is stated and proved, so
that \ref{thm:filtration-strict} rests on a statement and not on the shape of a released proof.
\textbf{What it is worth.} Monotonicity enters only in the last sentence: the identity of
exponents and the integrability conditions hold for every delay datum, and whether
$(b_0,k_d)$ is an origin datum is then the question whether $k_d$ is nonincreasing, which for the
witness family is \ref{lem:threshold-criterion}. \textbf{Not claimed}: that $k_d$ is
nonincreasing, which fails for the witness in high dimension; nor that $F_{\mathrm I}(\cdot^2/2)$ is
admissible or in any $\mathcal{A}_d$ when it is not; nor anything about a law: the lemma is an
identity between functions, and the law, where one is needed, is \ledger{A3}(iii) at the place it
is used. Provenance: Q-0316, from the blind register review of the isotropic paper, Appendix~A
flag~1.)}
\end{lemma}

\begin{proof}
\leanok
For $u > 0$ let $g_u(x) = (2\pi u)^{-d/2}e^{-|x|^2/2u}$ be the heat kernel of
\ref{lem:positive-transform-origin}. It is even and its transform is $e^{-u|\xi|^2/2}$, so
\[
  1 - e^{-u|\xi|^2/2} = \int_{\RR^d}\bigl(1 - \cos\langle\xi,x\rangle\bigr)\,g_u(x)\,dx .
\]
Integrate against $k_{\mathrm I}(u)\,du/u$ and exchange the integrals by Tonelli's theorem, which
applies because the integrand is nonnegative:
\[
  F_{\mathrm I}\bigl(\tfrac12|\xi|^2\bigr) - \tfrac12 b_0|\xi|^2
   = \int_{\RR^d}\bigl(1 - \cos\langle\xi,x\rangle\bigr)\,f(x)\,dx ,
  \qquad f(x) := \int_0^\infty g_u(x)\,k_{\mathrm I}(u)\,\frac{du}{u} ,
\]
with $f$ measurable and valued in $[0,\infty]$. Fix $x \ne 0$ and put $r = |x|$. Under the
substitution $u = r^2/2v$, $du/u$ becomes $dv/v$ with the orientation reversed, and
$g_u(x) = \pi^{-d/2}r^{-d}v^{d/2}e^{-v}$. So, by \ref{eq:polar-profile-d} and
$|S^{d-1}| = 2\pi^{d/2}/\Gamma(d/2)$,
\[
  f(x) = \pi^{-d/2}r^{-d}\int_0^\infty v^{d/2-1}e^{-v}\,k_{\mathrm I}\Bigl(\frac{r^2}{2v}\Bigr)dv
   = \frac{k_d(r)}{|S^{d-1}|\,r^d} .
\]
This is the identity of exponents. The value $k_d(r)$ is finite. For fixed $x \ne 0$ the factor
$g_u(x)/u$ is bounded for $u \in (0,1]$, because $e^{-r^2/2u}$ decays faster than every power of
$1/u$ grows, and $g_u(x) \le (2\pi)^{-d/2}$ for $u \ge 1$. The two conditions on $k_{\mathrm I}$
then bound the integral defining $f(x)$ over $(0,1]$ and over $[1,\infty)$. Measurability of $k_d$
is Tonelli's theorem again. For the integrability conditions, $\int|x|^2g_u(x)\,dx = du$ and
$\int g_u = 1$ give $\int(|x|^2 \wedge 1)\,g_u(x)\,dx \le \min\{du, 1\}$. By Tonelli's theorem
\[
  \int_{\RR^d}(|x|^2 \wedge 1)\,f(x)\,dx \le d\int_0^1 k_{\mathrm I}(u)\,du
   + \int_1^\infty k_{\mathrm I}(u)\,\frac{du}{u} < \infty ,
\]
and in polar coordinates the left side is $\int_0^1 k_d(r)\,r\,dr + \int_1^\infty k_d(r)\,dr/r$.
If $k_d$ is also nonincreasing, $(b_0,k_d)$ satisfies \ref{def:isotropic-origin-datum}. Its
\Levy{} measure has the density $f$, and its exponent \ref{eq:isotropic-exponent} is
$F_{\mathrm I}(\tfrac12|\cdot|^2)$ by the identity above.
\end{proof}

% draft: new material. Provenance: Q-0044, from P-0014: "If (c_d) is strictly decreasing then any
% c in (c_{d+1}, c_d] puts F^(c) in A_d minus A_{d+1}" (this module writes w and w_d for [B]'s c and
% c_d since Q-0286). The bound w <= 6 + 4 sqrt 2 in clause (1) is not needed for the mathematics --
% the criterion at d = 1 gives admissibility for every w <= w_1 -- but it keeps the admissibility of
% F^(w) inside what [B, Lem. 4.12(1)] states, and clause (2) shows it binds at no d >= 2.
% CHANGED Q-0281: typed (the statement skeleton). Clause (3)'s last identification A_infty = S is
% split to lem:filtration-top (ADR-0025; records/formalization/PRICING.md, finding 3): it is
% [B, Prop. 4.9(4)], whose inclusion A_infty in S is Schoenberg's radial theorem (entry A23 of [B]'s ledger),
% stated in no Lean of either module, while the strictness of the chain does not read it. The text is
% unchanged; the node now \uses the child, and its \lean tag names the child's declaration too.
% CHANGED Q-0281: \uses re-pointed to the split child it reads (eq:witness-polar and the criterion beside the monotonicity of w_d).
% CHANGED Q-0281: blind review R2-3/R5-5: \uses gains def:isotropic-origin-datum (the proof of (1)
% checks its integrability conditions directly) and lem:polar-decomposition (k^(w)_{d+1} = k_{d+1}
% a.e. "by polar coordinates").
% CHANGED Q-0316 (the blind register review of the isotropic paper, Appendix A flag 1): the
% membership and exclusion steps read lem:delay-polar-profile, the polar-profile computation stated
% without monotonicity, where they read an unstated feature of [B]'s proof; \uses gains it. The
% statement is unchanged.
\begin{theorem}[the dimension filtration is strict at every step]
\label{thm:filtration-strict}
% CHANGED Q-0281 wave 2: clauses (1)-(3) proved (filtration_witness, filtration_witness_exists,
% filtration_strict, AffineHemigroup.FiltrationStrict, guarded; A2-A4). The node stayed not ready
% because its tag names filtration_top (lem:filtration-top, then unscheduled; proved in wave N). \uses gains
% prop:isotropic-polar-classification, which the membership step now reads for admissibility; the
% statusA's sentence on [B, Lem. 4.12(1)] says so (blind review W2A-4: "in the membership step"
% became "beside the membership step (which, since Q-0281, reads the admissibility from
% prop:isotropic-polar-classification(2) instead)").
% CHANGED Q-0281 wave N (2026-10-07): green. Its fourth declaration, filtration_top
% (lem:filtration-top), is proved on ledger A10, so the guard prints A10 on this node; the status
% names A10 and \uses gains thm:cited-A10, and the Assignment says the entry is spent through the
% child. Annotation only: the statement and the proof are unchanged, and so is the isotropic paper's
% shared copy, which cites [B, Prop. 4.9(4)] for the identification.
\uses{def:filtration-thresholds,lem:threshold-monotone,prop:isotropic-polar-reading,prop:isotropic-polar-classification,lem:filtration-top,lem:threshold-criterion,def:isotropic-origin-datum,lem:polar-decomposition,lem:delay-polar-profile,thm:cited-A2,thm:cited-A3,thm:cited-A4,thm:cited-A10}
\lean{AffineHemigroup.filtration_witness, AffineHemigroup.filtration_witness_exists,
AffineHemigroup.filtration_strict, AffineHemigroup.filtration_top}\leanok
In the notation of \ref{def:filtration-thresholds}:
\begin{enumerate}
% CHANGED Q-0281 wave 2: the bound w <= 6+4 sqrt 2 dropped from clause (1), unused (the author's
% decision, 2026-10-05; a widening). Clause (2), a statement about the bound itself, keeps its
% content; its referent "the set of parameters w allowed in clause (1)" is written out as "the set of
% w in (w_{d+1}, w_d] with w <= 6+4 sqrt 2", which is what it meant before the drop. The statusA
% gains a sentence recording the drop.
\item \emph{(One witness per step.)} Let $d \ge 1$ be an integer and let
$w \in (w_{d+1},\,w_d\,]$. Then
$F^{(w)} \in \mathcal{A}_d \setminus \mathcal{A}_{d+1}$.
\item \emph{(There is one at every step.)} $w_2 < 6 + 4\sqrt2$. So the set of
$w \in (w_{d+1},\,w_d\,]$ with $w \le 6 + 4\sqrt2$ is $(w_{d+1}, w_d\,]$ itself for every
$d \ge 2$, and it is nonempty for every $d \ge 1$.
\item \emph{(The chain.)} $\mathcal{A}_1 \supsetneq \mathcal{A}_2 \supsetneq \mathcal{A}_3
\supsetneq \dots \supsetneq \mathcal{A}_\infty = \mathcal{S}$, every inclusion being strict, and
$\mathcal{A}_\infty \subsetneq \mathcal{A}_d$ for every $d$. In particular the filtration is
infinite: no $\mathcal{A}_D$ equals $\mathcal{A}_\infty$, so admissibility in a finite dimension
never forces the Brownian-subordination picture.
\end{enumerate}
\statusA\quad\ledger{A2}\quad\ledger{A3}\quad\ledger{A4}\quad\ledger{A10}\quad\emph{(\textbf{Assignment.} The
cited analysis is spent through \ref{prop:isotropic-polar-reading}(1), in both directions, and
through the \Levy--Khintchine representation directly. \ledger{A10}, Schoenberg's radial theorem,
is spent only through \ref{lem:filtration-top}, for the identification
$\mathcal{A}_\infty = \mathcal{S}$ in clause (3); the strictness of the chain does not read it. The converse direction of that clause puts
the witness's law into $\mathcal{A}_d$; it spends \ledger{A2}, and \ledger{A3}(iii) for the law of
the datum. The forward direction keeps the witness out of $\mathcal{A}_{d+1}$; it spends
\ledger{A4}, which makes the given self-decomposable law infinitely divisible, with \ledger{A3}
and \ledger{A2}. The exclusion also spends \ledger{A3} directly: (iii) for the law that
\ref{lem:delay-polar-profile} builds in dimension $d+1$, and the uniqueness (ii) to identify its
\Levy{} measure with the datum's. Everything else is \ref{lem:threshold-monotone},
\ref{lem:delay-polar-profile} (which cites nothing) and restated facts of [B], each named at the
place it is used: [B, Lem.~4.12(1)] beside the membership step (which, since Q-0281, reads the
admissibility from \ref{prop:isotropic-polar-classification}(2) instead), [B, Prop.~4.9(1)] in
clause (3), and (4), which \ref{lem:filtration-top} proves on \ledger{A10}. Until Q-0316 the membership and exclusion steps read [B, Prop.~4.9(2)]'s computation at a delay
profile that is not monotone, on the strength of how [B, App.~A] carries it out; that computation is
now \ref{lem:delay-polar-profile}, stated and proved without monotonicity. Until
Q-0286 the status named \ledger{A2} alone and the proof cited the polar reading as the
old clause (5) of \ref{thm:origin-boundedness-d}, dropped then as a duplicate of
\ref{prop:isotropic-polar-reading}(1); the forward direction had always spent \ledger{A4}, and the
exclusion the uniqueness of \ledger{A3}, inside that clause, so the declaration is corrected, not
widened. [B, Lem.~4.12(3)] reads its two parameter points through the same two cited facts, the
polar criterion and the \Levy--Khintchine representation. Until 2026-10-05 clause (1) also
required $w \le 6 + 4\sqrt2$, the bound up to which [B, Lem.~4.12(1)] states admissibility; since
Q-0281 no step reads it, and it was dropped (the author's decision), so clause (2) now says of that
bound only that it binds at no $d \ge 2$. \textbf{Not claimed}: any statement about
which \emph{other} members separate the steps (the witness family is one curve through the cone,
and the theorem says nothing about the shape of $\mathcal{A}_d \setminus \mathcal{A}_{d+1}$), and no
quantitative form of the last clause, for which see \ref{lem:threshold-limit}'s annotation.
\textbf{What the clauses are worth.} [B, Lem.~4.12(3)] has clause (1) at $d = 1$ and $d = 2$,
from the two parameter points $w = 8$ and $w = 4$, and [B, Rem.~4.10] says in a remark that the
thresholds keep falling. That is the state P-0014 objects to: the sentence explaining the lemma
rests on a computation. Clause (3) is the objection answered, and it is also the answer to the
question P-0014 calls the one with the larger payoff either way: a collapse
$\mathcal{A}_D = \mathcal{A}_\infty$ at a finite $D$ would have made admissibility in dimension $D$
force subordination, and there is none. Clause (3) also carries the strictness
$\mathcal{S} \subsetneq \mathcal{A}_d$ that \ref{prop:isotropic-bridge} stated as its clause (5)
until Q-0286.)}
\end{theorem}

\begin{proof}
\leanok
% CHANGED Q-0281 wave 2: clause (2)'s constants rewritten to the Lean route (coarser, and enough).
\emph{Clause (2).} By \ref{eq:threshold-constant}, $2M_2 w_2 = \Gamma(3/2) < 1$, since $\Gamma$
is below $1$ on $(1,2)$. And $M_2 > 1/15$: taking $z = \tfrac12$ in
\ref{eq:threshold-constant} at $d = 2$, where $m = \tfrac32$ and $2^{-d/2} = \tfrac12$, gives
$M_2 \ge 2^{-3/2}\bigl(e^{-1/2} - \tfrac12 e^{-1/4}\bigr)$, and
$e^{-1/2} > 3/5$ because $e < 25/9$, while $e^{-1/4} < 4/5$ because $e^{1/4} > 1 + \tfrac14$, so
the bracket exceeds $3/5 - 2/5 = 1/5$. And $2^{-3/2} = \tfrac12\sqrt{1/2} > 1/3$ because
$\sqrt{1/2} > 2/3$. Hence $M_2 > 1/3 \cdot 1/5 = 1/15$. Therefore
$w_2 < 1/(2 \cdot 1/15) = 15/2 < 10 \le 6 + 4\sqrt2$. Since
$(w_d)$ is strictly decreasing (\ref{lem:threshold-monotone}), $w_d \le w_2 < 6+4\sqrt2$ for every
$d \ge 2$, so for $d \ge 2$ the constraint $w \le 6+4\sqrt2$ removes nothing from
$(w_{d+1},w_d\,]$, which is nonempty. At $d = 1$ the set is
$(w_2,\,\min\{w_1, 6+4\sqrt2\}\,]$, nonempty because $w_2 < 6+4\sqrt2$ and $w_2 < w_1$.

% CHANGED Q-0281 wave 2: the membership step rewritten to the Lean route. Admissibility comes from
% prop:isotropic-polar-classification(2), not from [B, Lem. 4.12(1)], so the bound w <= 6 + 4 sqrt 2
% is not read; it stayed in the statement until the author dropped it (2026-10-05, a widening).
\emph{Clause (1), membership.} Let $w$ be as in the clause; then $w > w_{d+1} > 0$. The delay datum $D^{(w)}$ has
drift $0$ and the delay profile $k^{(w)}_{\mathrm I}$, which is not monotone, and
\ref{lem:delay-polar-profile} applies to it, since that lemma reads no monotonicity. It gives that
$k^{(w)}_d$ is finite and satisfies the two integrability conditions of
\ref{def:isotropic-origin-datum}, and that
\[
  F^{(w)}(|\xi|) = F^{(w)}_{\mathrm I}\bigl(\tfrac12|\xi|^2\bigr)
   = \int_{\RR^d}\bigl(1 - \cos\langle\xi,x\rangle\bigr)\,
     \frac{k^{(w)}_d(|x|)}{|S^{d-1}|\,|x|^d}\,dx , \qquad \xi \in \RR^d .
\]
Since $w \le w_d$, $k^{(w)}_d$ is nonincreasing on $(0,\infty)$ (\ref{lem:threshold-monotone}). So
$(0,k^{(w)}_d)$ is an isotropic origin datum in dimension $d$ whose exponent is
$\xi \mapsto F^{(w)}(|\xi|)$. By \ref{prop:isotropic-polar-classification}(2), which reads
\ref{prop:isotropic-polar-reading}(1) in its converse direction, $F^{(w)} \in \mathcal{A}_d$: it
is admissible, and $\xi \mapsto F^{(w)}(|\xi|)$ is the exponent of a self-decomposable law on
$\RR^d$. (For $w \le 6 + 4\sqrt2$ the admissibility is also \prog{B}{Lem.~4.12(1)}.)

\emph{Clause (1), exclusion.} Since $w > w_{d+1}$, \ref{eq:threshold-criterion} fails at some
$z_0 > 0$ in dimension $d+1$, so $k^{(w)}_{d+1}$ has a positive derivative at
$r_0 = \sqrt{2z_0}$ and is continuous and strictly increasing on an interval around $r_0$.
Therefore it agrees with no nonincreasing function almost everywhere. (Given $r < s$ in that
interval with $k^{(w)}_{d+1}(r) < k^{(w)}_{d+1}(s)$, points $u_n \downarrow r$ and
$v_n \uparrow s$ of the conull agreement set of a nonincreasing $g$ would give
$k^{(w)}_{d+1}(u_n) = g(u_n) \ge g(v_n) = k^{(w)}_{d+1}(v_n)$, and continuity would force
$k^{(w)}_{d+1}(r) \ge k^{(w)}_{d+1}(s)$.) Suppose $F^{(w)} \in \mathcal{A}_{d+1}$. By
\ref{prop:isotropic-polar-reading}(1), in its forward direction, $\xi \mapsto F^{(w)}(|\xi|)$ on
$\RR^{d+1}$ is the exponent of an isotropic origin datum $(a, k_{d+1})$ in dimension $d+1$, whose
profile $k_{d+1}$ is nonincreasing. On the other hand \ref{lem:delay-polar-profile}, in dimension
$d+1$, holds for every $w > 0$, since it reads no monotonicity. Let $\nu'$ be the rotation-invariant
measure on $\RR^{d+1} \setminus \{0\}$ with density $k^{(w)}_{d+1}(|x|)(|S^{d}|\,|x|^{d+1})^{-1}$.
By that lemma $\int(|x|^2 \wedge 1)\,\nu'(dx) < \infty$ and
$F^{(w)}(|\xi|) = \int(1 - \cos\langle\xi,x\rangle)\,\nu'(dx)$, so by \ledger{A3}(iii) the
infinitely divisible law with triplet $(0,\nu',0)$ has transform $e^{-F^{(w)}(|\xi|)}$: $\nu'$ is
symmetric, so the odd part of the integrand integrates to $0$. The datum's
law has the same transform, so the two laws are equal, and the uniqueness of the triplet,
\ledger{A3}(ii), gives $a = 0$ and $\nu' = \nu_{d+1}$. The two radial densities then agree almost
everywhere on $\RR^{d+1}$, so $k^{(w)}_{d+1} = k_{d+1}$ almost everywhere on $(0,\infty)$, by polar
coordinates. That makes $k^{(w)}_{d+1}$ agree almost everywhere with a nonincreasing function, which
the first part of this paragraph excludes, a contradiction. So $F^{(w)} \notin \mathcal{A}_{d+1}$.

\emph{Clause (3).} Clauses (1) and (2) give, for every integer $d \ge 1$, a member of
$\mathcal{A}_d \setminus \mathcal{A}_{d+1}$. The nesting $\mathcal{A}_{d+1} \subseteq
\mathcal{A}_d$ is \restatedc{B}{Prop.~4.9}{prop:line-filtration}{(1)}{(1)}; it holds because the image of a self-decomposable law on
$\RR^{d+1}$ with transform $e^{-F(|\xi|)}$ under the projection on the first $d$ coordinates is
self-decomposable on $\RR^d$, with transform $\xi \mapsto e^{-F(|(\xi,0)|)} = e^{-F(|\xi|)}$, and
the residuals project likewise.
% CHANGED Q-0281 wave 2: the nesting argued (the Lean proves it by the projection).
% CHANGED Q-0381: proof rewritten to the Lean route (the author's decision on #96). Until then
% the identification A_infty = S was cited from [B, Prop. 4.9(4)]; the Lean proves it as
% lem:filtration-top (filtration_top) on ledger A10, and proves the strictness (filtration_strict)
% without it. The proof now says so, outlines lem:filtration-top's route, and keeps [B] as the
% pointer for the classical statement. The isotropic paper prints the same text with the pointer
% to lem:filtration-top replaced, since it does not print that lemma, and without the closing
% paragraph on what is machine-checked, which it states in its section 1.1 (pinned).
So $\mathcal{A}_{d+1} \subsetneq \mathcal{A}_d$ for every $d \ge 1$. Since $\mathcal{A}_\infty$ is
the intersection of the classes $\mathcal{A}_d$, $\mathcal{A}_\infty \subseteq \mathcal{A}_{d+1}
\subsetneq \mathcal{A}_d$, so $\mathcal{A}_\infty \subsetneq \mathcal{A}_d$ for every $d$, and no
$\mathcal{A}_D$ with $D$ finite equals $\mathcal{A}_\infty$. None of this reads what
$\mathcal{A}_\infty$ is.

The identification $\mathcal{A}_\infty = \mathcal{S}$, which completes the displayed chain, is
\ref{lem:filtration-top}. It is the statement of \restatedc{B}{Prop.~4.9}{prop:line-filtration}{(4)}{(3)}, proved here from
Schoenberg's radial theorem, \ledger{A10} (\ref{thm:cited-A10}), by the following route. For
$\mathcal{S} \subseteq \mathcal{A}_\infty$, let $F = F_{\mathrm I}(\cdot^2/2)$ with $F_{\mathrm I}$
causally admissible. $F$ is admissible by \restated{B}{Lem.~4.2}{prop:line-bridge}, and since the delay profile of
$F_{\mathrm I}$ is nonincreasing, so is its polar profile \ref{eq:polar-profile-d} in every
dimension $d$; \ref{lem:delay-polar-profile} and \ref{prop:isotropic-polar-reading}(1), in its
converse direction, then put $F$ in $\mathcal{A}_d$. For $\mathcal{A}_\infty \subseteq \mathcal{S}$,
let $F \in \mathcal{A}_\infty$ and put $f(\sigma) = F(\sqrt{2\sigma})$. In every dimension $d$,
\ref{prop:isotropic-polar-reading}(1), in its forward direction, writes $\eta \mapsto F(|\eta|)$ in
\Levy--Khintchine form with a \Levy{} measure $\nu_d$ whose polar profile is nonincreasing, so that
the image of $\nu_d$ under $x \mapsto cx$ is at most $\nu_d$ for $0 < c \le 1$. Hence
$\eta \mapsto F(|\eta|)$ and $\eta \mapsto F(|\eta|) - F(c|\eta|)$ are continuous and negative
definite on $\RR^d$ for every $d$, and \ledger{A10} represents $f$ and each residual
$\sigma \mapsto f(\sigma) - f(c^2\sigma)$ in the form $a + b\sigma + \int(1 - e^{-\sigma t})\,m(dt)$.
Comparing the representation of $f$ with the sum of those of $f(c^2\,\cdot)$ and of the residual,
the uniqueness of such a representation shows that the image of the measure $m$ of $f$ under every
dilation $t \mapsto \lambda t$, $0 < \lambda \le 1$, is at most $m$. So $m(dt) = k(t)\,t^{-1}dt$
with $k$ nonincreasing, the constant term of $f$ is $0$ because $F$ is continuous with $F(0) = 0$,
and $F = F_{\mathrm I}(\cdot^2/2)$ for the causally admissible $F_{\mathrm I}$ with drift $b$ and
delay profile $k$. With the identification the chain is the displayed one, ending at
$\mathcal{S}$, with every step strict.

Every step is machine-checked. The cited interfaces are \ledger{A2}, \ledger{A3} and \ledger{A4},
through \ref{prop:isotropic-polar-reading}(1) and, in the exclusion step of clause (1), \ledger{A3}
directly; and, for the identification $\mathcal{A}_\infty = \mathcal{S}$ alone, \ledger{A10}, through
\ref{lem:filtration-top}. The strictness of the chain, clause (3) without that identification, does
not read \ledger{A10}.
\end{proof}

% ADDED Q-0281 wave N (2026-10-07): the first of the three children of lem:filtration-top (ADR-0025),
% split so that the proof of record reads in steps and each piece is priced alone. Written against
% the Lean statement (lkForm_sub_negDef): the comparison clause (alpha', L') is what the proof of
% lem:filtration-top needs for the residual F(|eta|) - F(c|eta|), whose Levy measure is the
% difference nu_d - (c.)_* nu_d; stating it as a comparison avoids writing that difference as a
% measure (Mathlib's Measure.sub needs a finite measure, and nu_d is not finite). Priced at ~150
% lines (120-180 in the design); 206 written, with the module docstring.
\begin{lemma}[the negative definiteness of a \Levy--Khintchine form]
\label{lem:lk-negative-definite}
\lean{AffineHemigroup.lkForm_sub_negDef, AffineHemigroup.lkForm_negDef}\leanok
Let $d \ge 0$ be an integer, let $\alpha' \le \alpha$ be real numbers, and let $L' \le L$ be
measures on $\RR^d$ such that $\int_{\RR^d}(1 - \cos\langle\eta,x\rangle)\,L(dx) < \infty$ for
every $\eta \in \RR^d$. Put
\[
  \psi(\eta) = \alpha|\eta|^2 + \int_{\RR^d}\bigl(1 - \cos\langle\eta,x\rangle\bigr)\,L(dx),
  \qquad
  \psi'(\eta) = \alpha'|\eta|^2 + \int_{\RR^d}\bigl(1 - \cos\langle\eta,x\rangle\bigr)\,L'(dx),
\]
and $\phi = \psi - \psi'$. Then $\phi$ is negative definite on $\RR^d$: for every $n$, all
$\eta_1, \dots, \eta_n \in \RR^d$ and all $c_1, \dots, c_n \in \CC$,
\[
  \sum_{j,k=1}^n \bigl(\phi(\eta_j) + \phi(\eta_k) - \phi(\eta_j - \eta_k)\bigr)\,c_j\,\overline{c_k}
  \ge 0 .
\]
In particular, if $\alpha \ge 0$, then $\psi$ is negative definite on $\RR^d$.
\statusT\quad\emph{(Cites nothing: a trigonometric identity, the polarization identity and the
monotonicity of the integral in the measure. ``Negative definite'' is the sense of
\ref{thm:cited-A10}, Schoenberg's, with the inequality in $\CC$: the sum is real and nonnegative.
\textbf{What it is for.} \ref{lem:filtration-top} applies it to the exponent of an isotropic origin
datum and to the residual $F(|\eta|) - F(c|\eta|)$, which supplies the hypothesis of
\ref{thm:cited-A10} in every dimension. The comparison form, with $L' \le L$, is what lets the
residual be handled without writing its \Levy{} measure as a difference of measures.
\textbf{Not claimed}: any converse (that a continuous negative definite function has such a form is
the \Levy--Khintchine formula on $\RR^d$, which is not used); and nothing about $\psi$ when the
integral diverges.)}
\end{lemma}

\begin{proof}
\leanok
First let the coefficients be real, $a_1, \dots, a_n \in \RR$. Write $\theta_j = \langle\eta_j,x\rangle$.
Since $\cos(\theta_j - \theta_k) = \cos\theta_j\cos\theta_k + \sin\theta_j\sin\theta_k$,
\[
  (1 - \cos\theta_j) + (1 - \cos\theta_k) - \bigl(1 - \cos(\theta_j - \theta_k)\bigr)
  = (1 - \cos\theta_j)(1 - \cos\theta_k) + \sin\theta_j\sin\theta_k ,
\]
and $|\eta_j|^2 + |\eta_k|^2 - |\eta_j - \eta_k|^2 = 2\langle\eta_j,\eta_k\rangle$. Each function
$x \mapsto 1 - \cos\langle\eta,x\rangle$ is integrable for $L$ by hypothesis, so the integrals may be
added, and
\[
  \sum_{j,k=1}^n a_j a_k\bigl(\psi(\eta_j) + \psi(\eta_k) - \psi(\eta_j - \eta_k)\bigr)
  = 2\alpha\Bigl|\sum_j a_j\eta_j\Bigr|^2 + \int_{\RR^d} P_a(x)\,L(dx),
\]
where
\[
  P_a(x) = \Bigl(\sum_j a_j\bigl(1 - \cos\langle\eta_j,x\rangle\bigr)\Bigr)^2
  + \Bigl(\sum_j a_j\sin\langle\eta_j,x\rangle\Bigr)^2 \ \ge\ 0 .
\]
$P_a$ is integrable for $L$, being a finite combination of the integrable functions above. Since
$L' \le L$, every function integrable for $L$ is integrable for $L'$, so the same identity holds
for $\psi'$, with $\alpha'$ and $L'$; and $\int P_a\,dL' \le \int P_a\,dL$, $P_a$ being nonnegative.
Subtracting,
\[
  \sum_{j,k} a_j a_k\bigl(\phi(\eta_j) + \phi(\eta_k) - \phi(\eta_j - \eta_k)\bigr)
  = 2(\alpha - \alpha')\Bigl|\sum_j a_j\eta_j\Bigr|^2
  + \int P_a\,dL - \int P_a\,dL' \ \ge\ 0 .
\]
Now let the $c_j$ be complex. The matrix $K_{jk} = \phi(\eta_j) + \phi(\eta_k) - \phi(\eta_j -
\eta_k)$ is real and symmetric, because $\phi$ is even. Writing $c_j = a_j + ib_j$,
\[
  \sum_{j,k} K_{jk}\,c_j\overline{c_k}
  = \sum_{j,k} K_{jk}(a_ja_k + b_jb_k) + i\sum_{j,k} K_{jk}(b_ja_k - a_jb_k) .
\]
The imaginary part vanishes, since exchanging $j$ and $k$ changes its sign and leaves $K$ unchanged,
and the real part is the sum of the two real cases just treated, so it is nonnegative. For the last
sentence take $\alpha' = 0$ and $L' = 0$, so that $\psi' = 0$.
\end{proof}

% ADDED Q-0281 wave N (2026-10-07): the second child of lem:filtration-top. It is the last sentence
% of [SSV, Thm. 3.2], which ledger A10 deliberately does not carry (the ledger is as narrow as the
% argument allows). Priced at ~150-200 lines; 159 written, with the module docstring. The Lean takes the uniqueness of the
% Laplace transform of a measure on (0,infty) from the shared library
% (ScaleSpace.laplace_uniqueness_locally_finite, Lean core, ScaleSpaceCore v0.4.0, moved there from
% Paper V's package by Q-0360; Paper V's until the switch to v0.4.0), a standard fact that [V] v0.1
% does not print as a numbered statement.
\begin{lemma}[the uniqueness of a Bernstein triplet]
\label{lem:bernstein-triplet-unique}
\lean{AffineHemigroup.bernstein_triplet_unique}\leanok
Let $a, b, a', b' \ge 0$ and let $m, m'$ be measures on $(0,\infty)$ such that, for every
$\sigma > 0$,
\[
  a + b\sigma + \int_{(0,\infty)}\bigl(1 - e^{-\sigma t}\bigr)\,m(dt)
  = a' + b'\sigma + \int_{(0,\infty)}\bigl(1 - e^{-\sigma t}\bigr)\,m'(dt) < \infty .
\]
Then $a = a'$, $b = b'$ and $m = m'$.
\statusT\quad\emph{(Cites nothing in the paper's sense: the uniqueness of the Laplace transform of a
measure on $(0,\infty)$ whose transform is finite on $(0,\infty)$, a standard fact, which the Lean
takes from the shared library, where it is proved. \textbf{What it is.} The last sentence of
\sscite{schilling2012bernstein}, Thm.~3.2, p.~21 (``the triplet $(a, b, \mu)$ determines $f$
uniquely''), which \ref{thm:cited-A10} does not carry. No integrability of $1 \wedge t$ is assumed:
the finiteness of both sides is all the proof reads. \textbf{Not claimed}: that every function of
this form is a Bernstein function, or anything about derivatives.)}
\end{lemma}

\begin{proof}
\leanok
Write $G(\sigma)$ for the common value. For $t > 0$ and $\sigma > 0$,
\[
  \bigl(1 - e^{-\sigma t}\bigr) + \bigl(1 - e^{-(\sigma+2)t}\bigr)
  + e^{-\sigma t}\bigl(1 - e^{-t}\bigr)^2 = 2\bigl(1 - e^{-(\sigma+1)t}\bigr),
\]
as expanding the square shows, and $b\sigma + b(\sigma + 2) = 2b(\sigma + 1)$. Integrating against
$m$, every term being nonnegative,
\[
  G(\sigma) + G(\sigma + 2) + \int_{(0,\infty)} e^{-\sigma t}\bigl(1 - e^{-t}\bigr)^2\,m(dt)
  = 2G(\sigma + 1),
\]
and the same with $a'$, $b'$ and $m'$. All the values of $G$ are finite, so the measures
$n(dt) = (1 - e^{-t})^2\,m(dt)$ and $n'(dt) = (1 - e^{-t})^2\,m'(dt)$ on $(0,\infty)$ have the same
Laplace transform $\int e^{-\sigma t}\,n(dt) = 2G(\sigma+1) - G(\sigma) - G(\sigma+2)$ at every
$\sigma > 0$, and it is finite. Such a transform determines the measure: $n_1(dt) := e^{-t}\,n(dt)$
is finite, its image under $t \mapsto e^{-t}$ is a finite measure on $(0,1)$ whose $k$-th moment is
$\int e^{-(k+1)t}\,n(dt)$, and likewise for $n'$; so the two images have the same moments, finite
measures on $[0,1]$ are determined by their moments (by the Weierstrass approximation theorem), and
$n_1 = n'_1$, so $n = n'$. The density $(1 - e^{-t})^2$ is positive and finite on
$(0,\infty)$, so dividing by it gives $m = m'$. Then the two integrals in the hypothesis are equal
and finite, so $a + b\sigma = a' + b'\sigma$ for every $\sigma > 0$; at $\sigma = 1$ and $\sigma = 2$
this gives $b = b'$ and $a = a'$.
\end{proof}

% ADDED Q-0281 wave N (2026-10-07): the third child of lem:filtration-top, restated from
% [V, Lem. 7.1], the step (1) => (3) of its proof from the dilation inequality eq:dilation-decrease
% on ([V] states it inside that proof, not as a numbered statement; checked at the line paper's tag
% v0.1, paper/07-characterization.tex). Its Lean is the trunk's ScaleSpace.exists_profile_of_dilate_le
% (ScaleSpaceCore v0.3.0, AnalysisDirection.lean, moved from Paper V's package by Q-0305; Lean core),
% tagged directly: its Levy pair carries a Gaussian coefficient that this statement does not read,
% and the proof of lem:filtration-top sets it to 0. No adapter was needed (priced at ~60 lines).
\begin{lemma}[a measure that decreases under dilation has a nonincreasing profile]
\label{lem:dilation-decreasing-profile}
\lean{ScaleSpace.exists_profile_of_dilate_le}\leanok
\restates{line:lem:selfdecomposable-exponents}
Let $\nu$ be a measure on $(0,\infty)$ with $\int_{(0,\infty)}(1 \wedge x^2)\,\nu(dx) < \infty$,
and suppose that the image of $\nu$ under $x \mapsto \lambda x$ is at most $\nu$ for every
$\lambda \in (0,1]$. Then $\nu(dx) = k(x)\,x^{-1}dx$ for a nonincreasing
$k : (0,\infty) \to [0,\infty)$ with $\int_0^1 x\,k(x)\,dx < \infty$ and
$\int_1^\infty k(x)\,x^{-1}dx < \infty$.
\statusT\quad\emph{(Restated from \prog{V}{Lem.~7.1}, the step from the dilation inequality to the
profile in its proof of (1) implies (3), and proved in the shared library, citing nothing. Measure
theory only: a change of variables, monotone convergence and continuity of a measure from above and
below. \textbf{What it is for.} \ref{lem:filtration-top} applies it to the \Levy{} measure of the
Bernstein function $\sigma \mapsto F(\sqrt{2\sigma})$, read through
\ref{lem:bernstein-triplet-unique}; it replaces the half-line characterization of self-decomposable
laws that [B]'s proof of Prop.~4.9(4) cites. \textbf{Not claimed}: a right-continuous choice of $k$,
which [V]'s lemma adds and which is not used here.)}
\end{lemma}

\begin{proof}
\leanok
Let $\tilde\nu$ be the image of $\nu$ under $\log$, a measure on $\RR$. Its rays are finite:
$\tilde\nu((\theta,\infty)) = \nu((e^\theta,\infty)) < \infty$, by $\int(1 \wedge x^2)\,\nu(dx) <
\infty$. Dilation by $\lambda = e^{-h}$, $h > 0$, becomes translation by $-h$, so the hypothesis
reads $\tilde\nu(A + h) \le \tilde\nu(A)$ for every Borel $A$ and every $h > 0$: $\tilde\nu$
decreases under every right shift. Put
\[
  f(\theta) := \sup_{n \ge 0}\, 2^n\,\tilde\nu\bigl((\theta,\theta + 2^{-n}]\bigr) \in [0,\infty] .
\]
Each term is nonincreasing in $\theta$, by the shift property, and the sequence is nondecreasing in
$n$, because splitting $(\theta,\theta+2\delta]$ at its midpoint and shifting the right half gives
$\tilde\nu((\theta,\theta+2\delta]) \le 2\,\tilde\nu((\theta,\theta+\delta])$; so $f$ is
nonincreasing. For $a < b$, monotone convergence and Tonelli's theorem give
$\int_a^b f = \lim_n 2^n\int|(a,b] \cap [x - 2^{-n}, x)|\,\tilde\nu(dx)$, which lies between
$\tilde\nu((a + 2^{-n}, b])$ and $\tilde\nu((a, b + 2^{-n}])$; both tend to $\tilde\nu((a,b])$, the
second because the rays are finite. So $\tilde\nu(d\theta) = f(\theta)\,d\theta$. And $f$ is finite
at every point: $f(\theta) = \infty$ would force $f = \infty$ on $(\theta - 1, \theta]$, $f$ being
nonincreasing, and so $\tilde\nu((\theta - 1,\theta]) = \infty$, which the finite rays exclude.
Returning to $x = e^\theta$, $\nu(dx) = k(x)\,x^{-1}dx$ with $k(x) := f(\log x)$, finite,
nonnegative and nonincreasing. The two integrability conditions are
$\int(1 \wedge x^2)\,\nu(dx) = \int_0^1 x\,k(x)\,dx + \int_1^\infty k(x)\,x^{-1}dx < \infty$.
\end{proof}

% CHANGED Q-0281: split from thm:filtration-strict (ADR-0025; records/formalization/PRICING.md,
% finding 3): the last identification of clause (3), A_infty = S, which is [B, Prop. 4.9(4)]. Its
% Lean statement is typed and not scheduled: the inclusion A_infty in S is Schoenberg's radial theorem
% with the representation of a Bernstein function (entry A23 of [B]'s ledger), which neither module states in
% Lean and this module's ledger does not carry. Proving it here would need a new interface, the
% author's call. The parent's text is unchanged and \uses this node.
% CHANGED Q-0281 wave N (2026-10-07): proved, on ledger A10, which the author admitted that day
% (AffineHemigroup.FiltrationTop: classAInf_of_classS, classS_of_classAInf and the bundle
% filtration_top, guarded). The statement is unchanged. The status becomes [A] with A2-A4 and A10,
% the proof of record is written to the Lean route, and three children are split off (ADR-0025):
% lem:lk-negative-definite, lem:bernstein-triplet-unique, lem:dilation-decreasing-profile. Priced at
% ~700-1,050 Lean lines for the node and ~150-200 for the uniqueness; written: 425 (FiltrationTop)
% + 206 + 159 = 790, docstrings included.
\begin{lemma}[the top of the filtration is the subordinated slice]
\label{lem:filtration-top}
\uses{def:filtration-thresholds,def:isotropic-origin-datum,prop:isotropic-polar-reading,lem:delay-polar-profile,lem:lk-negative-definite,lem:bernstein-triplet-unique,lem:dilation-decreasing-profile,thm:cited-A2,thm:cited-A3,thm:cited-A4,thm:cited-A10}
\lean{AffineHemigroup.classAInf_of_classS, AffineHemigroup.classS_of_classAInf,
AffineHemigroup.filtration_top}\leanok
In the notation of \ref{def:filtration-thresholds}, $\mathcal{A}_\infty = \mathcal{S}$.
\statusA\quad\ledger{A2}\quad\ledger{A3}\quad\ledger{A4}\quad\ledger{A10}\quad\emph{(This is
[B, Prop.~4.9(4)], restated, and proved here. \textbf{Assignment.} The inclusion
$\mathcal{A}_\infty \subseteq \mathcal{S}$ spends \ledger{A10} twice, for the function
$\sigma \mapsto F(\sqrt{2\sigma})$ and for each of its residuals, and \ledger{A2}, \ledger{A3}(i)--(ii)
and \ledger{A4} through \ref{prop:isotropic-polar-reading}(1) in its forward direction, which
supplies the \Levy--Khintchine form from which \ref{lem:lk-negative-definite} proves the hypothesis
of \ledger{A10} in every dimension. The inclusion $\mathcal{S} \subseteq \mathcal{A}_\infty$ spends
\ledger{A2} and \ledger{A3}(iii) through the converse direction of the same clause, and reads the
admissibility from [B, Lem.~4.2], restated. Everything else is \ref{lem:lk-negative-definite},
\ref{lem:bernstein-triplet-unique}, \ref{lem:dilation-decreasing-profile} and
\ref{lem:delay-polar-profile}, which cite nothing. \textbf{What the route saves.} [B]'s proof
of Prop.~4.9(4) spends, besides the radial theorem with the representation of a Bernstein function
(its ledger entry A23), the infinite divisibility of the residual of a self-decomposable law, the
complete monotonicity of $e^{-f}$ with Bernstein's theorem on completely monotone functions, and the
characterization of self-decomposable laws on the half-line. None is used here: the residual is
handled by comparing \Levy{} measures, $(c\,\cdot)_*\nu_d \le \nu_d$, which the nonincreasing polar
profile gives at once, and the half-line characterization is replaced by the uniqueness of a
Bernstein triplet and the dilation lemma. Self-decomposability in every dimension is read twice:
once to make $F(\sqrt{2\,\cdot})$ a Bernstein function, and once, through the monotone polar profile
in each dimension, to make each residual one; $\mathcal{A}_1$ alone could not do it, since
$\mathcal{S} \subsetneq \mathcal{A}_1$. Split from clause (3) of \ref{thm:filtration-strict} on
2026-10-04 (Q-0281, ADR-0025); the strictness of the chain, and $\mathcal{A}_\infty \subsetneq
\mathcal{A}_d$, do not read it.)}
\end{lemma}

\begin{proof}
\leanok
\emph{$\mathcal{S} \subseteq \mathcal{A}_\infty$.} Let $F = F_{\mathrm I}(\cdot^2/2)$ with
$F_{\mathrm I}$ causally admissible, of drift $b_0$ and nonincreasing delay profile
$k_{\mathrm I}$, and let $d \ge 1$. $F$ is admissible by \restated{B}{Lem.~4.2}{prop:line-bridge}, restated: \emph{for
causally admissible $F_{\mathrm I}$ with data $(b_0, k_{\mathrm I})$, $F_{\mathrm I}(\omega^2/2)$ is of
the form \ref{eq:admissible-line} with $a_{\mathrm V} = \tfrac12 b_0$.} The polar profile $k_d$
of $(b_0, k_{\mathrm I})$, \ref{eq:polar-profile-d}, is nonincreasing on $(0,\infty)$: for
$0 < r \le s$ and $v > 0$, $k_{\mathrm I}(r^2/2v) \ge k_{\mathrm I}(s^2/2v)$, and both integrands are
integrable by \ref{lem:delay-polar-profile}. So by that lemma $(b_0, k_d)$ is an isotropic origin
datum in dimension $d$ whose exponent is $\xi \mapsto F_{\mathrm I}(\tfrac12|\xi|^2) = F(|\xi|)$,
and by \ref{prop:isotropic-polar-reading}(1), in its converse direction, that is the exponent of a
self-decomposable law on $\RR^d$. So $F \in \mathcal{A}_d$ for every $d \ge 1$.

\emph{$\mathcal{A}_\infty \subseteq \mathcal{S}$.} Let $F \in \mathcal{A}_\infty$. Being admissible,
of the form \ref{eq:admissible-line}, $F$ is nonnegative on $[0,\infty)$ with $F(0) = 0$, and
continuous there by dominated convergence. Put
$f(\sigma) = F(\sqrt{2\sigma})$ for $\sigma \ge 0$, and $g_c(\sigma) = f(\sigma) - f(c^2\sigma)$ for
$0 < c \le 1$.

\emph{Step 1: the hypothesis of \ref{thm:cited-A10}.} Fix $d \ge 1$. By
\ref{prop:isotropic-polar-reading}(1), in its forward direction, there is an isotropic origin datum
$(a_d, k_d)$ in dimension $d$ with
$F(|\eta|) = \tfrac12 a_d|\eta|^2 + \int(1 - \cos\langle\eta,x\rangle)\,\nu_d(dx)$, where $\nu_d$ is
its \Levy{} measure, of density $k_d(|x|)/(|S^{d-1}|\,|x|^d)$. For $0 < c \le 1$ the image
$(c\,\cdot)_*\nu_d$ of $\nu_d$ under $x \mapsto cx$ has density $k_d(|x|/c)/(|S^{d-1}|\,|x|^d)$, the
Jacobian $c^{-d}$ cancelling against $|x/c|^{-d}$; so $(c\,\cdot)_*\nu_d \le \nu_d$, since
$|x|/c \ge |x|$ and $k_d$ is nonincreasing. And
$F(c|\eta|) = \tfrac12 a_dc^2|\eta|^2 + \int(1 - \cos\langle\eta,y\rangle)\,(c\,\cdot)_*\nu_d(dy)$.
So by \ref{lem:lk-negative-definite}, with $(\alpha, L) = (\tfrac12 a_d, \nu_d)$ and
$(\alpha', L') = (0, 0)$, respectively $(\tfrac12 a_dc^2, (c\,\cdot)_*\nu_d)$, the functions
$\eta \mapsto F(|\eta|)$ and $\eta \mapsto F(|\eta|) - F(c|\eta|)$ are negative definite on
$\RR^d$. Since $f(|\xi|^2) = F(|\sqrt2\,\xi|)$ and
$g_c(|\xi|^2) = F(|\sqrt2\,\xi|) - F(c|\sqrt2\,\xi|)$, the functions $\xi \mapsto f(|\xi|^2)$ and
$\xi \mapsto g_c(|\xi|^2)$ are negative definite on $\RR^d$ as well: their sums at points $\xi_j$
are those of the former at the points $\sqrt2\,\xi_j$. They are continuous, $F$ being continuous.
In dimension $0$ both vanish identically, as $f(0) = g_c(0) = 0$. Finally $f \ge 0$, and
$g_c \ge 0$: the sum with one point $\xi_1$, $|\xi_1|^2 = \sigma$, and coefficient $1$ reads
$2g_c(\sigma) - g_c(0) \ge 0$.

\emph{Step 2: the representations.} By \ref{thm:cited-A10} there are $a, b \ge 0$ and a measure $m$
on $(0,\infty)$ with $\int(1 \wedge t)\,m(dt) < \infty$ and
$f(\sigma) = a + b\sigma + \int(1 - e^{-\sigma t})\,m(dt)$ for $\sigma > 0$; and for each $c$ there
are $a_c, b_c \ge 0$ and $m_c$ with $g_c(\sigma) = a_c + b_c\sigma + \int(1 - e^{-\sigma t})\,m_c(dt)$.

\emph{Step 3: $a = 0$.} For $\sigma > 0$, $a \le f(\sigma)$, and $f(\sigma) \to F(0) = 0$ as
$\sigma \downarrow 0$.

\emph{Step 4: $m$ decreases under dilation.} Let $D_\lambda m$ be the image of $m$ under
$t \mapsto \lambda t$, so that
$\int(1 - e^{-c^2\sigma t})\,m(dt) = \int(1 - e^{-\sigma s})\,D_{c^2}m(ds)$. Writing
$f(\sigma) = f(c^2\sigma) + g_c(\sigma)$ with Step 2,
\[
  b\sigma + \int\bigl(1 - e^{-\sigma t}\bigr)\,m(dt)
  = a_c + (bc^2 + b_c)\,\sigma + \int\bigl(1 - e^{-\sigma t}\bigr)\,(D_{c^2}m + m_c)(dt)
  \qquad (\sigma > 0),
\]
both sides finite, and $D_{c^2}m + m_c$ is a measure on $(0,\infty)$. By
\ref{lem:bernstein-triplet-unique}, $m = D_{c^2}m + m_c \ge D_{c^2}m$. As $c$ ranges over $(0,1]$,
so does $c^2$: $D_\lambda m \le m$ for every $\lambda \in (0,1]$.

\emph{Step 5: the profile.} $\int(1 \wedge t^2)\,m(dt) \le \int(1 \wedge t)\,m(dt) < \infty$, so by
\ref{lem:dilation-decreasing-profile} $m(dt) = k(t)\,t^{-1}dt$ with $k \ge 0$ nonincreasing on
$(0,\infty)$.

\emph{Step 6: causal admissibility.} $\int_0^1 k(t)\,dt = \int_{(0,1)} t\,m(dt)$ and
$\int_1^\infty k(t)\,t^{-1}dt = m((1,\infty))$ are both at most $\int(1 \wedge t)\,m(dt) < \infty$,
so $F_{\mathrm I} := (b, k)$ is causally admissible, and for $\sigma > 0$
\[
  F_{\mathrm I}(\sigma) = b\sigma + \int_0^\infty\bigl(1 - e^{-\sigma t}\bigr)\,k(t)\,\frac{dt}{t}
  = f(\sigma) .
\]
For $r > 0$, $F(r) = f(r^2/2) = F_{\mathrm I}(r^2/2)$, and both sides vanish at $r = 0$. So
$F \in \mathcal{S}$.

Every step is machine-checked; the cited interfaces are \ref{thm:cited-A10}, in Step 2, and those of
\ref{prop:isotropic-polar-reading}(1), in Step 1 and in the first inclusion.
\end{proof}

% draft: new material. Provenance: Q-0044, from P-0014's second question, the cell [B]'s Figure 4
% leaves blank: where the Cin rays sit among A_2, A_3, .... [B, Prop. 4.9(3)] puts them outside
% A_infty, by the lattice-step argument; this node puts them outside A_2, which is as far out as
% the chain goes. The postdoc put the answer "no" at p = 0.6 and expected a Schoenberg-type
% computation in d = 2; the proof below is not that. It is a projection: a rotation-invariant Levy
% measure whose one-dimensional shadow is the ray's must live in the ball of radius tau, and a
% BOUNDED profile on that ball -- bounded because it is nonincreasing -- casts a shadow that dies
% at the edge, where the ray's shadow is 1/(2|y|) and does not. Nothing is computed with Bessel
% functions, and nothing is computed numerically.
% CHANGED Q-0281: \uses re-pointed to the split child it reads (eq:polar-decomposition).
\begin{theorem}[no $\Cin$ ray reaches the plane]
\label{thm:no-cin-ray-in-the-plane}
% CHANGED Q-0281 wave 2: proved (AffineHemigroup.NoCinRay, guarded; A2-A4). \uses drops
% lem:isotropic-exponent-log, which only the separate proof of a = 0 read; the Lean route gets a = 0
% from the uniqueness of the triplet and reads it nowhere. Blind review W2A-7: the Assignment names
% both clauses of A3 spent on the line, (iii) and (ii), where it named the uniqueness alone.
\uses{def:isotropic-origin-datum,def:filtration-thresholds,prop:isotropic-polar-reading,lem:polar-decomposition,thm:cited-A2,thm:cited-A3,thm:cited-A4}
\lean{AffineHemigroup.no_cin_ray_in_the_plane, AffineHemigroup.cin_ray_scaled}\leanok
Fix $\tau > 0$ and $b > 0$ and let $F(\omega) = b\,\Cin(\tau|\omega|)$ be the corresponding $\Cin$
ray: the admissible exponent with Gaussian coefficient $0$ and profile
$b\,\mathbf{1}_{(0,\tau)}$, of \restated{B}{Lem.~3.1}{prop:line-extreme-rays}. Then for every integer $d \ge 2$ the radial function
$\xi \mapsto F(|\xi|)$ on $\RR^d$ is \textbf{not} the exponent of a self-decomposable law on
$\RR^d$. That is, $F \notin \mathcal{A}_d$. In particular $\mathcal{A}_2$ contains no $\Cin$ ray.
\statusA\quad\ledger{A2}\quad\ledger{A3}\quad\ledger{A4}\quad\emph{(\textbf{Assignment.}
\ledger{A2}, \ledger{A3} and \ledger{A4} are spent through \ref{prop:isotropic-polar-reading}(1) in
its forward direction, which turns the hypothesis into an isotropic origin datum. \ledger{A3} is
spent once more, directly, on the line: (iii), for a law with the shadow's triplet, and (ii), the
uniqueness of the triplet, which together turn an identity between \emph{functions} into an
identity between \emph{measures} and identify the shadow of the datum's \Levy{} measure with the
ray's. Until Q-0316 the proof also gave the law a triplet by
\ledger{A4} and \ledger{A3}(i) and read it by uniqueness twice, which repeated what the polar reading
supplies (the blind register review of the isotropic paper, Appendix~A flag~17); the declaration is
unchanged. Until Q-0286 the status named
\ledger{A2} and \ledger{A3} alone and the proof cited the polar reading as the
old clause (5) of \ref{thm:origin-boundedness-d}, dropped then as a duplicate of
\ref{prop:isotropic-polar-reading}(1). The first step of the proof had always applied \ledger{A3}(i)
to a law given only as self-decomposable, which needs \ledger{A4}, and the identification of the
law's \Levy{} measure with the datum's was a fourth use of the uniqueness, inside that clause; the
declaration is corrected, not widened. The rest (the projection formula, the support argument and
the estimate at the edge) is Tonelli and one elementary integral. \textbf{Not claimed}: that
$\xi \mapsto \Cin(\tau|\xi|)$ fails to be the exponent of an infinitely divisible law on $\RR^d$.
The theorem assumes self-decomposability and contradicts it; whether the function is negative
definite on $\RR^d$ at all is a Schoenberg question the proof never asks, and P-0014's
expectation that the answer would come from such a computation is not what happened. Nor anything
about the isotropic data whose polar profile is a unit step $b\,\mathbf{1}_{(0,\tau)}$ in $\RR^d$
itself, with exponents $\int_0^\tau(1 - J_0(r\rho))\,dr/r$ in the plane and their
higher-dimensional analogues: those are members of $\mathcal{A}_d$ by
\ref{prop:isotropic-polar-reading}(1) and (2), and this theorem says only that the \emph{line's}
ray is not among them. \textbf{What it is worth.} [B, Prop.~3.3] makes the Gaussian ray and the
$\Cin$ rays the extreme rays of the admissible cone, and [B, Prop.~4.9(3)] excludes the $\Cin$ rays
from $\mathcal{A}_\infty$ by a lattice-step argument that says nothing about any finite dimension.
With this theorem the picture is complete and is a rigidity statement: of the extreme rays of the
line's cone, exactly one extends radially to the plane, the Gaussian, which lies in
$\mathcal{S} = \mathcal{A}_\infty$, being the delay datum with drift and no profile. Every other one
stops at the first step of the filtration. So the plane keeps none of the line's extreme structure
but the Gaussian, while by \ref{thm:filtration-strict} it keeps strictly more than the
subordinated slice $\mathcal{S}$: the two readings sit together, and the cell [B]'s Figure~4
leaves blank is filled.)}
\end{theorem}

\begin{proof}
\leanok
Let $d \ge 2$ and suppose $F \in \mathcal{A}_d$: then $\xi \mapsto F(|\xi|)$ is the exponent of a
self-decomposable law on $\RR^d$.

\emph{The datum.} By \ref{prop:isotropic-polar-reading}(1), in its forward direction,
$\xi \mapsto F(|\xi|)$ is the exponent $F_d$ of an isotropic origin datum $(a,k_d)$ in dimension
$d$. Write $\nu := \nu_d$ for its \Levy{} measure. By \ref{eq:isotropic-exponent},
\begin{equation}\label{eq:cin-ray-exponent}
  F(|\xi|) = \tfrac12 a|\xi|^2
   + \int_{\RR^d}\bigl(1 - \cos\langle\xi,x\rangle\bigr)\,\nu(dx), \qquad \xi \in \RR^d ,
\end{equation}
and $\nu$ is rotation invariant, with the radial density
$f_d(r) = k_d(r)\bigl(|S^{d-1}|r^d\bigr)^{-1}$, where $k_d$ is nonincreasing and finite at every
$r > 0$ (\ref{def:isotropic-origin-datum}).
% CHANGED Q-0281 wave 2: the proof rewritten to the Lean route. The Gaussian coefficient is not shown
% to vanish separately (the uniqueness of the triplet gives a = 0 with the rest, and nothing below
% reads it), and the support and the edge are read from the masses of slabs through
% \ref{eq:polar-decomposition}, not from the density of the shadow.

\emph{The shadow on a line.} Let $\Pi(x) = x_1$ and let $\bar\nu$ be the image of $\nu$ under
$\Pi$, restricted to $\RR \setminus \{0\}$. Since $|x_1| \le |x|$,
$\int(1 \wedge y^2)\,\bar\nu(dy) \le \int(1 \wedge |x|^2)\,\nu(dx) < \infty$, so $\bar\nu$ is a
\Levy{} measure on $\RR$. It is symmetric, because $\nu$ is invariant under $x \mapsto -x$. Take
$\xi = (\omega,0,\dots,0)$ in \ref{eq:cin-ray-exponent}. The integrand vanishes where $x_1 = 0$,
so
\[
  \tfrac12 a\omega^2 + \int_{\RR}\bigl(1 - \cos(\omega y)\bigr)\,\bar\nu(dy) = b\,\Cin(\tau|\omega|),
  \qquad \omega \in \RR .
\]
The right side is the exponent of the datum $(0, b\,\mathbf{1}_{(0,\tau)})$ in dimension $1$, whose
\Levy{} measure is $\nu_1(dy) := b\,\mathbf{1}_{\{0<|y|<\tau\}}\,(2|y|)^{-1}dy$, read through
\ref{eq:polar-decomposition} at $d = 1$: there $S^0 = \{\pm1\}$ and
$\int_{\RR}(1 - \cos\omega y)\,\nu_1(dy) = b\int_0^\tau(1-\cos\omega r)\,dr/r = b\,\Cin(\tau|\omega|)$,
which is \restated{B}{Lem.~3.1}{prop:line-extreme-rays} scaled by $b$. As $\bar\nu$ is symmetric, the exponent of the triplet
$(a,\bar\nu,0)$ reduces to the left side, so by \ledger{A3}(iii) the law with that triplet has
characteristic function $e^{-b\Cin(\tau|\omega|)}$, the same as the law of the datum, whose triplet
is $(0,\nu_1,0)$. By the uniqueness in \ledger{A3}(ii), $a = 0$ and $\bar\nu = \nu_1$.

\emph{The shadow determines the support.} As $\nu_1$ does not charge $\{|y| > \tau\}$,
$\nu(\{x : |x_1| > \tau\}) = \bar\nu(\{|y| > \tau\}) = 0$. A reflection carrying $e_1$ to a unit
vector $u$ preserves $\nu$, so $\nu(\{x : |\langle u,x\rangle| > \tau\}) = 0$ for every unit $u$.
Take a dense sequence $(u_n)$ in $S^{d-1}$. If $|x| > \tau$, some $u_n$ has
$|u_n - x/|x|| < (|x|-\tau)/|x|$, and then $\langle u_n,x\rangle \ge |x| - |u_n - x/|x||\,|x| >
\tau$. So $\{|x| > \tau\}$ is covered by countably many null slabs, and
$\nu(\{|x| > \tau\}) = 0$. By \ref{lem:polar-decomposition} the image of $\nu$ under
$x \mapsto |x|$ is $k_d(r)\,dr/r$ on $(0,\infty)$, so $k_d = 0$ almost everywhere on
$(\tau,\infty)$, and since $k_d$ is nonincreasing and nonnegative, $k_d(r) = 0$ for every
$r > \tau$.

\emph{The edge.} Let $0 < \varepsilon \le \tfrac12$ and consider the slab
$S_\varepsilon = \{x : \tau(1-\varepsilon) < |x_1| < \tau\}$. By \ref{eq:polar-decomposition},
\[
  \nu(S_\varepsilon) = \frac{1}{|S^{d-1}|}\int_{S^{d-1}}\sigma_{d-1}(d\eta)\int_0^\infty
    \mathbf 1_{S_\varepsilon}(r\eta)\,k_d(r)\,\frac{dr}{r} .
\]
A point $r\eta$ of $S_\varepsilon$ with $k_d(r) \ne 0$ has $r \le \tau$, by the previous step, so
$r \ge r|\eta_1| > \tau(1-\varepsilon)$, and $|\eta_1| > \tau(1-\varepsilon)/r \ge 1-\varepsilon$:
$\eta$ lies in the cap $K_\varepsilon = \{\eta : |\eta_1| > 1-\varepsilon\}$ and
$r \in (\tau(1-\varepsilon), \tau]$, an interval of length $\tau\varepsilon$ on which
$k_d(r)/r \le M := 2k_d(\tau/2)/\tau$. Hence
\[
  \nu(S_\varepsilon) \le \frac{\sigma_{d-1}(K_\varepsilon)}{|S^{d-1}|}\,M\,\tau\varepsilon .
\]
On the other side, $\nu(S_\varepsilon) \ge \bar\nu(\{\tau(1-\varepsilon) < |y| < \tau\}) =
\nu_1(\{\tau(1-\varepsilon) < |y| < \tau\}) = b\int_{\tau(1-\varepsilon)}^\tau dr/r \ge
(b/\tau)\,\tau\varepsilon$. Dividing by $\tau\varepsilon$,
$b/\tau \le M\,\sigma_{d-1}(K_\varepsilon)/|S^{d-1}|$ for every $\varepsilon \in (0,\tfrac12]$.
As $\varepsilon \downarrow 0$ the caps decrease to the two poles $\{|\eta_1| = 1\} = \{\pm e_1\}$,
which lie on the great subsphere $\{\eta_2 = 0\}$ because $d \ge 2$, and which are therefore
$\sigma_{d-1}$-null (a great subsphere is null, the cone over it lying in a hyperplane). So $\sigma_{d-1}(K_\varepsilon) \to 0$, and $b/\tau \le 0$, a contradiction.
Hence $F \notin \mathcal{A}_d$.
\end{proof}

% not in the draft: queue item Q-0425 (hub), decision D2 of the response to the external review of
% the isotropic paper, round 1 (records/affine/PLAN-review-response-paper-iso.md): the dictionary
% between A_d and the radial characteristic functions, with what follows from it. Four nodes:
%   lem:shadow-polar-profile          the polar profile of the shadow in R^d of a datum of R^{d+1}:
%                                     a formula, and it is continuous; cites nothing;
%   prop:radial-dictionary            the dictionary, for a finite profile at the origin;
%   cor:dictionary-filtration         the continuity criterion, strictness at every step with the
%                                     step datum as witness, and the Cin rays;
%   prop:separating-family-mixtures   the separating family: Gaussian variance mixtures in every
%                                     dimension, self-decomposable up to the critical one.
% thm:filtration-strict and thm:no-cin-ray-in-the-plane are not touched: their statements, proofs
% and Lean stand, and the corollary's annotation says which of their clauses it reproves. The
% derivations are the verification records' (records/affine/reviews/VERIFY-round1-prior-art.md
% section 4, VERIFY-round1-filtration-map-blind-check.md, VERIFY-round1-mathematics.md item 9),
% checked line by line; the departures are in records/affine/NODES-q0425.md. No new cited fact:
% Gneiting's and Schoenberg's theorems and the representation of a self-decomposable law by its
% driving Levy process are named as attributions and are not used. No \lean tag on any of the
% four.
%
% The form of this lemma, decided when the statement was written. It is stated on DATA and
% measures, not on classes or laws, so it cites nothing: the uniqueness of the datum (A3(ii),
% through prop:isotropic-polar-reading(1)) enters only in cor:dictionary-filtration(1). The
% normalisation C_d int (1+u^2)^{-(d+1)/2} du = 1 is proved from polar coordinates in the two
% dimensions, with no Beta integral.
% Two cautions from the attribution item Q-0426 (its record, section 2.3), for whoever cites the
% printed forms of this formula. (a) Fang-Kotz-Ng's remark 1 on p. 37 ("the marginal densities of
% dimensions <= n-1 are continuous") is printed without a hypothesis and fails for the
% (n-1)-marginal of a general density generator; continuity here comes from the monotone profile,
% and that remark is not to be cited for it. (b) Cambanis-Huang-Simons print "(s_1 = 1)" on
% p. 378, which is inconsistent with their own (20)-(22) in codimension one; the constant is to be
% taken from C_d = |S^{d-1}|/|S^d|, not from their (22).
% CHANGED Q-0425 (the blind check; records/formalization/blind-check-q0425-report-B.md, verdict
% "correct as stated" for the four nodes, five optional precision edits and two cautions on the
% attributions; all applied). This node: clause (3) and its proof say that the restriction to
% R^d minus the origin removes nothing (edit 5); the attribution says that Beta(d/2, 1/2) is the
% marginal one dimension down, the general parameter being (m/2, (n-m)/2) (second caution).
% prop:radial-dictionary: the integrability step of the existence part names the a.e. identity it
% uses (edit 2); the attribution says that clauses (1) and (2) are the compound Poisson case of
% the representation, and which laws are outside (first caution). cor:dictionary-filtration: the
% statement of clause (1) says that a datum in dimension d+1 exists and its last sentence has the
% logarithmic moment (edits 1 and 4), and the proofs of clauses (2) and (3) name
% prop:isotropic-polar-classification(2) for the nesting (edit 1).
% prop:separating-family-mixtures: the integrability of nu' is by polar coordinates for a
% measurable density, not by eq:polar-decomposition (edit 3). No clause changed in content.
% CHANGED Q-0429 (2026-10-12; the author's decision of that day: the shadow lemma and the
% consequences on data are formalized, the closed-form constant, the dictionary and the family as
% laws stay proved in print). The four nodes are six, so that what will carry a \leanok tag is
% exactly what is proved (ADR-0025):
%   lem:shadow-polar-profile                unchanged (statement and proof); to be formalized;
%   prop:profile-one-dimension-down         NEW label: the three clauses on data of the old
%                                           cor:dictionary-filtration; to be formalized;
%   prop:radial-dictionary                  unchanged; print-only;
%   cor:dictionary-filtration               the label kept, for the readings in the dictionary
%                                           that the old node carried in three sentences;
%                                           print-only;
%   thm:separating-family-critical-values   NEW: the first sentence of the old
%                                           prop:separating-family-mixtures(3) and its
%                                           consequences, with the strict decrease and the limit
%                                           of the critical values; to be formalized;
%   prop:separating-family-mixtures         clauses (1) and (2) unchanged, clause (3) reduced to
%                                           the reading in laws; print-only.
% The proofs are cut, not rewritten; the sentences that had to change are listed in
% records/affine/NODES-q0425.md, section "After the split (Q-0429)". This node: annotation only
% (two references re-pointed to prop:profile-one-dimension-down).
% CORRECTED after the blind check of Q-0429 (records/formalization/blind-check-q0429-report.md,
% finding F-B1). "Unchanged" and "word for word", in the list above and in the comments of the
% five other nodes, compare with the text as it stood before the split. That text already had the
% repairs of the blind check of Q-0425 (the CHANGED Q-0425 paragraph above), made after that check
% had read the nodes. So against the text that check refereed (blind-check-q0425-input.md) these
% places differ besides the sentences written for the split: this lemma, clause (3) of the
% statement (the sentence "The restriction removes nothing: ..." and "The exponent of (a, k-bar)"
% for "Its exponent") and three sentences of the proof of clause (3) (the null line);
% prop:radial-dictionary, the proof of clause (2), existence (two insertions: "and c G = k_d
% almost everywhere" and the integral of k_d(r+)); cor:dictionary-filtration(1), statement and
% proof ("with E log^+|X| < infty"); prop:profile-one-dimension-down, the proof of clause (2),
% the sentence on the inclusion; thm:separating-family-critical-values, "w > w_D excludes it",
% the sentence with the display of the integral of |x|^2 ^ 1. The blind check of Q-0429 read
% each of them and found each correct (its sections 2.1 and 2.3).
% CHANGED Q-0429, part 2: the three nodes marked "to be formalized" in the list above are typed
% and proved. This node: Formalization/AffineHemigroup/ShadowProfile.lean, Lean core, all three
% clauses (no clause was split off). The statement is unchanged. Added: the \lean and \leanok
% tags, the passage "Formalization" of the annotation, and the closing sentences of the proof,
% which said that nothing was machine-checked.
\begin{lemma}[the polar profile of a shadow]
\label{lem:shadow-polar-profile}
\uses{def:isotropic-origin-datum,lem:polar-decomposition}
\lean{AffineHemigroup.shadow_profile_mean, AffineHemigroup.shadow_profile_regular,
AffineHemigroup.shadow_datum, AffineHemigroup.lintegral_shadow,
AffineHemigroup.OriginDatum.shadow, AffineHemigroup.shadowProfile}\leanok
Let $d \ge 1$, let $(a, k')$ be an isotropic origin datum in dimension $d + 1$, with \Levy{}
measure $\nu'$ and exponent $F'$, and let $\Pi(x,y) := x$ be the projection of
$\RR^{d+1} = \RR^d \times \RR$ on the first $d$ coordinates. Put
\[
  C_d := \frac{|S^{d-1}|}{|S^{d}|}, \qquad
  \bar k(r) := C_d\int_{\RR} k'\bigl(r\sqrt{1+u^2}\bigr)\,(1+u^2)^{-(d+1)/2}\,du \quad (r > 0) .
\]
\begin{enumerate}
\item $C_d\int_{\RR}(1+u^2)^{-(d+1)/2}\,du = 1$. So $\bar k(r)$ is a mean of the values of $k'$
on $[r,\infty)$, and $0 \le \bar k(r) \le k'(r)$.
\item $\bar k$ is nonincreasing and continuous on $(0,\infty)$, and $\bar k(0{+}) = k'(0{+})$ in
$[0,\infty]$.
\item $(a, \bar k)$ is an isotropic origin datum in dimension $d$. Its \Levy{} measure $\bar\nu$
is the image of $\nu'$ under $\Pi$, restricted to $\RR^d \setminus \{0\}$: for every nonnegative
Borel $g$ on $\RR^d$ vanishing at the origin, $\int g(\Pi z)\,\nu'(dz) = \int g\,d\bar\nu$. The
restriction removes nothing: the image of $\nu'$ gives the origin of $\RR^d$ the mass
$\nu'(\{z \ne 0 : \Pi z = 0\}) = 0$. The exponent of $(a,\bar k)$ is
$\xi \mapsto F'\bigl((\xi,0)\bigr)$, $\xi \in \RR^d$.
\end{enumerate}
In particular, if $F'(\zeta) = F(|\zeta|)$ for a function $F$ on $[0,\infty)$, then $(a,\bar k)$
is an isotropic origin datum in dimension $d$ with exponent $\xi \mapsto F(|\xi|)$.
\statusT\quad\emph{(Cites nothing: Tonelli's theorem on $\RR^d \times \RR$, one substitution,
polar coordinates in the two dimensions, and dominated and monotone convergence. The statement is
about data, measures and functions; no law is named. \textbf{What it is.} The shadow of a
rotation-invariant \Levy{} measure one dimension down, with its polar profile in closed form. It
is the general form of the estimate at the edge in the proof of
\ref{thm:no-cin-ray-in-the-plane}, which projects from $\RR^d$ to the line, and of the short
projection argument of the external referee of the isotropic paper (round 1), verified in
\texttt{records/affine/reviews/VERIFY-round1-mathematics.md}, item~9, where this formula was
found. \textbf{What it is worth.} First, the nesting of the classes at the level of data: a datum
in dimension $d + 1$ gives a datum in dimension $d$ with the same radial exponent, the same
Gaussian coefficient and the same profile at the origin, as
\ref{cor:origin-profile-exponent}(3) requires. Second, the profile of a shadow is continuous,
whatever the jumps of $k'$: this is what \ref{prop:profile-one-dimension-down}(1) turns into a
criterion. Third, it computes. The shadow of the step $k' = b\,\mathbf 1_{(0,\tau)}$ of the plane
on the line is $\bar k(r) = (2b/\pi)\arccos(r/\tau)$ on $(0,\tau)$, which falls continuously from
$b$ to $0$. Applied $n$ times it gives the shadow in dimension $d$ of a datum in dimension
$d + n$. \textbf{Checks}
(\texttt{records/affine/reviews/verify-round1-scripts/item09\_projection\_of\_polar\_profile.py},
rerun for this node): the normalisation of clause (1) for $d = 1,\dots,8$; that the shadow of the
\Matern{} datum of dimension $d + 1$ is the \Matern{} datum of dimension $d$, for
$d = 1,\dots,4$; the shadow of a step. \textbf{Attribution (the anchors read at the page image
for the attribution item Q-0426; its record
\texttt{records/affine/PRIOR-ART-review-round1-additions.md}, section~2, and the reading note
\texttt{READING-cambanis1981theory-marginals.md}).} The formula is in print for
probability laws. The radius of the projection to $\RR^d$ of a rotation-invariant law of
$\RR^{d+1}$ is $R\sqrt{B}$, with $B$ of the $\mathrm{Beta}(d/2, 1/2)$ law and independent of the
radius $R$: \sscite{cambanis1981theory}, Lemma~2, p.~371, and Cor.~2 with (8), p.~372. This is
their statement for the marginal one dimension down, $m = d$ and $n = d + 1$. For the marginal
on $\RR^m$ of a rotation-invariant law of $\RR^n$ they have $B$ of the
$\mathrm{Beta}(m/2, (n-m)/2)$ law, so $(d/2, 1/2)$ is not the parameter of a shadow more than
one dimension down. Applied to
the radial measure $k'(s)\,ds/s$ it gives $\bar k(r) = \mathbb{E}\,k'(r/\sqrt{B})$, which is the
integral above with $B = 1/(1+u^2)$; the two forms were compared numerically for that record, for
$d = 1, 2, 3$. In the form of a density it is \sscite{fang1990symmetric}, Thm.~2.10 with (2.23),
p.~36. That the \Levy{} measure of a projection is the projection of the \Levy{} measure is
\sscite{sato1999levy}~Prop.~11.10, p.~65. None of these is used: the lemma is proved from
Tonelli's theorem, for a measure that need not be finite, and the constant is this chapter's
$C_d$. The continuity of clause (2) is not taken from those sources. It is proved from the
monotonicity of $k'$, without which it can fail. \textbf{Not claimed}: which profiles in dimension $d$ are
shadows, that is an inverse of the formula; anything about laws or about the classes
$\mathcal{A}_d$, for which the uniqueness of the datum is needed and which is
\ref{prop:profile-one-dimension-down}(1). \textbf{Formalization} (Q-0429). $d$ is a natural
number with $d \ge 1$. The projection $\Pi$ is the map \texttt{dropLast} of
\ref{thm:filtration-strict}'s Lean, the first $d$ coordinates of a point of $\RR^{d+1}$, and
$(\xi,0)$ is its \texttt{padZero}. $C_d$ is the quotient of the closed forms
$|S^{n-1}| = 2\pi^{n/2}/\Gamma(n/2)$, which \ref{def:isotropic-origin-datum}'s Lean identifies
with the total mass of the surface measure. $\bar k$ is the total function
\texttt{shadowProfile}, a Bochner integral over the line with a real power in the weight. It is
read at $r > 0$ only. Clause (1) is the equality of real numbers
$C_d\int_{\RR}(1+u^2)^{-(d+1)/2}du = 1$ and, for every datum and every $r > 0$, the two
inequalities $0 \le \bar k(r) \le k'(r)$. That $\bar k(r)$ is a mean is the formula with this
normalisation and has no declaration of its own. Clause (2) is monotonicity and continuity on
the set $(0,\infty)$, and the limit $\bar k(r) \to k'(0{+})$ as $r \downarrow 0$ in
$[0,\infty]$, where $k'(0{+})$ is the supremum that \ref{def:isotropic-origin-datum}'s Lean
proves to be the right limit of $k'$ at $0$. Clause (3) constructs the datum,
\texttt{OriginDatum.shadow}, with Gaussian coefficient $a$ and profile $\bar k$. What it says of
the \Levy{} measure is stated in three forms: the identity of integrals for every nonnegative
Borel $g$ on $\RR^d$, without the condition $g(0) = 0$; the equality of the image of $\nu'$
under $\Pi$ with $\bar\nu$ as measures on $\RR^d$, without a restriction; and
$\nu'(\{z : \Pi z = 0\}) = 0$. In Lean the \Levy{} measure of a datum is a measure on $\RR^d$
whose density vanishes at the origin, so the three together are the clause. The exponent is the
real-valued one, at every $\xi \in \RR^d$, and the last sentence of the statement is an
implication for every function $F$ on the line. The identity of the proof's first step is a
declaration as well, for every nonnegative Borel $h$ on the line. The properties of $\bar k$ are
proved for every function $k'$ that is nonnegative and nonincreasing on $(0,\infty)$, not only
for the profile of a datum. Witness: the step datum $(0, b\,\mathbf 1_{(0,\tau)})$ in dimension
$d + 1$, whose profile at the origin is $b > 0$. Probes: at $d = 1$ the normalisation reads
$\int_{\RR}(1+u^2)^{-1}du = \pi$, and a second probe proves the same equality from Mathlib's
value of that integral, without the lemma (the two probes have one type, and what is independent
is their proofs); the shadow of the step datum
has a continuous profile with the value $b$ at the origin, where the step is discontinuous; and
the lemma applied twice. Not in Lean: the closed form $(2b/\pi)\arccos(r/\tau)$ of the shadow of
the plane's step, and the comparison with the printed forms of the formula.)}
\end{lemma}

\begin{proof}
\leanok
\emph{The integral over the last coordinate.} Let $h : (0,\infty) \to [0,\infty)$ be Borel and
put $\bar h(r) := C_d\int_{\RR}h(r\sqrt{1+u^2})\,(1+u^2)^{-(d+1)/2}\,du \in [0,\infty]$ for
$r > 0$. Let $g \ge 0$ be Borel on $\RR^d$ with $g(0) = 0$, and write $z = (x,y)$ with
$x \in \RR^d$ and $y \in \RR$. By Tonelli's theorem,
\[
  \int_{\RR^{d+1}} g(x)\,\frac{h(|z|)}{|S^{d}|\,|z|^{d+1}}\,dz
   = \int_{\RR^d} g(x)\Bigl[\int_{\RR}
       \frac{h\bigl(\sqrt{|x|^2+y^2}\bigr)}{|S^{d}|\,(|x|^2+y^2)^{(d+1)/2}}\,dy\Bigr]dx .
\]
For $x \ne 0$ put $r = |x|$ and substitute $y = ru$. The inner integral becomes
$(|S^{d}|\,r^{d})^{-1}\int_{\RR}h(r\sqrt{1+u^2})\,(1+u^2)^{-(d+1)/2}\,du =
\bar h(r)\,(|S^{d-1}|\,r^{d})^{-1}$. As $g(0) = 0$,
\begin{equation}\label{eq:shadow-identity}
  \int_{\RR^{d+1}} g(\Pi z)\,\frac{h(|z|)}{|S^{d}|\,|z|^{d+1}}\,dz
   = \int_{\RR^d} g(x)\,\frac{\bar h(|x|)}{|S^{d-1}|\,|x|^{d}}\,dx ,
\end{equation}
both sides in $[0,\infty]$.

\emph{Clause (1).} Take $h = \mathbf 1_{(1,2)}$ and $g = 1$ off the origin in
\ref{eq:shadow-identity}. In polar coordinates of $\RR^{d+1}$ the left side is
$\int_1^2 dt/t = \log 2$; the set $\{\Pi z = 0\}$ is a line, which is null because $d + 1 \ge 2$.
In polar coordinates of $\RR^d$, and by Tonelli's theorem, the right side is
\[
  \int_0^\infty \bar h(r)\,\frac{dr}{r}
   = C_d\int_{\RR}(1+u^2)^{-(d+1)/2}\Bigl[\int_0^\infty
       \mathbf 1_{(1,2)}\bigl(r\sqrt{1+u^2}\bigr)\frac{dr}{r}\Bigr]du
   = C_d\int_{\RR}(1+u^2)^{-(d+1)/2}\,du\cdot\log 2 ,
\]
the inner integral being over an interval whose endpoints have the ratio $2$. So
$C_d\int_{\RR}(1+u^2)^{-(d+1)/2}du = 1$. As $k'$ is nonincreasing and nonnegative,
$0 \le k'(r\sqrt{1+u^2}) \le k'(r)$, so $0 \le \bar k(r) \le k'(r) < \infty$.

\emph{Clause (2).} For each $u$, $r \mapsto k'(r\sqrt{1+u^2})$ is nonincreasing, so $\bar k$ is.
Let $r_0 > 0$. A nonincreasing function has at most countably many discontinuities, and for each
of them, $t$, at most two real $u$ have $r_0\sqrt{1+u^2} = t$. So for all $u$ outside a countable
set, $k'(r\sqrt{1+u^2}) \to k'(r_0\sqrt{1+u^2})$ as $r \to r_0$. For $r \ge r_0/2$ the integrand
of $\bar k(r)$ is at most $k'(r_0/2)\,(1+u^2)^{-(d+1)/2}$, which is integrable. By dominated
convergence $\bar k(r) \to \bar k(r_0)$. As $r \downarrow 0$, $k'(r\sqrt{1+u^2})$ increases to
$k'(0{+})$ for every $u$, so by monotone convergence and clause (1),
$\bar k(r) \to k'(0{+})$, in $[0,\infty]$.

\emph{Clause (3).} By \ref{def:isotropic-origin-datum}, $\nu'$ has the density
$k'(|z|)\,(|S^{d}|\,|z|^{d+1})^{-1}$ on $\RR^{d+1} \setminus \{0\}$. Let $\bar\nu$ be the measure
on $\RR^d \setminus \{0\}$ with density $\bar k(|x|)\,(|S^{d-1}|\,|x|^{d})^{-1}$. Then
\ref{eq:shadow-identity} with $h = k'$ reads $\int g(\Pi z)\,\nu'(dz) = \int g\,d\bar\nu$ for
every nonnegative Borel $g$ vanishing at the origin. The set $\{z \ne 0 : \Pi z = 0\}$ lies in a
line, which is Lebesgue-null in $\RR^{d+1}$ because $d + 1 \ge 2$, and $\nu'$ has a density. So
$\nu'$ gives it the mass $0$, and the image of $\nu'$ under $\Pi$ gives the origin no mass. Take
$g(x) = |x|^2 \wedge 1$. As
$|\Pi z| \le |z|$,
\[
  \int_0^1 \bar k(r)\,r\,dr + \int_1^\infty \bar k(r)\,\frac{dr}{r}
   = \int_{\RR^d}\bigl(|x|^2 \wedge 1\bigr)\,\bar\nu(dx)
   \le \int_{\RR^{d+1}}\bigl(|z|^2 \wedge 1\bigr)\,\nu'(dz)
   = \int_0^1 k'(r)\,r\,dr + \int_1^\infty k'(r)\,\frac{dr}{r} < \infty ,
\]
the first equality by polar coordinates in $\RR^d$ and the last by
\ref{eq:polar-decomposition} in dimension $d + 1$. With clauses (1) and (2), $\bar k$ is finite,
nonnegative and nonincreasing and satisfies the two conditions of
\ref{def:isotropic-origin-datum}. So $(a,\bar k)$ is an isotropic origin datum in dimension $d$,
and its \Levy{} measure is $\bar\nu$. For $\xi \in \RR^d$ take $g(x) = 1 - \cos\langle\xi,x\rangle$.
Since $\langle(\xi,0),z\rangle = \langle\xi,\Pi z\rangle$ and $|(\xi,0)| = |\xi|$,
\[
  F'\bigl((\xi,0)\bigr) = \tfrac12 a|\xi|^2
     + \int_{\RR^{d+1}}\bigl(1 - \cos\langle\xi,\Pi z\rangle\bigr)\,\nu'(dz)
   = \tfrac12 a|\xi|^2 + \int_{\RR^d}\bigl(1 - \cos\langle\xi,x\rangle\bigr)\,\bar\nu(dx) ,
\]
which is the exponent \ref{eq:isotropic-exponent} of $(a,\bar k)$ at $\xi$. If
$F'(\zeta) = F(|\zeta|)$, the left side is $F(|\xi|)$.

No step cites a fact of the ledger, and every step is machine-checked. The Lean differs from
this proof in four places, none of which changes the route. It writes the two densities as
$h(|z|)$ times $(|S^{d}|\,|z|^{d+1})^{-1}$, with the convention $1/0 = 0$, and treats $x = 0$ as
a null set of $\RR^d$, so that \ref{eq:shadow-identity} holds for every nonnegative Borel $g$,
without $g(0) = 0$. In clause (1) it takes $g = 1$ and does not compute $\log 2$: by the
invariance of $dr/r$ under dilations both sides are multiples of $I = \int_1^2 dt/t$, and
$0 < I < \infty$. In clause (3) it obtains the two integrability conditions of the shadow from
$\bar k \le k'$, without the test function $|x|^2 \wedge 1$, and the mass
$\nu'(\{z : \Pi z = 0\}) = 0$ from the identity of the images and $\bar\nu(\{0\}) = 0$. The limit
at the origin is taken along $r = 1/(n+1)$.
\end{proof}

% not in the draft: Q-0425 and Q-0429. The three clauses on data of the corollary that Q-0425
% wrote as cor:dictionary-filtration, in a node of their own since Q-0429 (2026-10-12), because
% they will be formalized and the dictionary will not (ADR-0025: a node's Lean covers all of it,
% or the node is split). The statement names no law, no characteristic function and not the
% function Omega_d: it needs lem:shadow-polar-profile, the polar form of the classes and the
% uniqueness of a datum, and nothing of prop:radial-dictionary, which now follows it. The readings
% in the dictionary are cor:dictionary-filtration.
% What the cut changed, against the text the blind check of Q-0425 read (report B, "correct as
% stated"). Statement: clause (1) lost its last sentence (the reading Pr(|X| = r) = 0); clause (2)
% defines G as the radial exponent of the datum (0, b 1_(0,tau)), where it was defined by the
% formula b int_0^r (1 - Omega_d(tau u)) du/u and then identified with that exponent, and it lost
% the two readings in the dictionary; clause (3) is word for word. Proof: clause (1) is word for
% word without its last sentence; in clause (2) the membership, which was read from
% prop:radial-dictionary(1), is proved on the datum (four sentences, new), and the exclusion and
% the inclusion are word for word; in clause (3) the first sentence, which read Omega_1 = cos, is
% the computation at d = 1 from eq:polar-decomposition (the one in the proof of
% thm:no-cin-ray-in-the-plane), and the rest is word for word.
% CORRECTED after the blind check of Q-0429 (its report, F-B1): "word for word" is against the
% text before the split. Against the text the blind check of Q-0425 read, the sentence on the
% inclusion in the proof of clause (2) differs: it read "and the inclusion A_{d+1} in A_d is part
% of clause (1)" and since that check's repairs it says why, with
% prop:isotropic-polar-classification(2); and in clause (3) "with the datum that
% prop:isotropic-polar-classification(2) supplies" is from the same repairs. The blind check of
% Q-0429 read both and found them correct. After it the wording of the first was made exact
% (F-B5: the exponent is zeta -> F(|zeta|), not F), and "finite" was added to "nonnegative and
% radial" in the same proof (F-B4). F-B7 of the same check: the closing paragraph of the proof
% says that the node's Lean reads thm:no-cin-ray-in-the-plane's Lean (the datum on the line and
% its exponent), and that theorem is not in \uses. It stays out: \uses is the printed proof's,
% which names the theorem and does not use it as a step; the closing paragraph now says so.
% No \lean tag yet: the Lean is the second part of Q-0429. Priced at the target types
% (records/affine/NODES-q0425.md, "After the split"): about 270 lines after the shadow lemma, on
% A2, A3, A4. The library has classA_iff_originDatum, polar_reading_unique, exists_ae_eq_mem_Ioo,
% OriginDatum.exponent_eq_of_norm_eq, cinDatum with cinDatum_exponent, and classA_of_succ.
% CHANGED Q-0429, part 2: typed and proved
% (Formalization/AffineHemigroup/ProfileOneDimensionDown.lean, on A2, A3, A4). The statement is
% unchanged. Added: the \lean and \leanok tags, the passage "Formalization" of the annotation,
% and the closing sentences of the proof; "to be formalized" in the annotation is "formalized".
% CHANGED Q-0430 (2026-10-12): annotation only, one see-also in Attribution (Daley and Porcu,
% Prop. 11, p. 1821; records/affine/READING-daley2014dimension.md).
\begin{proposition}[the profile one dimension down is continuous; the step datum separates at every step; the $\Cin$ rays]
\label{prop:profile-one-dimension-down}
\uses{lem:shadow-polar-profile,def:isotropic-origin-datum,def:filtration-thresholds,lem:polar-decomposition,prop:isotropic-polar-reading,prop:isotropic-polar-classification,thm:cited-A2,thm:cited-A3,thm:cited-A4}
\lean{AffineHemigroup.profile_one_dimension_down,
AffineHemigroup.profile_one_dimension_down_of_classA,
AffineHemigroup.not_classA_succ_of_profile_discontinuous, AffineHemigroup.step_datum_separates,
AffineHemigroup.cin_ray_not_classA, AffineHemigroup.stepOriginDatum}\leanok
Let $d \ge 1$.
\begin{enumerate}
\item \emph{(The profile one dimension down is continuous.)} Let $F \in \mathcal{A}_{d+1}$, and
let $(a',k')$ be an isotropic origin datum in dimension $d + 1$ with exponent
$\zeta \mapsto F(|\zeta|)$; there is one, by \ref{prop:isotropic-polar-classification}(2). Then
$F \in \mathcal{A}_d$, and every isotropic origin datum
$(a,k_d)$ in dimension $d$ with exponent $\xi \mapsto F(|\xi|)$ has $a = a'$ and
\[
  k_d(r) = \frac{|S^{d-1}|}{|S^{d}|}\int_{\RR} k'\bigl(r\sqrt{1+u^2}\bigr)\,(1+u^2)^{-(d+1)/2}\,du
  \qquad\text{for every } r > 0 .
\]
In particular $k_d$ is continuous on $(0,\infty)$. So if $F \in \mathcal{A}_d$ and some isotropic
origin datum in dimension $d$ with exponent $\xi \mapsto F(|\xi|)$ has a polar profile with a
discontinuity, then $F \notin \mathcal{A}_{d+1}$.
\item \emph{(Strictness at every step, with a witness.)} Let $b > 0$ and $\tau > 0$. Then
$(0, b\,\mathbf 1_{(0,\tau)})$ is an isotropic origin datum in dimension $d$, and its exponent is
radial: it is $\xi \mapsto G(|\xi|)$ for a function $G$ on $[0,\infty)$. This $G$ belongs to
$\mathcal{A}_d \setminus \mathcal{A}_{d+1}$. Hence $\mathcal{A}_{d+1} \subsetneq \mathcal{A}_d$
for every $d \ge 1$.
\item \emph{(The $\Cin$ rays.)} At $d = 1$, $G(r) = b\,\Cin(\tau r)$. So no $\Cin$ ray belongs
to $\mathcal{A}_2$, nor to $\mathcal{A}_D$ for any $D \ge 2$.
\end{enumerate}
\statusA\quad\ledger{A2}\quad\ledger{A3}\quad\ledger{A4}\quad\emph{(\textbf{Assignment.} The node
cites nothing of its own. \ledger{A2}, \ledger{A3} and \ledger{A4} are spent through
\ref{prop:isotropic-polar-classification}(2), which places a function in a class or takes it
from one, and through the uniqueness of the datum in \ref{prop:isotropic-polar-reading}(1), which
is \ledger{A3}(ii). The formula of clause (1) is
\ref{lem:shadow-polar-profile}, which cites nothing. \textbf{What the citations carry and what
they do not.} They carry the polar form of $\mathcal{A}_d$ and the uniqueness of a datum. The
continuity of a shadow's profile, on which the three clauses turn, is proved.
\textbf{What it is worth.} Clause (1) is a test that needs one dimension only: a member of
$\mathcal{A}_d$ whose polar profile there has a jump goes no further. It also contains the
nesting $\mathcal{A}_{d+1} \subseteq \mathcal{A}_d$, with the profile of the lower dimension
computed from that of the higher. Clause (2) names a witness at every step whose profile is the
simplest one, a step. The three clauses are stated on the classes, on data and on their
profiles: neither the dictionary nor a random vector enters, and a reader decides on the datum
alone. \textbf{How it stands to the two theorems of this section, which are
unchanged.} Clause (2) reproves the part of \ref{thm:filtration-strict}(3) that follows from
strictness alone: every inclusion $\mathcal{A}_{d+1} \subsetneq \mathcal{A}_d$ is strict, hence
$\mathcal{A}_\infty \subsetneq \mathcal{A}_d$ for every $d$ and no $\mathcal{A}_D$ equals
$\mathcal{A}_\infty$. It does so with another witness. It gives neither clause (1) of that
theorem, the window $(w_{d+1}, w_d\,]$ of the one-parameter family $F^{(w)}$, nor its clause
(2), nor the identification $\mathcal{A}_\infty = \mathcal{S}$, which is
\ref{lem:filtration-top}. What that theorem adds to this node is the family: exponents with
smooth profiles, of the subordinated form $F_{\mathrm I}(r^2/2)$, that separate consecutive
classes at explicit critical values, with \ref{lem:threshold-monotone},
\ref{lem:threshold-limit}, \ref{lem:threshold-asymptotic},
\ref{thm:separating-family-critical-values} and
\ref{prop:separating-family-mixtures} about them. Clause (3) is the whole statement of
\ref{thm:no-cin-ray-in-the-plane}. That theorem's proof is the same projection carried out from
$\RR^d$ to the line in one step, without the nesting, and it is machine-checked. It also answers
what that theorem's annotation leaves aside, where the data with a step profile in dimension
$d$ sit: in $\mathcal{A}_d$ and not in $\mathcal{A}_{d+1}$. \textbf{The split.} Until Q-0429
these clauses stood in \ref{cor:dictionary-filtration} with their readings in the dictionary.
They are a node of their own because they are formalized and the dictionary is not. In
that node the exponent of clause (2) is written out,
$G(r) = b\int_0^r(1 - \Omega_d(\tau u))\,du/u$ with $\Omega_d$ the mean of the cosine over the
sphere. Here $G$ is the exponent of the datum, which is all the clause reads. \textbf{Readings.}
In clause (1) the equality of profiles holds at every $r > 0$, for every datum in dimension $d$
with that exponent: a nonincreasing function equal almost everywhere to a continuous one equals
it everywhere. The displayed integral depends on $k'$ only up to a Lebesgue-null set, so it is
the same for every datum in dimension $d + 1$. In clause (2), $G$ is determined on $[0,\infty)$
by the datum. \textbf{Attribution (the anchors read
at the page image for the attribution item Q-0426; its record
\texttt{records/affine/PRIOR-ART-review-round1-additions.md}, sections 1 and 2).} The
formula of clause (1) is in print for probability laws, with the sources given in the annotation
of \ref{lem:shadow-polar-profile}; its continuity is proved there. For radial characteristic
functions see also \sscite{daley2014dimension}, Prop.~11 with (3.8), p.~1821: the radius law of
such a function in a lower dimension is the law of its radius times an independent variable
with values in $(0,1)$ (read at the page image for Q-0430;
\texttt{records/affine/CITATIONS-primary-sources-B.md}, item 4). The projection argument for
the step datum is the external referee's of the isotropic paper (round 1). The $\Cin$ rays of
clause (3) are the elementary exponents of
the representation of a self-decomposable exponent that \sscite{jurek1983integral}, Thm.~7.2,
p.~259, calls the Urbanik representation: one jump size in one direction. None of these is used.
\textbf{Not claimed}: the converse of clause (1). A continuous profile in dimension $d$ does not
put $F$ in $\mathcal{A}_{d+1}$: for $w \in (w_{d+1}, w_d\,]$ the exponent $F^{(w)}$ of
\ref{thm:filtration-strict}(1) has a smooth profile in dimension $d$
(\ref{lem:threshold-criterion}) and is not in $\mathcal{A}_{d+1}$. Nor that
$\zeta \mapsto G(|\zeta|)$ fails to be the exponent of an infinitely divisible law on
$\RR^{d+1}$: the clause excludes a self-decomposable one. Nor anything about the shape of
$\mathcal{A}_d \setminus \mathcal{A}_{d+1}$ beyond these members. \textbf{Formalization}
(Q-0429). $d$ and $D$ are natural numbers, and $\mathcal{A}_d$ is the predicate of
\ref{def:filtration-thresholds}'s Lean on functions on the line, which reads a function on
$[0,\infty)$ only. In clause (1) the hypothesis of the Lean statement is the datum $(a',k')$ of
dimension $d + 1$ with exponent $\zeta \mapsto F(|\zeta|)$; that $F \in \mathcal{A}_{d+1}$
follows from it and is not asked. A second declaration starts from $F \in \mathcal{A}_{d+1}$ and
concludes that there is such a datum for which the conclusions hold: there the datum is an
existential. ``Every isotropic origin datum $(a,k_d)$ in dimension $d$ with exponent
$\xi \mapsto F(|\xi|)$'' is a universal quantifier over the data of dimension $d$ whose exponent
is $F(|\xi|)$ at every $\xi \in \RR^d$. So it covers every version of the profile, and the
equality with the displayed integral is an equality of real numbers at every real $r > 0$.
Continuity is continuity on the set $(0,\infty)$. ``Has a discontinuity'' is its negation, and the
last sentence of the clause does not ask $F \in \mathcal{A}_d$, which the datum gives. In clause
(2) the datum is a term, with Gaussian coefficient $0$ and the profile $b\,\mathbf 1_{(0,\tau)}$
as a function on the whole line. $G$ is not a definition: one conjunct says that a function $G$
with $G(|\xi|)$ equal to the exponent at every $\xi$ exists, and one that every such function is
in $\mathcal{A}_d$ and not in $\mathcal{A}_{d+1}$. The strict inclusion is between the two sets
of functions. In clause (3) the exponent at $d = 1$ is $b\,\Cin(\tau|\omega|)$ at every $\omega$
of the one-dimensional Euclidean space, with the shared library's $\Cin$. The conclusion, for
every natural number $D \ge 2$, is the proposition that \ref{thm:no-cin-ray-in-the-plane}'s Lean
proves. A probe holds the two proofs of it, this node's and that theorem's; its type is that one
proposition twice, so as a statement it adds nothing. Witnesses: the step datum of dimension
$d + 1$ for the hypothesis of clause (1), and the step datum on the line, the $\Cin$ ray, for its last
sentence. Probes: no datum in the plane has the exponent $b\,\Cin(\tau|\zeta|)$, which is
obtained from the continuity at the one point $\tau$ and would not follow from continuity almost
everywhere, the step being continuous at every other point; and the $\Cin$ ray is in
$\mathcal{A}_1$ and not in $\mathcal{A}_2$.)}
\end{proposition}

\begin{proof}
\leanok
\emph{Clause (1).} By \ref{lem:shadow-polar-profile}, with $\bar k$ the function displayed in
the clause, $(a',\bar k)$ is an isotropic origin datum in dimension $d$ with exponent
$\xi \mapsto F(|\xi|)$. So $F \in \mathcal{A}_d$, by
\ref{prop:isotropic-polar-classification}(2). Let $(a,k_d)$ be any isotropic origin datum in
dimension $d$ with that exponent. By the uniqueness in \ref{prop:isotropic-polar-reading}(1),
$a = a'$ and $k_d = \bar k$ almost everywhere on $(0,\infty)$. Let $r > 0$, and choose points
$r_n \downarrow r$ and $s_n \uparrow r$ at which $k_d = \bar k$. As $k_d$ is nonincreasing,
$\bar k(r_n) = k_d(r_n) \le k_d(r) \le k_d(s_n) = \bar k(s_n)$, and both ends tend to
$\bar k(r)$, because $\bar k$ is continuous (\ref{lem:shadow-polar-profile}(2)). So
$k_d(r) = \bar k(r)$ for every $r > 0$, and $k_d$ is continuous. The next sentence of the clause
is the contrapositive.

\emph{Clause (2).} The function $k := b\,\mathbf 1_{(0,\tau)}$ is nonnegative and nonincreasing
on $(0,\infty)$, with $\int_0^1 k(r)\,r\,dr \le b/2$ and
$\int_1^\infty k(r)\,dr/r = b\log^+\tau$. So $(0,k)$ is an isotropic origin datum in dimension
$d$ (\ref{def:isotropic-origin-datum}). Its exponent $F_d$ is finite and nonnegative, as that
of every datum, and radial: for
$Q \in \Orth(d)$ the substitution $x = Qy$ in \ref{eq:isotropic-exponent} gives
$F_d(Q\xi) = F_d(\xi)$, because $Q$ preserves Lebesgue measure and the radial density of the
\Levy{} measure. So $F_d(\xi) = G(|\xi|)$ with $G(r) := F_d(re_1)$ for $r \ge 0$, and
$G \in \mathcal{A}_d$ by \ref{prop:isotropic-polar-classification}(2). The profile $k$ has the
left limit $b$ and the right limit $0$ at $\tau$. By clause (1), $G \notin \mathcal{A}_{d+1}$.
So $\mathcal{A}_{d+1} \ne \mathcal{A}_d$. The inclusion
$\mathcal{A}_{d+1} \subseteq \mathcal{A}_d$ follows from clause (1), since for every
$F \in \mathcal{A}_{d+1}$, $\zeta \mapsto F(|\zeta|)$ is the exponent of an isotropic origin datum
in dimension $d + 1$, by \ref{prop:isotropic-polar-classification}(2).

\emph{Clause (3).} At $d = 1$ the sphere is $S^0 = \{\pm1\}$, with $\sigma_0$ the counting
measure. By \ref{eq:polar-decomposition} the exponent of the datum
$(0, b\,\mathbf 1_{(0,\tau)})$ at $\omega \in \RR$ is
$b\int_0^\tau(1 - \cos\omega r)\,dr/r$, which is $b\,\Cin(\tau|\omega|)$: for $\omega \ne 0$ by
the substitution $v = |\omega| r$, and both vanish at $\omega = 0$. So
$G(r) = b\,\Cin(\tau r)$.
This is the $\Cin$ ray of \ref{thm:no-cin-ray-in-the-plane}, the admissible exponent with
Gaussian coefficient $0$ and profile $b\,\mathbf 1_{(0,\tau)}$. By clause (2) it is not in
$\mathcal{A}_2$. By clause (1), with the datum that
\ref{prop:isotropic-polar-classification}(2) supplies,
$\mathcal{A}_{D+1} \subseteq \mathcal{A}_D$ for every $D \ge 1$, so
$\mathcal{A}_D \subseteq \mathcal{A}_2$ for $D \ge 2$, and the ray is in none of them.

The cited facts enter as the Assignment says, and every step is machine-checked. The Lean
follows this proof. Two points of form: the conditions of \ref{def:isotropic-origin-datum} for
the step datum do not read the dimension and are taken from the datum on the line of
\ref{thm:no-cin-ray-in-the-plane}'s Lean, and the exponent at $d = 1$ is read from that Lean,
which computes it through the admissible exponent of the line. These are two declarations
about the datum on the line, a definition and a computation. The printed proof names that
theorem and does not use it, so it is not among the dependencies of this node. The nesting used
in clauses (2)
and (3) is the one from clause (1), on data; \ref{thm:filtration-strict}'s Lean has the same
inclusion by the image of the law.
\end{proof}

% not in the draft: Q-0425, the dictionary (see the comment before lem:shadow-polar-profile). The
% form, decided when the statement was written.
% (i) The FINITE-c form, not the general one (F in A_d iff r F'(r) is the radial exponent of a
% rotation-invariant infinitely divisible law on R^d with a finite logarithmic moment). Three
% reasons. The finite case is where the driver is a compound Poisson process and its jump law a
% probability law, so the statement is about characteristic functions, the objects of Schoenberg's
% and Gneiting's chain, which is the connection the response plan wants drawn (its section 0). The
% general form is not proved in the records (the blind check argues it from two printed
% statements, one with its proof unread, and checks it numerically on one family). And stated with
% a logarithmic moment of the LAW it needs a fact relating the moments of an infinitely divisible
% law to its Levy measure, which the ledger does not carry; this item may add no cited fact.
% (ii) The law is named by a random vector X, and its radial characteristic function is written
% phi, not varphi: varphi is the cosine profile of lem:isotropic-exponent-log's proof.
% (iii) The class is described intrinsically, by F(r)/log r -> c, through
% cor:origin-profile-exponent, so that "no Gaussian part and a finite positive profile at the
% origin" is a property of F and not of a chosen datum.
% (iv) "No atom at the origin" is a normalisation: an atom of mass p changes (c, phi) to
% (c(1-p), the function of the law conditioned off the origin) and leaves F unchanged. It is what
% makes (c, law) a function of F.
% (v) Clause (3) is proved from clauses (1) and (2) in the two dimensions and
% cor:origin-profile-exponent; it does not read lem:shadow-polar-profile.
% CHANGED Q-0429 (2026-10-12): annotation only, one reference in "Examples" (the step profile is
% now cor:dictionary-filtration(2) and prop:profile-one-dimension-down(3)). The statement and the
% proof are unchanged. Proved in print, not formalized: the author's decision of that day.
% CORRECTED after the blind check of Q-0429 (its report, F-B1): unchanged by Q-0429. Against the
% text the blind check of Q-0425 read, the proof of clause (2), existence, has two insertions,
% from that check's repairs: "and c G = k_d almost everywhere" and the integral of k_d(r+) before
% that of k_d(r). The blind check of Q-0429 read them and found them correct.
% CHANGED Q-0430 (2026-10-12): annotation only, one sentence in Attribution (Sato and Yamazato
% 1983, Jurek 1982, Wolfe 1982; records/affine/CITATIONS-primary-sources-B.md, items 3 and 7).
\begin{proposition}[the dictionary between $\mathcal{A}_d$ and the radial characteristic functions]
\label{prop:radial-dictionary}
\uses{def:isotropic-origin-datum,def:filtration-thresholds,lem:polar-decomposition,lem:isotropic-exponent-log,cor:origin-profile-exponent,prop:isotropic-polar-reading,prop:isotropic-polar-classification,thm:cited-A2,thm:cited-A3,thm:cited-A4}
Let $d \ge 1$. For $t \ge 0$ put
\[
  \Omega_d(t) := \frac{1}{|S^{d-1}|}\int_{S^{d-1}}\cos\bigl(t\langle e,\eta\rangle\bigr)\,
     \sigma_{d-1}(d\eta) ,
\]
which does not depend on the unit vector $e$; $\Omega_1 = \cos$. Call a random vector $X$ in
$\RR^d$ \emph{rotation invariant} when $QX$ has the law of $X$ for every $Q \in \Orth(d)$. Its
characteristic function is then real and radial,
\[
  \mathbb{E}\,e^{i\langle\xi,X\rangle} = \phi(|\xi|), \qquad
  \phi(r) := \mathbb{E}\,\Omega_d\bigl(r|X|\bigr) \quad (r \ge 0) ,
\]
the function $\phi$ determines the law of $X$, and $X$ has the law of $R\,U$, where $R$ has the
law of $|X|$ and $U$ is uniformly distributed on $S^{d-1}$ and independent of $R$.
\begin{enumerate}
\item \emph{(From a law to an exponent.)} Let $X$ be a rotation-invariant random vector in
$\RR^d$ with $\Pr(X = 0) = 0$, let $\phi$ be as above, and let $0 < c < \infty$. Then
$\int_0^1 (1 - \phi(u))\,du/u < \infty$ if and only if $\mathbb{E}\log^+|X| < \infty$. In that
case the function
\[
  F(r) := c\int_0^r \bigl(1 - \phi(u)\bigr)\,\frac{du}{u} \qquad (r \ge 0)
\]
belongs to $\mathcal{A}_d$: $\xi \mapsto F(|\xi|)$ is the exponent of the isotropic origin datum
$(0,k_d)$ in dimension $d$ with $k_d(r) = c\,\Pr(|X| > r)$. Its profile at the origin is
$k_d(0{+}) = c$, and $F(r)/\log r \to c$ as $r \to \infty$.
\item \emph{(From an exponent to a law.)} Conversely, let $F \in \mathcal{A}_d$ and suppose that
$F(r)/\log r$ has a limit $c \in (0,\infty)$ as $r \to \infty$. Then $F$ is the function of
clause (1) for this $c$ and for a rotation-invariant $X$ with $\Pr(X = 0) = 0$ and
$\mathbb{E}\log^+|X| < \infty$, whose law is unique. Moreover $F$ is continuously differentiable
on $(0,\infty)$, with
\[
  \phi(r) = 1 - \frac{r\,F'(r)}{c} \quad (r > 0), \qquad
  \Pr(|X| > r) = \frac{k_d(r{+})}{c} \quad (r > 0)
\]
for every isotropic origin datum $(0,k_d)$ in dimension $d$ with exponent $\xi \mapsto F(|\xi|)$.
So $(c, \text{the law of } X) \mapsto F$ is a bijection from the pairs of clause (1) with
$\mathbb{E}\log^+|X| < \infty$ onto the set of the $F \in \mathcal{A}_d$ for which $F(r)/\log r$
has a limit in $(0,\infty)$. These are the members of $\mathcal{A}_d$ with Gaussian coefficient
$0$ and profile at the origin in $(0,\infty)$.
\item \emph{(The next dimension.)} Let $F$, $c$, $X$ and $\phi$ be as in clause (1), with
$\mathbb{E}\log^+|X| < \infty$. Then $F \in \mathcal{A}_{d+1}$ if and only if
$\zeta \mapsto \phi(|\zeta|)$, $\zeta \in \RR^{d+1}$, is the characteristic function of a random
vector $X'$ in $\RR^{d+1}$. Every such $X'$ is rotation invariant, with $\Pr(X' = 0) = 0$ and
$\mathbb{E}\log^+|X'| < \infty$, its first $d$ coordinates have the law of $X$, and $F$ is the
function of clause (1) in dimension $d + 1$ for the same $c$ and for $X'$.
\end{enumerate}
\statusA\quad\ledger{A2}\quad\ledger{A3}\quad\ledger{A4}\quad\emph{(\textbf{Assignment.} The node
cites nothing of its own. The classes $\mathcal{A}_d$ and $\mathcal{A}_{d+1}$ are read through
\ref{prop:isotropic-polar-classification}(2), in both directions, which spends \ledger{A2},
\ledger{A3} and \ledger{A4}. The identity behind the dictionary is an exchange of two integrals of
a nonnegative function and spends nothing: for a datum $(0,k_d)$ with a finite positive profile at
the origin it reads $F(\rho) = c\int_0^\rho(1 - \phi(v))\,dv/v$, and it is Step~A of the proof.
The limit $F(r)/\log r \to c$ and the equality of the constants $c$ in the two dimensions are
\ref{cor:origin-profile-exponent}, which cites nothing. \textbf{What the citations carry and what
they do not.} \ledger{A2}, \ledger{A3} and \ledger{A4} carry the polar form of the
self-decomposable laws, which is what makes $\mathcal{A}_d$ the set of the exponents of the data.
They carry nothing about characteristic functions of rotation-invariant laws. \textbf{What it is
worth.} A reader can decide membership of $\mathcal{A}_d$ on a familiar object. An exponent with
logarithmic growth is in $\mathcal{A}_d$ exactly when $1 - rF'(r)/c$ is the characteristic
function of a rotation-invariant law on $\RR^d$, read along a ray. The law is that of the jumps
of the driving compound Poisson process, $c$ is the rate of the jumps, and the polar profile is
$c$ times the tail of the jump size. Clause (3) explains why the classes shrink with the
dimension: $F$ survives in dimension $d + 1$ exactly when the same radial function is still a
characteristic function there, that is when the jump law is the shadow of a rotation-invariant
law of $\RR^{d+1}$. Membership of a radial function in the class of the characteristic functions
on $\RR^d$ is the subject of Schoenberg's and Gneiting's work, and through clause (3) every
function that separates two of their classes separates $\mathcal{A}_d$ from
$\mathcal{A}_{d+1}$. \textbf{Examples.} The step profile $k_d = b\,\mathbf 1_{(0,\tau)}$ is
$c = b$ and $X$ uniform on the sphere of radius $\tau$, $\phi(r) = \Omega_d(\tau r)$; at $d = 1$
this is the $\Cin$ ray, with $\phi(r) = \cos(\tau r)$
(\ref{cor:dictionary-filtration}(2), \ref{prop:profile-one-dimension-down}(3)). The \Matern{} exponent $\gamma\log(1 + r^2/2)$ has
$c = 2\gamma$ and $\phi(r) = (1 + r^2/2)^{-1}$ in every dimension. \textbf{Attribution (the
anchors read at the page image for the attribution item Q-0426; its record
\texttt{records/affine/PRIOR-ART-review-round1-additions.md}, section~1, and its reading
notes).} Clauses (1) and (2) are one case of a known theorem, in another notation: the case in
which the driving \Levy{} process is a compound Poisson process. They are not that theorem for
every rotation-invariant self-decomposable law. The theorem: a law on $\RR^d$ is
self-decomposable exactly when its exponent is $\int_0^\infty\psi(e^{-s}\xi)\,ds$ for the exponent
$\psi$ of an infinitely divisible law with a finite logarithmic moment, the law of the driving
\Levy{} process at time $1$, and $\psi$ is then unique: \sscite{jurek1983integral}, Thm.~3.2,
p.~252, and Thm.~4.1, p.~254; \sscite{sato1984operator}, Thm.~4.1, p.~86; and
\sscite{sato1999levy}~Thm.~17.5, pp.~108--109, whose Notes, pp.~117--118, credit Sato and
Yamazato. (In Thm.~17.5 the letter $c$ is the rate of the Ornstein--Uhlenbeck process, not the
$c$ of this node.) The same representation is stated in three further papers, read at the page
for Q-0430 (\texttt{records/affine/CITATIONS-primary-sources-B.md}, items 3 and 7): as the law
of a stationary process of Ornstein--Uhlenbeck type, with the proofs referred to the paper of
Sato and Yamazato just cited, in \sscite{sato1983stationary}, Thms.~5.1--5.3, pp.~546--547; for
operator-self-decomposable laws on a Banach space, as the limit in distribution of
$\int_{(0,t]}e^{-sQ}\,dY(s)$, in \sscite{jurek1982integral}, Thms.~2.1 and~3.1, pp.~388
and~390; and on the line, as the limit in law of $e^{-\gamma t}\int_0^t e^{\gamma u}\,dB(u)$
with $\gamma > 0$, in \sscite{wolfe1982continuous}, Thm.~1, p.~306, and Cor.~3, p.~307. The radial form is \sscite{jurek1992operator}, Cor.~1.7, pp.~1056--1057, with
its proof on pp.~1060--1061, restated as \sscite{jurek1993operator}, Prop.~3.9.6, pp.~180--182.
It is printed as a series in Bessel functions with a probability measure on the radii. In the
notation of this chapter it reads $F(r) = a_0 r^2 + \int_0^r\psi(u)\,du/u$, with $a_0 \ge 0$ and
$\psi$ the radial exponent of the driver; the two forms were compared numerically for that
record. The dictionary is the case $a_0 = 0$ with a compound Poisson driver. Its intensity is
$c$ and its jumps have the law of $X$, so that $\psi = c\,(1 - \phi)$ and the formula is that of
clause (1). The logarithmic moment $\mathbb{E}\log^+|X| < \infty$ is the condition under which
the integral converges: \sscite{sato1984operator}, (1.3), p.~74, with Thm.~4.1, p.~86, has it on
the \Levy{} measure of the driver, as the integrability of $\log|x|$ over $\{|x| > 1\}$, and
here that measure is $c$ times the law of $X$. A rotation-invariant self-decomposable law with a
Gaussian part, or with an infinite profile at the origin, has a driver that is not compound
Poisson and is outside the node: the rotation-invariant stable laws and the separating family of
\ref{thm:filtration-strict} are of the second kind. That the correspondence commutes with
linear maps, a projection included, is
\sscite{jurek1983integral}, Lemma~4.2, (4.2c), p.~254, where the map is any linear operator, and
\sscite{jurek1993operator}, Prop.~3.6.10(c), p.~128, with its proof on pp.~129--130. The classes
of the radial characteristic functions are those of \sscite{schoenberg1938metric}, Thm.~1,
p.~816, and \sscite{gneiting1999radial}. None of these is used: the proof below rests on this
chapter's statements, so that no cited fact is added. Clause (3) is proved here, and nothing is
attributed for it. \textbf{Readings.} $F$ is a
function on $[0,\infty)$, as the members of $\mathcal{A}_d$ are. The characteristic function of a
random vector on $\RR^{d+1}$ in clause (3) is that of any probability measure; rotation
invariance is a conclusion. \textbf{Not claimed}: the general form, without the restriction to a
finite $c$ and to no Gaussian part, in which $rF'(r)$ is the radial exponent of a
rotation-invariant infinitely divisible law with a finite logarithmic moment. It is in print, in
the sources above. It is not a clause here because this chapter's statements do not prove it:
with a logarithmic moment of the law it needs a fact about the moments of infinitely divisible
laws that the ledger does not carry, and it would then be a cited statement. The separating family of
\ref{thm:filtration-strict} has $c = \infty$ and falls outside this node; it is described in
\ref{prop:separating-family-mixtures}. Nor a description of the radial functions that are
characteristic functions on $\RR^{d+1}$; nor any use of Gneiting's separating functions, each of
which would give, through clause (3), a member of $\mathcal{A}_d \setminus \mathcal{A}_{d+1}$ on
one further cited fact.)}
\end{proposition}

\begin{proof}
\emph{The function $\Omega_d$.} For unit vectors $e, e'$ take $Q \in \Orth(d)$ with $Qe = e'$.
Then $\langle e',\eta\rangle = \langle e,Q^{\tp}\eta\rangle$, and $\sigma_{d-1}$ is invariant
under $Q^{\tp}$, so the integral defining $\Omega_d(t)$ is the same for $e'$ as for $e$. At
$d = 1$ the sphere is $\{\pm1\}$ and $\Omega_1(t) = \cos t$. $\Omega_d$ is continuous, by
dominated convergence, $|\Omega_d| \le 1$ and $\Omega_d(0) = 1$. For $x \in \RR^d$ and $t \ge 0$,
\begin{equation}\label{eq:sphere-mean-cos}
  \frac{1}{|S^{d-1}|}\int_{S^{d-1}}e^{it\langle x,\theta\rangle}\,\sigma_{d-1}(d\theta)
   = \Omega_d\bigl(t|x|\bigr) .
\end{equation}
Indeed the imaginary part vanishes by the symmetry $\theta \mapsto -\theta$ of $\sigma_{d-1}$.
The real part is the definition of $\Omega_d(t|x|)$ with $e = x/|x|$ when $x \ne 0$, and both
sides are $1$ when $x = 0$.

\emph{The preamble.} Let $X$ be rotation invariant. As $-I \in \Orth(d)$, $-X$ has the law of
$X$, so the characteristic function is real:
$\mathbb{E}\,e^{i\langle\xi,X\rangle} = \mathbb{E}\cos\langle\xi,X\rangle$. Let $\xi \in \RR^d$
and let $\theta$ be a unit vector. There is $Q \in \Orth(d)$ with $Q\xi = |\xi|\theta$, and
$Q^{\tp}X$ has the law of $X$, so
\[
  \mathbb{E}\cos\langle\xi,X\rangle = \mathbb{E}\cos\langle\xi,Q^{\tp}X\rangle
   = \mathbb{E}\cos\bigl(|\xi|\langle\theta,X\rangle\bigr) .
\]
The left side does not depend on $\theta$. Averaging over $\theta$ against
$\sigma_{d-1}/|S^{d-1}|$, exchanging the two integrals of a bounded continuous function, and using
\ref{eq:sphere-mean-cos} gives
$\mathbb{E}\,e^{i\langle\xi,X\rangle} = \mathbb{E}\,\Omega_d(|\xi||X|) = \phi(|\xi|)$. By the
uniqueness theorem for characteristic functions, $\phi$ determines the law of $X$. Let $R$ have
the law of $|X|$ and let $U$ be uniform on $S^{d-1}$, that is of law $\sigma_{d-1}/|S^{d-1}|$,
and independent of $R$. By \ref{eq:sphere-mean-cos} and Fubini's theorem,
$\mathbb{E}\,e^{i\langle\xi,RU\rangle} = \mathbb{E}\,\Omega_d(R|\xi|) = \phi(|\xi|)$. So $RU$ has
the law of $X$, by the same theorem.

\emph{Step A: the identity.} Let $R$ be a random variable with values in $(0,\infty)$, let
$0 < c < \infty$, and put $k(r) := c\,\Pr(R > r)$ for $r > 0$ and
$\phi(v) := \mathbb{E}\,\Omega_d(vR)$ for $v \ge 0$. Then for every $\rho > 0$,
\begin{equation}\label{eq:dictionary-identity}
  \int_0^\infty\bigl(1 - \Omega_d(\rho r)\bigr)\,k(r)\,\frac{dr}{r}
   = c\int_0^\rho\bigl(1 - \phi(v)\bigr)\,\frac{dv}{v} ,
\end{equation}
both sides in $[0,\infty]$. Indeed, as $1 - \Omega_d \ge 0$, Tonelli's theorem applies twice, with
the substitution $r = vR/\rho$ in between:
\[
  \int_0^\infty\bigl(1 - \Omega_d(\rho r)\bigr)\,c\,\Pr(R > r)\,\frac{dr}{r}
   = c\,\mathbb{E}\int_0^R\bigl(1 - \Omega_d(\rho r)\bigr)\frac{dr}{r}
   = c\,\mathbb{E}\int_0^\rho\bigl(1 - \Omega_d(vR)\bigr)\frac{dv}{v}
   = c\int_0^\rho\bigl(1 - \phi(v)\bigr)\frac{dv}{v} .
\]
Moreover, if $\nu$ is the measure on $\RR^d \setminus \{0\}$ with the density
$k(|x|)\,(|S^{d-1}|\,|x|^d)^{-1}$, then for $\xi \ne 0$, by polar coordinates, Tonelli's theorem
and \ref{eq:sphere-mean-cos},
\begin{equation}\label{eq:dictionary-polar}
  \int_{\RR^d}\bigl(1 - \cos\langle\xi,x\rangle\bigr)\,\nu(dx)
   = \int_0^\infty\bigl(1 - \Omega_d(|\xi| r)\bigr)\,k(r)\,\frac{dr}{r} .
\end{equation}
The same holds for every nonnegative Borel $k$ in place of this one.

\emph{Step B: the two integrability conditions.} With $R$, $c$ and $k$ as in Step~A, Tonelli's
theorem gives
\[
  \int_1^\infty k(r)\,\frac{dr}{r} = c\,\mathbb{E}\int_1^{\max\{R,1\}}\frac{dr}{r}
   = c\,\mathbb{E}\log^+R ,
\]
and $\int_0^1 k(r)\,r\,dr \le c/2$ because $k \le c$.

\emph{Step C: the logarithmic moment read on $\phi$.} Put
$H_d(s) := \int_0^s (1 - \Omega_d(v))\,dv/v$ for $s \ge 0$. It is nonnegative and nondecreasing.
The pair $(0,\mathbf 1_{(0,1)})$ is an isotropic origin datum in dimension $d$, with profile at
the origin $1$ and with $\int_1^\infty\mathbf 1_{(0,1)}(r)\,dr/r = 0$. By
\ref{eq:dictionary-polar} and \ref{eq:dictionary-identity} with $R \equiv 1$ and $c = 1$, its
exponent at a point of norm $s$ is $H_d(s)$. So \ref{lem:isotropic-exponent-log}(1) gives
$H_d(s) \le \log s + \tfrac{13}{4}$ for $s \ge 1$, and hence
$H_d(s) \le \log^+ s + \tfrac{13}{4}$ for every $s \ge 0$, $H_d$ being nondecreasing. And
\ref{cor:origin-profile-exponent}(2) gives $H_d(s)/\log s \to 1$ as $s \to \infty$, so there is
$s_0 > 1$ with $H_d(s) \ge \tfrac12\log s$ for $s \ge s_0$. Now let $R$ and $\phi$ be as in
Step~A. By Tonelli's theorem and the substitution $v \mapsto vR$,
\[
  \int_0^1\bigl(1 - \phi(v)\bigr)\frac{dv}{v}
   = \mathbb{E}\int_0^1\bigl(1 - \Omega_d(vR)\bigr)\frac{dv}{v} = \mathbb{E}\,H_d(R) .
\]
By the two bounds on $H_d$,
$\tfrac12\,\mathbb{E}\log^+R - \tfrac12\log s_0 \le \mathbb{E}\,H_d(R) \le
\mathbb{E}\log^+R + \tfrac{13}{4}$. So $\int_0^1(1 - \phi(v))\,dv/v$ is finite if and only if
$\mathbb{E}\log^+R$ is.

\emph{Clause (1).} Let $X$ be as in the clause and $R := |X|$, which is positive almost surely.
By the preamble $\phi(v) = \mathbb{E}\,\Omega_d(vR)$, and the equivalence is Step~C. Suppose
$\mathbb{E}\log^+|X| < \infty$ and put $k_d(r) := c\,\Pr(|X| > r)$. It is nonincreasing with
values in $[0,c]$, and $k_d(0{+}) = c\,\Pr(|X| > 0) = c$. By Step~B it satisfies the two
conditions of \ref{def:isotropic-origin-datum}. So $(0,k_d)$ is an isotropic origin datum in
dimension $d$. By \ref{eq:dictionary-polar} and \ref{eq:dictionary-identity} its exponent at
$\xi \ne 0$ is $c\int_0^{|\xi|}(1 - \phi(v))\,dv/v = F(|\xi|)$, and both vanish at $\xi = 0$. By
\ref{prop:isotropic-polar-classification}(2), $F \in \mathcal{A}_d$. By
\ref{cor:origin-profile-exponent}(2) for this datum, read along a ray, $F(r)/\log r \to c$.

\emph{Clause (2), existence.} Let $F \in \mathcal{A}_d$ with $F(r)/\log r \to c \in (0,\infty)$.
By \ref{prop:isotropic-polar-classification}(2), $\xi \mapsto F(|\xi|)$ is the exponent of an
isotropic origin datum $(a,k_d)$ in dimension $d$. By \ref{cor:origin-profile-exponent}(1),
$F(r)/r^2 \to \tfrac12 a$; as $F(r)/\log r$ has a finite limit, $F(r)/r^2 \to 0$, so $a = 0$.
Then \ref{cor:origin-profile-exponent}(2) gives $F(r)/\log r \to k_d(0{+})$, so $k_d(0{+}) = c$.
As in the proof of \ref{lem:isotropic-exponent-log}, $k_d(r) \to 0$ as $r \to \infty$. Put
$G(r) := k_d(r{+})/c$ for $r > 0$. It is nonincreasing and right-continuous, it tends to $1$ as
$r \downarrow 0$ and to $0$ as $r \to \infty$, and $c\,G = k_d$ off the countable set of the
discontinuities of $k_d$. So there is a random variable $R$ with values in $(0,\infty)$ and
$\Pr(R > r) = G(r)$ for every $r > 0$. Let $U$ be uniform on $S^{d-1}$ and independent of $R$,
and $X := RU$. Then $X$ is rotation invariant, $|X| = R$, $\Pr(X = 0) = 0$, and by Step~B and
$c\,G = k_d$ almost everywhere,
$c\,\mathbb{E}\log^+|X| = \int_1^\infty k_d(r{+})\,dr/r = \int_1^\infty k_d(r)\,dr/r < \infty$.
Its function $\phi$ is that of
Step~A. Let $\rho > 0$ and let $e$ be a unit vector. By \ref{eq:dictionary-polar} for $k_d$, the
equality $k_d = c\,\Pr(R > \cdot)$ almost everywhere, and \ref{eq:dictionary-identity},
\[
  F(\rho) = \int_0^\infty\bigl(1 - \Omega_d(\rho r)\bigr)\,k_d(r)\,\frac{dr}{r}
   = c\int_0^\rho\bigl(1 - \phi(v)\bigr)\frac{dv}{v} .
\]
So $F$ is the function of clause (1) for $c$ and $X$. As $\phi$ is continuous, the integrand
$(1 - \phi(v))/v$ is continuous on $(0,\infty)$, so $F$ is continuously differentiable there, with
$rF'(r) = c\,(1 - \phi(r))$. The datum $(0,k_d)$ was any datum with exponent $F(|\cdot|)$, and
$\Pr(|X| > r) = k_d(r{+})/c$ by construction.

\emph{Clause (2), uniqueness.} Suppose $F$ is also the function of clause (1) for $c_1$ and
$X_1$, with $\mathbb{E}\log^+|X_1| < \infty$ and radial function $\phi_1$. By clause (1),
$F(r)/\log r \to c_1$, so $c_1 = c$. Both integrands being continuous, differentiation gives
$c\,(1 - \phi_1(r)) = rF'(r) = c\,(1 - \phi(r))$ for every $r > 0$, and
$\phi_1(0) = \phi(0) = 1$. So $\phi_1 = \phi$, and $X_1$ has the law of $X$ by the preamble. In
particular $\Pr(|X| > r) = k_d(r{+})/c$ holds for the one law and every datum, since the
construction above can start from any datum. \emph{The bijection.} By clause (1) the map sends a
pair $(c, \text{law of } X)$ into the set named; by the existence part it is onto, and by the
uniqueness part it is one to one. The set is that of the members of $\mathcal{A}_d$ with
Gaussian coefficient $0$ and profile at the origin in $(0,\infty)$: one inclusion was shown at
the start of the existence part, and the other is \ref{cor:origin-profile-exponent}(2).

\emph{Clause (3), if.} Let $X'$ be a random vector in $\RR^{d+1}$ with characteristic function
$\zeta \mapsto \phi(|\zeta|)$. For $Q \in \Orth(d+1)$ the characteristic function of $QX'$ is
$\zeta \mapsto \phi(|Q^{\tp}\zeta|) = \phi(|\zeta|)$, so $QX'$ has the law of $X'$: $X'$ is
rotation invariant. Let $\Pi$ be the projection on the first $d$ coordinates. The characteristic
function of $\Pi X'$ at $\xi \in \RR^d$ is that of $X'$ at $(\xi,0)$, that is $\phi(|\xi|)$. So
$\Pi X'$ has the law of $X$, and $\Pr(X' = 0) \le \Pr(\Pi X' = 0) = \Pr(X = 0) = 0$. By the
preamble in dimension $d + 1$, the radial function of $X'$ is $\phi$. As
$\int_0^1(1 - \phi(u))\,du/u = F(1)/c < \infty$, clause (1) in dimension $d + 1$ gives
$\mathbb{E}\log^+|X'| < \infty$ and $F \in \mathcal{A}_{d+1}$, with $F$ the function of that
clause for $c$ and $X'$.

\emph{Clause (3), only if.} Let $F \in \mathcal{A}_{d+1}$. By clause (1) in dimension $d$,
$F(r)/\log r \to c$. So clause (2) applies in dimension $d + 1$: there is a rotation-invariant
$X'$ in $\RR^{d+1}$ with $\Pr(X' = 0) = 0$ and $\mathbb{E}\log^+|X'| < \infty$ whose radial
function is $r \mapsto 1 - rF'(r)/c$ on $(0,\infty)$. By clause (2) in dimension $d$ this
function is $\phi$. Both radial functions equal $1$ at $0$. So the characteristic function of
$X'$ is $\zeta \mapsto \phi(|\zeta|)$.

The cited facts \ledger{A2}, \ledger{A3} and \ledger{A4} enter through
\ref{prop:isotropic-polar-classification}(2) only, each time a function is placed in a class
$\mathcal{A}_d$ or taken from one. Nothing here is machine-checked yet.
\end{proof}

% not in the draft: Q-0425, what follows from the dictionary and the shadow. The environment is
% `theorem` because the blueprint has no `corollary` environment, as for
% cor:origin-profile-exponent; the isotropic paper would print it as a corollary.
% CHANGED Q-0429 (2026-10-12): SPLIT (ADR-0025; the author's decision of that day: the
% consequences on data are formalized, the dictionary is not). Until then this node had three
% clauses: (1) the continuity of the profile one dimension down, with its reading in the
% dictionary as a last sentence; (2) the step datum as a witness at every step, with two readings
% in the dictionary; (3) the Cin rays. The statements on data are now
% prop:profile-one-dimension-down, which precedes prop:radial-dictionary and does not read it.
% What stays under this label is what needs the dictionary, and the label still says what the
% node is: a corollary of the dictionary about the filtration.
%   clause (1) here = the last sentence of the old clause (1), word for word (CORRECTED after
%                     the blind check of Q-0429, its report, F-B1: word for word against the
%                     text before the split; against the text the blind check of Q-0425 read it
%                     has "with E log^+|X| < infty", in the statement and in the proof, from
%                     that check's repairs; read by the blind check of Q-0429 and correct);
%   clause (2) here = the formula for G, the identification with the step datum and the two
%                     readings of the old clause (2). The old clause said "Then G in A_d minus
%                     A_{d+1}. Here ... Hence A_{d+1} strictly inside A_d for every d >= 1, and
%                     ..."; the membership, the exclusion and the strictness are
%                     prop:profile-one-dimension-down(2), and this clause refers to it;
%   the old clause (3) is prop:profile-one-dimension-down(3) and has no reading here.
% Proof: the last three sentences of the old proof of clause (1), with "For the last one, let"
% become "Let" and "by prop:profile-one-dimension-down(1)" added where the old proof said "so its
% profile is continuous" of what stood three lines above; the first two sentences and the last
% sentence of the old proof of clause (2), word for word, joined by one new sentence that takes
% the exclusion from prop:profile-one-dimension-down(2). No \lean tag, and none is scheduled.
% CHANGED Q-0430 (2026-10-12): annotation only, one see-also in Attribution (Daley and Porcu,
% Lemma 12, p. 1822; records/affine/READING-daley2014dimension.md).
\begin{theorem}[the consequences on data, read in the dictionary]
\label{cor:dictionary-filtration}
\uses{prop:radial-dictionary,prop:profile-one-dimension-down,def:isotropic-origin-datum,def:filtration-thresholds,thm:cited-A2,thm:cited-A3,thm:cited-A4}
Let $d \ge 1$, and let $\Omega_d$ be the function of \ref{prop:radial-dictionary}.
\begin{enumerate}
\item \emph{(The jump size has no atom.)} In the terms of
\ref{prop:radial-dictionary}: if $F$ and $X$ are as in its clause (1), with
$\mathbb{E}\log^+|X| < \infty$, and $F \in \mathcal{A}_{d+1}$, then $\Pr(|X| = r) = 0$ for every
$r > 0$.
\item \emph{(The step datum.)} Let $b > 0$ and $\tau > 0$, and put
\[
  G(r) := b\int_0^r \bigl(1 - \Omega_d(\tau u)\bigr)\,\frac{du}{u} \qquad (r \ge 0) .
\]
Then $\xi \mapsto G(|\xi|)$ is the exponent of the datum $(0, b\,\mathbf 1_{(0,\tau)})$ in
dimension $d$, so that $G$ is the function of \ref{prop:profile-one-dimension-down}(2), and in
\ref{prop:radial-dictionary} it is the function of $c = b$ and of $X$ uniformly distributed on
the sphere of radius $\tau$. Hence $\zeta \mapsto \Omega_d(\tau|\zeta|)$ is not a characteristic
function on $\RR^{d+1}$.
\end{enumerate}
\statusA\quad\ledger{A2}\quad\ledger{A3}\quad\ledger{A4}\quad\emph{(\textbf{Assignment.} The node
cites nothing of its own. \ledger{A2}, \ledger{A3} and \ledger{A4} are spent through
\ref{prop:radial-dictionary} and through \ref{prop:profile-one-dimension-down}, each of which
reads the classes through \ref{prop:isotropic-polar-classification}(2), and the second also the
uniqueness of the datum, \ledger{A3}(ii). \textbf{What the citations carry and what they do
not.} They carry the polar form of $\mathcal{A}_d$ and the uniqueness of a datum. They carry
nothing about characteristic functions of rotation-invariant laws: the dictionary is proved from
this chapter's statements. \textbf{What it is.} The two readings of
\ref{prop:profile-one-dimension-down} in the terms of \ref{prop:radial-dictionary}. Clause (1):
the profile is $c$ times the tail of the jump size $|X|$, so a continuous profile is a jump size
without atoms, and the jump law of a member of $\mathcal{A}_{d+1}$, read in dimension $d$,
charges no sphere. Clause (2): the step datum is the uniform law on a sphere, whose radius is an
atom, and its radial function $\Omega_d(\tau\,\cdot)$ is a characteristic function on $\RR^d$
that is not one on $\RR^{d+1}$. \textbf{Proved in print, not formalized.} The node has no Lean
and none is scheduled (the author's decision of 2026-10-12, with \ref{prop:radial-dictionary},
whose Lean it would need). What a reader uses to decide whether an exponent is admitted in a
dimension is \ref{prop:profile-one-dimension-down}, which is stated on data and is
machine-checked. Until Q-0429 the two nodes were one, under this label. \textbf{Attribution (the
anchors read
at the page image for the attribution item Q-0426; its record
\texttt{records/affine/PRIOR-ART-review-round1-additions.md}, sections 1 and 2).} For a finite
profile at the origin clause (1) is \sscite{cambanis1981theory}, Cor.~2, p.~372, read
through \ref{prop:radial-dictionary}: a radial characteristic function on $\RR^{d+1}$ has, as a
radial characteristic function on $\RR^d$, a radius law that is absolutely continuous on
$(0,\infty)$. Applied to the uniform law on a sphere, whose radius is an atom, it gives the last
assertion of clause (2). The inclusions of the classes of radial characteristic functions
are \sscite{schoenberg1938metric}, p.~816, and \sscite{gneiting1999radial}, p.~94, has their
strictness with another function. See also \sscite{daley2014dimension}, Lemma~12, p.~1822,
where $\Omega_d$ is shown not to be a radial characteristic function on $\RR^{d+1}$, from the
function of Gneiting's paper (read at the page image for Q-0430;
\texttt{records/affine/CITATIONS-primary-sources-B.md}, item 4). None of these is used.
\textbf{Not claimed}: the converse of clause (1), that a jump size without atoms puts $F$ in
$\mathcal{A}_{d+1}$. Nor that $\zeta \mapsto G(|\zeta|)$ fails to be the exponent of an
infinitely divisible law on $\RR^{d+1}$. Nor a description of the radial functions that are
characteristic functions on $\RR^{d+1}$.)}
\end{theorem}

\begin{proof}
\emph{Clause (1).} Let $F$ and $X$ be as in
\ref{prop:radial-dictionary}(1), with $\mathbb{E}\log^+|X| < \infty$, and
$F \in \mathcal{A}_{d+1}$. By that clause
$(0, c\,\Pr(|X| > \cdot))$ is a datum in dimension $d$ with exponent $F(|\cdot|)$, so its profile
is continuous, by \ref{prop:profile-one-dimension-down}(1), and
$\Pr(|X| = r) = \Pr(|X| \ge r) - \Pr(|X| > r)$ is $1/c$ times the jump of
that profile at $r$, which is $0$.

\emph{Clause (2).} Let $U$ be uniformly distributed on $S^{d-1}$ and $X := \tau U$. It is
rotation invariant, $\Pr(X = 0) = 0$, $\mathbb{E}\log^+|X| = \log^+\tau < \infty$, and its radial
function is $\phi(r) = \Omega_d(\tau r)$. By \ref{prop:radial-dictionary}(1) with $c = b$, the
function $G$ belongs to $\mathcal{A}_d$, and $\xi \mapsto G(|\xi|)$ is the exponent of the datum
$(0,k_d)$ with $k_d(r) = b\,\Pr(\tau > r) = b\,\mathbf 1_{(0,\tau)}(r)$. So $G$ is the function
of \ref{prop:profile-one-dimension-down}(2), and $G \notin \mathcal{A}_{d+1}$ by that clause. By
\ref{prop:radial-dictionary}(3), $G \notin \mathcal{A}_{d+1}$ says that
$\zeta \mapsto \phi(|\zeta|) = \Omega_d(\tau|\zeta|)$ is not a characteristic function on
$\RR^{d+1}$.

The cited facts enter as the Assignment says. Nothing here is machine-checked.
\end{proof}

% not in the draft: queue item Q-0434 (hub), from the second external round of the isotropic paper
% (records/affine/PLAN-review-response-paper-iso-round2.md; the verification record
% records/affine/reviews/VERIFY-round2-mathematics.md, item 13). The paper's section 4.1 asked
% whether the exponent of a step datum, which is not that of a self-decomposable law in the next
% dimension (prop:profile-one-dimension-down(2)), is still the exponent of an infinitely
% divisible law there. It is, in the next two dimensions and in no further one.
% Proved in print, no Lean: the statement is in terms of Omega_d, which is defined in
% prop:radial-dictionary, a node proved in print (the author's decision of 2026-10-12). It stands
% with the readings in the dictionary for that reason.
% What differs from the verification record's draft, and why (records/affine/NODES-q0434.md).
% (i) The record's clause (1), "of a self-decomposable law iff D <= d", is
% prop:profile-one-dimension-down(2) with the nesting, and is recalled before the clauses, not
% stated as one.
% (ii) The identity G_d = G_{d+2} + (b/d)(1 - Omega_{d+2}(tau .)) is proved from
% lem:shadow-polar-profile applied twice, with one differentiation in tau. The record proves it
% from Omega_d = Omega_{d+2} + (t/d) Omega_{d+2}', by the density of one coordinate of a uniform
% point of a sphere and an integration by parts. That density is not a statement of this chapter;
% the shadow lemma is, and it is machine-checked. The two identities are equivalent.
% (iii) The profile in dimension d + 1 is verified by eq:shadow-identity, read for a profile that
% is not monotone, in place of the record's computation of the image of the surface measure of a
% sphere under a projection. Only the first of the record's two closed forms is stated.
% (iv) In "only if" the Haar average of the record is replaced by the decomposition X = R U of
% prop:radial-dictionary's preamble, and the null set of the sphere by polar coordinates.
% (v) The record's remark, that for D >= d + 3 and a small height b the function e^{-G} is not a
% characteristic function at all, is not a clause: it needs the continuity theorem in the form
% that produces the limit law, which is not among the chapter's cited facts. It is described
% under "Not claimed".
% Numerical checks: records/affine/reviews/verify-round2-scripts/item13_step_exponent_dimensions.py
% (the record's), and records/affine/q0434-scripts/check_q0434_nodes.py for the identities of (ii)
% and (iii) in the form the proof uses.
% CHANGED after the blind check of Q-0434 (records/formalization/blind-check-q0434-report-A.md:
% correct, the "only if" included, for every b > 0). Proof only. F-A5: the transform of the law
% of a symmetric triplet is a fourth fact, (d), cited where the proof passed from fact (c) to a
% transform. F-A6: the proof of prop:isotropic-polar-reading is cited from its third sentence on,
% and the proof says that only the modulus of the representation is used.
\begin{proposition}[the step exponent one and two dimensions up]
\label{prop:step-exponent-dimensions}
\uses{prop:profile-one-dimension-down,lem:shadow-polar-profile,prop:radial-dictionary,cor:dictionary-filtration,def:isotropic-origin-datum,lem:polar-decomposition,prop:isotropic-polar-reading,def:filtration-thresholds,thm:cited-A2,thm:cited-A3,thm:cited-A4}
Let $d \ge 1$, $b > 0$ and $\tau > 0$. For an integer $n \ge 1$ let $\Omega_n$ be the function of
\ref{prop:radial-dictionary} in dimension $n$, and put
\[
  G_n(\rho) := b\int_0^\tau \bigl(1 - \Omega_n(r\rho)\bigr)\,\frac{dr}{r} \qquad (\rho \ge 0) ,
\]
so that $\xi \mapsto G_n(|\xi|)$, $\xi \in \RR^n$, is the exponent of the step datum
$(0, b\,\mathbf 1_{(0,\tau)})$ in dimension $n$ (\ref{cor:dictionary-filtration}(2), in whose
form the substitution $r = \tau u/\rho$ puts the integral). Write
$G := G_d$. It is the function of \ref{prop:profile-one-dimension-down}(2), and by that clause
with the nesting of the classes, $G \in \mathcal{A}_D$ exactly when $D \le d$. Say that a
function $\psi \ge 0$ on $\RR^D$ is \emph{the exponent of an infinitely divisible law} when
$e^{-\psi}$ is the characteristic function of an infinitely divisible probability measure on
$\RR^D$.
\begin{enumerate}
\item \emph{(Exactly two dimensions more.)} Let $D \ge 1$ be an integer. The function
$\zeta \mapsto G(|\zeta|)$, $\zeta \in \RR^D$, is the exponent of an infinitely divisible law on
$\RR^D$ if and only if $D \le d + 2$.
\item \emph{(Dimension $d + 2$.)} For every $\rho \ge 0$,
\begin{equation}\label{eq:step-exponent-two-up}
  G_d(\rho) = G_{d+2}(\rho) + \frac{b}{d}\bigl(1 - \Omega_{d+2}(\tau\rho)\bigr) .
\end{equation}
So the infinitely divisible law on $\RR^{d+2}$ with exponent $\zeta \mapsto G(|\zeta|)$ is the
convolution of the law of the step datum $(0, b\,\mathbf 1_{(0,\tau)})$ of dimension $d + 2$ with
the compound Poisson law $e^{-b/d}\sum_{n \ge 0}(b/d)^n\sigma_\tau^{*n}/n!$, where $\sigma_\tau$
is the uniform probability measure on the sphere of radius $\tau$ of $\RR^{d+2}$: jumps at the
rate $b/d$, uniformly distributed on that sphere. Its generating triplet is
$(0,\ \nu^{\mathrm{step}}_{d+2} + \tfrac{b}{d}\,\sigma_\tau,\ 0)$, with $\nu^{\mathrm{step}}_{d+2}$
the \Levy{} measure of that step datum.
\item \emph{(Dimension $d + 1$.)} The infinitely divisible law on $\RR^{d+1}$ with exponent
$\zeta \mapsto G(|\zeta|)$ has the generating triplet $(0, \nu', 0)$, where $\nu'$ is the
rotation-invariant measure on $\RR^{d+1} \setminus \{0\}$ with the density
$k'(|z|)\,(|S^{d}|\,|z|^{d+1})^{-1}$ and
\[
  k'(s) := 2b\,C_{d+1}\Bigl[J_{d+2}\Bigl(\frac{\tau}{s}\Bigr)
     + \frac{1}{d}\Bigl(\frac{s}{\tau}\Bigr)^{d}\frac{s}{\sqrt{\tau^2 - s^2}}\Bigr]
  \quad (0 < s < \tau), \qquad k'(s) := 0 \quad (s \ge \tau) ,
\]
with $J_n(x) := \int_1^x v^{1-n}(v^2-1)^{-1/2}\,dv$ for $x \ge 1$ and $C_n = |S^{n-1}|/|S^{n}|$
as in \ref{lem:shadow-polar-profile}. The function $k'$ is positive on $(0,\tau)$,
$k'(s) \to b$ as $s \downarrow 0$, and $k'(s) \to \infty$ as $s \uparrow \tau$. For $d = 1$,
$k'(s) = b\,\tau/\sqrt{\tau^2 - s^2}$, and for $d = 2$,
$k'(s) = (2b/\pi)\bigl[\arccos(s/\tau) + s/\sqrt{\tau^2 - s^2}\bigr]$, on $(0,\tau)$.
\end{enumerate}
In particular a $\Cin$ ray, $G(r) = b\,\Cin(\tau r)$, which is the case $d = 1$
(\ref{prop:profile-one-dimension-down}(3)), is the exponent of an infinitely divisible law in the
plane and in space, of a self-decomposable law in neither, and of no infinitely divisible law on
$\RR^4$.
\statusA\quad\ledger{A2}\quad\ledger{A3}\quad\ledger{A4}\quad\emph{(\textbf{Assignment.}
\ledger{A3} is spent in its three clauses: (iii) for the law of a triplet, in clauses (2) and (3)
and in both directions of clause (1); (i) and (ii) in the ``only if'' of clause (1), for the
triplet of a law given infinitely divisible and for the identification of a \Levy{} measure; (ii)
again for ``its generating triplet is'' in clauses (2) and (3). \ledger{A2} and \ledger{A4} are
not spent by the three clauses. They enter through \ref{prop:profile-one-dimension-down}, in the
sentence recalled before the clauses and in the words ``of a self-decomposable law in neither''
of the last sentence. \textbf{What the citations carry and what they do not.} \ledger{A3}
carries the correspondence between infinitely divisible laws and triplets. The identity
\ref{eq:step-exponent-two-up}, the profile $k'$ and the exclusion of dimension $d + 3$ are
proved here, from \ref{lem:shadow-polar-profile} and the preamble of
\ref{prop:radial-dictionary}. \textbf{What it is.} The step exponent of dimension $d$ is still
the exponent of an infinitely divisible law one and two dimensions up, where it is no longer
that of a self-decomposable one: in those two dimensions $e^{-G(|\zeta|)}$ is a probability
kernel, with every convolution root a probability kernel, and what is lost is the cascade. Three
dimensions up it is the exponent of no infinitely divisible law. The reason is visible in clause
(2): two dimensions up the \Levy{} measure has a part carried by a sphere, and a
rotation-invariant measure one dimension further up cannot project onto a measure that charges a
sphere. \textbf{What the identity says.} With \ref{eq:step-exponent-two-up} the proof obtains
$1 - \Omega_{d+2}(t) = d\int_0^1(1 - \Omega_d(tu))\,u^{d-1}\,du$. By the preamble of
\ref{prop:radial-dictionary} the right side belongs to the uniform law on the unit ball of
$\RR^d$: that law has the radial characteristic function $\Omega_{d+2}$, the one of the uniform
law on the unit sphere of $\RR^{d+2}$. For $d = 1$ this is $\int_0^1\cos(tu)\,du = \sin t/t$.
The identity is equivalent to $\Omega_d(t) = \Omega_{d+2}(t) + (t/d)\,\Omega_{d+2}'(t)$, the form
in which the verification record has it. \textbf{For $d = 1$} the plane and the space were
observed numerically in the first review round
(\texttt{records/affine/reviews/VERIFY-round1-filtration-map-blind-check.md}, \S 5): $\Cin(|\zeta|)$
is the exponent of the \Levy{} measure with polar profile $(1 - s^2)^{-1/2}$ on the unit disc,
and of the radial measure $\mathbf 1_{(0,1)}\,dr/r$ with a unit mass on the unit sphere in
$\RR^3$. These are clauses (3) and (2) at $d = 1$, $b = \tau = 1$. \textbf{Checks.}
\texttt{records/affine/reviews/verify-round2-scripts/item13\_step\_exponent\_dimensions.py}:
the identity, the two closed forms, the positivity of $k'$, its shadow, and two tests of
dimension $d + 3$ that do not use the proof. \texttt{records/affine/q0434-scripts/check\_q0434\_nodes.py}:
the shadow two dimensions down of the step, the shadow of the second term of $k'$, and the
identity in the form the proof uses, for $d = 1,\dots,4$. \textbf{Proved in print.} The node has
no Lean: $\Omega_d$ is defined in \ref{prop:radial-dictionary}, which is proved in print and not
formalized, and the node stands with the readings in the dictionary for that reason. \textbf{Not
claimed}: that $\zeta \mapsto e^{-G(|\zeta|)}$ fails to be a characteristic function on $\RR^D$
for $D \ge d + 3$. Clause (1) excludes an infinitely divisible law with that transform, not every
law. For every sufficiently small height $b$ the function is not a characteristic function there:
this follows from clause (1) and the continuity theorem for characteristic functions in the form
that produces the limit law, a theorem that is not among this chapter's cited facts, so it is not
a clause (the argument is in the verification record, item 13). Whether it is a characteristic
function for some large $b$ is open. Nor is it claimed that $k'$ is monotone, nor anything for a
datum whose profile is not a step.)}
\end{proposition}

\begin{proof}
\emph{Four facts.} (a) For $n \ge 1$ and $\zeta \in \RR^n$, \ref{eq:sphere-mean-cos} with
$t = \tau$ and $x = \zeta$ says that $\zeta \mapsto \Omega_n(\tau|\zeta|)$ is the characteristic
function of the uniform probability measure on the sphere of radius $\tau$ of $\RR^n$, the image
of $\sigma_{n-1}/|S^{n-1}|$ under $\theta \mapsto \tau\theta$. (b) $\Omega_n$ is continuous and
$0 \le 1 - \Omega_n(t) \le t^2/2$, because $1 - \cos u \le u^2/2$ and
$|\langle e,\eta\rangle| \le 1$ on the sphere. (c) Let $(0,k)$ be an isotropic origin datum in
dimension $n$. By \ref{eq:polar-decomposition} with $g(x) = 1 - \cos\langle\xi,x\rangle$, by
Tonelli's theorem for a nonnegative integrand, and by the definition of $\Omega_n$ with
$e = \xi/|\xi|$, its exponent at $\xi \ne 0$ is
\[
  \frac{1}{|S^{n-1}|}\int_{S^{n-1}}\sigma_{n-1}(d\eta)\int_0^\infty
     \bigl(1 - \cos(r\langle\xi,\eta\rangle)\bigr)\,k(r)\,\frac{dr}{r}
   = \int_0^\infty \bigl(1 - \Omega_n(r|\xi|)\bigr)\,k(r)\,\frac{dr}{r} ,
\]
and both sides vanish at $\xi = 0$. For $k = b\,\mathbf 1_{(0,\tau)}$ this is $G_n(|\xi|)$, which
is finite by (b). (d) Let $\nu$ be a symmetric \Levy{} measure on $\RR^n$ and $v \ge 0$. The
infinitely divisible law with triplet $(vI, \nu, 0)$ that \ledger{A3}(iii) supplies has the
transform $\exp[-\tfrac12 v|\zeta|^2 - \int(1 - \cos\langle\zeta,z\rangle)\,\nu(dz)]$. Indeed
the integrand $e^{i\langle\zeta,z\rangle} - 1 - i\langle\zeta,z\rangle\mathbf 1_D(z)$ of the
representation has the real part $\cos\langle\zeta,z\rangle - 1$, and its imaginary part is odd
in $z$ and integrable against $\nu$, being at most $|\zeta|^2|z|^2/2$ in modulus on $D$ (as
$|\sin u - u| \le u^2/2$ for every real $u$) and at most $1$ off $D$; so the integral of the
imaginary part against the symmetric measure $\nu$ is $0$.

\emph{The shadow of the step two dimensions down.} The step datum
$(0, b\,\mathbf 1_{(0,\tau)})$ of dimension $d + 2$ has the exponent
$\zeta \mapsto G_{d+2}(|\zeta|)$, by (c). By \ref{lem:shadow-polar-profile} with $d + 1$ in place
of $d$, its shadow $(0,k_1)$ is an isotropic origin datum in dimension $d + 1$ with exponent
$\zeta \mapsto G_{d+2}(|\zeta|)$, $\zeta \in \RR^{d+1}$, and
\[
  k_1(s) = C_{d+1}\int_{\RR} b\,\mathbf 1_{\{s\sqrt{1+w^2} < \tau\}}\,(1+w^2)^{-(d+2)/2}\,dw
         = 2b\,C_{d+1}\,J_{d+2}\Bigl(\frac{\tau}{s}\Bigr) \quad (0 < s < \tau),
\]
by the substitution $v = \sqrt{1+w^2}$ on $w > 0$, for which $dw = v\,(v^2-1)^{-1/2}dv$; and
$k_1(s) = 0$ for $s \ge \tau$. By the same lemma, the shadow $(0,k_2)$ of $(0,k_1)$ is an
isotropic origin datum in dimension $d$ with exponent $\xi \mapsto G_{d+2}(|\xi|)$,
$\xi \in \RR^d$, and
\[
  k_2(r) = C_d\,C_{d+1}\,b\int_{\RR}\int_{\RR}
     \mathbf 1_{\{r\sqrt{(1+u^2)(1+w^2)} < \tau\}}\,(1+w^2)^{-(d+2)/2}\,(1+u^2)^{-(d+1)/2}\,dw\,du .
\]
Substitute $w = y/\sqrt{1+u^2}$ in the inner integral. Then $(1+u^2)(1+w^2) = 1 + u^2 + y^2$ and
$(1+w^2)^{-(d+2)/2}(1+u^2)^{-(d+1)/2}\,dw = (1+u^2+y^2)^{-(d+2)/2}\,dy$, so with $q = (u,y)$
\[
  k_2(r) = C_d\,C_{d+1}\,b\int_{\RR^2}\mathbf 1_{\{r\sqrt{1+|q|^2} < \tau\}}\,
     (1+|q|^2)^{-(d+2)/2}\,dq .
\]
The same computation with the indicator replaced by $1$, and the normalisation of
\ref{lem:shadow-polar-profile}(1) in the dimensions $d$ and $d + 1$, give
$C_d\,C_{d+1}\int_{\RR^2}(1+|q|^2)^{-(d+2)/2}\,dq = 1$. In polar coordinates of the plane, with
$v = \sqrt{1+|q|^2}$, the integral of $\mathbf 1_{\{v < x\}}(1+|q|^2)^{-(d+2)/2}$ is
$2\pi\int_1^x v^{-d-1}\,dv = (2\pi/d)(1 - x^{-d})$ for $x \ge 1$, and $0$ for $x < 1$. At
$x = \infty$ this gives
\begin{equation}\label{eq:shadow-constants}
  2\pi\,C_d\,C_{d+1} = d ,
\end{equation}
and at $x = \tau/r$ it gives
\[
  k_2(r) = b\Bigl(1 - \Bigl(\frac{r}{\tau}\Bigr)^{d}\Bigr) \quad (0 < r < \tau), \qquad
  k_2(r) = 0 \quad (r \ge \tau) .
\]

\emph{Clause (2), the identity.} Fix $\rho \ge 0$ and put $h(r) := 1 - \Omega_d(r\rho)$ and
$h_2(r) := 1 - \Omega_{d+2}(r\rho)$. By (c) for the datum $(0,k_2)$, whose exponent is
$G_{d+2}(|\cdot|)$, and by the definition of $G_{d+2}$,
\begin{equation}\label{eq:step-two-forms}
  \int_0^\tau h_2(r)\,\frac{dr}{r} = \frac{G_{d+2}(\rho)}{b}
   = \int_0^\tau h(r)\,\frac{dr}{r} - \tau^{-d}\int_0^\tau h(r)\,r^{d-1}\,dr .
\end{equation}
This holds for every $\tau > 0$, the parameter $\tau$ having been arbitrary. By (b) the
functions $h$ and $h_2$ are continuous with $0 \le h(r), h_2(r) \le r^2\rho^2/2$, so the three
integrals are continuously differentiable functions of $\tau > 0$. Differentiating
\ref{eq:step-two-forms} in $\tau$,
\[
  \frac{h_2(\tau)}{\tau} = \frac{h(\tau)}{\tau} + d\,\tau^{-d-1}\int_0^\tau h(r)\,r^{d-1}\,dr
     - \frac{h(\tau)}{\tau} ,
  \qquad\text{that is}\qquad
  \tau^{-d}\int_0^\tau h(r)\,r^{d-1}\,dr = \frac{1}{d}\bigl(1 - \Omega_{d+2}(\tau\rho)\bigr) .
\]
By the definition of $G_d$ and \ref{eq:step-two-forms},
$G_d(\rho) - G_{d+2}(\rho) = b\,\tau^{-d}\int_0^\tau h(r)\,r^{d-1}\,dr$, which is
\ref{eq:step-exponent-two-up}.

\emph{Clause (2), the law.} Let $\mu^{\mathrm{step}}$ be the law of the step datum of dimension
$d + 2$, the infinitely divisible law with triplet $(0,\nu^{\mathrm{step}}_{d+2},0)$ that
\ledger{A3}(iii) supplies; its transform is $e^{-G_{d+2}(|\zeta|)}$, by (c) and (d). Let
$\pi := e^{-b/d}\sum_{n \ge 0}(b/d)^n\sigma_\tau^{*n}/n!$. It is a probability measure, and by
(a) its transform is
$e^{-b/d}\sum_n\bigl((b/d)\,\Omega_{d+2}(\tau|\zeta|)\bigr)^n/n!
= \exp\bigl[-\tfrac{b}{d}(1 - \Omega_{d+2}(\tau|\zeta|))\bigr]$. It is infinitely divisible: for
$m \ge 1$ the law defined in the same way with $b/(dm)$ in place of $b/d$ has the transform
$\widehat\pi^{1/m}$, so its $m$-fold convolution is $\pi$. The convolution of two infinitely
divisible laws is infinitely divisible, the convolution of their roots being a root. So
$\mu^{\mathrm{step}} * \pi$ is infinitely divisible, and by \ref{eq:step-exponent-two-up} its
transform is $e^{-G_d(|\zeta|)}$. Put
$\nu^{(d+2)} := \nu^{\mathrm{step}}_{d+2} + \tfrac{b}{d}\,\sigma_\tau$. It is the sum of a
\Levy{} measure and a finite measure that gives the origin no mass, so it is a \Levy{} measure,
and it is symmetric. By (a) and (c),
$\int(1 - \cos\langle\zeta,z\rangle)\,\nu^{(d+2)}(dz) = G_{d+2}(|\zeta|)
+ \tfrac{b}{d}(1 - \Omega_{d+2}(\tau|\zeta|)) = G_d(|\zeta|)$. The infinitely divisible law with
triplet $(0,\nu^{(d+2)},0)$, which \ledger{A3}(iii) supplies, therefore has the transform
$e^{-G_d(|\zeta|)}$, by (d). So it is $\mu^{\mathrm{step}} * \pi$, and the triplet of a law is
unique (\ledger{A3}(ii)).

\emph{Clause (3).} Write $k' = k_1 + k_2'$ on $(0,\infty)$, with $k_1$ as above and
$k_2'(s) := (2b\,C_{d+1}/d)\,(s/\tau)^d\,s\,(\tau^2 - s^2)^{-1/2}$ for $0 < s < \tau$,
$k_2'(s) := 0$ for $s \ge \tau$. Both are nonnegative and Borel. The identity
\ref{eq:shadow-identity} holds for every nonnegative Borel profile $h$, monotone or not. For
$h = k_1$ its function $\bar h$ is $k_2$. For $h = k_2'$ and $0 < r < \tau$ put
$A := \sqrt{\tau^2 - r^2}$. The integrand of $\bar h(r)$ vanishes unless $|u| < A/r$, and there
$k_2'(r\sqrt{1+u^2})\,(1+u^2)^{-(d+1)/2} = (2b\,C_{d+1}/d)\,r^{d+1}\tau^{-d}\,(A^2 - r^2u^2)^{-1/2}$.
Since $\int_{-A/r}^{A/r}(A^2 - r^2u^2)^{-1/2}\,du = \pi/r$, and by \ref{eq:shadow-constants},
\[
  \bar h(r) = C_d\cdot\frac{2b\,C_{d+1}}{d}\cdot\frac{r^{d+1}}{\tau^{d}}\cdot\frac{\pi}{r}
            = b\Bigl(\frac{r}{\tau}\Bigr)^{d} \quad (0 < r < \tau), \qquad
  \bar h(r) = 0 \quad (r \ge \tau) .
\]
Adding the two, \ref{eq:shadow-identity} for $h = k'$ has $\bar h = k_2 + b\,(r/\tau)^d
\mathbf 1_{(0,\tau)} = b\,\mathbf 1_{(0,\tau)}$: for every nonnegative Borel $g$ on $\RR^d$
vanishing at the origin,
\begin{equation}\label{eq:step-one-up-image}
  \int_{\RR^{d+1}} g(\Pi z)\,\nu'(dz) = \int_{\RR^d} g\,d\nu^{\mathrm{step}}_d ,
\end{equation}
with $\Pi$ the projection on the first $d$ coordinates and $\nu^{\mathrm{step}}_d$ the \Levy{}
measure of the step datum of dimension $d$. The measure $\nu'$ is rotation invariant, having a
radial density, and gives the origin no mass. It is a \Levy{} measure: in polar coordinates of
$\RR^{d+1}$, $\int(|z|^2 \wedge 1)\,\nu'(dz) \le \int_0^\tau k'(s)\,s\,ds$, and $k'$ is bounded
on $(0,\tau/2]$ and at most a constant times $(\tau - s)^{-1/2}$ on $[\tau/2,\tau)$. Put
$\psi(\zeta) := \int(1 - \cos\langle\zeta,z\rangle)\,\nu'(dz)$ for $\zeta \in \RR^{d+1}$, which
is finite. It depends on $|\zeta|$ only, by the invariance of $\nu'$. For $\xi \in \RR^d$,
\ref{eq:step-one-up-image} with $g(x) = 1 - \cos\langle\xi,x\rangle$ gives
$\psi((\xi,0)) = G_d(|\xi|)$, by (c). So $\psi(\zeta) = G_d(|\zeta|)$ for every
$\zeta \in \RR^{d+1}$. The infinitely divisible law with triplet $(0,\nu',0)$, which
\ledger{A3}(iii) supplies, has the transform $e^{-\psi}$, by (d), and its triplet is unique (\ledger{A3}(ii)). The properties of $k'$: both terms are
positive on $(0,\tau)$, since $J_{d+2}(x) > 0$ for $x > 1$. As $s \downarrow 0$ the second term
tends to $0$, and $2C_{d+1}J_{d+2}(\tau/s) \to 2C_{d+1}J_{d+2}(\infty)
= C_{d+1}\int_{\RR}(1+w^2)^{-(d+2)/2}\,dw = 1$, by \ref{lem:shadow-polar-profile}(1) in dimension
$d + 1$. As $s \uparrow \tau$ the first term tends to $0$ and the second to $\infty$. The closed
forms. By \ref{lem:shadow-polar-profile}(1), $C_2 = (\int_{\RR}(1+u^2)^{-3/2}du)^{-1} = \tfrac12$
and $C_3 = (\int_{\RR}(1+u^2)^{-2}du)^{-1} = 2/\pi$. Differentiation shows
$J_3(x) = \sqrt{x^2-1}/x$ and $J_4(x) = \tfrac12\bigl[\sqrt{x^2-1}/x^2 + \arccos(1/x)\bigr]$,
both sides vanishing at $x = 1$. For $d = 1$ this gives
$k'(s) = b\bigl[\sqrt{\tau^2-s^2}/\tau + s^2/(\tau\sqrt{\tau^2-s^2})\bigr]
= b\,\tau/\sqrt{\tau^2-s^2}$, and for $d = 2$,
$k'(s) = (2b/\pi)\bigl[s\sqrt{\tau^2-s^2}/\tau^2 + \arccos(s/\tau)
+ s^3/(\tau^2\sqrt{\tau^2-s^2})\bigr] = (2b/\pi)\bigl[\arccos(s/\tau) + s/\sqrt{\tau^2-s^2}\bigr]$.

\emph{Clause (1), ``if''.} Let $D \le d + 2$, let $\mu$ be the law of clause (2), and let $P$ be
the projection of $\RR^{d+2}$ on its first $D$ coordinates. The image $P_*\mu$ has the transform
$\zeta \mapsto \widehat\mu((\zeta,0)) = e^{-G_d(|\zeta|)}$, $\zeta \in \RR^D$. It is infinitely
divisible: the image of a convolution under a linear map is the convolution of the images, so
the image of an $m$-th convolution root of $\mu$ is one of $P_*\mu$.

\emph{Clause (1), ``only if''.} Let $D \ge d + 3$ and suppose that an infinitely divisible law on
$\RR^D$ has the transform $e^{-G_d(|\zeta|)}$. Its image under the projection on the first
$d + 3$ coordinates is infinitely divisible with the transform $e^{-G_d(|\zeta|)}$ on
$\RR^{d+3}$, as in the previous paragraph. So let $D = d + 3$, and let $\mu'$ be such a law. By
\ledger{A3}(i) it has a triplet $(A,\nu',\gamma)$. Its transform is invariant under $\Orth(D)$,
so, as in the first paragraph of the proof of \ref{prop:isotropic-polar-reading} from its
third sentence on, which reads of
the law only that it is infinitely divisible with a radial transform, \ledger{A3}(ii) gives
$\langle\zeta,A\zeta\rangle = \alpha|\zeta|^2$ with $\alpha \ge 0$ and $Q_*\nu' = \nu'$ for every
$Q \in \Orth(D)$, and
\[
  G_d(|\zeta|) = \tfrac12\alpha|\zeta|^2
     + \int_{\RR^{d+3}}\bigl(1 - \cos\langle\zeta,z\rangle\bigr)\,\nu'(dz), \qquad \zeta \in \RR^{d+3} .
\]
Only the modulus of the representation is used for this display: $\mu'$ is not assumed
symmetric, and $-\log|\widehat{\mu'}(\zeta)| = G_d(|\zeta|)$ because $e^{-G_d}$ is real and
positive.
Let $\Pi$ be the projection of $\RR^{d+3}$ on its first $d + 2$ coordinates, and let $\bar\nu$ be
the image of $\nu'$ under $\Pi$, restricted to $\RR^{d+2} \setminus \{0\}$. It is symmetric, and
it is a \Levy{} measure, because $|\Pi z| \le |z|$ gives
$\int(|y|^2 \wedge 1)\,\bar\nu(dy) \le \int(|z|^2 \wedge 1)\,\nu'(dz) < \infty$. For
$\eta \in \RR^{d+2}$, $\langle(\eta,0),z\rangle = \langle\eta,\Pi z\rangle$, and the integrand
vanishes where $\Pi z = 0$, so
$G_d(|\eta|) = \tfrac12\alpha|\eta|^2 + \int(1 - \cos\langle\eta,y\rangle)\,\bar\nu(dy)$. By
\ledger{A3}(iii) there is an infinitely divisible law on $\RR^{d+2}$ with triplet
$(\alpha I,\bar\nu,0)$, and by (d) its transform is $e^{-G_d(|\eta|)}$. The law of clause (2) has the
same transform and the triplet $(0,\nu^{(d+2)},0)$. The two laws are equal, and by
\ledger{A3}(ii) $\alpha = 0$ and $\bar\nu = \nu^{(d+2)}$. Let
$E := \{y \in \RR^{d+2} : |y| = \tau\}$. Then
$\bar\nu(E) = \nu^{(d+2)}(E) \ge \tfrac{b}{d}\,\sigma_\tau(E) = \tfrac{b}{d} > 0$.

On the other hand $\bar\nu(E) = \nu'(\Pi^{-1}E)$, and
$\Pi^{-1}E = \{z \in \RR^{d+3} : |\Pi z| = \tau\}$ lies in $\{|z| \ge \tau\}$. Let $m$ be the
restriction of $\nu'$ to $\{|z| \ge \tau\}$. It is a finite measure, because
$\int(|z|^2 \wedge 1)\,d\nu' < \infty$, it is invariant under $\Orth(d+3)$, and
$m(\Pi^{-1}E) = \bar\nu(E) > 0$, so $m \ne 0$. Let $Z$ be a random vector in $\RR^{d+3}$ with the
law $m/m(\RR^{d+3})$. It is rotation invariant, so by the preamble of
\ref{prop:radial-dictionary} it has the law of $R\,U$, with $R$ of the law of $|Z|$, so that
$R \ge \tau$, and $U$ uniformly distributed on $S^{d+2}$ and independent of $R$. Hence
\[
  \frac{m(\Pi^{-1}E)}{m(\RR^{d+3})} = \Pr\bigl(R\,|\Pi U| = \tau\bigr)
   = \int \Pr\bigl(|\Pi U| = \tau/r\bigr)\,\Pr(R \in dr) .
\]
For every $w \in [0,1]$, $\Pr(|\Pi U| = w) = 0$. Indeed $|\Pi u|^2 = 1 - u_{d+3}^2$ on the
sphere, so the event is $\{U \in K\}$ with $K := \{z \in \RR^{d+3} : z_{d+3}^2 = (1-w^2)|z|^2\}$,
a cone. The cone is Lebesgue-null: if $w > 0$ it is the set
$\{w^2z_{d+3}^2 = (1-w^2)|\Pi z|^2\}$, which meets every line parallel to the last axis in at
most two points, and if $w = 0$ it is the last axis. In polar coordinates the Lebesgue measure
of $K \cap \{1 < |z| < 2\}$ is $\sigma_{d+2}(K \cap S^{d+2})\int_1^2 r^{d+2}\,dr$, so
$\sigma_{d+2}(K \cap S^{d+2}) = 0$. Therefore $m(\Pi^{-1}E) = 0$, a contradiction.

\emph{The last sentence.} For $d = 1$, $G(r) = b\,\Cin(\tau r)$ and the ray is in no
$\mathcal{A}_D$ with $D \ge 2$, by \ref{prop:profile-one-dimension-down}(3). The rest is clause
(1) with $d + 2 = 3$.

Nothing in this proof is machine-checked. It uses \ref{lem:shadow-polar-profile} and the identity
\ref{eq:shadow-identity} of its proof, which are machine-checked, the latter for every
nonnegative Borel profile; \ref{eq:polar-decomposition}; the preamble of
\ref{prop:radial-dictionary}, which is proved in print; and \ledger{A3} in its three clauses.
\end{proof}

% not in the draft: Q-0429 (hub), first part (2026-10-12). One theorem for the separating
% family. Provenance: the blind register review of section 3.4 of the isotropic paper of that day
% (flags 3, 10 and 33) found that the paper's prose, abstract and table quote, as the content of
% thm:filtration-strict, two statements which that theorem does not make: "F^(w) is admitted in
% dimension D exactly when w <= w_D", which was the first sentence of
% prop:separating-family-mixtures(3), and "sqrt(d) w_d -> sqrt(pi)", which is
% lem:threshold-asymptotic(3). This node states them together, with the strict decrease
% (lem:threshold-monotone) and the dimensions in which a member is admitted.
% thm:filtration-strict and thm:no-cin-ray-in-the-plane are not touched: they are machine-checked,
% the papers print them, and other nodes use them.
% The membership statement is MERGED into this theorem, not kept as a node of its own (the brief
% of Q-0429 left the choice). Reason: a separate proposition "F^(w) in A_D iff w <= w_D, hence for
% w in (w_{d+1}, w_d] exactly the dimensions D <= d, for w > w_1 none" would be clauses (1) and
% (4) of this theorem stated twice, which is the duplication the review flags between
% thm:no-cin-ray-in-the-plane and the corollary's third clause. Merged, the criterion has one
% home, the theorem's proof of record is the refereed proof of it, and the Lean is one assembly
% either way.
% The proof is cut from the proof of prop:separating-family-mixtures(3) as the blind check of
% Q-0425 read it (report B, "correct as stated"): the paragraphs "w <= w_D gives membership" and
% "w > w_D excludes it" word for word, with "Clause (3)" become "Clause (1)" in their headings,
% and the paragraph "the windows" word for word in clause (4). New text: the two lines for
% clauses (2) and (3), the three closing sentences of clause (4), and the closing paragraph.
% CORRECTED after the blind check of Q-0429 (its report, F-B1): the paragraph "w > w_D excludes
% it" is word for word against the text before the split, not against the text the blind check
% of Q-0425 read. There the finiteness of the integral of |x|^2 ^ 1 for nu' was drawn from "the
% display and the two integrability conditions"; since that check's repairs the sentence has the
% computation in polar coordinates, with the display of the two integrals. The blind check of
% Q-0429 read it and found it correct. The same check (F-B3, F-B6): the last sentence of clause
% (4) cited [B, Prop. 4.9](1) "proved in the proof of thm:filtration-strict(3)"; it now cites the
% statement of thm:filtration-strict(3), which has the inclusion and which the node uses.
% No \lean tag yet: the Lean is the second part of Q-0429, an assembly. The library has
% witness_classA and witness_not_classA (FiltrationStrict.lean: the two directions of clause (1),
% each with one inequality), threshold_strictAnti (clause (2)), threshold_asymptotic (clause (3),
% along the real parameter), threshold_limit (w_d -> 0, for "exactly one interval") and
% filtration_strict (the strict inclusion). Priced at the target types
% (records/affine/NODES-q0425.md, "After the split"): about 140 lines, clauses (1) and (4) on A2,
% A3, A4, clauses (2) and (3) on Lean core.
% CHANGED Q-0429, part 2: typed and proved (Formalization/AffineHemigroup/SeparatingFamily.lean,
% an assembly; clauses (1) and (4) on A2, A3, A4, clauses (2) and (3) and the sentence on the
% intervals on Lean core). The statement is unchanged. Added: the \lean and \leanok tags, the
% passage "Formalization" of the annotation, and the closing paragraph of the proof.
\begin{theorem}[the separating family and its critical values]
\label{thm:separating-family-critical-values}
\uses{def:filtration-thresholds,def:isotropic-origin-datum,lem:threshold-monotone,lem:threshold-criterion,lem:threshold-limit,lem:threshold-asymptotic,lem:delay-polar-profile,prop:isotropic-polar-classification,thm:filtration-strict,thm:cited-A2,thm:cited-A3,thm:cited-A4}
\lean{AffineHemigroup.separating_family_critical_values, AffineHemigroup.witness_classA_iff,
AffineHemigroup.criticalValues_strictAnti, AffineHemigroup.criticalValues_tendsto,
AffineHemigroup.witness_classA_iff_le, AffineHemigroup.witness_not_classA_of_gt,
AffineHemigroup.criticalValues_existsUnique_window,
AffineHemigroup.separating_family_strict}\leanok
Let $F^{(w)}$, $w > 0$, be the separating family of \ref{def:filtration-thresholds} and
$w_1, w_2, \dots$ its critical values.
\begin{enumerate}
\item \emph{(Membership.)} For every integer $D \ge 1$ and every $w > 0$,
$F^{(w)} \in \mathcal{A}_D$ if and only if $w \le w_D$.
\item \emph{(The critical values decrease strictly.)} $w_1 > w_2 > w_3 > \dots > 0$.
\item \emph{(Their limit.)} $\sqrt D\,w_D \to \sqrt\pi$ as $D \to \infty$.
\item \emph{(The dimensions in which a member is admitted.)} Let $w > 0$. If
$w \in (w_{d+1}, w_d\,]$ for an integer $d \ge 1$, then $F^{(w)} \in \mathcal{A}_D$ for
$D = 1, \dots, d$ and for no $D > d$. If $w > w_1$, then $F^{(w)}$ belongs to no
$\mathcal{A}_D$. Every $w \in (0, w_1]$ lies in exactly one of the intervals
$(w_{d+1}, w_d\,]$. Each of these intervals is nonempty, so
$\mathcal{A}_d \setminus \mathcal{A}_{d+1}$ contains a member of the family, and
$\mathcal{A}_{d+1} \subsetneq \mathcal{A}_d$, for every integer $d \ge 1$.
\end{enumerate}
\statusA\quad\ledger{A2}\quad\ledger{A3}\quad\ledger{A4}\quad\emph{(\textbf{Assignment.} Clause
(1) spends \ledger{A2}, \ledger{A3} and
\ledger{A4} through \ref{prop:isotropic-polar-classification}(2), in both directions, and
\ledger{A3}(iii) and (ii) directly, exactly as \ref{thm:filtration-strict}(1) does: its two
directions are the two steps of that theorem's proof, each of which reads one inequality.
Clause (4) spends them through clause (1) only. Clauses (2) and (3) cite nothing: they are
\ref{lem:threshold-monotone} and \ref{lem:threshold-asymptotic}(3). \textbf{What the citations
carry and what they do not.} They carry the polar form of
$\mathcal{A}_D$ and the uniqueness of a \Levy{} measure. The criterion $w \le w_D$, the strict
decrease and the limit are proved. \textbf{What it adds.} One statement for the family, from
which a reader decides in which dimensions a member is admitted: the position of $w$ among the
critical values. \ref{thm:filtration-strict}(1) has the two dimensions $d$ and $d + 1$ for
$w \in (w_{d+1}, w_d\,]$. Clause (1) is the uniform form with one inequality. It adds the case
$w > w_1$, where $F^{(w)}$ is not admissible at all, $\mathcal{A}_1$ being the whole admissible
cone [B, Prop.~4.9]; and with clause (4) it says that the dimensions in which $F^{(w)}$ is admitted are an initial
segment $1, \dots, d$, possibly empty, and never all of them. Clauses (2) and (3) say where the
critical values are. \textbf{How it stands to the other nodes.} \ref{thm:filtration-strict} is
unchanged. Its clause (1) is clause (1) here at $D = d$ and $D = d + 1$. Its clause (3) has the
last assertion of clause (4), and adds the comparison with $\mathcal{A}_\infty$ and the
identification $\mathcal{A}_\infty = \mathcal{S}$. Its clause (2), on the bound $6 + 4\sqrt2$ of
[B, Lem.~4.12(1)], is not read here. Until Q-0429 clause (1) was the first sentence of
\ref{prop:separating-family-mixtures}(3), and the first three sentences of clause (4) were its
consequences, stated there for the laws. That proposition keeps the description of the family as
laws and reads this theorem. \ref{prop:profile-one-dimension-down}(2) separates the same classes
with another witness, the step datum, and without critical values. \textbf{Numbers (computed,
not claimed, and not used).} $w_1 = 12.606$, $w_2 = 5.242$, $w_3 = 3.206$
(\texttt{records/affine/reviews/verify-round1-scripts/a\_critical\_values.py}; [B, Rem.~4.10]
reports $12.61$, $5.24$, $3.21$). With these values $w = 8$ is admitted in dimension $1$ only
and $w = 4$ in dimensions $1$ and $2$ only, in agreement with [B, Lem.~4.12(3)], which has these
two parameter points. The limit of clause
(3) is approached slowly and not monotonically, and at small $D$ the asymptote $\sqrt{\pi/D}$ is
not an approximation of $w_D$: $\sqrt D\,w_D$ is $12.61$, $7.41$ and $5.55$ at $D = 1, 2, 3$,
against $\sqrt\pi = 1.772$, so that $w_2$ and $w_3$ are $4.2$ and $3.1$ times
$\sqrt{\pi/D}$. It is $2.056$ at $D = 20$ and $1.778$ at $D = 50$, falls below $\sqrt\pi$
between $D = 55$ and $D = 56$, is least at $D = 81$, where it is $1.7661$, and returns to
$\sqrt\pi$ from below, as far as computed; the annotation of \ref{lem:threshold-asymptotic} has
the table. Bounds that hold at each $D$ are \ref{lem:threshold-limit} and
\ref{lem:threshold-asymptotic}(2). \textbf{Readings.} $D$ and $d$ are integers, and $w > 0$ is
the range for which \ref{def:filtration-thresholds} defines the family. The limit of clause (3)
is along the integers; \ref{lem:threshold-asymptotic}(3) has it along the real parameter of
\ref{lem:threshold-monotone}. The first sentence of clause (4) says that for every integer
$D \ge 1$, $F^{(w)} \in \mathcal{A}_D$ if and only if $D \le d$. \textbf{Not claimed}: the value
of any critical value; a rate of convergence in clause (3), or any monotonicity of
$\sqrt D\,w_D$; which other members of $\mathcal{A}_d$ separate it from $\mathcal{A}_{d+1}$;
anything about the laws with transform $e^{-F^{(w)}(|\xi|)}$, which is
\ref{prop:separating-family-mixtures}. \textbf{Formalization} (Q-0429). $F^{(w)}$ and
$\mathcal{A}_D$ are those of \ref{thm:filtration-strict}'s Lean. $D$ and $d$ are natural numbers
with $D, d \ge 1$. $w_D$ is the function of the real parameter of \ref{lem:threshold-monotone}
at the integer $D$, and $w_{d+1}$ is that function at the integer $d + 1$. Clause (1) is one
equivalence, for every $w > 0$. Clause (2) says that $n \mapsto w_n$ is strictly decreasing on
the integers $n \ge 1$ and positive there. Clause (3) is a limit along the natural numbers.
Clause (4) is four declarations. The first is the equivalence that the Readings give: for
$w_{d+1} < w \le w_d$ and every integer $D \ge 1$, $F^{(w)} \in \mathcal{A}_D$ if and only if
$D \le d$. The second is for $w > w_1$ and every integer $D \ge 1$. The third says that for
$0 < w \le w_1$ there is exactly one integer $d \ge 1$ with $w_{d+1} < w \le w_d$. The fourth
has, for every $d \ge 1$, the interval nonempty, a parameter $w > 0$ with
$F^{(w)} \in \mathcal{A}_d \setminus \mathcal{A}_{d+1}$, and the strict inclusion of the two sets
of functions. One further declaration states the four clauses as a single proposition. Probes:
clause (3) read as $D\,w_D^2 \to \pi$; $w_3 < 4$, proved from $M_3 \ge h_3(1)$ with $h_3$ the
function of \ref{lem:threshold-monotone}, so that $F^{(4)}$ is in no $\mathcal{A}_D$ with
$D \ge 3$; and the member $F^{(w_d)}$ at the right end of each window, admitted in the
dimensions $1, \dots, d$ only. The other half of the parameter point $w = 4$, that $4 \le w_2$,
needs an upper bound of $M_2$, which the Lean does not have. So of the computed values above
only $w_3 < 4$ is machine-checked, beside $w_2 < 6 + 4\sqrt2$ of \ref{thm:filtration-strict}(2)
and the bounds of \ref{lem:threshold-limit} and \ref{lem:threshold-asymptotic}(2), which hold at
each $D$.)}
\end{theorem}

\begin{proof}
\leanok
\emph{Clause (1), $w \le w_D$ gives membership.} \ref{lem:delay-polar-profile} applies to
$D^{(w)}$ in dimension $D$, since it reads no monotonicity. It gives that the polar profile
$k^{(w)}_D$ is finite and satisfies the two conditions of \ref{def:isotropic-origin-datum}, and
that
\[
  F^{(w)}(|\xi|) = \int_{\RR^D}\bigl(1 - \cos\langle\xi,x\rangle\bigr)\,
     \frac{k^{(w)}_D(|x|)}{|S^{D-1}|\,|x|^D}\,dx , \qquad \xi \in \RR^D .
\]
If $w \le w_D$, $k^{(w)}_D$ is nonincreasing (\ref{lem:threshold-criterion}). So
$(0,k^{(w)}_D)$ is an isotropic origin datum in dimension $D$ with exponent
$\xi \mapsto F^{(w)}(|\xi|)$, and $F^{(w)} \in \mathcal{A}_D$ by
\ref{prop:isotropic-polar-classification}(2).

\emph{Clause (1), $w > w_D$ excludes it.} Let $w > w_D$ and suppose
$F^{(w)} \in \mathcal{A}_D$. By \ref{prop:isotropic-polar-classification}(2),
$\xi \mapsto F^{(w)}(|\xi|)$ is the exponent of an isotropic origin datum $(a,k_D)$ in dimension
$D$, with \Levy{} measure $\nu_D$ and $k_D$ nonincreasing. Let $\nu'$ be the measure on
$\RR^D \setminus \{0\}$ with density $k^{(w)}_D(|x|)\,(|S^{D-1}|\,|x|^D)^{-1}$. By polar
coordinates in $\RR^D$, which need no monotonicity of $k^{(w)}_D$, and the two integrability
conditions,
\[
  \int\bigl(|x|^2 \wedge 1\bigr)\,\nu'(dx)
   = \int_0^1 k^{(w)}_D(r)\,r\,dr + \int_1^\infty k^{(w)}_D(r)\,\frac{dr}{r} < \infty ,
\]
and by the display $F^{(w)}(|\xi|) = \int(1 - \cos\langle\xi,x\rangle)\,\nu'(dx)$. By
\ledger{A3}(iii) there are
infinitely divisible laws with the triplets $(0,\nu',0)$ and $(aI,\nu_D,0)$. Both \Levy{}
measures are symmetric, so the odd parts of the integrands integrate to $0$, and both laws have
the characteristic function $e^{-F^{(w)}(|\xi|)}$. So they are one law, and by the uniqueness
of the triplet, \ledger{A3}(ii), $a = 0$ and $\nu' = \nu_D$. The two radial densities then agree
almost everywhere on $\RR^D$, so $k^{(w)}_D = k_D$ almost everywhere on $(0,\infty)$, by polar
coordinates. On the other hand $M_D$ is attained at some $z^* > 0$
(\ref{lem:threshold-monotone}), and $w > w_D$ means $w\,M_D > \tfrac12\Gamma(m)$. So the
criterion of \ref{lem:threshold-criterion} fails at $z^*$: the smooth function $k^{(w)}_D$ has a
positive derivative at $r_0 = \sqrt{2z^*}$, and is continuous and strictly increasing on an
interval around $r_0$. Take $r < s$ in that interval. Points $u_n \downarrow r$ and
$v_n \uparrow s$ at which $k^{(w)}_D = k_D$ give
$k^{(w)}_D(u_n) = k_D(u_n) \ge k_D(v_n) = k^{(w)}_D(v_n)$, and by continuity
$k^{(w)}_D(r) \ge k^{(w)}_D(s)$, a contradiction. So $F^{(w)} \notin \mathcal{A}_D$.

\emph{Clause (2).} This is \ref{lem:threshold-monotone}: $w_{d+1} < w_d$ for every $d \ge 1$,
and every $w_d$ is positive.

\emph{Clause (3).} By \ref{lem:threshold-asymptotic}(3), $\sqrt d\,w_d \to \sqrt\pi$ as
$d \to \infty$ through the real numbers, hence through the integers.

\emph{Clause (4).} The critical values decrease strictly
(\ref{lem:threshold-monotone}). Let $w \in (w_{d+1}, w_d\,]$. For $D \le d$,
$w_D \ge w_d \ge w$. For $D \ge d + 1$, $w_D \le w_{d+1} < w$. If $w > w_1$, then $w > w_D$ for
every $D$. Finally $w_D \to 0$ (\ref{lem:threshold-limit}), so for $w \in (0,w_1]$ there is a
largest integer $d \ge 1$ with $w_d \ge w$, and it is the one $d$ with
$w \in (w_{d+1}, w_d\,]$.
With clause (1), this gives the first three sentences of the clause. For the last, let $d \ge 1$
be an integer. The interval $(w_{d+1}, w_d\,]$ is nonempty by clause (2), and for $w$ in it
$F^{(w)} \in \mathcal{A}_d \setminus \mathcal{A}_{d+1}$. The inclusion
$\mathcal{A}_{d+1} \subseteq \mathcal{A}_d$ is part of \ref{thm:filtration-strict}(3).

The cited facts \ledger{A2}, \ledger{A3} and \ledger{A4} enter in clause (1), and in clause (4)
through it; clauses (2) and (3) cite nothing. Every step is machine-checked, and the theorem is
one statement in Lean, assembled from the Lean of the nodes named. The two directions of clause
(1) are the two steps of \ref{thm:filtration-strict}(1)'s Lean, each of which reads one
inequality. Clause (2) is \ref{lem:threshold-monotone}'s. Clause (3) is
\ref{lem:threshold-asymptotic}(3)'s, composed with the inclusion of the integers in the reals.
The nesting in clause (4) is \ref{thm:filtration-strict}(3)'s. For the sentence on the intervals
the Lean takes the least integer $n \ge 1$ with $w_n < w$, which exists because $w_D \to 0$, and
$d = n - 1$, which is the largest $d$ with $w_d \ge w$.
\end{proof}

% not in the draft: Q-0425, the description of the separating family that the blind check of the
% dictionary found (records/affine/reviews/VERIFY-round1-filtration-map-blind-check.md, section 7:
% "the witnesses are Gaussian variance mixtures: infinitely divisible in every dimension,
% self-decomposable exactly up to dimension d. Supported by the proof; not yet stated"). Checked
% before it was stated, and stated as a proposition because it can be proved from what the chapter
% has. What the check found.
% (i) "Gaussian variance mixture" needs a mixing law, a probability measure on [0,infty) with
% Laplace transform exp(-F_I^{(w)}). The record takes it from Bernstein's theorem, which this
% module's ledger does not carry ([B, Rem. 4.3] is a hypothesis wherever this chapter uses it:
% thm:origin-boundedness-d(3), prop:isotropic-bridge). For THIS family the mixing law is explicit,
% a one-sided stable law of index 1/2 convolved with a compound Poisson law, so clause (2)
% constructs it and cites nothing.
% (ii) Infinite divisibility then needs no Levy-Khintchine theorem either: the same construction
% with the parameters divided by n gives the n-th convolution root.
% (iii) "In A_D exactly for D up to the critical one" is thm:filtration-strict(1) with the nesting
% for w <= w_1, and for w > w_1 it needs the exclusion step at D = 1, which that theorem does not
% state. Clause (3) is the uniform form, F^{(w)} in A_D iff w <= w_D; its two directions are the
% two steps of that theorem's proof, each of which reads one inequality, and the Lean has them in
% that form already (witness_classA, witness_not_classA in FiltrationStrict.lean).
% CHANGED Q-0429 (2026-10-12): SPLIT (ADR-0025; the author's decision of that day: the family as
% laws stays proved in print). Point (iii) above describes the node as it was. The criterion
% "F^{(w)} in A_D iff w <= w_D" and its consequences on exponents, which the Lean nearly has, are
% now thm:separating-family-critical-values(1) and (4), with their proof (the paragraphs
% "w <= w_D gives membership", "w > w_D excludes it" and "the windows" of this node's proof, moved
% word for word). What stays here: clauses (1) and (2), statement and proof unchanged, and of
% clause (3) the reading in laws, which needs the law mu^{(w)}_D of clause (2).
%   clause (3), statement: its first sentence said "F^{(w)} in A_D if and only if w <= w_D, and
%     this holds if and only if mu^{(w)}_D is self-decomposable"; it now says that mu^{(w)}_D is
%     self-decomposable if and only if F^{(w)} in A_D, and cites the theorem for "if and only if
%     w <= w_D". The second sentence ("Consequently ... in no dimension") is word for word. The
%     third ("Every w in (0, w_1] lies in exactly one of the intervals") moved to the theorem.
%   proof: the paragraph "the reading in laws" is word for word; one new paragraph of three lines
%     takes the criterion and the windows from the theorem; the closing sentence no longer points
%     to thm:filtration-strict for the criterion.
% \uses loses what only the moved paragraphs read and gains the theorem. No \lean tag, and none is
% scheduled.
\begin{proposition}[the separating family: Gaussian variance mixtures, self-decomposable up to the critical dimension]
\label{prop:separating-family-mixtures}
\uses{def:filtration-thresholds,thm:separating-family-critical-values,prop:isotropic-polar-classification,thm:cited-A2,thm:cited-A3,thm:cited-A4}
Let $w > 0$, and let $F^{(w)}_{\mathrm I}$, $F^{(w)}$ and the critical values $w_D$ be as in
\ref{def:filtration-thresholds}.
\begin{enumerate}
\item \emph{(Closed form.)} For $\sigma \ge 0$ and $r \ge 0$,
\[
  F^{(w)}_{\mathrm I}(\sigma) = 2\sqrt{\pi\sigma} + w\int_1^2\bigl(1 - e^{-\sigma u}\bigr)\frac{du}{u},
  \qquad
  F^{(w)}(r) = \sqrt{2\pi}\,r + w\int_1^2\bigl(1 - e^{-u r^2/2}\bigr)\frac{du}{u} .
\]
\item \emph{(A Gaussian variance mixture in every dimension.)} Let $Z_0$ be a standard Gaussian
real variable, $N$ a Poisson variable of mean $w\log 2$, and $U_1, U_2, \dots$ variables with the
density $(u\log 2)^{-1}$ on $[1,2]$, all independent, and put
\[
  W := \frac{2\pi}{Z_0^2} + \sum_{j=1}^{N}U_j .
\]
Then $W > 0$ almost surely and $\mathbb{E}\,e^{-\sigma W} = e^{-F^{(w)}_{\mathrm I}(\sigma)}$ for
every $\sigma \ge 0$. For every integer $D \ge 1$, if $Z$ is a standard Gaussian vector of
$\RR^D$ independent of $W$, then $\sqrt W\,Z$ has the characteristic function
$\xi \mapsto e^{-F^{(w)}(|\xi|)}$ on $\RR^D$. Its law $\mu^{(w)}_D$ is rotation invariant and
infinitely divisible. It is the convolution of the law with characteristic function
$e^{-\sqrt{2\pi}|\xi|}$, a rotation-invariant Cauchy law, and of a compound Poisson law, with
characteristic function $\exp\bigl[-w\int_1^2(1 - e^{-u|\xi|^2/2})\,du/u\bigr]$.
\item \emph{(Self-decomposable up to the critical dimension.)} For every integer $D \ge 1$,
$\mu^{(w)}_D$ is self-decomposable if and only if $F^{(w)} \in \mathcal{A}_D$, that is, by
\ref{thm:separating-family-critical-values}(1), if and only if $w \le w_D$. Consequently, if
$w \in (w_{d+1}, w_d\,]$ for an integer
$d \ge 1$, then $\mu^{(w)}_D$ is self-decomposable exactly for $D \le d$, and if $w > w_1$ it is
self-decomposable in no dimension.
\end{enumerate}
\statusA\quad\ledger{A2}\quad\ledger{A3}\quad\ledger{A4}\quad\emph{(\textbf{Assignment.} Clauses
(1) and (2) cite nothing: one Gamma integral, one Gaussian integral, the Laplace transform of a
compound Poisson law, and the uniqueness theorem for characteristic functions. The mixing law is
written down, so Bernstein's theorem, which this module's ledger does not carry, is not needed
for this family; and the infinite divisibility is shown by exhibiting the convolution roots, so
\ledger{A3}(iii) is not needed for it either. Clause (3) spends \ledger{A2}, \ledger{A3} and
\ledger{A4} through \ref{thm:separating-family-critical-values}(1) and (4), which carry the
criterion, and through \ref{prop:isotropic-polar-classification}(2), which makes $F^{(w)}$
admissible when $\mu^{(w)}_D$ is self-decomposable.
\textbf{What the citations carry and what they do not.} They carry the polar form of
$\mathcal{A}_D$ and the uniqueness of a \Levy{} measure. The closed form and the mixture are
proved here, and the criterion $w \le w_D$ in \ref{thm:separating-family-critical-values}.
\textbf{Proved in print, not formalized.} The node has no Lean and none is scheduled (the
author's decision of 2026-10-12). What a reader uses to decide whether $F^{(w)}$ is admitted in
a dimension is \ref{thm:separating-family-critical-values}, which is stated on exponents and is
machine-checked. Until Q-0429 its clause (1) was the first sentence of clause (3) here.
\textbf{What it is worth.}
\ref{thm:filtration-strict}(1) says that $F^{(w)}$ separates $\mathcal{A}_d$ from
$\mathcal{A}_{d+1}$. This node says what the separating exponents are. In every dimension
$e^{-F^{(w)}(|\xi|)}$ is the transform of a positive rotation-invariant kernel, a Cauchy kernel
convolved with a compound Poisson mixture of Gaussians, and the kernel is a Gaussian variance
mixture and infinitely divisible. So each $F^{(w)}$ generates a convolution semigroup of such
kernels in every dimension, and it is the exponent of an isotropic cascade in dimension $D$
exactly when $w \le w_D$. The separation of $\mathcal{A}_d$ from $\mathcal{A}_{d+1}$ therefore
already takes place among the exponents of the form $f(r^2/2)$ with $f$ a Bernstein function,
which give infinitely divisible laws in every dimension, and what is lost at dimension $d + 1$ is
self-decomposability alone. \textbf{How it stands to
\ref{thm:separating-family-critical-values} and \ref{thm:filtration-strict}.} The first has the
criterion $F^{(w)} \in \mathcal{A}_D$ exactly when $w \le w_D$, in the uniform form with one
inequality and with the case $w > w_1$; at $D = d$ and $D = d + 1$ it is clause (1) of the
second. What clause (3) adds to them is the reading in laws. \textbf{A reading, not claimed.}
Differentiating clause (1), $r\,(F^{(w)})'(r) = \sqrt{2\pi}\,r + 2w\,(e^{-r^2/2} - e^{-r^2})$, a
Cauchy exponent plus $w$ times a difference of two Gaussian transforms. In the general form of
the dictionary, which \ref{prop:radial-dictionary} does not state, this would be the exponent of
the driver, and the criterion $w \le w_D$ the nonnegativity of its \Levy{} density on $\RR^D$.
The verification record checks this numerically for $D \le 8$
(\texttt{records/affine/reviews/VERIFY-round1-filtration-map-blind-check.md}, section~7); it is
not proved here. \textbf{Not claimed}: that the step witnesses of
\ref{prop:profile-one-dimension-down}(2) are not of this kind, that is that their exponents are not
Bernstein functions of $r^2/2$, or that they fail to be infinitely divisible one dimension up;
the record argues the first and computes examples for the second, and neither is stated. Nor a
density of $\mu^{(w)}_D$, nor its behaviour at the origin: the profile at the origin of
$F^{(w)}$ is infinite, so \ref{thm:origin-boundedness-d}(2) makes the density smooth where
$\mu^{(w)}_D$ is self-decomposable.)}
\end{proposition}

\begin{proof}
\emph{Clause (1).} By \ref{def:filtration-thresholds} the delay datum $D^{(w)}$ has drift $0$
and profile $k^{(w)}_{\mathrm I}(u) = u^{-1/2} + w\,\mathbf 1_{[1,2]}(u)$, so
$F^{(w)}_{\mathrm I}(\sigma) = \int_0^\infty(1 - e^{-\sigma u})\,u^{-3/2}\,du +
w\int_1^2(1 - e^{-\sigma u})\,du/u$. Integrating by parts, for $\sigma > 0$,
\[
  \int_0^\infty\bigl(1 - e^{-\sigma u}\bigr)\,u^{-3/2}\,du
   = \Bigl[-2u^{-1/2}\bigl(1 - e^{-\sigma u}\bigr)\Bigr]_0^\infty
     + 2\sigma\int_0^\infty u^{-1/2}e^{-\sigma u}\,du
   = 2\sigma\cdot\Gamma(\tfrac12)\,\sigma^{-1/2} = 2\sqrt{\pi\sigma} ,
\]
the bracket vanishing at both ends. At $\sigma = 0$ both sides are $0$. The second formula is the
first at $\sigma = r^2/2$.

\emph{Clause (2), a Gaussian integral.} For $\beta \ge 0$,
\begin{equation}\label{eq:gaussian-inverse-square}
  \int_0^\infty e^{-x^2 - \beta^2/x^2}\,dx = \frac{\sqrt\pi}{2}\,e^{-2\beta} .
\end{equation}
For $\beta = 0$ this is the Gaussian integral. Let $\beta > 0$. Then
$x^2 + \beta^2/x^2 = (x - \beta/x)^2 + 2\beta$. Put
$J := \int_0^\infty e^{-(x - \beta/x)^2}dx$. The substitution $x = \beta/y$ gives
$J = \int_0^\infty e^{-(y - \beta/y)^2}\,\beta y^{-2}\,dy$. Adding the two expressions,
$2J = \int_0^\infty e^{-(x - \beta/x)^2}(1 + \beta/x^2)\,dx = \int_{\RR}e^{-t^2}dt = \sqrt\pi$,
by the substitution $t = x - \beta/x$, an increasing bijection of $(0,\infty)$ onto $\RR$ with
$dt = (1 + \beta/x^2)\,dx$.

\emph{Clause (2), the Laplace transform of $W$.} Let $\varkappa > 0$ and $\sigma \ge 0$. With
$z = \sqrt2\,x$ and \ref{eq:gaussian-inverse-square} at $\beta^2 = \sigma\varkappa/2$,
\[
  \mathbb{E}\,e^{-\sigma\varkappa/Z_0^2}
   = \frac{2}{\sqrt{2\pi}}\int_0^\infty e^{-z^2/2 - \sigma\varkappa/z^2}\,dz
   = \frac{2}{\sqrt\pi}\int_0^\infty e^{-x^2 - \sigma\varkappa/(2x^2)}\,dx
   = e^{-\sqrt{2\sigma\varkappa}} .
\]
At $\varkappa = 2\pi$ this is $e^{-2\sqrt{\pi\sigma}}$. For the sum, conditioning on $N$,
\[
  \mathbb{E}\,e^{-\sigma\sum_{j \le N}U_j}
   = \sum_{n \ge 0}e^{-w\log 2}\frac{(w\log 2)^n}{n!}\bigl(\mathbb{E}\,e^{-\sigma U_1}\bigr)^n
   = \exp\bigl[-w\log 2\,\bigl(1 - \mathbb{E}\,e^{-\sigma U_1}\bigr)\bigr]
   = \exp\Bigl[-w\int_1^2\bigl(1 - e^{-\sigma u}\bigr)\frac{du}{u}\Bigr] ,
\]
because $\log 2\,(1 - \mathbb{E}\,e^{-\sigma U_1}) = \int_1^2(1 - e^{-\sigma u})\,du/u$. By
independence and clause (1), $\mathbb{E}\,e^{-\sigma W} = e^{-F^{(w)}_{\mathrm I}(\sigma)}$. And
$W \ge 2\pi/Z_0^2 > 0$ almost surely.

\emph{Clause (2), the mixture.} Let $D \ge 1$ and let $Z$ be a standard Gaussian vector of
$\RR^D$ independent of $W$. Conditioning on $W$, and using that $\sqrt u\,Z$ has the
characteristic function $e^{-u|\xi|^2/2}$,
\[
  \mathbb{E}\,e^{i\langle\xi,\sqrt W Z\rangle} = \mathbb{E}\,e^{-W|\xi|^2/2}
   = e^{-F^{(w)}_{\mathrm I}(|\xi|^2/2)} = e^{-F^{(w)}(|\xi|)} .
\]
The characteristic function is radial, so the law $\mu^{(w)}_D$ of $\sqrt W\,Z$ is invariant
under $\Orth(D)$, by the uniqueness theorem for characteristic functions. By clause (1) it
factors as $e^{-\sqrt{2\pi}|\xi|}$ times
$\exp[-w\int_1^2(1 - e^{-u|\xi|^2/2})\,du/u]$. The same computation with $W$ replaced by each
of its two summands shows that the first factor is the characteristic function of
$\sqrt{2\pi}\,Z/|Z_0|$ and the second that of $\sqrt{\sum_{j\le N}U_j}\,Z$. The second is also
the characteristic function of the compound Poisson sum $\sum_{j \le N}\sqrt{U_j}\,Z_j$, with
$Z_1, Z_2, \dots$ independent standard Gaussian vectors independent of the rest, by the
computation made above for the Laplace transform. So $\mu^{(w)}_D$ is the convolution of the
two laws.

\emph{Clause (2), infinite divisibility.} Let $n \ge 1$, let $N_n$ be a Poisson variable of mean
$(w\log 2)/n$, independent of the rest, and put
$W_n := 2\pi/(n^2Z_0^2) + \sum_{j \le N_n}U_j$. The two computations above, at
$\varkappa = 2\pi/n^2$ and with $w/n$ in place of $w$, give
$\mathbb{E}\,e^{-\sigma W_n} = e^{-F^{(w)}_{\mathrm I}(\sigma)/n}$. So the law $\mu_n$ of
$\sqrt{W_n}\,Z$ has the characteristic function $e^{-F^{(w)}(|\xi|)/n}$, and its $n$-fold
convolution has the characteristic function of $\mu^{(w)}_D$. Hence
$\mu^{(w)}_D = \mu_n^{\ast n}$ for every $n$: it is infinitely divisible.

\emph{Clause (3), the reading in laws.} $F^{(w)} \in \mathcal{A}_D$ means that $F^{(w)}$ is
admissible and that $\xi \mapsto F^{(w)}(|\xi|)$ is the exponent of a self-decomposable law on
$\RR^D$. A law with the characteristic function $e^{-F^{(w)}(|\xi|)}$ is $\mu^{(w)}_D$, by
clause (2) and the uniqueness theorem. And if $\mu^{(w)}_D$ is self-decomposable, $F^{(w)}$ is
admissible, by \ref{prop:isotropic-polar-classification}(2). So $F^{(w)} \in \mathcal{A}_D$
exactly when $\mu^{(w)}_D$ is self-decomposable.

\emph{Clause (3), the criterion and the windows.} By
\ref{thm:separating-family-critical-values}(1), $F^{(w)} \in \mathcal{A}_D$ exactly when
$w \le w_D$. By its clause (4) this holds exactly for $D \le d$ when $w \in (w_{d+1}, w_d\,]$,
and for no $D$ when $w > w_1$. With the reading in laws these are the assertions of the clause.

The cited facts enter in clause (3) only, as the Assignment says. Nothing here is
machine-checked.
\end{proof}

\section{The isotropic classification in dimension $d$}
\label{sec:isotropic-classification}

This section is the draft's \S 2.1 proper: the classification of the isotropic families of
$\RR^d$ and the bridge from the causal theory. Both are reduced to the line. A
rotation-invariant law on $\RR^d$ is determined by its law along one line, and the marginals of
an isotropic family along a line form a family of the line paper's kind. So the line's theorem
[V, Thm.~7.3] applies to them unchanged, and the only $d$-dimensional content left is the class
$\mathcal{A}_d$ of \ref{def:filtration-thresholds}: membership in it comes from the cascade, and
reading it as an explicit polar form is one cited step, \ref{prop:isotropic-polar-reading}.

% draft: section 2.1, Definition 2.1, in the operator form in which the draft states it. Not a node
% until Q-0406 (hub; the author's review of the isotropic paper's section 3,
% records/affine/REVIEW-author-paper-iso.md): the axioms stood inline in
% lem:operators-to-kernels(5) and were cited from there by lem:kernels-to-operators(3), while the
% Lean had the structure all along (AffineHemigroup.IsIsotropicOperatorFamily, typed in Q-0281 and
% re-typed literally in wave M, phase 1). The author asked that the isotropic paper open its
% section 3 with the axioms on operators, as the line papers do, so the definition is now a shared
% statement. No statement of another node changes: lem:operators-to-kernels(5),
% lem:operators-to-kernels-families(5) and lem:kernels-to-operators(3) keep their text, which lists
% these axioms, and gain a \uses edge to this node. Nothing is proved here and nothing new is typed:
% the tag names the existing structure and the single-operator structure it is built on.
% What writing it down found. (i) The line's nine axioms keep their numbers, and rotation invariance
% takes the number (A3): at d = 1, O(1) = {1, -1} and R_{-1} is the line's reflection R, so (A3)
% here is [V]'s (A3) there and the definition at d = 1 is [V, Def. 3.1] word for word. The clause
% lem:operators-to-kernels(5) lists the same axiom without a number, and
% lem:operators-to-kernels(1) says "(A3)" of [V]'s reflection axiom only; the two usages agree.
% (ii) (A8) is existential in the text, as (I4) is, and the Lean structure takes the family
% (S_lambda) as a parameter S, as IsIsotropicCascadeFamily does; the passages between the two
% forms keep the same S.
\begin{definition}[isotropic cascade measurement family, on operators]
\label{def:isotropic-operator-family}
\lean{AffineHemigroup.IsIsotropicOperatorFamily, AffineHemigroup.IsSingleOperator}\leanok
Fix an integer $d \ge 1$. On the real space $L^1(\RR^d)$ write $T_af := f(\cdot - a)$ for the
translation by $a \in \RR^d$, $R_Qf := f(Q^{-1}\,\cdot)$ for the action of an orthogonal matrix
$Q \in \Orth(d)$, and $D_\lambda f := \lambda^{-d}f(\cdot/\lambda)$ for the dilation by
$\lambda > 0$. A family $(\Phi_{s,t})_{0 \le s \le t}$ of linear operators on $L^1(\RR^d)$ is an
\emph{isotropic cascade measurement family in dimension $d$} if:
\begin{enumerate}
\item[(A1)] each $\Phi_{s,t}$ is a bounded operator on $L^1(\RR^d)$;
\item[(A2)] \emph{(translation invariance)} $\Phi_{s,t}T_a = T_a\Phi_{s,t}$ for every
$a \in \RR^d$;
\item[(A3)] \emph{(rotation invariance)} $\Phi_{s,t}R_Q = R_Q\Phi_{s,t}$ for every
$Q \in \Orth(d)$;
\item[(A4)] \emph{(positivity)} $\Phi_{s,t}$ maps nonnegative functions to nonnegative functions;
\item[(A5)] \emph{(mass)} $\int\Phi_{s,t}f = \int f$ for every nonnegative $f \in L^1(\RR^d)$;
\item[(A6)] \emph{(cascade)} $\Phi_{t,t} = \mathrm{Id}$, and $\Phi_{s,t}\Phi_{r,s} = \Phi_{r,t}$
whenever $r \le s \le t$;
\item[(A7)] \emph{(continuity)} for every $f \in L^1(\RR^d)$ the map $(s,t) \mapsto \Phi_{s,t}f$
is continuous from $\{0 \le s \le t\}$ to $L^1(\RR^d)$;
\item[(A8)] \emph{(dilation covariance)} there is a family $(S_\lambda)_{\lambda > 0}$ of
increasing bijections of $[0,\infty)$ with
$D_\lambda\Phi_{s,t} = \Phi_{S_\lambda s,\,S_\lambda t}D_\lambda$ for all $0 \le s \le t$ and all
$\lambda > 0$;
\item[(ND)] \emph{(nondegeneracy)} $\Phi_{s,t} \ne \mathrm{Id}$ whenever $s < t$.
\end{enumerate}
\statusT\quad\emph{(Vocabulary, and no claim. The draft's Definition~2.1 as the draft states it,
on operators: the axioms that \ref{lem:operators-to-kernels}(5) lists and
\ref{lem:kernels-to-operators}(3) cites from it, given a name and a node of their own so that the
isotropic paper can open with them. \textbf{Readings.} The numbers are those of [V, Def.~3.1], and
rotation invariance is numbered (A3) because at $d = 1$ it is [V]'s reflection axiom:
$\Orth(1) = \{\pm1\}$ and $R_{-1}f = f(-\,\cdot)$. So at $d = 1$ the definition is [V, Def.~3.1]
itself. (A1), (A2), (A4) and (A5) concern one operator and are the single-operator axioms of
\ref{lem:operators-to-kernels}(1). In the Lean, (A1) is the type of $\Phi_{s,t}$, a bounded
linear operator; $R_Q$ is $D_Q$, since $|\det Q| = 1$, and $D_\lambda$ is $D_{\lambda I}$; and the
family $(S_\lambda)$ of (A8) is a parameter of the structure, as it is in
\ref{def:isotropic-kernel-family}. \textbf{The passage to the kernel form} is
\ref{lem:operators-to-kernels}(5), and the passage back is \ref{lem:kernels-to-operators}(3); both
keep the same $(S_\lambda)$. \textbf{Not claimed}: anything about operators on $L^p(\RR^d)$ for
$p > 1$.)}
\end{definition}

% draft: section 2.1, Definition 2.1, restated on kernels, on the pattern of
% def:affine-cascade-family. The operator form is the status annotation's business; the passage
% between the two is lem:operators-to-kernels(5) (Q-0207).
% CHANGED Q-0406: the operator form is a node of its own, def:isotropic-operator-family, placed
% before this one; the annotation names it. The statement is unchanged.
% CHANGED Q-0281: typed (the statement skeleton). The displayed identity hat{A_* mu}(xi) =
% hat{mu}(A^T xi) is lem:gw-scheme-claims(3) (charFun_push), which the \lean tag names; so is the
% lemma identifying the Lean spelling of a gauge (Paper V's conjunction) with (I4)'s "increasing
% bijections of [0,infty)". The text is unchanged.
% Q-0281 wave 1, blind review W1E-5: isGauge_iff_isIncreasingBijection is proved
% (AffineHemigroup.Defs.IsotropicFamily, Lean core, guarded). Every declaration of the tag is now
% defined in the proved library or proved (charFun_push in wave 1A), so the node is \leanok.
% CHANGED Q-0365: the tag names the shared library's ScaleSpace.IsSelfDecomposable (the measure form
% mu = (c .)_* mu * rho_c, ScaleSpaceCore v0.3.0) and ScaleSpace.isSelfDecomposable_iff_charFun, which
% makes it eq:self-decomposable for a probability measure on R^d; Paper VII's own copy, the
% transform form verbatim, is deleted. The text is unchanged.
\begin{definition}[isotropic cascade family, on kernels]
\label{def:isotropic-kernel-family}
\lean{AffineHemigroup.IsIsotropicCascadeFamily, AffineHemigroup.IsIncreasingBijection,
AffineHemigroup.IsGauge, ScaleSpace.IsSelfDecomposable, ScaleSpace.isSelfDecomposable_iff_charFun,
AffineHemigroup.push, AffineHemigroup.charFun_push,
AffineHemigroup.isGauge_iff_isIncreasingBijection}\leanok
Fix an integer $d \ge 1$. For $A \in \GL(d)$ write $A_*\mu := \mu \circ A^{-1}$ for the image of a
measure $\mu$ on $\RR^d$, and $\widehat\mu(\xi) = \int e^{i\langle\xi,x\rangle}\mu(dx)$, so that
$\widehat{A_*\mu}(\xi) = \widehat\mu(A^{\tp}\xi)$. An \emph{isotropic cascade family in dimension
$d$} is a family $(\mu_{s,t})_{0 \le s \le t}$ of probability measures on $\RR^d$ such that
\begin{enumerate}
\item[(I1)] \emph{(cascade)} $\mu_{t,t} = \delta_0$, and $\mu_{r,t} = \mu_{s,t} * \mu_{r,s}$
whenever $r \le s \le t$;
\item[(I2)] \emph{(continuity)} for every $\xi \in \RR^d$ the map
$(s,t) \mapsto \widehat\mu_{s,t}(\xi)$ is continuous on $\{0 \le s \le t\}$;
\item[(I3)] \emph{(rotation invariance)} $Q_*\mu_{s,t} = \mu_{s,t}$ for every $Q \in \Orth(d)$
and every pair;
\item[(I4)] \emph{(dilation covariance)} there is a family $(S_\lambda)_{\lambda > 0}$ of
increasing bijections of $[0,\infty)$ with $(\lambda I)_*\mu_{s,t} = \mu_{S_\lambda s,\,S_\lambda t}$
for every $\lambda > 0$ and every pair;
\item[(ND)] \emph{(nondegeneracy)} $\mu_{s,t} \ne \delta_0$ whenever $s < t$.
\end{enumerate}
\statusT\quad\emph{(Vocabulary, and no claim. \textbf{The words the chapter uses with it},
restated from the released line paper and the source so that the nodes stand alone, and kept in
this annotation rather than in the statement because they are the line's vocabulary, which the
article restates once, in its section on the line papers. A \emph{gauge} is an increasing
bijection $\chi$ of $[0,\infty)$. A probability measure $\mu$ on $\RR^d$ is
\emph{self-decomposable} if for every $c \in (0,1)$ there is a probability measure $\rho_c$ on
$\RR^d$, a \emph{residual}, with
\begin{equation}\label{eq:self-decomposable}
  \widehat\mu(\xi) = \widehat\mu(c\,\xi)\,\widehat\rho_c(\xi), \qquad \xi \in \RR^d .
\end{equation}
This is \sscite{sato1999levy}'s Definition~15.1, p.~90, with $c = b^{-1}$, and at $d = 1$ it is
[V, Def.~2.12]; the nodes of chapters~4 and~5 that cite \ref{eq:self-decomposable} read it as this
definition (vocabulary, as a definition is; the contraction is written $c$ as in [V], and is not the
profile at the origin); it is read at the page in
\texttt{records/affine/READING-sato1999levy-prop-15-5.md} \S 1, and no ledger entry carries it. A
function $F : [0,\infty) \to [0,\infty)$ is \emph{admissible} if it has the form [V, (2.3)],
\begin{equation}\label{eq:admissible-line}
  F(r) = a_{\mathrm V}\,r^2 + \int_0^\infty \bigl(1 - \cos rx\bigr)\,k(x)\,\frac{dx}{x}, \qquad r \ge 0,
\end{equation}
with $a_{\mathrm V} \ge 0$ and $k : (0,\infty) \to [0,\infty)$ nonincreasing,
$\int_0^1 x\,k(x)\,dx < \infty$ and $\int_1^\infty k(x)\,x^{-1}dx < \infty$. It writes
$a_{\mathrm V}$ for [V]'s Gaussian coefficient, so that the letter $a$ stays the Gaussian
coefficient of \ref{def:isotropic-origin-datum}; the two are related in
\ref{prop:isotropic-polar-reading}. The class $\mathcal{A}_d$ is defined once, in
\ref{def:filtration-thresholds}. Until Q-0286 this vocabulary, and a sentence restating
$\mathcal{A}_d$ (``So the class $\mathcal{A}_d$ \dots{} is the set of admissible $F$ \dots''),
stood in the statement; the restatement was a second definition of the class, and the rest
re-defined what the article restates once, so the statement was slimmed to the axioms, with no
change to them. \textbf{The operator form.} The draft's Definition~2.1 states the axioms on
operators $\Phi_{s,t}$ on $L^1(\RR^d)$: (A1) and (A4)--(A7) of [V, Def.~3.1], (A2) for every
$a \in \RR^d$, rotation invariance $\Phi_{s,t}R_Q = R_Q\Phi_{s,t}$ with $R_Qf = f(Q^{-1}\,\cdot)$,
(A8) with $D_\lambda f = \lambda^{-d}f(\cdot/\lambda)$, and (ND). That form is
\ref{def:isotropic-operator-family}. This node is the form every proof
below reads, and the passage from the operator form to it is
\ref{lem:operators-to-kernels}(5), the node that also serves \ref{def:affine-cascade-family}:
[V, Lem.~4.1] proved in $\RR^d$, with $R_Q(\mu * f) = (Q_*\mu) * R_Qf$ and
$D_\lambda(\mu * f) = ((\lambda I)_*\mu) * D_\lambda f$ turning rotation invariance into (I3) and
(A8) into (I4). Reflection symmetry, [V]'s (A3), is not listed: it is (I3) at $Q = -I$. At $d = 1$,
$\Orth(1) = \{\pm1\}$ and the definition is the kernel form of [V, Def.~3.1]. In the
$(G, W)$ scheme of chapter~4 this definition is the slice $\{sI\}$ of the corner at the similarities
and the rays: (ND) at every pair follows there from (K4), and the free gauge $S_\lambda$ becomes the
scheme's $s \mapsto \lambda^2 s$ after a change of the scale coordinate unique up to a constant
(\ref{prop:isotropic-corner}).)}
\end{definition}

% draft: section 2.1, the observation the draft's port does not make (Q-0210, section 1). It
% replaces the draft's "no-lattice lemma is not needed in d >= 2" port note as well: (2) below is
% what the reduction needs of the zero set, and it holds in every dimension.
\begin{lemma}[a rotation-invariant law is determined by its marginal]
\label{lem:radial-marginal}
\uses{def:isotropic-kernel-family}
\lean{AffineHemigroup.radial_marginal, AffineHemigroup.radial_marginal_eq_dirac,
AffineHemigroup.radial_marginal_ops, AffineHemigroup.marginal}\leanok
Let $d \ge 1$, let $e \in \RR^d$ be a unit vector and $\pi(x) := \langle e, x\rangle$, and for a
probability measure $\mu$ on $\RR^d$ write $\pi_*\mu$ for its image on $\RR$.
\begin{enumerate}
\item If $Q_*\mu = \mu$ for every $Q \in \Orth(d)$, then $\kappa := \pi_*\mu$ is a symmetric
probability measure on $\RR$, and
\[
  \widehat\mu(\xi) = \widehat\kappa(|\xi|), \qquad \xi \in \RR^d .
\]
In particular $\widehat\mu$ is real and radial, and $\kappa$ does not depend on $e$.
\item Under the hypothesis of (1), $\pi_*\mu = \delta_0$ if and only if $\mu = \delta_0$.
\item For all probability measures $\mu, \mu'$ on $\RR^d$ and every $\lambda > 0$,
$\pi_*(\mu * \mu') = \pi_*\mu * \pi_*\mu'$, $\pi_*\bigl((\lambda I)_*\mu\bigr) =
(\lambda I)_*(\pi_*\mu)$ on $\RR$, and $\pi_*\delta_0 = \delta_0$.
\end{enumerate}
\statusT\quad\emph{(Cites nothing. It reads the uniqueness of the Fourier transform of a probability
measure on $\RR^d$, as the chapter's earlier nodes do, and the transitivity of $\Orth(d)$ on each
sphere. \textbf{What it replaces.} The draft's port notes need, in $d \ge 2$, that the zero set
$\{\xi : \widehat\mu(\xi) = 1\}$ of a rotation-invariant law is all of $\RR^d$ unless
$\mu = \delta_0$, and argue it by sums of sphere points. Clause (2) is the fact the reduction needs,
in every dimension including $d = 1$, and clause (1) is why it is immediate: the zero set is the
line's, spread over spheres.)}
\end{lemma}

% CHANGED Q-0281 wave 1: proof of (1) rewritten to the Lean route: the transform of the marginal,
% hat{kappa}(omega) = hat{mu}(omega e), is computed once, and the orthogonal matrix is named (the
% reflection carrying one point of a sphere to another), which covers xi = 0 and the independence of
% e without the evenness of hat{kappa}. Clauses (2) and (3) are the Lean route as they stood.
\begin{proof}
\leanok
\emph{(1).} $\pi_*\mu$ is a probability measure because $\pi$ is measurable, and for $\omega \in \RR$
\[
  \widehat\kappa(\omega) = \int_{\RR} e^{i\omega y}\,\kappa(dy)
   = \int e^{i\omega\langle e, x\rangle}\,\mu(dx) = \widehat\mu(\omega e) .
\]
A rotation-invariant law has a radial transform: if $|\eta| = |\eta'|$, let $Q$ be the identity when
$\eta = \eta'$ and otherwise the reflection in the hyperplane orthogonal to $\eta - \eta'$, an
element of $\Orth(d)$ with $Q^{\tp} = Q$ and $Q\eta = \eta'$; then
$\widehat\mu(\eta) = \widehat{Q_*\mu}(\eta) = \widehat\mu(Q^{\tp}\eta) = \widehat\mu(\eta')$.
Since $|\,|\xi|e\,| = |\xi|$, this gives $\widehat\mu(\xi) = \widehat\mu(|\xi|e) =
\widehat\kappa(|\xi|)$ for every $\xi \in \RR^d$. Since $-I \in \Orth(d)$ and $\pi(-x) = -\pi(x)$,
the image of $\kappa$ under $\omega \mapsto -\omega$ is $\pi_*((-I)_*\mu) = \kappa$, so $\kappa$ is
symmetric and $\widehat\kappa$ is real; hence so is $\widehat\mu$. For another unit vector $e'$ with
marginal $\kappa'$, $|\omega e'| = |\omega e|$ gives $\widehat{\kappa'}(\omega) = \widehat\mu(\omega e')
= \widehat\mu(\omega e) = \widehat\kappa(\omega)$ for every $\omega \in \RR$, so $\kappa' = \kappa$
by uniqueness of the transform on $\RR$.

\emph{(2).} If $\mu = \delta_0$ then $\pi_*\mu = \delta_0$. Conversely, if $\kappa = \delta_0$ then
$\widehat\kappa \equiv 1$, so $\widehat\mu \equiv 1 = \widehat{\delta_0}$ by (1), and
$\mu = \delta_0$ by uniqueness of the transform on $\RR^d$.

\emph{(3).} $\mu * \mu'$ is the image of $\mu \otimes \mu'$ under $(x,y) \mapsto x + y$, and
$\pi(x+y) = \pi(x) + \pi(y)$, so $\pi_*(\mu * \mu')$ is the image of $\pi_*\mu \otimes \pi_*\mu'$
under addition on $\RR$, which is $\pi_*\mu * \pi_*\mu'$. The other two identities are
$\pi(\lambda x) = \lambda\pi(x)$ and $\pi(0) = 0$.
\end{proof}

% draft: section 2.1, port notes, replaced. This is the statement that [V, Thm. 7.3] applies to
% the marginal family AS IT STANDS (Q-0210, section 1; in [V]'s Lean it is the observation that the
% marginals are a CascadeData). [V, Def. 3.1] is restated in the node so that nothing here \uses a
% node of [V]'s blueprint.
% CHANGED Q-0365 (the switch to ScaleSpaceCore v0.4.0, referee V04-1): tag only, no statement text;
% the declarations return the trunk's ScaleSpace.CascadeData / ScaleSpace.CascadeFamily.
% CHANGED Q-0376 (fidelity review, row R1): the \lean tag names the shared library's
% ScaleSpace.CascadeFamily and ScaleSpace.CascadeData, the structures the two declarations of the tag
% state their conclusions in since Q-0365; it named Paper V's SpatialLine.CascadeFamily and
% SpatialLine.CascadeData, which still resolve in the required package but are not the ones the Lean
% statement uses. Tag only; the text is unchanged.
\begin{lemma}[the marginals of an isotropic family form a family on the line]
\label{lem:isotropic-marginal-family}
\uses{def:isotropic-kernel-family,lem:radial-marginal}
\lean{AffineHemigroup.isotropic_marginal_family, AffineHemigroup.isotropic_marginal_cascadeData,
ScaleSpace.CascadeFamily, ScaleSpace.CascadeData}\leanok
\statedin{line:def:cascade-family}\statedin{line:thm:main-characterization}
Let $(\mu_{s,t})$ be an isotropic cascade family in dimension $d$ with dilation action
$(S_\lambda)$, let $e$ and $\pi$ be as in \ref{lem:radial-marginal}, and put
\[
  \kappa_{s,t} := \pi_*\mu_{s,t}, \qquad \Phi_{s,t}f := \kappa_{s,t} * f
  \quad (f \in L^1(\RR)), \qquad 0 \le s \le t .
\]
Then $\widehat\mu_{s,t}(\xi) = \widehat\kappa_{s,t}(|\xi|)$ for every $\xi \in \RR^d$, each
$\kappa_{s,t}$ is a symmetric probability measure, and $(\Phi_{s,t})$ is a symmetric cascade
measurement family in the sense of \restated{V}{Def.~3.1}{def:line-family}, with the same $(S_\lambda)$. That is, restated:
\begin{enumerate}
\item[(A1)] each $\Phi_{s,t}$ is a bounded operator on $L^1(\RR)$;
\item[(A2)] $\Phi_{s,t}\tau_a = \tau_a\Phi_{s,t}$ for every $a \in \RR$, where
$\tau_af = f(\cdot - a)$;
\item[(A3)] $\Phi_{s,t}R = R\,\Phi_{s,t}$, where $Rf = f(-\,\cdot)$;
\item[(A4)] $\Phi_{s,t}$ maps nonnegative functions to nonnegative functions;
\item[(A5)] $\int\Phi_{s,t}f = \int f$ for every nonnegative $f \in L^1(\RR)$;
\item[(A6)] $\Phi_{t,t} = \mathrm{Id}$, and $\Phi_{s,t}\Phi_{r,s} = \Phi_{r,t}$ whenever
$r \le s \le t$;
\item[(A7)] for every $f \in L^1(\RR)$ the map $(s,t) \mapsto \Phi_{s,t}f$ is continuous from
$\{0 \le s \le t\}$ to $L^1(\RR)$;
\item[(A8)] $D_\lambda\Phi_{s,t} = \Phi_{S_\lambda s,\,S_\lambda t}D_\lambda$ for every
$\lambda > 0$, where $D_\lambda f = \lambda^{-1}f(\cdot/\lambda)$;
\item[(ND)] $\Phi_{s,t} \ne \mathrm{Id}$ whenever $s < t$.
\end{enumerate}
\statusT\quad\emph{(Cites one released fact of [V], \Levy's continuity theorem on the line,
[V, Prop.~2.8], restated in the proof where it is used; it is a machine-checked node of [V] with no
ledger entry. The argument for (A7) is the one in [V, Thm.~7.3]'s sufficiency proof, written out
here in full rather than cited. \textbf{Not claimed}: anything
in dimension $d$ beyond the marginal. The lemma says nothing about the operator form of
\ref{def:isotropic-kernel-family}; it builds operators on the \emph{line} from the kernels, and
that direction --- convolution by a probability measure satisfies (A1), (A2), (A4) and (A5) --- is
elementary, being the converse half of [V, Lem.~4.1]. \textbf{What it is worth.} This is the
sentence ``the marginals form a family of [V]'s kind'' of \texttt{notes/PLAN-d-dimensional-port.md}
\S 1, and in [V]'s Lean it is the statement that the marginals are a \texttt{CascadeData}, the
data of [V, Thm.~7.3] (the structure \texttt{CascadeFamily} is [V, Def.~3.1]); the
reduction's formalization would begin here.)}
\end{lemma}

% CHANGED Q-0281 wave 1: proved; the last paragraph, saying which steps are machine-checked where,
% is added. The argument is unchanged.
\begin{proof}
\leanok
\emph{The kernels.} By (I3) each $\mu_{s,t}$ is rotation invariant, so by
\ref{lem:radial-marginal}(1) each $\kappa_{s,t}$ is a symmetric probability measure with
$\widehat\mu_{s,t}(\xi) = \widehat\kappa_{s,t}(|\xi|)$.

\emph{(A1), (A2), (A4), (A5).} For a probability measure $\kappa$ on $\RR$ and $f \in L^1(\RR)$,
$(\kappa * f)(x) = \int f(x-y)\,\kappa(dy)$. Tonelli gives $\|\kappa * f\|_1 \le \|f\|_1$, so the
operator is bounded, and Fubini gives $\int \kappa * f = \int f$. A nonnegative $f$ gives a
nonnegative $\kappa * f$, and $\tau_a(\kappa * f) = \kappa * \tau_af$ is a change of variables.

\emph{(A3).} As $\kappa_{s,t}$ is symmetric, $R(\kappa_{s,t} * f)(x) = \int f(-x-y)\,\kappa_{s,t}(dy)
= \int f(-x+y)\,\kappa_{s,t}(dy) = (\kappa_{s,t} * Rf)(x)$.

\emph{(A6).} By (I1) and \ref{lem:radial-marginal}(3), $\kappa_{t,t} = \pi_*\delta_0 = \delta_0$
and $\kappa_{r,t} = \pi_*(\mu_{s,t} * \mu_{r,s}) = \kappa_{s,t} * \kappa_{r,s}$. So
$\Phi_{t,t} = \mathrm{Id}$, and $\Phi_{s,t}\Phi_{r,s}f = \kappa_{s,t} * (\kappa_{r,s} * f)
= (\kappa_{s,t} * \kappa_{r,s}) * f = \Phi_{r,t}f$, by Fubini in the middle step.

\emph{(A8).} By (I4) and \ref{lem:radial-marginal}(3),
$(\lambda I)_*\kappa_{s,t} = \pi_*\bigl((\lambda I)_*\mu_{s,t}\bigr) = \kappa_{S_\lambda s,S_\lambda t}$.
For a probability measure $\kappa$ on $\RR$, substituting $z = \lambda y$,
\[
  D_\lambda(\kappa * f)(x) = \int \lambda^{-1}f\Bigl(\frac{x - \lambda y}{\lambda}\Bigr)\kappa(dy)
   = \int (D_\lambda f)(x - z)\,\bigl((\lambda I)_*\kappa\bigr)(dz) ,
\]
so $D_\lambda\Phi_{s,t}f = \kappa_{S_\lambda s,S_\lambda t} * D_\lambda f
= \Phi_{S_\lambda s,S_\lambda t}D_\lambda f$.

\emph{(ND).} Let $s < t$. By (ND) of the family $\mu_{s,t} \ne \delta_0$, so
$\kappa_{s,t} \ne \delta_0$ by \ref{lem:radial-marginal}(2). If $\kappa_{s,t} * f = f$ for every
$f \in L^1(\RR)$, take $f$ the standard Gaussian density: its transform is positive, and the
transform of the identity reads $\widehat\kappa_{s,t}\widehat f = \widehat f$, so
$\widehat\kappa_{s,t} \equiv 1$ and $\kappa_{s,t} = \delta_0$, a contradiction. So
$\Phi_{s,t} \ne \mathrm{Id}$.

\emph{(A7), the collapsing increment.} For $f \in L^1(\RR)$ put
$\Theta_f(y) := \|\tau_yf - f\|_1$, a function that is bounded by $2\|f\|_1$, continuous, and zero
at $y = 0$, because translation acts continuously on $L^1(\RR)$. For a probability measure $\kappa$,
$\kappa * f - f = \int(\tau_yf - f)\,\kappa(dy)$, so Minkowski's integral inequality gives
\begin{equation}\label{eq:collapsing-increment}
  \|\kappa * f - f\|_1 \le \int \Theta_f\,d\kappa .
\end{equation}
Moreover $\Theta_{\kappa' * g} \le \Theta_g$ for every probability measure $\kappa'$ and every
$g \in L^1(\RR)$, because $\tau_y(\kappa' * g) - \kappa' * g = \kappa' * (\tau_yg - g)$ and
convolution by $\kappa'$ does not increase the norm. Now let pairs $(s_n,t_n)$ with
$s_n \le t_n$ converge to a diagonal point $(t_0,t_0)$. By (I2) at $\xi = \omega e$,
$\widehat\kappa_{s_n,t_n}(\omega) = \widehat\mu_{s_n,t_n}(\omega e) \to \widehat\mu_{t_0,t_0}(\omega e) = 1$
for every $\omega \in \RR$. \Levy's continuity theorem on the line \restatedp{V}{Prop.~2.8}{sec:line-facts} says that if
the transforms of probability measures on $\RR$ converge at every point to the transform of a
probability measure, the measures converge weakly to it. So $\kappa_{s_n,t_n} \to \delta_0$
weakly. Since $\Theta_g$ is bounded and continuous, $\int\Theta_g\,d\kappa_{s_n,t_n} \to
\Theta_g(0) = 0$ for every $g$.

\emph{(A7), the general pair.} Let $(s_n,t_n) \to (s,t)$ with $s_n \le t_n$. If $s = t$ then
$\Phi_{s,t} = \mathrm{Id}$ and \ref{eq:collapsing-increment} with the previous paragraph gives
$\Phi_{s_n,t_n}f \to f$. If $s < t$ put $u_n := \max\{s_n,s\}$ and $v_n := \min\{t_n,t\}$, so
that $u_n \to s$, $v_n \to t$ and, for all large $n$, $s_n \le u_n \le v_n \le t_n$ and
$s \le u_n \le v_n \le t$. By (A6),
\[
  \Phi_{s_n,t_n}f = \Phi_{v_n,t_n}\,\Phi_{u_n,v_n}\,\Phi_{s_n,u_n}f ,
  \qquad
  \Phi_{s,t}f = \Phi_{v_n,t}\,\Phi_{u_n,v_n}\,\Phi_{s,u_n}f .
\]
The left endpoints are moved first, where contraction absorbs them: with
$g_n := \Phi_{u_n,v_n}\Phi_{s_n,u_n}f$ and $h_n := \Phi_{u_n,v_n}\Phi_{s,u_n}f$,
\[
  \|g_n - h_n\|_1 \le \|\Phi_{s_n,u_n}f - f\|_1 + \|f - \Phi_{s,u_n}f\|_1 \to 0 ,
\]
both increments collapsing onto $(s,s)$. The right endpoints are moved last, where
\ref{eq:collapsing-increment} is read at $h_n$ and $\Theta_{h_n} \le \Theta_f$:
\[
  \|\Phi_{v_n,t_n}h_n - \Phi_{v_n,t}h_n\|_1
   \le \int\Theta_f\,d\kappa_{v_n,t_n} + \int\Theta_f\,d\kappa_{v_n,t} \to 0 ,
\]
both increments collapsing onto $(t,t)$. Hence
$\|\Phi_{s_n,t_n}f - \Phi_{s,t}f\|_1 \le \|g_n - h_n\|_1
+ \|\Phi_{v_n,t_n}h_n - \Phi_{v_n,t}h_n\|_1 \to 0$.

\emph{What is machine-checked where.} The checked proof verifies, for the marginals, exactly the
facts about kernels the paragraphs above start from: each $\kappa_{s,t}$ is a symmetric
probability measure, $\kappa_{t,t} = \delta_0$, $\kappa_{r,t} = \kappa_{s,t} * \kappa_{r,s}$, the map
$(s,t) \mapsto \widehat\kappa_{s,t}(\omega)$ is continuous for every $\omega$,
$\kappa_{s,t} \ne \delta_0$ for $s < t$, and $(\lambda I)_*\kappa_{s,t} =
\kappa_{S_\lambda s,S_\lambda t}$. The passage from these to (A1)--(A8) and (ND) for convolution
by the kernels, the operator-level steps of the paragraphs above, is the released line paper's own
machine-checked construction, which makes the same arguments once for every family of kernels with
these properties and cites no ledger entry of \progpaper{V}.
\end{proof}

% draft: section 2.1, Theorem 2.2, in the kernel form and with the class A_d in place of the
% explicit (2.1); (2.1) is prop:isotropic-polar-reading. Provenance: Q-0208, on Q-0205 (c) by
% reduction. [V, Thm. 7.3] is restated in the proof, cited by its v0.1 number.
% CHANGED Q-0281: typed (the statement skeleton). The uniqueness clause is NARROWED: "a real function
% F' on [0,infty)" becomes "a real function F' on [0,infty) with F'(0) = 0". As written the clause was
% false: only differences of F' enter eq:isotropic-classification, so F' := F + 1 satisfies it and is
% not F(c^{-1} .); the proof's "taking logarithms of the kernels from the origin" reads F'(0) = 0. The
% proof's sentence says so now; nothing else changed.
% CHANGED Q-0281 wave 1: blind review of group E. W1E-6: "unique up to ...: if" now reads "unique up
% to ..., in the sense that if", so the clause is not read as also asserting the converse, which the
% Lean (isotropic_classification_unique) does not state. W1E-2: the annotation's sentence on what
% sufficiency reads is reworded to the rewritten proof (continuity of F, F(0) = 0, and F not
% identically 0 through [V, Lem. 6.3]); the law and residuals of A_d are read by the split child
% lem:isotropic-realisation only. The sentence "In the Lean, ..." was edited in wave 1E without a
% marker; this line marks it.
% CHANGED Q-0365: the Lean switched to the shared library's ScaleSpace.main_analysis, which takes the
% uniqueness clause as a hypothesis; AffineHemigroup.line_levy_unique (prop:line-levy-unique, now
% proved from this module's A3 at d = 1) discharges it, so isotropic_classification prints this
% module's A3(ii)-(iii) and no interface of Paper V's. The comment inside the environment is the
% wave-1 record and is kept (the papers carry it verbatim); the status sentence "In the Lean, ..."
% is updated to match. The text of the statement is unchanged.
% CHANGED Q-0417 (the freeze of the isotropic paper; REPOSITIONING-paper-iso.md section 19): "the
% kernels from the origin" in the statement is "the kernels from scale zero", the papers' wording
% for the kernels mu_{0,t} (in this chapter "origin" is the centre of R^d). Wording only; both
% papers' copies are re-transcribed.
\begin{theorem}[the isotropic classification in dimension $d$]
\label{thm:isotropic-classification}
\uses{def:isotropic-kernel-family,lem:radial-marginal,lem:isotropic-marginal-family,def:filtration-thresholds,prop:line-levy-unique,thm:cited-A3}
% CHANGED 2026-10-08, the author's decision: regraded [T] -> [A] on ledger A3, since Q-0365 the
% necessity direction's uniqueness clause is prop:line-levy-unique, proved from A3(ii)-(iii) at d = 1;
% \uses gains thm:cited-A3 (the chapter's convention for an [A] node). The statement is unchanged.
% Q-0281 wave 1 (group E): proved. AffineHemigroup.isotropic_classification prints, besides Lean
% core, Paper V's cited interface SpatialLine.fourier_toolbox_levy_unique ([V]'s ledger A3, Sato
% Thm. 8.1, uniqueness of the Levy pair), through SpatialLine.main_analysis; it is in this module's
% trust base by reference (CLAUDE.md, "Trust"). It is listed in CIAxiomGuard.lean, V's interface
% admitted through the include line of blueprint/trust-boundary.txt (V's boundary at the manifest's
% pin; before article-kit read an include, Q-0298, it was proved and unguarded by the author's
% decision of 2026-10-05, blind review W1E-3). The other two declarations of the tag, and the
% sufficiency direction, are Lean core and guarded.
\lean{AffineHemigroup.isotropic_classification, AffineHemigroup.isotropic_classification_origin,
AffineHemigroup.isotropic_classification_unique}\leanok
Let $d \ge 1$. A family $(\mu_{s,t})_{0 \le s \le t}$ of probability measures on $\RR^d$ is an
isotropic cascade family in dimension $d$ if and only if there are a gauge $\chi$ and an
$F \in \mathcal{A}_d$ with $F \not\equiv 0$ such that
\begin{equation}\label{eq:isotropic-classification}
  \widehat\mu_{s,t}(\xi) = \exp\bigl[-\bigl(F(\chi(t)|\xi|) - F(\chi(s)|\xi|)\bigr)\bigr],
  \qquad 0 \le s \le t,\ \xi \in \RR^d .
\end{equation}
In particular the kernels from scale zero are $\widehat\mu_{0,t}(\xi) = e^{-F(\chi(t)|\xi|)}$. The
pair is unique up to $(\chi,F) \mapsto (c\,\chi,\,F(c^{-1}\,\cdot))$ with $c > 0$, in the sense that if a gauge $\chi'$
and a real function $F'$ on $[0,\infty)$ with $F'(0) = 0$ also satisfy
\ref{eq:isotropic-classification}, then
$\chi' = c\,\chi$ and $F' = F(c^{-1}\,\cdot)$ for some $c > 0$, and $c = 1$ when
$\chi(1) = \chi'(1) = 1$.
\statusA\quad\ledger{A3}\quad\emph{(\textbf{Assignment.} \ledger{A3}(ii)--(iii), at $d = 1$,
carries the one analytic fact the necessity direction reads: the uniqueness of a symmetric \Levy{}
pair, which \ref{prop:line-levy-unique} proves from it. Everything else is the released line theorem
[V, Thm.~7.3], restated in the proof where it is used, its structural part reproved through the
shared library's form of it, [V, Lem.~6.3], proved there in three lines, and nodes of this module.
Graded \texttt{[A]} on 2026-10-08 by the author's decision; it was \texttt{[T]} while that clause was
[V]'s cited interface. \textbf{Where the analysis is.} Necessity is [V, Thm.~7.3] applied to the marginal family of
\ref{lem:isotropic-marginal-family}; sufficiency needs neither that theorem nor Bochner's: it reads
only that $F$ is continuous with $F(0) = 0$, and $F \not\equiv 0$ through [V, Lem.~6.3]. The law
in $\mathcal{A}_d$ and its residuals are read only by \ref{lem:isotropic-realisation}, which builds
a family from each pair. In the prose the trust passes through [V]'s ledger along that theorem's
necessity direction, cited by release and not copied here. In the Lean, [V]'s necessity direction from kernels, \texttt{main\_analysis}
applied to the marginals, reads only the uniqueness clause of [V]'s A3
(\texttt{notes/PLAN-d-dimensional-port.md}, S11), which the shared library's form of that theorem
takes as a hypothesis and \ref{prop:line-levy-unique} proves from this module's A3 at $d = 1$; the
Lean spends no interface of [V]'s. \textbf{Against the draft's Theorem~2.2}, each point accepted by the author as the statement of
record (Q-0222, 2026-09-30). (i) It is stated on kernels, as
\ref{def:isotropic-kernel-family} is; the operator form is Q-0207. (ii) It names the class
$\mathcal{A}_d$ where the draft writes the explicit polar form (2.1); the two are the same class by
\ref{prop:isotropic-polar-reading}, which is \texttt{[A]} on ledger A2--A4, so this node carries no
entry; the draft's (2.1) is that node's clause (3), and its normalisation is recorded in that
node's annotation. The word \emph{gauge} and the condition \ref{eq:self-decomposable} are restated
in \ref{def:isotropic-kernel-family}'s annotation. (iii) The draft's port
notes --- the Wendel step, the null array with $|x|^2$, the no-lattice argument, the gauge, Sato's
polar form --- are not used: the proof consumes [V, Thm.~7.3] whole. (iv) The uniqueness clause is
added; the draft does not state one. Its normalisation $F'(0) = 0$ was added on 2026-10-04 (Q-0281,
typing the statement): without it the clause is false, $F' = F + 1$ satisfying
\ref{eq:isotropic-classification}, whose exponent reads only differences of $F'$. At $d = 1$ the theorem is [V, Thm.~7.3] read on kernels,
$\mathcal{A}_1$ being the whole admissible cone ([B, Prop.~4.9]).)}
\end{theorem}

% CHANGED Q-0281 wave 1: proof rewritten to the Lean route. Necessity restates [V, Thm. 7.3] in the
% form of its necessity direction that reads the representing kernels as given, so the
% identification of kernels becomes the remark that the two forms agree; the residual identity of
% F in A_d is read off eq:isotropic-classification. Sufficiency verifies the axioms for the GIVEN
% family, which is what the equivalence asserts; the construction of a family from (chi, F), which
% the previous proof used to the same end, is kept as the last paragraph, marked as not needed for
% the statement and not machine-checked. Uniqueness is unchanged.
% CHANGED Q-0281 wave 1: blind review W1E-1: that last paragraph, "Every pair is realised", is moved
% out of this \leanok proof into the child node lem:isotropic-realisation, whose proof it is and
% which is machine-checked there.
\begin{proof}
\leanok
\emph{Necessity.} Let $(\mu_{s,t})$ be an isotropic cascade family and let $\kappa_{s,t}$ and
$\Phi_{s,t}$ be as in \ref{lem:isotropic-marginal-family}, which says that $(\Phi_{s,t})$
satisfies (A1)--(A8) and (ND) of \restated{V}{Def.~3.1}{def:line-family} and is convolution by the symmetric probability
measures $\kappa_{s,t}$. \restated{V}{Thm.~7.3}{thm:line-classification}, restated in the form of its necessity direction that this
proof reads: \emph{if a family of operators on $L^1(\RR)$ satisfying (A1)--(A8) and (ND) is of the
form $\Phi_{s,t}f = \nu_{s,t} * f$ for probability measures $\nu_{s,t}$ on $\RR$, then
$\widehat\nu_{s,t}(\omega) = \exp[-(F(\chi(t)\omega) - F(\chi(s)\omega))]$ for a gauge $\chi$ and
an $F \not\equiv 0$ of the form \ref{eq:admissible-line}.} (The theorem as \progpaper{V} states it also
produces the kernels $\nu_{s,t}$; the two forms agree because a family has at most one kernel at
each pair: two probability measures $\kappa, \nu$ on $\RR$ with $\kappa * f = \nu * f$ for every
$f \in L^1(\RR)$ are equal, since at $f$ the standard Gaussian density, whose transform is
positive, $\widehat\kappa\widehat f = \widehat\nu\widehat f$ gives $\widehat\kappa = \widehat\nu$.)
Applied with $\nu_{s,t} = \kappa_{s,t}$, and with $\widehat\mu_{s,t}(\xi) =
\widehat\kappa_{s,t}(|\xi|)$ from \ref{lem:isotropic-marginal-family}, it gives
\ref{eq:isotropic-classification}.

It remains to see $F \in \mathcal{A}_d$. $F$ is admissible, since it has the form
\ref{eq:admissible-line}, and $F(0) = 0$. Put $t_1 := \chi^{-1}(1)$ and $\mu := \mu_{0,t_1}$, so
that $\widehat\mu(\xi) = e^{-F(|\xi|)}$, because $\chi(0) = 0$. For $\lambda \in (0,1)$ put
$s_\lambda := \chi^{-1}(\lambda) \in (0,t_1)$. By \ref{eq:isotropic-classification},
$\widehat\mu(\lambda\xi) = e^{-F(\lambda|\xi|)}$ and
$\widehat\mu_{s_\lambda,t_1}(\xi) = e^{-(F(|\xi|) - F(\lambda|\xi|))}$, so
\[
  \widehat\mu(\xi) = \widehat\mu(\lambda\xi)\,\widehat\mu_{s_\lambda,t_1}(\xi) ,
\]
which is \ref{eq:self-decomposable} with residual $\rho_\lambda = \mu_{s_\lambda,t_1}$. So $\mu$ is
self-decomposable, $\xi \mapsto F(|\xi|)$ is its exponent, and $F \in \mathcal{A}_d$.

\emph{Sufficiency.} Let $\chi$ be a gauge and $F \in \mathcal{A}_d$ with $F \not\equiv 0$, and let
$(\mu_{s,t})$ be a family of probability measures on $\RR^d$ satisfying
\ref{eq:isotropic-classification}. First, $F$ is
continuous on $[0,\infty)$ with $F(0) = 0$: on $0 \le r \le \Omega$ the integrand of
\ref{eq:admissible-line} is at most $\max\{2,\Omega^2/2\}\,(1 \wedge x^2)\,k(x)/x$, which is
integrable by the two conditions on $k$, so dominated convergence applies. The axioms are read off
the transform, a probability measure on $\RR^d$ being determined by its
transform. (I1): the exponents of $\mu_{s,t}$ and $\mu_{r,s}$ add up to that of $\mu_{r,t}$, and
the exponent of $\mu_{t,t}$ is $0$. (I2): $F$ is continuous, and so is $\chi$, an increasing
bijection of $[0,\infty)$. (I3): $\widehat{Q_*\mu_{s,t}}(\xi) = \widehat\mu_{s,t}(Q^{\tp}\xi)$ and
$|Q^{\tp}\xi| = |\xi|$. (I4): with $S_\lambda := \chi^{-1}(\lambda\,\chi(\cdot))$, an increasing
bijection of $[0,\infty)$,
$\widehat{(\lambda I)_*\mu_{s,t}}(\xi) = \widehat\mu_{s,t}(\lambda\xi)
= \widehat\mu_{S_\lambda s,S_\lambda t}(\xi)$. (ND): let $s < t$ and suppose
$\mu_{s,t} = \delta_0$, so that $F(\chi(t)r) = F(\chi(s)r)$ for every $r \ge 0$. If $s = 0$ this
says $F(\chi(t)r) = 0$ for every $r$, and $\chi(t) > 0$ gives $F \equiv 0$. If $s > 0$ it says
$F(r) = F(\varkappa r)$ for every $r$, with $\varkappa := \chi(s)/\chi(t) \in (0,1)$. Iterating,
$F(r) = F(\varkappa^nr) \to F(0) = 0$, so again $F \equiv 0$ (\restated{V}{Lem.~6.3}{sec:line-facts}). Both contradict
$F \not\equiv 0$.

\emph{Uniqueness.} Let $(\chi',F')$ also satisfy \ref{eq:isotropic-classification}. Both exponents
are real, so taking logarithms of the kernels from scale zero, $F(\chi(t)r) = F'(\chi'(t)r)$ for
all $t, r \ge 0$, since $\chi(0) = \chi'(0) = 0$ and $F(0) = F'(0) = 0$. At $t' := \chi'^{-1}(1)$ this reads $F'(r) = F(c^{-1}r)$ with
$c^{-1} := \chi(t') > 0$. Hence $F(\chi(t)r) = F(c^{-1}\chi'(t)r)$ for all $t$ and $r$. If
$\chi(t) \ne c^{-1}\chi'(t)$ for some $t > 0$, then $F(r) = F(\varkappa r)$ for every $r$ with
$\varkappa := c^{-1}\chi'(t)/\chi(t) \ne 1$, and $F \equiv 0$ as in the previous paragraph (with
$\varkappa^{-1}$ in place of $\varkappa$ if $\varkappa > 1$), a contradiction. So
$\chi' = c\,\chi$. When $\chi(1) = \chi'(1) = 1$ this gives $c = 1$.
\end{proof}

% not in the draft as a statement: queue item Q-0407 (hub), the author's decision of 2026-10-10 on
% pull request #100 ("The 'classification on operators' should be a Lean corollary"). Since Q-0406
% the isotropic paper opens with the axioms on operators (def:isotropic-operator-family) and said
% after this theorem, in two sentences of prose, that the theorem classifies the operators too; no
% node stated it, and its counterpart in chapter 4, thm:affine-rigidity-operators, has been a node
% since Q-0314. This node is that sentence, on the model of thm:affine-rigidity-operators. No
% statement of another node changes.
% The environment is `theorem` because the blueprint has no `corollary` environment (the
% framework's theorems.tex declares four kinds); the label says what the statement is, and the
% isotropic paper prints it as a corollary.
% Typed and proved in the same item, the declarations going straight into the proved library
% (AffineHemigroup.IsotropicOperators): assembly, priced when the target type was written at about
% 100 Lean lines, which is what it took. What writing the type down found.
% (i) The measures have to be named. The paper's sentence said "for probability measures mu_{s,t} of
% the form (eq:isotropic-classification)", which puts the measures under the same existential
% quantifier as the pair; the uniqueness clause must then take a second family mu' with its own
% pair, and it concludes mu' = mu before it can use the theorem's clause, by the uniqueness of
% lem:operators-to-kernels(1). So the clause is not the theorem's clause alone, and it says so.
% (ii) "Phi_{s,t} f = mu_{s,t} * f" is typed as the hypothesis form the two passage lemmas already
% take (for every pair, mu_{s,t} is a probability measure and Phi_{s,t} = mconvL1E mu_{s,t}); no
% new definition was needed.
% (iii) The necessity half returns a normalised gauge, chi(1) = 1, as the theorem's does
% (isotropic_classification_operators_necessity); the statement does not claim it, as the theorem's
% does not.
% Axioms, printed before this annotation was written: isotropic_classification_operators prints
% what isotropic_classification prints (propext, Classical.choice, Quot.sound,
% levy_khintchine_unique, levy_khintchine_converse) and isotropic_classification_operators_unique
% prints Lean core. Both are pinned in CIAxiomGuard.lean and are headline declarations of
% linkage.toml [boundary].
\begin{theorem}[the isotropic classification on operators]
\label{cor:isotropic-classification-operators}
\uses{thm:isotropic-classification,def:isotropic-operator-family,def:isotropic-kernel-family,def:filtration-thresholds,lem:operators-to-kernels,lem:operators-to-kernels-families,lem:kernels-to-operators,thm:cited-A3}
\lean{AffineHemigroup.isotropic_classification_operators,
AffineHemigroup.isotropic_classification_operators_unique}\leanok
Let $d \ge 1$. A family $(\Phi_{s,t})_{0 \le s \le t}$ of linear operators on $L^1(\RR^d)$ is an
isotropic cascade measurement family in dimension $d$ (\ref{def:isotropic-operator-family}) if and
only if there are a gauge $\chi$, an $F \in \mathcal{A}_d$ with $F \not\equiv 0$, and probability
measures $\mu_{s,t}$ on $\RR^d$ satisfying \ref{eq:isotropic-classification}, such that
\[
  \Phi_{s,t}f = \mu_{s,t} * f, \qquad 0 \le s \le t,\ f \in L^1(\RR^d) .
\]
The measures are then unique, and the pair is unique up to
$(\chi,F) \mapsto (c\,\chi,\,F(c^{-1}\,\cdot))$ with $c > 0$, in the sense that if a gauge $\chi'$,
a real function $F'$ on $[0,\infty)$ with $F'(0) = 0$, and probability measures $\mu'_{s,t}$ on
$\RR^d$ satisfying \ref{eq:isotropic-classification} with $(\chi', F')$ in place of $(\chi, F)$
also give $\Phi_{s,t}f = \mu'_{s,t} * f$ for all $0 \le s \le t$ and all $f \in L^1(\RR^d)$, then
$\mu'_{s,t} = \mu_{s,t}$ for all $0 \le s \le t$, and $\chi' = c\,\chi$ and
$F' = F(c^{-1}\,\cdot)$ for some $c > 0$, and $c = 1$ when $\chi(1) = \chi'(1) = 1$.
\statusA\quad\ledger{A3}\quad\emph{(\textbf{Assignment.} Nothing is spent here that
\ref{thm:isotropic-classification} does not spend: \ledger{A3}(ii)--(iii), at $d = 1$, through that
theorem's necessity direction. The node itself cites nothing. It assembles
\ref{lem:operators-to-kernels}(1) and~(5), \ref{thm:isotropic-classification} and
\ref{lem:kernels-to-operators}(3), and is \ref{thm:isotropic-classification} read on the operator
axioms of \ref{def:isotropic-operator-family}, as \ref{thm:affine-rigidity-operators} is
\ref{thm:affine-rigidity} read on those of the case $(\GL(d), \text{Loewner})$. The sufficiency
direction and the uniqueness clause spend no entry. \textbf{Readings.} In the Lean, (A1) is the
type of $\Phi_{s,t}$, a bounded linear operator, as in \ref{def:isotropic-operator-family}; this
loses nothing, since an isotropic cascade measurement family consists of bounded operators by (A1)
and an operator $f \mapsto \mu * f$ is bounded. The families $(\Phi_{s,t})$ and $(\mu_{s,t})$ are
read on the pairs $0 \le s \le t$ only, and the family $(S_\lambda)$ of (A8) is under an
existential quantifier, as in the statement. The uniqueness clause names a second family
$(\mu'_{s,t})$ because the measures stand under the same existential quantifier as the pair; that
they are determined by the operators is the uniqueness clause of
\ref{lem:operators-to-kernels}(1), and the rest is the uniqueness clause of
\ref{thm:isotropic-classification}, whose normalisation $F'(0) = 0$ is kept. As there, the converse
of the clause is not asserted. \textbf{At $d = 1$} the statement is the classification on the line,
\restated{V}{Thm.~7.3}{thm:line-classification}: the definition is then
\restatedp{V}{Def.~3.1}{def:line-family}, since $\Orth(1) = \{\pm1\}$,
$\mathcal{A}_1$ is the whole admissible cone
\restatedc{B}{Prop.~4.9}{prop:line-filtration}{(1)}{(1)}, and a measure satisfying
\ref{eq:isotropic-classification} is symmetric. \textbf{Not claimed}: the kernels from scale zero,
which are those of \ref{thm:isotropic-classification}; that every pair $(\chi, F)$ occurs, which is
\ref{lem:isotropic-realisation} followed by \ref{lem:kernels-to-operators}(3); and anything about
operators on $L^p(\RR^d)$ for $p > 1$.)}
\end{theorem}

\begin{proof}
\leanok
\emph{Necessity.} Let $(\Phi_{s,t})$ be an isotropic cascade measurement family in dimension $d$.
By \ref{lem:operators-to-kernels}(1) each $\Phi_{s,t}$ is convolution by a probability measure
$\mu_{s,t}$ on $\RR^d$, and by \ref{lem:operators-to-kernels}(5) the family $(\mu_{s,t})$ is an
isotropic cascade family in dimension $d$. The necessity direction of
\ref{thm:isotropic-classification} gives a gauge $\chi$ and an $F \in \mathcal{A}_d$ with
$F \not\equiv 0$ satisfying \ref{eq:isotropic-classification}.

\emph{Sufficiency.} Let $\chi$, $F$ and $(\mu_{s,t})$ be as in the statement. By the sufficiency
direction of \ref{thm:isotropic-classification} the family $(\mu_{s,t})$ is an isotropic cascade
family in dimension $d$. By \ref{lem:kernels-to-operators}(3) the operators
$f \mapsto \mu_{s,t} * f$, which are the $\Phi_{s,t}$, form an isotropic cascade measurement family
in dimension $d$.

\emph{Uniqueness.} Let $\chi'$, $F'$ and $(\mu'_{s,t})$ be as in the last sentence of the
statement. For $0 \le s \le t$ the probability measures $\mu_{s,t}$ and $\mu'_{s,t}$ give the same
operator $\Phi_{s,t}$, so they are equal, by the uniqueness clause of
\ref{lem:operators-to-kernels}(1). Hence $(\chi', F')$ satisfies
\ref{eq:isotropic-classification} for the family $(\mu_{s,t})$, and the uniqueness clause of
\ref{thm:isotropic-classification} gives $c$.

Every step is machine-checked. The necessity direction spends what that of
\ref{thm:isotropic-classification} spends, and the other two directions spend no cited fact.
\end{proof}

% not in the draft as a statement: the last paragraph of the proof of thm:isotropic-classification
% until the blind review of wave 1E, where it stood inside a \leanok proof without being
% machine-checked.
% CHANGED Q-0281 wave 1: split from thm:isotropic-classification's proof (blind review W1E-1). The
% statement is what that paragraph proved; the proof is the paragraph, with the case t = 0 folded
% into the one formula (chi(0) = 0 makes the pair (0, 0) read rho_0 pushed by the zero map, which
% is delta_0), as the Lean has it. Proved: AffineHemigroup.isotropic_realisation, Lean core only,
% guarded.
\begin{lemma}[every pair is realised]
\label{lem:isotropic-realisation}
\uses{thm:isotropic-classification,def:isotropic-kernel-family,def:filtration-thresholds}
\lean{AffineHemigroup.isotropic_realisation}\leanok
Let $d \ge 1$, let $\chi$ be a gauge, and let $F \in \mathcal{A}_d$ with $F \not\equiv 0$. Then
there is an isotropic cascade family $(\mu_{s,t})_{0 \le s \le t}$ in dimension $d$ that satisfies
\ref{eq:isotropic-classification} with the pair $(\chi, F)$.
\statusT\quad\emph{(Cites nothing beyond the uniqueness of the Fourier transform and the
sufficiency direction of \ref{thm:isotropic-classification}; no ledger entry, and the Lean is Lean
core: it does not pass through [V]'s classification. \textbf{What it is worth.} With
\ref{thm:isotropic-classification} it makes the pairs $(\chi, F)$, $\chi$ a gauge and
$F \in \mathcal{A}_d$ with $F \not\equiv 0$, a parametrisation of the isotropic cascade families in
dimension $d$, up to the rescaling of that theorem's uniqueness clause; the classification alone
says only that every family has such a pair.)}
\end{lemma}

\begin{proof}
\leanok
Let $\mu$ be a self-decomposable probability measure on $\RR^d$ with
$\widehat\mu(\xi) = e^{-F(|\xi|)}$, which the definition of $\mathcal{A}_d$ provides. For
$c \in (0,1)$ let $\rho_c$ be a residual of $\mu$. As $\widehat\mu(c\,\xi) = e^{-F(c|\xi|)} \ne 0$,
\ref{eq:self-decomposable} forces $\widehat\rho_c(\xi) = e^{-(F(|\xi|) - F(c|\xi|))}$; put
$\rho_0 := \mu$ and $\rho_1 := \delta_0$, so that this formula holds for every $c \in [0,1]$, as
$F(0) = 0$. For $0 \le s \le t$ put
\[
  \mu_{s,t} := \bigl(\chi(t)I\bigr)_*\,\rho_{\chi(s)/\chi(t)} ,
\]
with $\chi(s)/\chi(t) := 0$ when $\chi(t) = 0$, that is, when $s = t = 0$; then $\mu_{0,0}$ is
the image of $\mu$ under the zero map, $\delta_0$. Since
$0 \le \chi(s) \le \chi(t)$, the index lies in $[0,1]$, and
\[
  \widehat\mu_{s,t}(\xi) = \widehat\rho_{\chi(s)/\chi(t)}\bigl(\chi(t)\xi\bigr)
  = \exp\bigl[-\bigl(F(\chi(t)|\xi|) - F(\chi(s)|\xi|)\bigr)\bigr] ,
\]
because $\chi(s)/\chi(t)\cdot\chi(t)|\xi| = \chi(s)|\xi|$, both sides being $0$ when
$\chi(t) = 0$. Each $\mu_{s,t}$ is a probability measure, the image of one, so by the sufficiency
direction of \ref{thm:isotropic-classification} $(\mu_{s,t})$ is an isotropic cascade family, and
it satisfies \ref{eq:isotropic-classification} with $(\chi, F)$ by construction. All of this is
machine-checked.
\end{proof}

% not in the draft as a statement: queue item Q-0434 (hub), from the second external round of the
% isotropic paper (records/affine/PLAN-review-response-paper-iso-round2.md; the verification record
% records/affine/reviews/VERIFY-round2-mathematics.md, item 5a). Two reviewers asked that the main
% result, which the paper's section 1 states in a sentence ("the rotation-invariant
% self-decomposable laws on R^d are exactly the kernels from scale zero of the isotropic scale
% spaces that the axioms define"), be a numbered statement. This node is that sentence, with the
% qualifications it needs: the point mass at the origin is excluded, and the scale is positive.
% The environment is `theorem` because the blueprint has no `corollary` environment; the label
% says what the statement is, as for cor:isotropic-classification-operators.
% Typed and proved in the same item (Formalization/AffineHemigroup/ScaleZeroKernels.lean, about
% 320 lines against the estimate of 100 to 150; the difference is clause (3), which the estimate
% did not count). What writing the type down found.
% (i) The step the verification record calls unwritten, that the transform of a rotation-invariant
% self-decomposable law is e^{-F(|xi|)} with F >= 0, was in the library already:
% selfDecomposable_charFun_ne_zero (prop:gw-ray-families(1), the shared library's doubling
% inequality, Lean core) and exists_exp_of_rotationInvariant (written for
% prop:shape-adapted-apertures(3)). So the Lean does not read A4 with A3(i) for it. The printed
% proof does, because those are statements this chapter's reader has; the cited facts of the node
% are the same either way, admissibility of F being prop:isotropic-polar-classification(2).
% (ii) "The kernel from scale zero to some scale" has to exclude the scale 0, where the kernel is
% delta_0 by (I1): the statement says t > 0.
% (iii) "The family with that kernel is unique up to the relabelling of the scales" needs a
% normalisation to be an equivalence. A family has many pairs (chi, F), related by the rescaling of
% thm:isotropic-classification, and fixing mu_{0,t_1} = mu fixes F(chi(t_1) .). With F the exponent
% of mu itself the condition is chi(t_1) = 1, and then the families with the kernel mu at the scale
% t_1 and the gauges with chi(t_1) = 1 correspond one to one. Clause (3) is stated so.
% (iv) The operator clause is the equivalence only. The families of operators correspond to the
% families of kernels by cor:isotropic-classification-operators, which clause (3) does not repeat.
% Axioms, printed before this annotation was written: scale_zero_kernels, scale_zero_operators,
% scale_zero_exponent and scale_zero_families print propext, Classical.choice, Quot.sound,
% sato_polar_criterion, levy_khintchine_exists, levy_khintchine_unique, levy_khintchine_converse
% and selfDecomposable_infinitelyDivisible (A2, A3(i)-(iii), A4, what polar_reading_admissible
% prints); scale_zero_kernel_of_family and scale_zero_family_iff print levy_khintchine_unique and
% levy_khintchine_converse (what isotropic_classification prints); scale_zero_exponent_unique,
% scale_zero_dilation, scale_zero_gauge_unique and scale_zero_family_unique print Lean core.
% CHANGED after the blind check of Q-0434 (records/formalization/blind-check-q0434-report-B.md;
% no error, one gap, six notes). F-B2, the gap: the closing sentence of the statement said "up
% to the relabelling of the scales" and no clause, proof step or Lean conjunct stated it; the
% annotation had it as a reading. It is now clause (4), with a proof paragraph and its Lean:
% scale_zero_relabelling (two isotropic cascade families with the same kernel at t_1; prints A2,
% A3, A4, as scale_zero_exponent) and scale_zero_relabelling_display (theta with chi o theta =
% chi', Lean core), which is also the last conjunct of scale_zero_families. No cited fact is
% added. F-B3: scale_zero_gauge_unique asks t_1 > 0 (at t_1 = 0 it was vacuous). F-B4, F-B5,
% F-B6, F-B7: wording, applied; the title says "other than delta_0". F-B8 is not applied: the
% environment stays `theorem`, because the framework's theorems.tex declares four kinds and no
% corollary, and that file is article-kit's; the paper prints the statement as a corollary, as
% it does cor:isotropic-classification-operators.
% CHANGED Q-0447 (the third pass of the blind check, on clause (4): faithful; four notes). F-B9:
% the docstring of scale_zero_families is headed "(3) and the second sentence of (4)" and
% describes its last conjunct. F-B10: scale_zero_relabelling_display asks t_1 > 0 (at t_1 = 0 its
% premises chi(0) = chi'(0) = 1 are false for a gauge and it was vacuous); its two callers had
% t_1 > 0 already. F-B11, F-B12: two sentences of the annotation ("clauses (3) and (4)"; the
% history sentence on the relabelling is gone). No statement, proof or axiom report changes.
\begin{theorem}[the kernels from scale zero are the rotation-invariant self-decomposable laws other than $\delta_0$]
\label{cor:isotropic-scale-zero-kernels}
\uses{thm:isotropic-classification,lem:isotropic-realisation,prop:isotropic-polar-classification,def:isotropic-kernel-family,def:isotropic-operator-family,def:filtration-thresholds,lem:radial-marginal,lem:operators-to-kernels,lem:operators-to-kernels-families,lem:kernels-to-operators,thm:cited-A2,thm:cited-A3,thm:cited-A4}
\lean{AffineHemigroup.scale_zero_kernels, AffineHemigroup.scale_zero_operators,
AffineHemigroup.scale_zero_families, AffineHemigroup.scale_zero_relabelling}\leanok
Let $d \ge 1$ and let $\mu$ be a probability measure on $\RR^d$. Consider the condition
\begin{enumerate}
\item[(S)] $\mu$ is rotation invariant, that is $Q_*\mu = \mu$ for every $Q \in \Orth(d)$, it is
self-decomposable, and $\mu \ne \delta_0$.
\end{enumerate}
\begin{enumerate}
\item \emph{(On kernels.)} (S) holds if and only if there are an isotropic cascade family
$(\mu_{s,t})_{0 \le s \le t}$ in dimension $d$ and a scale $t > 0$ with $\mu_{0,t} = \mu$.
\item \emph{(On operators.)} (S) holds if and only if there are an isotropic cascade measurement
family $(\Phi_{s,t})_{0 \le s \le t}$ in dimension $d$ (\ref{def:isotropic-operator-family}) and
a scale $t > 0$ with $\Phi_{0,t}f = \mu * f$ for every $f \in L^1(\RR^d)$.
\item \emph{(The families with this kernel.)} Suppose (S). Then there is an $F \in \mathcal{A}_d$
with $F \not\equiv 0$ and $\widehat\mu(\xi) = e^{-F(|\xi|)}$ for every $\xi \in \RR^d$, and $\mu$
determines $F$ on $[0,\infty)$. Let $t_1 > 0$. A family $(\mu_{s,t})_{0 \le s \le t}$ of
probability measures on $\RR^d$ is an isotropic cascade family in dimension $d$ with
$\mu_{0,t_1} = \mu$ if and only if it satisfies \ref{eq:isotropic-classification} with this $F$
and a gauge $\chi$ with $\chi(t_1) = 1$. The gauge is then unique. For every gauge $\chi$ with
$\chi(t_1) = 1$ there is exactly one such family, and its kernels from scale zero are the
dilations of $\mu$: $\mu_{0,t} = (\chi(t)I)_*\mu$ for every $t > 0$.
\item \emph{(Up to the relabelling of the scales.)} Let $t_1 > 0$, and let
$(\mu_{s,t})_{0 \le s \le t}$ and $(\mu'_{s,t})_{0 \le s \le t}$ be isotropic cascade families
in dimension $d$ with $\mu'_{0,t_1} = \mu_{0,t_1}$. Then there is an increasing bijection
$\theta$ of $[0,\infty)$ with $\theta(t_1) = t_1$ and
$\mu'_{s,t} = \mu_{\theta(s),\theta(t)}$ for all $0 \le s \le t$. If $\chi$ and $\chi'$ are the
gauges that clause (3) gives the two families, for the law $\mu = \mu_{0,t_1}$, then
$\theta := \chi^{-1} \circ \chi'$ is such a bijection.
\end{enumerate}
So the rotation-invariant self-decomposable laws on $\RR^d$ other than the point mass at the
origin are exactly the kernels from scale zero to a positive scale of the isotropic cascade
families in dimension $d$, and a family with a given kernel from scale zero to a given scale is
determined by the gauge (clause (3)), that is up to the relabelling of the scales (clause (4)).
\statusA\quad\ledger{A2}\quad\ledger{A3}\quad\ledger{A4}\quad\emph{(\textbf{Assignment.} The node
cites nothing of its own. \ledger{A3}(ii)--(iii), at $d = 1$, is spent through the necessity
direction of \ref{thm:isotropic-classification}, in the direction from a family to its kernel.
\ledger{A2}, \ledger{A3} and \ledger{A4} are spent through
\ref{prop:isotropic-polar-classification}(2), which makes the function $F$ of a law with (S)
admissible, in the direction from the law to a family. The printed proof spends \ledger{A4} and
\ledger{A3}(i) once more, for the fact that the transform of a self-decomposable law has no zero.
\textbf{What the citations carry and what they do not.} They carry the classification on the
line and the polar form of the self-decomposable laws. The equivalences are assembled here from
\ref{thm:isotropic-classification}, \ref{lem:isotropic-realisation} and
\ref{prop:isotropic-polar-classification}(2), and no source states them. \textbf{What it is.}
The main result of the isotropic paper as one statement. \ref{thm:isotropic-classification} says
that every isotropic family has a pair $(\chi, F)$ with $F \in \mathcal{A}_d$, and
\ref{lem:isotropic-realisation} that every pair occurs. Neither says which laws are kernels.
The step between them is that a law with (S) has the transform $e^{-F(|\cdot|)}$ with $F \ge 0$,
for which the law must be self-decomposable as well as rotation invariant: the transform of a
rotation-invariant law is real and radial, and it is positive because it has no zero.
\textbf{What it is worth.} A designer may take as aperture any rotation-invariant
self-decomposable law other than $\delta_0$, and no other law. Once the law is chosen, at the scale $t_1$, what is left
is the gauge: how the scales are labelled. The cascade is then determined, stages included.
\textbf{Readings.} In clauses (1) and (2) the scale is positive: at $t = 0$ the kernel is
$\delta_0$, by (I1). In clause (3) ``$\mu$ determines $F$ on $[0,\infty)$'' says that a function
$F'$ with $\widehat\mu(\xi) = e^{-F'(|\xi|)}$ for every $\xi$ equals $F$ at every $r \ge 0$; the
members of $\mathcal{A}_d$ are functions on $[0,\infty)$. ``The gauge is then unique'' is the
uniqueness clause of \ref{thm:isotropic-classification} with the constant fixed by
$\chi(t_1) = 1$: two gauges with the value $1$ at $t_1$ that satisfy
\ref{eq:isotropic-classification} with $F$ for one family agree on $[0,\infty)$. ``Exactly one
such family'': one exists by \ref{lem:isotropic-realisation}, and two families of probability
measures that satisfy \ref{eq:isotropic-classification} with the same pair agree at every pair
$0 \le s \le t$. In the last sentence of the clause $t > 0$, so that $\chi(t)I \in \GL(d)$; at
$t = 0$ the kernel is $\delta_0$. Clause (4) is what ``up to the relabelling of the scales'' says: two
isotropic families with the same kernel at the scale $t_1$ run through the same kernels, at
scales matched by $\theta$.
\textbf{At $d = 1$} a law on the line is rotation invariant when it is symmetric, and the
statement is the corresponding one for the symmetric families of
\restated{V}{Thm.~7.3}{thm:line-classification}. \textbf{Not claimed}: that the stages
$\mu_{s,t}$ with $s > 0$ satisfy (S); they are in general not self-decomposable
(\ref{prop:isotropic-stage-atom}); nor that a law which is a kernel from scale zero in dimension
$d$ is one in another dimension, which is the subject of the filtration
(\ref{thm:filtration-strict}); nor an operator form of clauses (3) and (4), the families of operators
corresponding to the families of kernels by \ref{cor:isotropic-classification-operators}.
\textbf{Formalization} (Q-0434). $\mu$ is a measure on the Euclidean space of dimension $d$
with the instance of a probability measure. Rotation invariance is invariance under the image by
every orthogonal matrix, self-decomposability is the shared library's, and the families are those
of \ref{def:isotropic-kernel-family}'s and \ref{def:isotropic-operator-family}'s Lean, with the
family $(S_\lambda)$ under an existential quantifier. In clause (2), $\Phi_{0,t}f = \mu * f$ for
every $f$ is the equality of $\Phi_{0,t}$ with the operator of convolution by $\mu$. Clause (3) is
one declaration: there is an $F$ in the class, not identically $0$ on $[0,\infty)$, with the
transform of $\mu$, such that every function with that transform equals $F$ on $[0,\infty)$ and,
for every $t_1 > 0$ and every family of probability measures on the pairs $0 \le s \le t$, the
equivalence holds, two gauges as above agree on $[0,\infty)$, two families with the display of
the same function $\chi$, a gauge or not, agree at every pair, the kernels from scale zero of a
family with the display of any gauge, normalised at $t_1$ or not, are the images
$(\chi(t)I)_*\mu$ for every $t \ge 0$ (at $t = 0$ the image under the zero
map, which is $\delta_0$), every gauge with $\chi(t_1) = 1$ has a family of probability
measures with the display, and two families of probability measures with the displays of two
gauges $\chi$ and $\chi'$ that take the value $1$ at $t_1$ have a gauge $\theta$ with
$\theta(t_1) = t_1$, $\chi(\theta(t)) = \chi'(t)$ for every $t \ge 0$ and
$\mu'_{s,t} = \mu_{\theta(s),\theta(t)}$ at every pair. Clause (4) is that last conjunct, which
is its second sentence, and a second declaration, which is its first: it starts from two
isotropic cascade families with $\mu'_{0,t_1} = \mu_{0,t_1}$ and concludes that there is a
gauge $\theta$ with $\theta(t_1) = t_1$ and the equality at every pair. A gauge is an increasing
bijection of $[0,\infty)$; $\theta$ is read on $[0,\infty)$ only. Witness: the standard Gaussian
law satisfies (S) in every dimension
$d \ge 1$. Probes: the Gaussian law is the kernel from scale zero to scale $1$ of the heat
family; $\delta_0$ is the kernel from scale zero to no positive scale of any isotropic
family; and the heat family run at the scales $t^2$ is the heat family relabelled by a $\theta$
with $\theta(t) = t^2$.)}
\end{theorem}

\begin{proof}
\leanok
\emph{A kernel from scale zero satisfies (S).} Let $(\mu_{s,t})$ be an isotropic cascade family in
dimension $d$ and let $t > 0$. Then $\mu_{0,t}$ is rotation invariant by (I3), and
$\mu_{0,t} \ne \delta_0$ by (ND). By \ref{thm:isotropic-classification} there are a gauge $\chi$
and an $F \in \mathcal{A}_d$ with \ref{eq:isotropic-classification}, so that
$\widehat\mu_{0,t}(\xi) = e^{-F(\chi(t)|\xi|)}$. By the definition of $\mathcal{A}_d$
(\ref{def:filtration-thresholds}) there is a self-decomposable probability measure $\mu'$ on
$\RR^d$ with $\widehat{\mu'}(\xi) = e^{-F(|\xi|)}$. Put $\lambda := \chi(t) > 0$. Then
$\widehat{(\lambda I)_*\mu'}(\xi) = \widehat{\mu'}(\lambda\xi) = e^{-F(\lambda|\xi|)}
= \widehat\mu_{0,t}(\xi)$, so $\mu_{0,t} = (\lambda I)_*\mu'$, a probability measure on $\RR^d$
being determined by its transform. The image of a self-decomposable law under $\lambda I$ is
self-decomposable: if $\rho_c$ is a residual of $\mu'$ for the ratio $c \in (0,1)$, then
$\widehat{\mu'}(\lambda\xi) = \widehat{\mu'}(c\lambda\xi)\,\widehat\rho_c(\lambda\xi)$ is
\ref{eq:self-decomposable} for $(\lambda I)_*\mu'$ with the residual $(\lambda I)_*\rho_c$. So
$\mu_{0,t}$ satisfies (S).

\emph{The exponent of a law with (S).} Let $\mu$ satisfy (S). By \ref{lem:radial-marginal}(1)
its transform is real and radial. It has no zero: $\mu$ is infinitely divisible by \ledger{A4},
so by \ledger{A3}(i) it has a generating triplet, and its transform is the exponential of a
complex-valued function. A real continuous function on $\RR^d$ that has no zero and takes the
value $1$ at the origin is positive. Fix a unit vector $e$ and put
$F(r) := -\log\widehat\mu(re)$ for $r \ge 0$. Then $F(0) = 0$, $F \ge 0$ because
$|\widehat\mu| \le 1$, and $\widehat\mu(\xi) = \widehat\mu(|\xi|e) = e^{-F(|\xi|)}$ for every
$\xi$, the transform being radial. So $\xi \mapsto F(|\xi|)$ is the exponent of the
self-decomposable law $\mu$, and $F$ is admissible by
\ref{prop:isotropic-polar-classification}(2). Hence $F \in \mathcal{A}_d$. If $F$ vanished on
$[0,\infty)$ the transform of $\mu$ would be $1$ and $\mu$ would be $\delta_0$, so
$F \not\equiv 0$. A function $F'$ with $\widehat\mu(\xi) = e^{-F'(|\xi|)}$ for every $\xi$ has
$F'(r) = -\log\widehat\mu(re) = F(r)$ for every $r \ge 0$.

\emph{Clause (3), the equivalence.} Let $\mu$ satisfy (S), let $F$ be its exponent, let
$t_1 > 0$, and let $(\mu_{s,t})$ be a family of probability measures on $\RR^d$. Suppose first
that it satisfies \ref{eq:isotropic-classification} with $F$ and a gauge $\chi$ with
$\chi(t_1) = 1$. By the sufficiency direction of \ref{thm:isotropic-classification}, $F$ being in
$\mathcal{A}_d$ and not identically $0$, it is an
isotropic cascade family, and
$\widehat\mu_{0,t_1}(\xi) = e^{-F(\chi(t_1)|\xi|)} = e^{-F(|\xi|)} = \widehat\mu(\xi)$, because
$\chi(0) = 0$ and $F(0) = 0$. So $\mu_{0,t_1} = \mu$. Conversely let $(\mu_{s,t})$ be an
isotropic cascade family with $\mu_{0,t_1} = \mu$. By the necessity direction of
\ref{thm:isotropic-classification} there are a gauge $\chi'$ and an $F' \in \mathcal{A}_d$ with
\ref{eq:isotropic-classification}. Put $c := \chi'(t_1) > 0$. The kernel from scale zero to
scale $t_1$ gives $e^{-F'(c|\xi|)} = \widehat\mu(\xi) = e^{-F(|\xi|)}$, that is $F'(cr) = F(r)$
for every $r \ge 0$. Then $\chi := c^{-1}\chi'$ is a gauge with $\chi(t_1) = 1$, and
$F'(\chi'(t)|\xi|) = F'(c\,\chi(t)|\xi|) = F(\chi(t)|\xi|)$ for every $t \ge 0$. So the family
satisfies \ref{eq:isotropic-classification} with $(\chi, F)$.

\emph{Clause (3), the rest.} Let two gauges $\chi$ and $\chi''$ with
$\chi(t_1) = \chi''(t_1) = 1$ both satisfy \ref{eq:isotropic-classification} with $F$ for one
family. By the uniqueness clause of \ref{thm:isotropic-classification}, applied with $F' = F$,
which vanishes at $0$, there is $c > 0$ with $\chi'' = c\,\chi$, and at $t_1$ this reads $c = 1$.
Let $\chi$ be any gauge with $\chi(t_1) = 1$. By \ref{lem:isotropic-realisation} there is a
family of probability measures that satisfies \ref{eq:isotropic-classification} with
$(\chi, F)$, and by the equivalence it is an isotropic cascade family with $\mu_{0,t_1} = \mu$.
It is the only one, because \ref{eq:isotropic-classification} gives the transform of every
$\mu_{s,t}$. For $t > 0$,
$\widehat{(\chi(t)I)_*\mu}(\xi) = \widehat\mu(\chi(t)\xi) = e^{-F(\chi(t)|\xi|)}
= \widehat\mu_{0,t}(\xi)$, so $\mu_{0,t} = (\chi(t)I)_*\mu$.

\emph{Clause (1).} One direction is the first paragraph. For the other let $\mu$ satisfy (S).
Clause (3) with $t_1 = 1$ and the gauge $\chi = \mathrm{id}$ gives an isotropic cascade family
with $\mu_{0,1} = \mu$.

\emph{Clause (2).} Let $(\Phi_{s,t})$ be an isotropic cascade measurement family in dimension
$d$ and let $t > 0$ with $\Phi_{0,t}f = \mu * f$ for every $f$. By
\ref{lem:operators-to-kernels}(1) each $\Phi_{s,t}$ is convolution by a probability measure
$\mu_{s,t}$, which it determines, and by \ref{lem:operators-to-kernels}(5) the family
$(\mu_{s,t})$ is an isotropic cascade family in dimension $d$. As $\mu_{0,t}$ and $\mu$ give the
same operator, $\mu_{0,t} = \mu$, and (S) holds by clause (1). Conversely let $\mu$ satisfy (S).
By clause (1) there are an isotropic cascade family $(\mu_{s,t})$ and a $t > 0$ with
$\mu_{0,t} = \mu$. By \ref{lem:kernels-to-operators}(3) the operators
$\Phi_{s,t}f := \mu_{s,t} * f$ form an isotropic cascade measurement family in dimension $d$,
and $\Phi_{0,t}f = \mu * f$.

\emph{Clause (4).} Let the two families be as in the clause and put $\mu := \mu_{0,t_1}$. By
the first paragraph $\mu$ satisfies (S). Let $F$ be its exponent, and let $\chi$ and $\chi'$ be
the gauges of the two families by clause (3), so that $\chi(t_1) = \chi'(t_1) = 1$. Then
$\theta := \chi^{-1} \circ \chi'$ is an increasing bijection of $[0,\infty)$ with
$\theta(t_1) = t_1$ and $\chi \circ \theta = \chi'$. Let $0 \le s \le t$, so that
$0 \le \theta(s) \le \theta(t)$. By \ref{eq:isotropic-classification} the transforms of
$\mu'_{s,t}$ and of $\mu_{\theta(s),\theta(t)}$ are both
$\exp[-(F(\chi'(t)|\xi|) - F(\chi'(s)|\xi|))]$, so $\mu'_{s,t} = \mu_{\theta(s),\theta(t)}$.

Every step is machine-checked. In clause (4) the Lean takes $\theta(t)$, for $t \ge 0$, to be
the point of $[0,\infty)$ at which $\chi$ has the value $\chi'(t)$. The Lean follows this proof,
with one difference. That the
transform of a self-decomposable law has no zero is taken from the shared library, which proves
it without a cited fact, by a doubling inequality for the transform of a symmetrized residual; the printed proof reads
it from \ledger{A4} and \ledger{A3}(i). The cited facts that the declarations report are the
same: \ledger{A2}, \ledger{A3} and \ledger{A4}, through
\ref{prop:isotropic-polar-classification}(2) and the necessity direction of
\ref{thm:isotropic-classification}.
\end{proof}

% draft: section 2.1, the explicit form (2.1) of Theorem 2.2, and its bracket "[Sato, Thm. 15.10,
% p. 95, with lambda the uniform measure on the sphere and k_xi = k by rotation invariance]".
% Provenance: Q-0208, step 2 of the plan's section 5; ledger A4 admitted by Q-0205.
% CHANGED Q-0281: typed (the statement skeleton). Clauses (2)-(3) are split to
% prop:isotropic-polar-classification (ADR-0025; records/formalization/PRICING.md: the reading (1),
% which thm:filtration-strict and thm:no-cin-ray-in-the-plane read, against what follows from it); the
% text is unchanged and the \lean tag names the child's declarations too.
% CHANGED Q-0281: \uses re-pointed to the split child it reads (eq:polar-decomposition).
% CHANGED Q-0281 wave 2: the annotation's "Against the draft" names the checked route, "the radial
% image of the Levy measure and the fact that a rotation-invariant measure is determined by its radial
% image", where it said "an average over O(d)" (blind review W2A-4).
\begin{proposition}[the polar form of $\mathcal{A}_d$]
\label{prop:isotropic-polar-reading}
\uses{def:isotropic-origin-datum,def:filtration-thresholds,def:isotropic-kernel-family,thm:isotropic-classification,lem:polar-decomposition,thm:cited-A2,thm:cited-A3,thm:cited-A4}
\lean{AffineHemigroup.polar_reading, AffineHemigroup.polar_reading_unique,
AffineHemigroup.polar_reading_admissible, AffineHemigroup.classA_iff_originDatum,
AffineHemigroup.polar_classification}\leanok
Let $d \ge 1$ and $F : [0,\infty) \to [0,\infty)$.
\begin{enumerate}
\item \emph{(The reading.)} $\xi \mapsto F(|\xi|)$ is the exponent of a self-decomposable law on
$\RR^d$ if and only if it is the exponent $F_d$ of an isotropic origin datum $(a,k_d)$ in
dimension $d$ (\ref{def:isotropic-origin-datum}):
\[
  F(|\xi|) = \tfrac12 a|\xi|^2 + \int_{\RR^d}\bigl(1 - \cos\langle\xi,x\rangle\bigr)\,
    \frac{k_d(|x|)}{|S^{d-1}|\,|x|^d}\,dx , \qquad \xi \in \RR^d .
\]
The datum is then unique, $k_d$ up to modification on a Lebesgue-null set.
\item \emph{(The class.)} If this holds, $F$ is admissible, of the form
\ref{eq:admissible-line} with $a_{\mathrm V} = \tfrac12 a'$ and $k = k_1$ for its datum
$(a',k_1)$ in dimension $1$. So $\mathcal{A}_d$ is the set of $F$ for which $F(|\cdot|)$ is the
exponent of an isotropic origin datum in dimension $d$.
\item \emph{(The classification in polar form.)} A family $(\mu_{s,t})_{0 \le s \le t}$ of
probability measures on $\RR^d$ is an isotropic cascade family in dimension $d$ if and only if
there are a gauge $\chi$ and an isotropic origin datum $(a,k_d)$ in dimension $d$, not both trivial
($a > 0$, or $k_d$ not almost everywhere $0$), with \ref{eq:isotropic-classification} for
\[
  F(r) = \tfrac12 a\,r^2 + \int_{\RR^d}\bigl(1 - \cos(r\langle e_1,x\rangle)\bigr)\,
    \frac{k_d(|x|)}{|S^{d-1}|\,|x|^d}\,dx , \qquad r \ge 0 .
\]
\end{enumerate}
\statusA\quad\ledger{A2}\quad\ledger{A3}\quad\ledger{A4}\quad\emph{(\textbf{Assignment.}
\ledger{A4}, \sscite{sato1999levy}~Prop.~15.5, p.~93, is spent once, in clause (1)'s forward
direction: the law is given self-decomposable, and the two theorems that read it are stated for
infinitely divisible laws. \ledger{A3}, \sscite{sato1999levy}~Thm.~8.1, pp.~37--38, gives that law a
triplet (i), makes its \Levy{} measure rotation invariant and its Gaussian form a multiple of
$|\xi|^2$ by uniqueness (ii), and in the converse direction supplies the law of a datum (iii);
uniqueness is spent again for the uniqueness of the datum and in clause (3). \ledger{A2},
\sscite{sato1999levy}~Thm.~15.10, p.~95, is the polar criterion in both directions, for a law that
may have a Gaussian part, which the source's ``no restriction on $A$'' allows. Clause (2) is
clause (1) at $d = 1$ applied to a marginal, and clause (3) is clause (2) combined with
\ref{thm:isotropic-classification}; neither spends anything further. \textbf{Who reads clause
(1).} \ref{thm:filtration-strict} in both directions and \ref{thm:no-cin-ray-in-the-plane} in the
forward one, each with one further use of the uniqueness of \ledger{A3}(ii) to identify a law's
\Levy{} measure with its datum's. Until Q-0286 they read it as the old clause (5) of
\ref{thm:origin-boundedness-d}, the case $a = 0$ of clause (1) stated a second time with its own proof; that clause was dropped as a
duplicate, after both citers were re-pointed here. \textbf{The normalisation}, fixed here once: this
module's $(a, k_d)$ is \ref{def:isotropic-origin-datum}'s, with Gaussian term
$\tfrac12 a|\xi|^2$ and $k_d$ the sphere-total of \sscite{sato1999levy}'s $k_\eta$. The draft's
(2.1) is clause (3) in the draft's normalisation,
$F(r) = a_{\mathrm{dr}}r^2 + \int_{\RR^d}(1 - \cos(r\langle e_1,x\rangle))\,k(|x|)\,|x|^{-d}dx$
with $a_{\mathrm{dr}} = \tfrac12 a$ and $k = k_d/|S^{d-1}|$; [V] takes the folded profile with
$a_{\mathrm V}\omega^2$, and at $d = 1$ clause (2) with $a' = a$ identifies $k_1$ with [V]'s folded
profile and $a_{\mathrm V}$ with $\tfrac12 a$. Until Q-0286 the draft's normalisation was a sentence
of clause (3), whose title was the draft's ``Theorem~2.2 with (2.1)''; both moved here. Clause (3)
is kept, not dropped as a duplicate of \ref{thm:isotropic-classification}: it adds the explicit
polar form and the equivalence of $F \equiv 0$ with a trivial datum, which the uniqueness of
clause (1) supplies. \textbf{Against the draft.} The draft reads ``$k_\xi = k$ by rotation
invariance'' in (15.8); that is not literal, (15.8) not determining $(\lambda,k_\xi)$, and the
proof replaces it by the radial image of the \Levy{} measure and the fact that a
rotation-invariant measure is determined by its radial image. \textbf{Not claimed}: anything about $F$ beyond
membership, in particular no explicit profile for a given $F$, which for a subordinated $F$ is
\ref{eq:polar-profile-d} and is the business of \ref{prop:isotropic-bridge}.)}
\end{proposition}

\begin{proof}
\leanok
\emph{(1), forward.} Let $\mu$ be self-decomposable on $\RR^d$ with
$\widehat\mu(\xi) = e^{-F(|\xi|)}$. By \ledger{A4} it is infinitely divisible, and by \ledger{A3}(i)
it has a generating triplet $(A,\nu,\gamma)$. For $Q \in \Orth(d)$,
$\widehat{Q_*\mu}(\xi) = \widehat\mu(Q^{\tp}\xi) = \widehat\mu(\xi)$, so $Q_*\mu = \mu$. The
unit ball is $Q$-invariant, so substituting $x = Q^{\tp}y$ in the representation shows that
$Q_*\mu$ has triplet $(QAQ^{\tp}, Q_*\nu, Q\gamma)$. By uniqueness, \ledger{A3}(ii),
$QAQ^{\tp} = A$ and $Q_*\nu = \nu$ for every $Q$. For unit vectors $e, e'$ choose $Q$ with
$Qe = e'$. Then $\langle e',Ae'\rangle = \langle e, Q^{\tp}AQ\,e\rangle = \langle e,Ae\rangle$, so
$\langle\xi,A\xi\rangle = a|\xi|^2$ with $a := \langle e,Ae\rangle \ge 0$. The real part of
$e^{i\theta} - 1 - i\theta\mathbf 1_D$ is $\cos\theta - 1$, so taking moduli in the representation,
\[
  F(|\xi|) = -\log\bigl|\widehat\mu(\xi)\bigr|
   = \tfrac12 a|\xi|^2 + \int_{\RR^d}\bigl(1 - \cos\langle\xi,x\rangle\bigr)\,\nu(dx) .
\]
Since $\mu$ is self-decomposable, \ledger{A2} gives $\nu(B) = \int_{S}\lambda(d\eta)\int_0^\infty
\mathbf 1_B(r\eta)\,k_\eta(r)\,dr/r$ with $\lambda$ a finite measure on $S = S^{d-1}$ and
$k_\eta(r)$ measurable in $\eta$ and nonincreasing in $r$.
% CHANGED Q-0281 wave 2: proof rewritten to the Lean route (the radial image of nu and the
% uniqueness of a rotation-invariant measure with a given radial image, in place of the Haar
% average over the orthogonal group). Blind review W2A-6: the set {|x| > eps} is U_eps, not S_eps,
% S being the sphere here.
A nonincreasing $k_\eta$ differs from its right limit
$k_\eta(r+) = \sup\{k_\eta(q) : q \in \mathbb Q,\ q > r\}$ at countably many $r$ only, a
Lebesgue-null set. Replacing $k_\eta(r)$ by $k_\eta(r+)$ therefore changes no integral in the
display, and makes the kernel jointly measurable in $(\eta,r)$, as a countable supremum of
measurable functions. Taking $B = \{x : |x| \in C\}$ for Borel $C \subset (0,\infty)$ and
exchanging the integrals, the image of $\nu$ under $x \mapsto |x|$ is
\[
  \nu\bigl(\{x : |x| \in C\}\bigr) = \int_C \bar k(r)\,\frac{dr}{r} ,
  \qquad \bar k(r) := \int_S k_\eta(r+)\,\lambda(d\eta) ,
\]
and $\bar k$ is nonincreasing with values in $[0,\infty]$. Integrating $|x|^2 \wedge 1$ against
$\nu$ through this image gives
$\int_0^\infty(r^2 \wedge 1)\bar k(r)\,dr/r = \int(|x|^2 \wedge 1)\,\nu(dx) < \infty$, which
contains the pair of conditions of \ref{def:isotropic-origin-datum}. They force $\bar k < \infty$
on $(0,\infty)$, because a nonincreasing function infinite at some $r_0$ is infinite on $(0,r_0]$.
So $(a, \bar k)$ is an isotropic origin datum. Let $\nu_d$ be its \Levy{} measure, the right side
of \ref{eq:polar-decomposition}. By \ref{lem:polar-decomposition} the image of $\nu_d$ under
$x \mapsto |x|$ is $\bar k(r)\,dr/r$ on $(0,\infty)$, the same as that of $\nu$.

It remains to see that $\nu_d = \nu$. Both are invariant under every $Q \in \Orth(d)$ ($\nu$ by the
first paragraph, $\nu_d$ by the invariance of $\sigma_{d-1}$), both vanish at the origin, and both
are finite on each set $U_\varepsilon = \{x : |x| > \varepsilon\}$, $\varepsilon > 0$. A finite
measure $m$ invariant under $\Orth(d)$ has a transform that depends only on its radial image:
$\widehat m(Q\xi) = \widehat m(\xi)$ for every $Q$, a reflection carries $\xi$ to $|\xi|\theta$ for
any unit $\theta$, and averaging over $\theta$ against $\sigma_{d-1}/|S^{d-1}|$ and exchanging the
integrals gives
\[
  \widehat m(\xi) = \frac{1}{|S^{d-1}|}\int_{[0,\infty)}\Phi(|\xi|,r)\,(|\cdot|_*m)(dr) ,
  \qquad
  \Phi(s,r) := \int_S e^{i s r\langle e_1,\theta\rangle}\,\sigma_{d-1}(d\theta) ,
\]
the inner sphere integral $\int_S e^{i s\langle x,\theta\rangle}\sigma_{d-1}(d\theta)$ depending on
$x$ only through $|x|$, again by the invariance of $\sigma_{d-1}$. Applied to the restrictions of
$\nu$ and $\nu_d$ to $U_\varepsilon$, which are invariant and have the same radial image, this gives
equal transforms, so the restrictions are equal. Letting $\varepsilon \downarrow 0$ along
$1/(n+1)$ gives $\nu = \nu_d$ off the origin, hence everywhere. The display for $F$ is therefore
$F(|\xi|) = \tfrac12 a|\xi|^2 + \int(1 - \cos\langle\xi,x\rangle)\,\nu_d(dx)$, the exponent of the
datum $(a, \bar k)$.

\emph{(1), converse.} Let $(a,k_d)$ be a datum with $F_d = F(|\cdot|)$. Its \Levy{} measure
$\nu_d$ satisfies $\int(|x|^2 \wedge 1)\,d\nu_d < \infty$ by the two conditions, so by
\ledger{A3}(iii) there is an infinitely divisible $\mu$ with triplet $(aI, \nu_d, 0)$. As $\nu_d$
is symmetric, the odd part of the integrand integrates to $0$ and
$\widehat\mu = e^{-F_d}$: $\mu$ is the datum's law. Its \Levy{} measure has the form of
\ledger{A2} with $\lambda = \sigma_{d-1}$ and $k_\eta = k_d/|S^{d-1}|$ for every $\eta$,
nonincreasing in $r$. So $\mu$ is self-decomposable by \ledger{A2}.

\emph{(1), uniqueness.} Two data with the same exponent have the same law, hence by \ledger{A3}(ii)
the same Gaussian matrix $aI$ and the same \Levy{} measure. The images of the two \Levy{}
measures under $x \mapsto |x|$ are then equal; by \ref{lem:polar-decomposition} they are
$k_d(r)\,dr/r$ and $k_d'(r)\,dr/r$ on $(0,\infty)$, so $k_d = k_d'$ almost everywhere.
% CHANGED Q-0281 wave 2: the last sentence rewritten to the Lean route (radial images).

\emph{(2).} Let $\mu$ be the law of clause (1), $e_1$ the first basis vector and $\pi(x) :=
\langle e_1,x\rangle$. Then $\pi_*\mu$ is self-decomposable on $\RR$. If $\rho_\varkappa$ is a
residual of $\mu$ for the ratio $\varkappa \in (0,1)$, then $\pi_*\rho_\varkappa$ is one of
$\pi_*\mu$:
\[
  \widehat{\pi_*\mu}(\omega) = \widehat\mu(\omega e_1)
   = \widehat\mu(\varkappa\omega e_1)\,\widehat\rho_\varkappa(\omega e_1)
   = \widehat{\pi_*\mu}(\varkappa\omega)\,\widehat{\pi_*\rho_\varkappa}(\omega) .
\]
Its exponent is $F(|\omega|)$.
Clause (1) at $d = 1$ gives a datum $(a',k_1)$ in dimension $1$ with, as $|S^0| = 2$ and the
integrand is even,
\[
  F(|\omega|) = \tfrac12 a'\omega^2 + \int_{\RR}\bigl(1 - \cos\omega x\bigr)\frac{k_1(|x|)}{2|x|}\,dx
   = \tfrac12 a'\omega^2 + \int_0^\infty\bigl(1 - \cos\omega x\bigr)\,k_1(x)\,\frac{dx}{x} ,
\]
which is \ref{eq:admissible-line} with $a_{\mathrm V} = \tfrac12 a'$ and $k = k_1$. The two
integrability conditions are the same pair. With the definition of $\mathcal{A}_d$ in
\ref{def:filtration-thresholds} this is the description of the class.

\emph{(3).} By \ref{thm:isotropic-classification} and clause (2), the family is isotropic exactly
when \ref{eq:isotropic-classification} holds for a gauge and an $F \not\equiv 0$ with $F(|\cdot|)$
the exponent $F_d$ of a datum. The displayed $F(r)$ is $F_d(re_1)$, and $F_d$ is radial. Finally
$F \equiv 0$ exactly when the datum is trivial. If $a = 0$ and $k_d = 0$ almost everywhere then
$F_d \equiv 0$. Conversely $F_d \equiv 0$ makes the law $\delta_0$, whose triplet is $(0,0,0)$,
so $a = 0$ and $k_d = 0$ almost everywhere by the uniqueness of clause (1).
\end{proof}

% CHANGED Q-0281: split from prop:isotropic-polar-reading (ADR-0025; records/formalization/PRICING.md):
% the parent's clauses (2) and (3), copied verbatim with the parent's clause numbers and its preamble.
% The parent's text and numbering are unchanged.
\begin{proposition}[the class $\mathcal{A}_d$ and the classification in polar form]
\label{prop:isotropic-polar-classification}
\uses{prop:isotropic-polar-reading,def:isotropic-origin-datum,def:filtration-thresholds,def:isotropic-kernel-family,thm:isotropic-classification,thm:cited-A2,thm:cited-A3,thm:cited-A4}
\lean{AffineHemigroup.polar_reading_admissible, AffineHemigroup.classA_iff_originDatum,
AffineHemigroup.polar_classification}\leanok
Let $d \ge 1$ and $F : [0,\infty) \to [0,\infty)$.
\begin{enumerate}
\item[(2)] \emph{(The class.)} If $\xi \mapsto F(|\xi|)$ is the exponent of a self-decomposable law
on $\RR^d$, then $F$ is admissible, of the form \ref{eq:admissible-line} with
$a_{\mathrm V} = \tfrac12 a'$ and $k = k_1$ for its datum $(a',k_1)$ in dimension $1$. So
$\mathcal{A}_d$ is the set of $F$ for which $F(|\cdot|)$ is the exponent of an isotropic origin
datum in dimension $d$.
\item[(3)] \emph{(The classification in polar form.)} A family $(\mu_{s,t})_{0 \le s \le t}$ of
probability measures on $\RR^d$ is an isotropic cascade family in dimension $d$ if and only if
there are a gauge $\chi$ and an isotropic origin datum $(a,k_d)$ in dimension $d$, not both trivial
($a > 0$, or $k_d$ not almost everywhere $0$), with \ref{eq:isotropic-classification} for
\[
  F(r) = \tfrac12 a\,r^2 + \int_{\RR^d}\bigl(1 - \cos(r\langle e_1,x\rangle)\bigr)\,
    \frac{k_d(|x|)}{|S^{d-1}|\,|x|^d}\,dx , \qquad r \ge 0 .
\]
\end{enumerate}
\statusA\quad\ledger{A2}\quad\ledger{A3}\quad\ledger{A4}\quad\emph{(\textbf{Assignment.} Nothing is
spent here that \ref{prop:isotropic-polar-reading}(1) does not spend: clause (2) is that clause at
$d = 1$ applied to a marginal, and clause (3) is clause (2) combined with
\ref{thm:isotropic-classification} and the uniqueness of the datum, \ledger{A3}(ii) through clause
(1). Clauses (2) and (3) of \ref{prop:isotropic-polar-reading}, split off on 2026-10-04 (Q-0281,
ADR-0025), with the parent's clause numbers; the hypothesis of (2), ``if this holds'' in the parent,
is written out. The parent's text is unchanged.)}
\end{proposition}

\begin{proof}
\leanok
% polar_classification (clause (3)) prints Paper V's interface through isotropic_classification,
% admitted through the include line of blueprint/trust-boundary.txt; all of it is guarded.
This is the part of the proof of \ref{prop:isotropic-polar-reading} headed (2) and (3).
\end{proof}

% not in the draft: queue item Q-0434 (hub), from the second external round of the isotropic paper
% (records/affine/PLAN-review-response-paper-iso-round2.md, "The stages"; the verification record
% records/affine/reviews/VERIFY-round2-mathematics.md, item 12). The paper's section 4.1 listed
% the stages mu_{s,t}, s > 0, as an open problem: which have a density, which an atom at the
% centre. A reviewer answered it from the paper's own statements. The answer is three nodes:
%   lem:isotropic-stage-levy-measure   the Levy measure of a stage, its mass, the transform of
%                                      the stage through it; cites nothing;
%   prop:isotropic-stage-atom          no Gaussian part and a finite c: a point mass at the
%                                      centre plus an absolutely continuous part; cites nothing;
%   prop:isotropic-stage-density       the stage as an infinitely divisible law with its triplet
%                                      (A3(iii)); a density when c is infinite (A3(iii)) or
%                                      when there is a Gaussian part (A1).
% Why three (ADR-0025: what is tagged is exactly what is formalized). The author has not decided
% whether these statements are formalized (the plan's decision 2). The cut is by what the proofs
% read, so that each of the three outcomes is a matter of tags: all in print (no tag, as now);
% the atom in Lean (the first two nodes, on Lean's logical axioms alone); all in Lean (the third
% as well, on A3(iii) and A1). The first node is separate from the second because the third
% reads it too. The record's draft has one proposition with four clauses; its clause (1) is the
% first node together with clause (1) of the third, its clause (2) the second node, its clauses
% (3) and (4) the clauses (2) and (3) of the third. A paper may print the three as one
% proposition.
% None of the three has a \lean tag. Their statements are typed, with `sorry`, in a scratch file
% that nothing imports (Formalization/scratch/StageSkeleton.lean, kept as
% records/formalization/skeleton-q0434-stages.lean), compiled on 2026-10-13, and priced beside
% the types in records/affine/NODES-q0434.md.
% What writing the types down found.
% (i) The first node is a statement about a datum and two numbers 0 < l < L, with no law in it;
% the family enters in its last sentence only. So the mass and the transform are stated once and
% both of the other nodes read them.
% (ii) The library has the mass for a finite c (OriginDatum.lintegral_hb, the identity behind
% lem:isotropic-exponent-log(2)) and not for c infinite, which the third node needs.
% (iii) The smoothing part is named as a measure, lambda^c sum_{n >= 1} nu^{*n}/n!, not by its
% density: the formula for the density through the convolution powers of a function would be a
% second object to formalize and adds nothing to the measure.
% (iv) "Not self-decomposable" has a proof that reads nothing of the classification: a law with
% an atom of weight in (0,1) at the origin and an absolutely continuous remainder is not
% self-decomposable. The record's proof goes through the polar reading and the uniqueness of the
% triplet (A2, A3, A4); this one cites nothing, and it is the one below.
% (v) The smooth density under a Gaussian part is A1 read at every order, from the bound
% e^{-v|xi|^2/2} on the transform; the record's proof by a Gaussian convolution needs A3(iii)
% and a differentiation under the integral, and its Lean would be the longer one.
% CHANGED after the blind check of Q-0434 (records/formalization/blind-check-q0434-report-A.md:
% the four nodes in print correct, six notes on proofs, one on the skeleton). This node, proof
% only: F-A1 (the passage to the limit holds for n >= (L/l)^{1/2}, where the condition under which
% its formula was derived holds) and F-A2 (why the integral of the dilated exponent is finite).
% prop:isotropic-stage-atom: annotation only (F-A7; the skeleton now types "c > 0" and the upper
% bound of the weight). prop:isotropic-stage-density: F-A3 in the proof (the bound on the
% imaginary part) and F-A4 in clause (1) of the STATEMENT, "is infinitely divisible and has the
% generating triplet" for "is infinitely divisible, with the generating triplet": what A3(iii)
% and the transform give is that the stage has the triplet; that a law has only one is A3(ii),
% which the clause does not read. A narrowing of the wording, no change of content.
% CHANGED Q-0447 (the author's decision of 2026-10-13: the stages are machine-checked): the three
% nodes are proved in Lean and tagged. No statement of the three changes. The proofs keep their
% text; the closing paragraph of each says what is machine-checked and where the Lean takes
% another step. Modules: AffineHemigroup/StageLevyMeasure.lean (456 lines against the estimate of
% 260 to 340), StageAtom.lean (392 against 300 to 380), StageDensity.lean (358 against 380 to
% 590). Axioms, printed before these annotations were written: every declaration of the first
% two modules prints propext, Classical.choice, Quot.sound and nothing else;
% isotropic_stage_triplet and isotropic_stage_absolutelyContinuous add levy_khintchine_converse
% (A3(iii)); isotropic_stage_smooth_density adds density_of_integrable_charFun (A1);
% isotropic_stage_dichotomy adds both.
% What proving them found.
% (i) ONE TYPE OF THE SKELETON WAS FALSE. not_selfDecomposable_of_atom, "the fact behind clause
% (4)" of the second node, was typed for every d, and at d = 0 it is false: R^0 is a point,
% Lebesgue measure on it is delta_0, so every measure is absolutely continuous, and
% delta_0 = p delta_0 + (1 - p) delta_0 is self-decomposable. The declaration asks d >= 1. The
% node's statement is unaffected (it has d >= 1); the sentence of its proof that states the
% general fact now says "on R^d, d >= 1".
% (ii) The transform of a symmetric triplet in real form was in the library
% (LevyTriplet.lkExponent_of_symm, written for the line), as were the compound Poisson law with
% its transform, the continuous densities of one measure being one function, and "a continuous
% function that tends to 0 at infinity is bounded". The estimate of the third node counted 80 to
% 150 lines for the first of these; it is why that node came in under its estimate.
% (iii) The dilated datum (a beta^2, k_d(./beta)) is not formed. Its Levy measure is taken as the
% image of nu_d under x -> beta x, whose density is computed (an identity of measures; the
% library had the inequality), and its exponent is read from the library's
% exponent_smul_eq_lkForm. So the two integrability conditions of the dilated profile are never
% needed.
% (iv) The mass for c infinite is obtained from the library's identity for a finite c
% (OriginDatum.lintegral_hb), applied to the capped profiles min(k_d, M): the integrand of the
% capped profile is below that of k_d, its integral is M log(L/l), and M is arbitrary. The
% printed proof treats both cases by one truncation; it stays, and the closing paragraph says
% what the Lean does.
% (v) "e^{-m} = lambda^c" in clause (1) of the second node had no conjunct in the skeleton (the
% weight was written as lambda^c throughout); isotropic_stage_exp_neg_mass states it. The last
% sentence of the third node ("a point mass exactly when ...") is isotropic_stage_dichotomy. The
% last sentence of this node's clause (3) is stageMeasure_univ_ne_top_iff, and the mass of the
% stage is isotropic_stage_mass.
% CHANGED after the blind check of Q-0447 (records/formalization/blind-check-q0447-report.md: no
% error, two gaps, six lesser findings; one line each in records/affine/NODES-q0447.md). No
% statement and no proof text changes. F-1, a gap: "in particular it has no point mass" of
% clause (2) of the third node had no conjunct; isotropic_stage_no_point_mass states it for every
% point, and is in the tag (A3(iii), as the absolute continuity it follows from). F-2, a gap:
% clause (3) of the second node names a density p and the Lean had the operator of convolution
% by the measure eta only; isotropic_stage_operator gives p >= 0, integrable, of integral
% 1 - lambda^c, a density of the series lambda^c sum_{n >= 1} nu^{*n}/n!, with
% mu * f = lambda^c f + p * f almost everywhere for f in L^1, and is in the tag (Lean core; p is
% the Radon-Nikodym derivative of the series, read in R). It also names eta by its series
% (F-3). F-4: the annotation of this node said the wrong surplus (what the last sentence does
% not ask is that the measures are probability measures). F-6: isotropic_stage_dichotomy does
% not ask F not identically 0. F-7: isotropic_stage_atom has the conjunct 0 < lambda^c. F-5 was
% in the docstring already. F-8: the finiteness of the integral behind the exponent is a lemma of
% the library (OriginDatum.exponentL_ne_top), in no tag; it is not added to the tag of
% def:isotropic-origin-datum, which the isotropic paper prints with its list of declarations,
% and the annotation of this node says where it is. The \uses of this node gains
% lem:isotropic-exponent-log, whose Frullani identity the Lean of clause (3) reads. After these
% repairs StageAtom.lean has 470 lines and StageDensity.lean 379, 1,305 in all against the
% estimate of 940 to 1,310.
\begin{lemma}[the \Levy{} measure of a stage]
\label{lem:isotropic-stage-levy-measure}
\uses{def:isotropic-origin-datum,lem:polar-decomposition,lem:isotropic-exponent-log,thm:isotropic-classification,def:isotropic-kernel-family}
\lean{AffineHemigroup.OriginDatum.stageProfile, AffineHemigroup.OriginDatum.stageMeasure,
AffineHemigroup.OriginDatum.stageProfile_nonneg, AffineHemigroup.OriginDatum.stageMeasure_zero,
AffineHemigroup.OriginDatum.stageMeasure_absolutelyContinuous,
AffineHemigroup.OriginDatum.stageMeasure_rotation,
AffineHemigroup.OriginDatum.stageMeasure_integrable,
AffineHemigroup.OriginDatum.exponent_smul_sub_smul, AffineHemigroup.OriginDatum.stageMeasure_univ,
AffineHemigroup.OriginDatum.stageMeasure_univ_ne_top_iff, AffineHemigroup.isotropic_stage_transform,
AffineHemigroup.isotropic_stage_mass}\leanok
Let $d \ge 1$, let $(a, k_d)$ be an isotropic origin datum in dimension $d$ with exponent $F_d$
and profile at the origin $c = k_d(0{+}) \in [0,\infty]$, and let $0 < \ell < L$. Put
\[
  \kappa(r) := k_d\Bigl(\frac{r}{L}\Bigr) - k_d\Bigl(\frac{r}{\ell}\Bigr) \qquad (r > 0) ,
\]
and let $\nu$ be the measure on $\RR^d \setminus \{0\}$ with the density
$\kappa(|x|)\,\bigl(|S^{d-1}|\,|x|^d\bigr)^{-1}$.
\begin{enumerate}
\item $\kappa \ge 0$. The measure $\nu$ is rotation invariant, and
$\int_{\RR^d}(|x|^2 \wedge 1)\,\nu(dx) < \infty$.
\item For every $\xi \in \RR^d$,
\[
  F_d(L\xi) - F_d(\ell\xi) = \tfrac12 a\,(L^2 - \ell^2)\,|\xi|^2
     + \int_{\RR^d}\bigl(1 - \cos\langle\xi,x\rangle\bigr)\,\nu(dx) .
\]
\item $\nu(\RR^d \setminus \{0\}) = c\,\log(L/\ell)$, in $[0,\infty]$. So $\nu$ is a finite
measure exactly when $c < \infty$.
\end{enumerate}
In particular, let $\chi$ be a gauge and $F$ a function on $[0,\infty)$ with
$F(|\xi|) = F_d(\xi)$ for every $\xi \in \RR^d$, let $(\mu_{s,t})_{0 \le s \le t}$ be probability
measures on $\RR^d$ that satisfy \ref{eq:isotropic-classification} with $(\chi, F)$, and let
$0 < s < t$. Write $\nu_{s,t}$ for the measure $\nu$ of $\ell = \chi(s)$ and $L = \chi(t)$, the
\emph{\Levy{} measure of the stage}, with polar profile
$\kappa_{s,t}(r) = k_d(r/\chi(t)) - k_d(r/\chi(s))$. Then
\begin{equation}\label{eq:stage-transform}
  \widehat\mu_{s,t}(\xi) = \exp\Bigl[-\tfrac12 a\bigl(\chi(t)^2 - \chi(s)^2\bigr)|\xi|^2
     - \int_{\RR^d}\bigl(1 - \cos\langle\xi,x\rangle\bigr)\,\nu_{s,t}(dx)\Bigr],
  \qquad \xi \in \RR^d ,
\end{equation}
and $\nu_{s,t}(\RR^d \setminus \{0\}) = c\,\log\bigl(\chi(t)/\chi(s)\bigr)$.
\statusT\quad\emph{(Cites nothing: a change of variable, the subtraction of two finite
integrals, and \ref{eq:polar-decomposition} on a shell. The statement is about a datum, a measure
and a function; no law is named before the last sentence, where the measures $\mu_{s,t}$ are any
probability measures with the transforms \ref{eq:isotropic-classification}. \textbf{What it is.}
The exponent of a datum read at two dilations, and the difference written as the exponent of one
measure. The dilation of a datum by $\beta$ has the profile $k_d(\cdot/\beta)$, and since the
profile is nonincreasing the larger dilation has the larger profile at every radius: the
difference $\kappa$ is the profile of a measure. \textbf{The direction.} In
\ref{eq:isotropic-classification} the kernel $\mu_{s,t}$ has the exponent
$F(\chi(t)|\xi|) - F(\chi(s)|\xi|)$, with $\chi(s) < \chi(t)$, so the profile of the stage is
$k_d(r/\chi(t)) - k_d(r/\chi(s))$, the profile at the coarser scale minus the profile at the
finer one. \textbf{The mass} is the identity behind \ref{lem:isotropic-exponent-log}(2), which
has it for a finite $c$ in the form $\int_0^\infty(k_d(u/b) - k_d(u))\,du/u = c\log b$; here it
is stated for $c = \infty$ as well, where it says that the stage has infinitely many small
jumps. \textbf{Independence of the pair.} A family has many pairs $(\chi, F)$, related by
$(\chi, F) \mapsto (\beta\chi, F(\beta^{-1}\,\cdot))$. The datum of $F(\beta^{-1}\,\cdot)$ is
$(a\beta^{-2}, k_d(\beta\,\cdot))$, so $a(\chi(t)^2 - \chi(s)^2)$, the profile $\kappa_{s,t}$,
the number $c$ and the ratio $\chi(t)/\chi(s)$ do not depend on the pair. \textbf{Not claimed}:
that $\mu_{s,t}$ is infinitely divisible with the triplet that \ref{eq:stage-transform}
displays, which is \ref{prop:isotropic-stage-density}(1) and reads \ledger{A3}(iii); nor that
$\kappa$ is monotone, which it is not in general, so that $(0,\kappa)$ is not an isotropic origin
datum and a stage is not, in general, self-decomposable. \textbf{Formalization} (Q-0447).
Machine-checked, on Lean's logical axioms alone. $\kappa$ and $\nu$ are definitions on a datum and
two reals, total functions of $\ell$ and $L$ that are read for $0 < \ell < L$; $\nu$ is a measure
on $\RR^d$ with the density written as the \Levy{} measure of a datum is, which gives the origin
no mass. Clause (1) is five declarations: $\kappa(r) \ge 0$ for $r > 0$ and
$0 < \ell \le L$; $\nu(\{0\}) = 0$; $\nu$ is absolutely continuous; $Q_*\nu = \nu$ for every
orthogonal matrix $Q$; and the integral of $|x|^2 \wedge 1$ is finite. The second, third and
fourth hold for every pair of reals $\ell, L$, the density being a function of $|x|$ whatever
they are. Clause (2) is an equality of real numbers together with the integrability of
$x \mapsto 1 - \cos\langle\xi,x\rangle$ against $\nu$. Clause (3) is an equality in $[0,\infty]$,
with its ``so'' as an equivalence. The last sentence is two declarations, for any family of
measures with the transforms \ref{eq:isotropic-classification} at the pairs $0 \le s \le t$, a
gauge $\chi$, and a function $F$ whose radial form is the exponent of the datum; that the
measures are probability measures is not asked, which is more than the print states. The
exponent of a datum is a real number read off an integral with values in $[0,\infty]$; that
integral is finite for every datum in a dimension $d \ge 1$, by a lemma of the library that is
in no tag and that this node's Lean uses.)}
\end{lemma}

\begin{proof}
\leanok
\emph{The dilation of a datum.} For $\beta > 0$ put $k^{(\beta)}(r) := k_d(r/\beta)$. It is
nonnegative and nonincreasing, and
\[
  \int_0^1 k^{(\beta)}(r)\,r\,dr = \beta^2\int_0^{1/\beta}k_d(u)\,u\,du < \infty, \qquad
  \int_1^\infty k^{(\beta)}(r)\,\frac{dr}{r} = \int_{1/\beta}^\infty k_d(u)\,\frac{du}{u} < \infty ,
\]
by the two conditions on $k_d$ and because $k_d$ is bounded on the interval between $1$ and
$1/\beta$. So $(a\beta^2, k^{(\beta)})$ is an isotropic origin datum in dimension $d$. Let
$\nu^{(\beta)}$ be its \Levy{} measure, of density $k_d(|x|/\beta)\,(|S^{d-1}||x|^d)^{-1}$, and
$\nu_d$ that of $(a,k_d)$. The substitution $y = \beta x$, $dy = \beta^d dx$, gives
$\int g(\beta x)\,\nu_d(dx) = \int g\,d\nu^{(\beta)}$ for every nonnegative Borel $g$. With
$g(y) = 1 - \cos\langle\xi,y\rangle$,
\begin{equation}\label{eq:dilated-exponent}
  F_d(\beta\xi) = \tfrac12 a\beta^2|\xi|^2
     + \int_{\RR^d}\bigl(1 - \cos\langle\xi,y\rangle\bigr)\,\nu^{(\beta)}(dy) ,
\end{equation}
so the exponent of $(a\beta^2, k^{(\beta)})$ is $\xi \mapsto F_d(\beta\xi)$, and the integral is
finite, as $\int(1 - \cos\langle\xi,y\rangle)\,d\nu^{(\beta)}
\le \int(|\xi|^2|y|^2/2 \wedge 2)\,d\nu^{(\beta)} < \infty$ by \ref{eq:polar-decomposition} and
the two conditions on $k^{(\beta)}$.

\emph{Clause (1).} For $r > 0$, $r/L < r/\ell$, and $k_d$ is nonincreasing, so
$\kappa(r) = k^{(L)}(r) - k^{(\ell)}(r) \ge 0$. The density of $\nu$ is the difference of the
densities of $\nu^{(L)}$ and $\nu^{(\ell)}$, so
\begin{equation}\label{eq:stage-measure-sum}
  \nu + \nu^{(\ell)} = \nu^{(L)}
\end{equation}
as measures on $\RR^d \setminus \{0\}$. In particular $\nu \le \nu^{(L)}$, and
$\int(|x|^2 \wedge 1)\,d\nu \le \int(|x|^2 \wedge 1)\,d\nu^{(L)}$, which is finite: by
\ref{eq:polar-decomposition} it is $\int_0^1 k^{(L)}(r)\,r\,dr + \int_1^\infty k^{(L)}(r)\,dr/r$.
The density of $\nu$ is a function of $|x|$ and Lebesgue measure is invariant under
$\Orth(d)$, so $\nu$ is.

\emph{Clause (2).} Let $g(x) := 1 - \cos\langle\xi,x\rangle \ge 0$. By
\ref{eq:dilated-exponent} the integrals of $g$ against $\nu^{(L)}$ and $\nu^{(\ell)}$ are finite,
so \ref{eq:stage-measure-sum} gives
$\int g\,d\nu = \int g\,d\nu^{(L)} - \int g\,d\nu^{(\ell)}$. By \ref{eq:dilated-exponent} at
$\beta = L$ and at $\beta = \ell$ this is
$F_d(L\xi) - F_d(\ell\xi) - \tfrac12 a(L^2 - \ell^2)|\xi|^2$.

\emph{Clause (3).} Let $0 < \varepsilon < A$ with $\varepsilon/\ell \le A/L$, and
$B := \{x : \varepsilon < |x| \le A\}$. By \ref{eq:polar-decomposition} for the datum
$(a\beta^2, k^{(\beta)})$ and $g = \mathbf 1_B$, and the substitution $u = r/\beta$,
$\nu^{(\beta)}(B) = \int_\varepsilon^A k_d(r/\beta)\,dr/r
= \int_{\varepsilon/\beta}^{A/\beta}k_d(u)\,du/u$, which is finite. By
\ref{eq:stage-measure-sum}, and since the intervals $[\varepsilon/L, A/L]$ and
$[\varepsilon/\ell, A/\ell]$ overlap in $[\varepsilon/\ell, A/L]$,
\[
  \nu(B) = \int_{\varepsilon/L}^{A/L}k_d(u)\,\frac{du}{u} - \int_{\varepsilon/\ell}^{A/\ell}k_d(u)\,\frac{du}{u}
         = \int_{\varepsilon/L}^{\varepsilon/\ell}k_d(u)\,\frac{du}{u}
           - \int_{A/L}^{A/\ell}k_d(u)\,\frac{du}{u} .
\]
Each of the last two integrals is over an interval whose endpoints have the ratio $L/\ell$, so
$\int du/u = \log(L/\ell)$ over it. As $k_d$ is nonincreasing, the first lies between
$k_d(\varepsilon/\ell)\log(L/\ell)$ and $k_d(\varepsilon/L)\log(L/\ell)$, and both bounds tend
to $c\log(L/\ell)$ as $\varepsilon \downarrow 0$, also when $c = \infty$. The second is at most
$k_d(A/L)\log(L/\ell)$, which tends to $0$ as $A \to \infty$: a nonincreasing nonnegative
function with $\int_1^\infty k_d(r)\,dr/r < \infty$ tends to $0$ at infinity. With
$\varepsilon = 1/n$ and $A = n$ for $n \ge (L/\ell)^{1/2}$, so that
$\varepsilon/\ell \le A/L$, the sets $B$ increase to $\RR^d \setminus \{0\}$, so
$\nu(\RR^d \setminus \{0\}) = \lim_n\nu(B) = c\log(L/\ell)$. Since $\log(L/\ell) > 0$, this is
finite exactly when $c$ is.

\emph{The last sentence.} A gauge is increasing with $\chi(0) = 0$, so
$0 < \chi(s) < \chi(t)$. By \ref{eq:isotropic-classification} and $F(|\zeta|) = F_d(\zeta)$,
$-\log\widehat\mu_{s,t}(\xi) = F_d(\chi(t)\xi) - F_d(\chi(s)\xi)$, and clauses (2) and (3) with
$\ell = \chi(s)$, $L = \chi(t)$ are the two assertions.

Every clause is machine-checked, and no cited fact is used. The Lean takes two steps another
way. It does not form the datum $(a\beta^2, k^{(\beta)})$: it takes $\nu^{(\beta)}$ as the image
of $\nu_d$ under $x \mapsto \beta x$, proves that this image has the density
$k_d(|x|/\beta)\,(|S^{d-1}||x|^d)^{-1}$, and reads \ref{eq:dilated-exponent} from the definition
of the exponent, so that the finiteness of the integral is that of $F_d(\beta\xi)$ and the bound
of clause (1) is
$\int(|x|^2 \wedge 1)\,d\nu^{(L)} \le \max(L^2,1)\int(|x|^2 \wedge 1)\,d\nu_d$. And for clause (3)
it computes $\nu(\RR^d \setminus \{0\}) = \int_0^\infty\kappa(r)\,dr/r$ by polar coordinates,
substitutes $r = \ell u$, and takes $\int_0^\infty(k_d(u/b) - k_d(u))\,du/u = c\log b$, with
$b = L/\ell$, from the proof of \ref{lem:isotropic-exponent-log}(2) when $c < \infty$, where it is
obtained by the truncation above; when $c = \infty$ it applies that identity to the profiles
$\min(k_d, M)$, $M > 0$, whose integrand is at most that of $k_d$ and whose integral is
$M\log b$, and lets $M \to \infty$.
\end{proof}

\begin{proposition}[a stage of a cascade with no Gaussian part and a finite profile at the origin: a point mass at the centre]
\label{prop:isotropic-stage-atom}
\uses{lem:isotropic-stage-levy-measure,thm:isotropic-classification,def:isotropic-kernel-family,def:isotropic-origin-datum,def:filtration-thresholds}
\lean{AffineHemigroup.isotropic_stage_compoundPoisson, AffineHemigroup.isotropic_stage_exp_neg_mass,
AffineHemigroup.profileAtOrigin_pos, AffineHemigroup.isotropic_stage_atom,
AffineHemigroup.isotropic_stage_operator, AffineHemigroup.isotropic_stage_not_selfDecomposable}\leanok
Let $d \ge 1$, let $(\mu_{s,t})_{0 \le s \le t}$ be an isotropic cascade family in dimension $d$,
and let $\chi$ be a gauge and $F \in \mathcal{A}_d$ with $F \not\equiv 0$ that satisfy
\ref{eq:isotropic-classification} (\ref{thm:isotropic-classification}). Let $(a, k_d)$ be an
isotropic origin datum in dimension $d$ with exponent $\xi \mapsto F(|\xi|)$, and suppose that
$a = 0$ and that $c := k_d(0{+})$ is finite. Let $0 < s < t$, put
$\lambda := \chi(s)/\chi(t) \in (0,1)$, and let $\nu_{s,t}$ be the \Levy{} measure of the stage
(\ref{lem:isotropic-stage-levy-measure}), a finite measure of mass
$m := c\,\log(\chi(t)/\chi(s))$ with a density. Then $c > 0$, and:
\begin{enumerate}
\item \emph{(The stage is compound Poisson.)}
$\mu_{s,t} = e^{-m}\sum_{n \ge 0}\nu_{s,t}^{*n}/n!$, with $\nu_{s,t}^{*0} = \delta_0$, and
$e^{-m} = \lambda^c$.
\item \emph{(A point mass at the centre and a smoothing part.)} The measure
$\eta_{s,t} := \lambda^c\sum_{n \ge 1}\nu_{s,t}^{*n}/n!$ is absolutely continuous with respect
to Lebesgue measure, of total mass $1 - \lambda^c$, and
\[
  \mu_{s,t} = \lambda^c\,\delta_0 + \eta_{s,t} .
\]
So $\mu_{s,t}(\{0\}) = (\chi(s)/\chi(t))^c \in (0,1)$, and $\mu_{s,t}$ has no density.
\item \emph{(The operator.)} With $p_{s,t} \ge 0$ a density of $\eta_{s,t}$, an integrable
function of integral $1 - \lambda^c$, the operator of the stage is
\[
  \mu_{s,t} * f = \Bigl(\frac{\chi(s)}{\chi(t)}\Bigr)^{c} f + p_{s,t} * f ,
  \qquad f \in L^1(\RR^d) .
\]
\item \emph{(Not self-decomposable.)} $\mu_{s,t}$ is not self-decomposable.
\end{enumerate}
\statusT\quad\emph{(Cites nothing: the exponential series of a finite measure, the uniqueness of
the transform of a probability measure, and the fact that the convolution of an absolutely
continuous finite measure with a finite measure is absolutely continuous. The datum is given in
the statement; that an $F \in \mathcal{A}_d$ has one is
\ref{prop:isotropic-polar-classification}(2), which the node does not read. \textbf{In plain
words, for a filter.} A stage of such a cascade passes the fraction $(\chi(s)/\chi(t))^c$ of the
record at scale $s$ through unchanged and adds to it a smoothed copy that carries the rest.
\textbf{The \Matern{} kernels.} For the \Matern{} kernels of index $\gamma$,
\ref{prop:isotropic-bridge}(2), the gauge is $\chi(t) = t$ and $F(r) = \gamma\log(1 + r^2/2)$, so
$a = 0$ and $c = 2\gamma$ by \ref{cor:origin-profile-exponent}(1) and~(2). The \Levy{} measure
of the stage from $s$ to $t$ has the mass $2\gamma\log(t/s)$, and the point mass has the weight
$(s/t)^{2\gamma}$: a quarter for $\gamma = 1$ and one octave. At $\gamma = 1$ the stage is a
mixture of the point mass and the kernel from scale zero,
$\mu_{s,t} = (s/t)^2\,\delta_0 + (1 - (s/t)^2)\,\mu_{0,t}$, in every dimension, because
$(1 + s^2x)/(1 + t^2x) = (s/t)^2 + (1 - (s/t)^2)/(1 + t^2x)$ at $x = |\xi|^2/2$.
\textbf{What the earlier papers have.} The cone paper has the case of the line and the
\Matern{} kernels of index $1$, as a realization: the mixture \prog{B}{(7.2)} of its
Proposition~7.2, in which the point mass has the weight $r^2$ with $r$ the ratio of the two
scales (read at the page image for the verification record of the second review round, item
12). The general statement, for every profile with a finite $c$ and in every dimension, is in
none of the three earlier papers. \textbf{Why the kernels from scale zero have no point mass.}
As $s \downarrow 0$ the weight $(\chi(s)/\chi(t))^c$ tends to $0$. \textbf{The smoothing part}
is rotation invariant, being $\mu_{s,t} - \lambda^c\delta_0$, and its density is
$p_{s,t} = \lambda^c\sum_{n \ge 1}f^{*n}/n!$ almost everywhere, with $f$ the density of
$\nu_{s,t}$; the statement names the measure. \textbf{Not claimed}: whether the smoothing part
is bounded near the centre, that is whether $p_{s,t}$ has a representative bounded on a
neighbourhood of the origin, and what number decides it; nor anything about the case $c = 0$,
which the hypotheses exclude ($a = 0$ and $c = 0$ make $F \equiv 0$). \textbf{Formalization}
(Q-0447). Machine-checked, on Lean's logical axioms alone. The family is any
family of probability measures with the transforms \ref{eq:isotropic-classification} for a
gauge $\chi$ and a function $F$ whose radial form is the exponent of the datum; no cascade
property is asked and $F \in \mathcal{A}_d$ is not asked, which is more than the print states.
Clause (1) is two declarations: $\nu_{s,t}$ is a finite measure and $\mu_{s,t}$ equals the
compound Poisson law of the library, which is $e^{-\nu(\RR^d)}\sum_{n \ge 0}\nu^{*n}/n!$ by
definition; and $\nu_{s,t}(\RR^d) = c\log(\chi(t)/\chi(s))$ as a real number with
$e^{-m} = (\chi(s)/\chi(t))^c$. That $c > 0$ when $F \not\equiv 0$ is a declaration of its own.
Clauses (2) and (3) are one
declaration, with $\eta_{s,t}$ under an existential quantifier, named by its properties: a
finite measure, absolutely continuous, with $\mu_{s,t} = \lambda^c\delta_0 + \eta_{s,t}$, of mass
$1 - \lambda^c$; then $\mu_{s,t}(\{0\}) = \lambda^c$, $0 < \lambda^c < 1$, $\mu_{s,t}$ is not
absolutely continuous, and the operator of convolution by $\mu_{s,t}$ on $L^1(\RR^d)$ is
$\lambda^c$ times the identity plus the operator of convolution by $\eta_{s,t}$. A second
declaration names $\eta_{s,t}$ by its series and gives the density of clause (3): the measure
$\lambda^c\sum_{n \ge 1}\nu^{*n}/n!$ equals the measure with a density $p \ge 0$ that is
integrable and has the integral $1 - \lambda^c$; $\mu_{s,t}$ is $\lambda^c\delta_0$ plus that
measure; and for every $f \in L^1(\RR^d)$ the function $\mu_{s,t} * f$ equals
$x \mapsto \lambda^c f(x) + \int p(y)\,f(x - y)\,dy$ almost everywhere. The function $p$ is the
Radon--Nikodym density of the series, taken with real values.
The first declaration of clauses (2) and (3) asks $F \not\equiv 0$ for the inequality
$\lambda^c < 1$, and that of clause (4) asks it too; clause (1) and the second declaration do
not. The general fact behind clause (4) is a
declaration on $\RR^d$ with $d \ge 1$; at $d = 0$ it is false, the space being a point.)}
\end{proposition}

\begin{proof}
\leanok
Write $\nu := \nu_{s,t}$, $\ell := \chi(s)$ and $L := \chi(t)$. By
\ref{lem:isotropic-stage-levy-measure}, $\nu$ is a finite measure of mass
$m = c\log(L/\ell)$, rotation invariant, hence symmetric, with a density, and by
\ref{eq:stage-transform} with $a = 0$,
$\widehat\mu_{s,t}(\xi) = \exp[-\int(1 - \cos\langle\xi,x\rangle)\,\nu(dx)]$. If $c$ were $0$,
the nonincreasing nonnegative function $k_d$ would vanish on $(0,\infty)$, the exponent of the
datum $(0,k_d)$ would be $0$, and $F$ would vanish on $[0,\infty)$. So $c > 0$, $m > 0$ and
$0 < e^{-m} < 1$, and $e^{-m} = (\ell/L)^c = \lambda^c$.

\emph{Clause (1).} Put $\pi := e^{-m}\sum_{n \ge 0}\nu^{*n}/n!$. The measure $\nu^{*n}$ has the
mass $m^n$, so $\pi$ is a probability measure. The transform of $\nu^{*n}$ is $\widehat\nu^{\,n}$
and $|\widehat\nu| \le m$, so $\widehat\pi = e^{-m}\sum_n\widehat\nu^{\,n}/n!
= \exp(\widehat\nu - m)$. As $\nu$ is symmetric, $\widehat\nu(\xi) = \int\cos\langle\xi,x\rangle
\,\nu(dx)$, and $\widehat\nu(\xi) - m = -\int(1 - \cos\langle\xi,x\rangle)\,\nu(dx)$. So
$\widehat\pi = \widehat\mu_{s,t}$, and $\mu_{s,t} = \pi$, a probability measure being determined
by its transform.

\emph{Clause (2).} Let $M$ be a finite measure and $N$ a Lebesgue-null Borel set. Then
$(\nu * M)(N) = \int\nu(N - y)\,M(dy) = 0$, because $N - y$ is Lebesgue-null and $\nu$ has a
density. So $\nu^{*n} = \nu * \nu^{*(n-1)}$ is absolutely continuous for every $n \ge 1$, and so
is $\eta_{s,t}$, a countable sum of such measures. Its mass is
$e^{-m}\sum_{n \ge 1}m^n/n! = 1 - e^{-m}$. The term $n = 0$ of clause (1) is $e^{-m}\delta_0$.
So $\mu_{s,t} = \lambda^c\delta_0 + \eta_{s,t}$, and
$\mu_{s,t}(\{0\}) = \lambda^c + \eta_{s,t}(\{0\}) = \lambda^c$, the origin being a Lebesgue-null
set. A measure with a density gives the origin no mass, and $\lambda^c > 0$, so $\mu_{s,t}$ has
no density.

\emph{Clause (3).} The convolution $\mu * f$ is additive in the measure $\mu$, and
$\delta_0 * f = f$. With clause (2),
$\mu_{s,t} * f = \lambda^c f + \eta_{s,t} * f$, and $\eta_{s,t} * f = p_{s,t} * f$.

\emph{Clause (4).} Suppose that $\mu_{s,t}$ is self-decomposable, fix $q \in (0,1)$, and let
$\rho$ be a residual: a probability measure with
$\widehat\mu_{s,t}(\xi) = \widehat\mu_{s,t}(q\xi)\,\widehat\rho(\xi)$ for every $\xi$, that is
$\mu_{s,t} = (qI)_*\mu_{s,t} * \rho$. By clause (2),
$(qI)_*\mu_{s,t} = \lambda^c\delta_0 + (qI)_*\eta_{s,t}$, and $(qI)_*\eta_{s,t}$ is absolutely
continuous. As in clause (2), $((qI)_*\eta_{s,t} * \rho)(\{0\}) = 0$. So the mass of the origin
under the right side is $\lambda^c\rho(\{0\})$, and under the left side it is $\lambda^c$. As
$\lambda^c > 0$, $\rho(\{0\}) = 1$, that is $\rho = \delta_0$. Then
$\widehat\mu_{s,t}(\xi) = \widehat\mu_{s,t}(q\xi)$ for every $\xi$, and by iteration
$\widehat\mu_{s,t}(\xi) = \widehat\mu_{s,t}(q^n\xi) \to \widehat\mu_{s,t}(0) = 1$, the transform
being continuous. So $\mu_{s,t} = \delta_0$, which contradicts $\mu_{s,t}(\{0\}) = \lambda^c < 1$.

Every clause is machine-checked, along this proof, and no cited fact is used. The Lean takes
$q = \tfrac12$ in clause (4), and it takes from Mathlib the fact of clause (2), that a
convolution with an absolutely continuous factor is absolutely continuous. The argument of
clause (4) shows more than the clause: on $\RR^d$, $d \ge 1$, a probability measure that is the
sum of a point mass of weight in $(0,1)$ at the origin and an absolutely continuous measure is
not self-decomposable.
\end{proof}

\begin{proposition}[a stage is infinitely divisible; the stages with a density]
\label{prop:isotropic-stage-density}
\uses{lem:isotropic-stage-levy-measure,prop:isotropic-stage-atom,thm:isotropic-classification,def:isotropic-kernel-family,def:isotropic-origin-datum,def:filtration-thresholds,thm:cited-A1,thm:cited-A3}
\lean{AffineHemigroup.isotropic_stage_triplet, AffineHemigroup.isotropic_stage_absolutelyContinuous,
AffineHemigroup.isotropic_stage_no_point_mass, AffineHemigroup.isotropic_stage_smooth_density,
AffineHemigroup.isotropic_stage_dichotomy}\leanok
Let $d \ge 1$, let $(\mu_{s,t})_{0 \le s \le t}$ be an isotropic cascade family in dimension $d$,
and let $\chi$ be a gauge and $F \in \mathcal{A}_d$ with $F \not\equiv 0$ that satisfy
\ref{eq:isotropic-classification}. Let $(a, k_d)$ be an isotropic origin datum in dimension $d$
with exponent $\xi \mapsto F(|\xi|)$, put $c := k_d(0{+}) \in [0,\infty]$, let $0 < s < t$, and
let $\nu_{s,t}$ be the \Levy{} measure of the stage (\ref{lem:isotropic-stage-levy-measure}).
\begin{enumerate}
\item \emph{(The stage is infinitely divisible.)} $\mu_{s,t}$ is infinitely divisible and has the
generating triplet $\bigl(a(\chi(t)^2 - \chi(s)^2)\,I,\ \nu_{s,t},\ 0\bigr)$: its Gaussian part
has the variance $a(\chi(t)^2 - \chi(s)^2)$ in every direction, and its \Levy{} measure has the
polar profile $k_d(r/\chi(t)) - k_d(r/\chi(s))$ and the mass $c\,\log(\chi(t)/\chi(s))$.
\item \emph{(No Gaussian part, an infinite profile at the origin: a density.)} If $a = 0$ and
$c = \infty$, then $\mu_{s,t}$ is absolutely continuous with respect to Lebesgue measure. In
particular it has no point mass.
\item \emph{(A Gaussian part: a smooth density.)} If $a > 0$, then $\mu_{s,t}$ has a bounded
density of class $C^\infty$, all of whose partial derivatives are bounded.
\end{enumerate}
So, with \ref{prop:isotropic-stage-atom}: a stage $\mu_{s,t}$ with $0 < s < t$ has a point mass
at the centre exactly when the datum has no Gaussian part and a finite profile at the origin,
and it has a density in the two other cases.
\statusA\quad\ledger{A1}\quad\ledger{A3}\quad\emph{(\textbf{Assignment.} \ledger{A3}(iii)
supplies the infinitely divisible law of a triplet: of the stage's triplet in clause (1), and of
the part of the \Levy{} measure carried by a small ball in clause (2). \ledger{A1}, at every
order $n$, gives the smooth density of clause (3) from the Gaussian bound on the transform.
Clause (3) does not read \ledger{A3}, and clauses (1) and (2) do not read \ledger{A1}.
\textbf{What the citations carry and what they do not.} \ledger{A3}(iii) carries the existence
of a law with a given triplet; that the triplet is the one displayed is
\ref{lem:isotropic-stage-levy-measure}, which cites nothing. \ledger{A1} carries the passage
from an integrable transform to a density. The absolute continuity of clause (2) is proved
here, by removing the jumps larger than $\varepsilon$ and letting $\varepsilon$ tend to $0$; no
cited fact on the absolute continuity of infinitely divisible laws is used. \textbf{What it
is.} With \ref{prop:isotropic-stage-atom} it answers which stages have a density and which a
point mass. The Gaussian kernels and the stable kernels of \ref{prop:isotropic-bridge}(3) have
stages with a density: the first by clause (3), the second by clause (2), their profile being
infinite at the origin. \textbf{Another route to clause (1).} That every kernel of an isotropic
cascade family is infinitely divisible is also
\ref{thm:cascade-infinitely-divisible}, by Khintchine's theorem on null arrays, ledger entry A6;
this node does not use it, and obtains the triplet as well. \textbf{Readings.} Clause (1) says that the stage
has the displayed triplet, which is what \ledger{A3}(iii) and \ref{eq:stage-transform} give.
That a law has one triplet only is \ledger{A3}(ii), which the clause does not read. In the last
sentence ``no Gaussian part'' is $a = 0$ and ``the two other cases'' are $a = 0$ with
$c = \infty$, and $a > 0$. The numbers $a(\chi(t)^2 - \chi(s)^2)$ and $c$ and the measure
$\nu_{s,t}$ do not depend on the pair $(\chi, F)$ chosen for the family
(\ref{lem:isotropic-stage-levy-measure}, annotation). \textbf{Not claimed}: that the density of
clause (2) is bounded, or continuous; nor any rate at the centre for a stage.
\textbf{Formalization} (Q-0447). Machine-checked: clauses (1) and (2) on \ledger{A3}(iii) and
Lean's logical axioms, clause (3) on \ledger{A1} and Lean's logical axioms, the last sentence on
both. The family is any family of probability measures with the transforms
\ref{eq:isotropic-classification} for a gauge $\chi$ and a function $F$ whose radial form is the
exponent of the datum; no cascade property is asked and $F \in \mathcal{A}_d$ is not asked, which
is more than the print states, and $F \not\equiv 0$ is asked nowhere. Clause
(1) is infinite
divisibility together with a generating triplet whose three components are the displayed ones;
the polar profile and the mass of $\nu_{s,t}$ are its definition and
\ref{lem:isotropic-stage-levy-measure}(3). Clause (2) is two declarations: absolute continuity
with respect to Lebesgue measure, and $\mu_{s,t}(\{x\}) = 0$ for every $x \in \RR^d$. Clause (3)
is a nonnegative
function that is $C^n$ for every $n$, has a bounded derivative of every order, and is a density
of $\mu_{s,t}$; ``bounded'' is the case $n = 0$, and the derivative of order $n$ is the
$n$-th derivative as a multilinear map, whose norm bounds every partial derivative of order
$n$. The last
sentence is one declaration: $\mu_{s,t}(\{0\}) > 0$ if and only if $a = 0$ and $c < \infty$, and
otherwise $\mu_{s,t}$ is absolutely continuous.)}
\end{proposition}

\begin{proof}
\leanok
Write $\nu := \nu_{s,t}$, $\ell := \chi(s)$, $L := \chi(t)$ and $v := a(L^2 - \ell^2) \ge 0$. By
\ref{lem:isotropic-stage-levy-measure}(1), $\nu$ gives the origin no mass, is symmetric, and has
$\int(|x|^2 \wedge 1)\,d\nu < \infty$: it is a symmetric \Levy{} measure.

\emph{A symmetric triplet.} Let $\nu_0$ be a symmetric \Levy{} measure on $\RR^d$ and
$v_0 \ge 0$. By \ledger{A3}(iii) there is an infinitely divisible law with the triplet
$(v_0I, \nu_0, 0)$, and its transform is
\begin{equation}\label{eq:symmetric-triplet}
  \exp\Bigl[-\tfrac12 v_0|\xi|^2 - \int_{\RR^d}\bigl(1 - \cos\langle\xi,x\rangle\bigr)\,\nu_0(dx)\Bigr] .
\end{equation}
Indeed the integrand of the representation, $e^{i\langle\xi,x\rangle} - 1
- i\langle\xi,x\rangle\mathbf 1_D(x)$, has the real part $\cos\langle\xi,x\rangle - 1$ and an
imaginary part that is an odd function of $x$, integrable against $\nu_0$ because it is at most
$|\xi|^2|x|^2/2$ in modulus on $D$ (as $|\sin u - u| \le u^2/2$ for every real $u$) and at most
$1$ off $D$; its integral against the symmetric
measure $\nu_0$ is $0$.

\emph{Clause (1).} By \ref{eq:symmetric-triplet} with $(v_0,\nu_0) = (v,\nu)$ and
\ref{eq:stage-transform}, the infinitely divisible law with the triplet $(vI, \nu, 0)$ has the
transform of $\mu_{s,t}$, so it is $\mu_{s,t}$. The polar profile and the mass of $\nu$ are
those of \ref{lem:isotropic-stage-levy-measure}.

\emph{Clause (2).} Let $a = 0$ and $c = \infty$, so that $\nu$ has infinite mass
(\ref{lem:isotropic-stage-levy-measure}(3)). For $\varepsilon > 0$ let $\nu^{>}_\varepsilon$ and
$\nu^{\le}_\varepsilon$ be the restrictions of $\nu$ to $\{|x| > \varepsilon\}$ and to
$\{0 < |x| \le \varepsilon\}$. The first is a finite measure, because
$\int(|x|^2 \wedge 1)\,d\nu < \infty$; it is symmetric and has a density; and its mass
$m_\varepsilon$ increases to $\nu(\RR^d \setminus \{0\}) = \infty$ as $\varepsilon \downarrow 0$.
The second is a symmetric \Levy{} measure. Let $\mu^\varepsilon$ be the law of
\ref{eq:symmetric-triplet} for $(0, \nu^{\le}_\varepsilon)$, and put
$\pi_\varepsilon := e^{-m_\varepsilon}\sum_{n \ge 0}(\nu^{>}_\varepsilon)^{*n}/n!$. As in the
proof of \ref{prop:isotropic-stage-atom}(1) and~(2), which uses of the measure only that it is
finite, symmetric and has a density, $\pi_\varepsilon$ is a probability measure with the
transform $\exp[-\int(1 - \cos\langle\xi,x\rangle)\,\nu^{>}_\varepsilon(dx)]$, and
$\pi_\varepsilon = e^{-m_\varepsilon}\delta_0 + \eta_\varepsilon$ with $\eta_\varepsilon$
absolutely continuous. The transforms of $\pi_\varepsilon$ and $\mu^\varepsilon$ multiply to
\ref{eq:stage-transform} with $a = 0$, so $\mu_{s,t} = \pi_\varepsilon * \mu^\varepsilon$. Let
$N$ be a Lebesgue-null Borel set. Then
$(\eta_\varepsilon * \mu^\varepsilon)(N) = \int\eta_\varepsilon(N - y)\,\mu^\varepsilon(dy) = 0$,
so
\[
  \mu_{s,t}(N) = e^{-m_\varepsilon}\mu^\varepsilon(N) \le e^{-m_\varepsilon} .
\]
Letting $\varepsilon \downarrow 0$ gives $\mu_{s,t}(N) = 0$. So $\mu_{s,t}$ is absolutely
continuous, and the origin is a Lebesgue-null set.

\emph{Clause (3).} Let $a > 0$. A gauge is increasing, so $v = a(L^2 - \ell^2) > 0$. By
\ref{eq:stage-transform}, $0 < \widehat\mu_{s,t}(\xi) \le e^{-v|\xi|^2/2}$, so
$\int|\widehat\mu_{s,t}(\xi)|\,|\xi|^n\,d\xi < \infty$ for every integer $n \ge 0$. By \ledger{A1}
at $n$, $\mu_{s,t}$ has a density $g_n$ of class $C^n$ whose partial derivatives of orders
$0,\dots,n$ tend to $0$ at infinity. Two continuous densities of one measure are equal, so the
$g_n$ are one function $p$. It is of class $C^n$ for every $n$, and each of its partial
derivatives is continuous and tends to $0$ at infinity, hence is bounded.

\emph{The last sentence.} If $a = 0$ and $c < \infty$, \ref{prop:isotropic-stage-atom}(2) gives
the point mass. If $a = 0$ and $c = \infty$, or $a > 0$, clauses (2) and (3) give a density,
which excludes a point mass.

Every clause is machine-checked, along this proof. The cited facts are \ledger{A3}(iii), in
clauses (1) and (2), and \ledger{A1}, in clause (3); they are the interfaces of the Lean, which
proves neither. In clause (2) the Lean takes $\varepsilon = 1/(n+1)$, and in clause (3) it bounds
the moments of the transform by $r^n \le n!\,e^r$ and $r \le 1/v + (v/4)r^2$.
\end{proof}

% draft: section 2.1, Proposition 2.3, in the reduction form of the plan's S13: the family clause
% quantifies over the causal kernels, as [B, Prop. 4.5] does; the Matern and stable corners are
% exponent identities at |xi|; the Student-t corner is conditional, as [B] made it (R177; ledger
% A18 in the line repository); strictness is thm:filtration-strict(3), not the d-dimensional Cin
% rays (it was this node's clause (5) until Q-0286, a duplicate of that clause, dropped).
% CHANGED Q-0281: typed (the statement skeleton). Clause (4), the conditional Student-t corner, is split
% to prop:isotropic-bridge-student (ADR-0025; records/formalization/PRICING.md); the text is unchanged
% and the \lean tag names the child's declarations too. The causal kernels are a hypothesis
% (IsCausalKernels), as the node takes them. Reading: clause (4)'s transform is stated at xi != 0,
% where K_vartheta(t|xi|) is finite; at xi = 0 the display is read as its limit, 1.
% CHANGED Q-0281: blind review R2-2: the tag names measurable_gaussVar, which the bridge family
% eq:isotropic-bridge (a Measure.bind in Lean) presupposes.
% CHANGED Q-0417 (the freeze of the isotropic paper; REPOSITIONING-paper-iso.md section 19, and
% the register reviews' flags on "corner" and "causal family", REVIEW-register-paper-iso-2026-10-11-
% fixes.md rows 36, 10 and 19): the statement in the paper's words. The title is the paper's ("and
% three named families" for "; the corners"); the clause titles say "kernels" for "corner" (the
% Matern kernels, the stable kernels, the Student-t kernels, as the paper's lead names them);
% clause (4) says "a time-causal family as above" for "a causal family as above, the causal Bessel
% family". The last is [B]'s name for the family whose delay kernels would be the inverse-gamma
% laws; the statement no longer uses the name, since "as above" and the citation of
% [B, Prop. 4.5(3)] identify it. No mathematical content changes. The child
% prop:isotropic-bridge-student, a verbatim copy of clause (4), follows. The annotation's "the
% draft's corners" is about the draft's wording and stays.
% CHANGED Q-0431 (the author's section review of 2026-10-12, records/affine/
% REVIEW-author-paper-iso.md: a statement that rests on a published theorem is presented as
% resting on it, never as a hypothesis): the wording of clause (4), here and in the child. The
% clause title loses "conditional", "If" its emphasis, and the parenthesis "of a time-causal
% family as above, whose existence [B, Prop. 4.5(3)] leaves as a hypothesis" becomes "of the
% family". The logical form is unchanged and is the Lean's: `isotropic_bridge_student_family`
% has the hypothesis `hT : muI 0 1 = ScaleSpace.inverseGammaLaw theta` beside
% `IsCausalKernels FI muI`, read for this item (Formalization/AffineHemigroup/
% IsotropicBridge.lean). Where the family's existence is in print is said in the annotation,
% not in the clause: a citation inside the clause ("they are, by Halgreen's theorem") would be an
% assertion of the node that no declaration states, against ADR-0025, and clauses (2) and (3)
% say nothing of the existence of their families either. The proof's "Under the hypothesis" is
% "When mu^I_{0,1} is the law of T_1". Point (ii) of "against the draft" keeps the word
% "conditional": it compares two wordings and is not printed in a paper.
\begin{proposition}[the bridge in dimension $d$, and three named families]
\label{prop:isotropic-bridge}
\uses{def:isotropic-kernel-family,thm:isotropic-classification,def:filtration-thresholds}
\lean{AffineHemigroup.isotropic_bridge_family, AffineHemigroup.isotropic_bridge_matern,
AffineHemigroup.isotropic_bridge_stable, AffineHemigroup.isotropic_bridge_student_transform,
AffineHemigroup.isotropic_bridge_student_family, AffineHemigroup.IsCausalKernels,
AffineHemigroup.bridgeFamily, AffineHemigroup.gaussVar,
AffineHemigroup.measurable_gaussVar}\leanok
Let $F_{\mathrm I}$ be causally admissible in the sense of \restated{B}{Def.~4.1}{def:line-causal}, restated:
$F_{\mathrm I}(\sigma) = b_0\sigma + \int_0^\infty(1 - e^{-\sigma u})\,k_{\mathrm I}(u)\,du/u$ with
$b_0 \ge 0$ and $k_{\mathrm I} : (0,\infty) \to [0,\infty)$ nonincreasing,
$\int_0^1 k_{\mathrm I}(u)\,du < \infty$ and $\int_1^\infty k_{\mathrm I}(u)\,u^{-1}du < \infty$.
Suppose $F_{\mathrm I} \not\equiv 0$. Let $(\mu^{\mathrm I}_{x,y})_{0 \le x \le y}$ be probability
measures on $[0,\infty)$ with
\[
  \int_{[0,\infty)} e^{-\sigma u}\,\mu^{\mathrm I}_{x,y}(du)
   = \exp\bigl[-\bigl(F_{\mathrm I}(y\sigma) - F_{\mathrm I}(x\sigma)\bigr)\bigr], \qquad \sigma \ge 0 ,
\]
the kernels of a time-causal family in its canonical gauge, as \restated{B}{Prop.~4.5}{sec:line-bridge} takes them. Fix
$d \ge 1$, let $g^{(d)}_u$ be the centred Gaussian law on $\RR^d$ with covariance $uI$ for $u > 0$
and $g^{(d)}_0 := \delta_0$, and put
\begin{equation}\label{eq:isotropic-bridge}
  \mu_{s,t} := \int_{[0,\infty)} g^{(d)}_u\;\mu^{\mathrm I}_{s^2,t^2}(du), \qquad 0 \le s \le t,
  \qquad F(r) := F_{\mathrm I}\bigl(\tfrac12 r^2\bigr) .
\end{equation}
\begin{enumerate}
\item \emph{(The family.)} $(\mu_{s,t})$ is an isotropic cascade family in dimension $d$ with
$S_\lambda t = \lambda t$, and
$\widehat\mu_{s,t}(\xi) = \exp[-(F(t|\xi|) - F(s|\xi|))]$: it is the family of
\ref{thm:isotropic-classification} with gauge $\chi = \mathrm{id}$ and exponent $F$, and
$F \in \mathcal{A}_d$, for every $d \ge 1$.
\item \emph{(The \Matern{} kernels.)} If $F_{\mathrm I}(\sigma) = \gamma\log(1+\sigma)$ with
$\gamma > 0$, then $F(r) = \gamma\log(1 + r^2/2)$ and
$\widehat\mu_{0,t}(\xi) = (1 + t^2|\xi|^2/2)^{-\gamma}$.
\item \emph{(The stable kernels.)} If $F_{\mathrm I}(\sigma) = \sigma^\alpha$ with
$0 < \alpha \le 1$, then $F(r) = 2^{-\alpha}r^{2\alpha}$ and
$\widehat\mu_{0,t}(\xi) = \exp(-2^{-\alpha}t^{2\alpha}|\xi|^{2\alpha})$, of index
$2\alpha \in (0,2]$. At $\alpha = 1$, $\mu_{s,t} = g^{(d)}_{t^2-s^2}$, the Gaussian.
\item \emph{(The Student-t kernels.)} Let $\vartheta > 0$ and let $T_1$ be the scaled
inverse-gamma law of shape $\vartheta$ of \prog{B}{Prop.~5.10}, with density
$u^{-\vartheta-1}e^{-1/(2u)}/(2^\vartheta\Gamma(\vartheta))$. For every
$t > 0$ the Gaussian variance mixture $\int g^{(d)}_u\,P(t^2T_1 \in du)$ is a probability measure
on $\RR^d$ with transform
\[
  \frac{2^{1-\vartheta}}{\Gamma(\vartheta)}\,(t|\xi|)^\vartheta\,K_\vartheta(t|\xi|) ,
\]
$K_\vartheta$ being the modified Bessel function of the second kind of order $\vartheta$.
If $\mu^{\mathrm I}_{0,1}$ is the law of $T_1$, that is, if the inverse-gamma laws are the
delay kernels of the family, then these mixtures are the kernels $\mu_{0,t}$ of
clause (1), and their exponent is in $\mathcal{A}_d$.
\end{enumerate}
\statusT\quad\emph{(Cites released statements of [B] (Def.~4.1, Lem.~4.2, Prop.~4.5 and
Prop.~5.10(2)), restated where they are used, and no ledger entry. \textbf{What is quantified
over, and where its existence is in print.} The statement is about the kernels of a time-causal
family and takes them as given, as [B, Prop.~4.5] does: in Lean, the hypothesis
\texttt{IsCausalKernels}. That a causally admissible $F_{\mathrm I} \not\equiv 0$ has them is the
classification of the time-causal paper (\sscite{fagerstrom2026hemigroup}, version~1.1.0,
Thm.~7.3), and [B, Rem.~4.3] has the delay law, with Laplace transform $e^{-F_{\mathrm I}}$, on
two cited facts about Bernstein functions. Clause (4)'s family sentence is about the family whose
delay law $\mu^{\mathrm I}_{0,1}$ is the inverse-gamma law: in Lean, the hypothesis
\texttt{hT} of \texttt{isotropic\_bridge\_student\_family}, which says exactly that. There is
such a family. The inverse-gamma laws are generalized gamma convolutions
(\sscite{bondesson1978infinite}, \S4, pp.~51--52), a class of self-decomposable laws on the
half-line (\sscite{halgreen1979self}, pp.~13--15, where $T_1$ is the law
$\mathrm{GIGD}(-\vartheta,1,0)$), and the time-causal paper then has the family
(\sscite{fagerstrom2026hemigroup}, version~1.1.0, Prop.~8.17, by its Thm.~7.3). Both are results
in print, read for this module (\texttt{records/affine/READING-halgreen1979self.md},
\texttt{READING-bondesson1978infinite.md}, and for the time-causal paper
\texttt{records/affine/reviews/VERIFY-round1-mathematics.md}, item 10). They are not ledger
entries: no declaration reads them, the declarations taking the family as given, so the tag
covers each clause as stated. The transform sentence of clause (4) needs no family.
\textbf{Notation.} [B, Prop.~5.10] writes $a$ for the shape parameter; this node
writes $\vartheta$, because $a$ is the Gaussian coefficient of \ref{def:isotropic-origin-datum}
(Q-0286, which also glossed $K_\vartheta$). \textbf{Strictness} is \ref{thm:filtration-strict}(3),
$\mathcal{S} = \mathcal{A}_\infty \subsetneq \mathcal{A}_d$ for every $d$, with $\mathcal{S}$ the
subordinated slice of \ref{def:filtration-thresholds}, the exponents $F$ of clause (1). Until
Q-0286 this node stated it again as its clause (5), with a third definition of $\mathcal{S}$ and a
proof by forward reference to that theorem; the clause was a duplicate and was dropped, nothing
citing it.
\textbf{Against the draft's Proposition~2.3}, each point accepted by the author as the statement
of record (Q-0222, 2026-09-30). (i) The draft states the family unconditionally, as
``$d$-dimensional Brownian motions at the subordinator times''; here the causal kernels are
hypothesised, and the probabilistic reading is \ref{eq:isotropic-bridge} as a mixture, with no
process asserted. (ii) The draft's Student-t corner is unconditional; here it is conditional.
(iii) The draft's corners are written in other gauges: its \Matern{} member
$(1 + t^2|\omega|^2)^{-\gamma}$ is clause (2) in the gauge $\chi(t) = \sqrt2\,t$, and its stable
member $e^{-t^\alpha|\omega|^\alpha}$ of index $\alpha \in (0,2]$ is clause (3) with the index
written $2\alpha$ and the gauge rescaled. The names ``Bessel potentials of $\RR^d$'' and
``elliptical Student-t with $2\vartheta$ degrees of freedom'' are the draft's; no density is
computed here. (iv) The draft witnesses strictness by the $d$-dimensional $\Cin$ rays, ``whose
symbol vanishes on spheres''; that claim is unverified and is not used. Strictness is
\ref{thm:filtration-strict}(3), whose witness is [B]'s family $F^{(w)}$, and the $\Cin$ rays of the
line are \ref{thm:no-cin-ray-in-the-plane}. (v) The draft's explicit profile
$k(r) = r^d\int_0^\infty(2\pi u)^{-d/2}e^{-r^2/2u}k_{\mathrm I}(u)\,du/u$ is not a clause: it is
[B, Prop.~4.9(2)]'s polar profile, recorded in \ref{def:isotropic-origin-datum} as
\ref{eq:polar-profile-d}, and the two agree, the substitution $u = r^2/2v$ in the draft's integral
giving $\pi^{-d/2}\int_0^\infty v^{d/2-1}e^{-v}k_{\mathrm I}(r^2/2v)\,dv$, which is
\ref{eq:polar-profile-d} divided by $|S^{d-1}| = 2\pi^{d/2}/\Gamma(d/2)$. At $d = 1$ clause (1) is
[B, Prop.~4.5]'s family clause read on kernels.)}
\end{proposition}

\begin{proof}
\leanok
\emph{(1).} Each $\mu_{s,t}$ is a mixture of probability measures, and $u \mapsto g^{(d)}_u(B)$
is measurable for every Borel $B$, so it is a probability measure. By Fubini and
$\widehat{g^{(d)}_u}(\xi) = e^{-u|\xi|^2/2}$,
\[
  \widehat\mu_{s,t}(\xi) = \int e^{-u|\xi|^2/2}\,\mu^{\mathrm I}_{s^2,t^2}(du)
   = \exp\bigl[-\bigl(F_{\mathrm I}(t^2|\xi|^2/2) - F_{\mathrm I}(s^2|\xi|^2/2)\bigr)\bigr] ,
\]
which is $\exp[-(F(t|\xi|) - F(s|\xi|))]$.
The axioms are read off this transform exactly as in the sufficiency half of
\ref{thm:isotropic-classification}, with $\chi = \mathrm{id}$ and $S_\lambda t = \lambda t$. That
half used three properties of $F$, and they hold here. $F$ is continuous on $[0,\infty)$, because
$F_{\mathrm I}$ is continuous by dominated convergence: its integrand is at most
$\max\{1,\Sigma\}(1 \wedge u)\,k_{\mathrm I}(u)/u$ on $\sigma \le \Sigma$, integrable by the two
conditions. $F(0) = F_{\mathrm I}(0) = 0$. And $F \not\equiv 0$, because $F_{\mathrm I} \not\equiv 0$
and $r^2/2$ runs over $[0,\infty)$. So $(\mu_{s,t})$ is an
isotropic cascade family, and \ref{thm:isotropic-classification} applies. Its conclusion is met by
$(\mathrm{id},F)$ once $F \in \mathcal{A}_d$, and the argument of its necessity half gives that.
$F$ is admissible by \restated{B}{Lem.~4.2}{prop:line-bridge}, restated: \emph{for causally admissible $F_{\mathrm I}$ with
data $(b_0,k_{\mathrm I})$, $F_{\mathrm I}(\omega^2/2)$ is of the form \ref{eq:admissible-line}
with $a_{\mathrm V} = \tfrac12 b_0$.} And $\mu_{0,1}$ has transform $e^{-F(|\cdot|)}$ and, by
(I1), residuals $\mu_{\lambda,1}$ in \ref{eq:self-decomposable}, as $\widehat\mu_{0,\lambda}(\xi)
= e^{-F(\lambda|\xi|)} = \widehat\mu_{0,1}(\lambda\xi)$. So $F \in \mathcal{A}_d$, and $d$ was
arbitrary.

\emph{(2) and (3)} are the exponent identities $\gamma\log(1 + r^2/2)$ and
$(r^2/2)^\alpha = 2^{-\alpha}r^{2\alpha}$, read at $r = t|\xi|$ through (1). At $\alpha = 1$,
$\widehat\mu_{s,t}(\xi) = e^{-(t^2-s^2)|\xi|^2/2}$, the transform of $g^{(d)}_{t^2-s^2}$.

\emph{(4).} The mixture is a probability measure as in (1), and by Fubini its transform is
$E\,e^{-t^2T_1|\xi|^2/2}$. On the line, \prog{B}{Prop.~5.10(2)} says that the law of the mixture of
$g^{(1)}_u$ against $T_1$ has transform
$\tfrac{2^{1-\vartheta}}{\Gamma(\vartheta)}|\omega|^\vartheta K_\vartheta(|\omega|)$. By the
same Fubini that transform is $E\,e^{-T_1\omega^2/2}$. Reading it at $\omega = t|\xi|$ gives the
display. When $\mu^{\mathrm I}_{0,1}$ is the law of $T_1$, by (1) and its Laplace transform at
$t^2|\xi|^2/2$, $\widehat\mu_{0,t}(\xi) = e^{-F_{\mathrm I}(t^2|\xi|^2/2)} = E\,e^{-t^2T_1|\xi|^2/2}$,
which is the transform of the mixture. Two probability measures on $\RR^d$ with the same
transform are equal, so $\mu_{0,t}$ is the mixture, and its exponent is in $\mathcal{A}_d$ by (1).
% CHANGED Q-0281 wave 2: the last step rewritten to the Lean route (the transforms on R^d are
% compared directly; no uniqueness of the Laplace transform on [0, infinity) is needed).
\end{proof}

% CHANGED Q-0281: split from prop:isotropic-bridge (ADR-0025; records/formalization/PRICING.md): the
% conditional Student-t clause (4), copied verbatim with the parent's clause number; the setting is
% the parent's. Its transform sentence reads the shared library's student_transform and the mixture
% identity bridge_families_bessel (ScaleSpaceCore v0.4.0, Lean core; Paper V's at v0.1 until the
% switch). The parent's text and numbering are unchanged.
% CHANGED Q-0281: blind review R2-2: the tag names measurable_gaussVar (the mixture over the inverse
% Gamma law is a Measure.bind in Lean).
% CHANGED Q-0417: the title, the clause title and "a time-causal family as above" follow the
% parent's clause (4) as reworded there (wording only).
% CHANGED Q-0431: the title and the clause follow the parent's clause (4) as reworded there
% (wording only; "conditional" is gone from both titles, and the annotation points to the
% parent's for where the family's existence is in print).
\begin{proposition}[the Student-t kernels of the bridge]
\label{prop:isotropic-bridge-student}
\uses{prop:isotropic-bridge,def:filtration-thresholds}
\lean{AffineHemigroup.isotropic_bridge_student_transform,
AffineHemigroup.isotropic_bridge_student_family, AffineHemigroup.measurable_gaussVar}\leanok
In the setting of \ref{prop:isotropic-bridge}:
\begin{enumerate}
\item[(4)] \emph{(The Student-t kernels.)} Let $\vartheta > 0$ and let $T_1$ be the scaled
inverse-gamma law of shape $\vartheta$ of [B, Prop.~5.10], with density
$u^{-\vartheta-1}e^{-1/(2u)}/(2^\vartheta\Gamma(\vartheta))$. For every
$t > 0$ the Gaussian variance mixture $\int g^{(d)}_u\,P(t^2T_1 \in du)$ is a probability measure
on $\RR^d$ with transform
\[
  \frac{2^{1-\vartheta}}{\Gamma(\vartheta)}\,(t|\xi|)^\vartheta\,K_\vartheta(t|\xi|) ,
\]
$K_\vartheta$ being the modified Bessel function of the second kind of order $\vartheta$.
If $\mu^{\mathrm I}_{0,1}$ is the law of $T_1$, that is, if the inverse-gamma laws are the
delay kernels of the family, then these mixtures are the kernels $\mu_{0,t}$ of
clause (1), and their exponent is in $\mathcal{A}_d$.
\end{enumerate}
\statusT\quad\emph{(Cites [B, Prop.~5.10(2)], restated in the proof of
\ref{prop:isotropic-bridge}, and no ledger entry. The family sentence is about the time-causal
family whose delay law is the inverse-gamma law, which it takes as given; the parent's annotation
says where the existence of that family is in print. Reading: the transform is stated at $\xi \ne 0$, where $K_\vartheta(t|\xi|)$ is finite; at
$\xi = 0$ the display is read as its limit, $1$. Clause (4) of \ref{prop:isotropic-bridge}, split off
on 2026-10-04 (Q-0281, ADR-0025), with the parent's clause number.)}
\end{proposition}

\begin{proof}
\leanok
This is the part of the proof of \ref{prop:isotropic-bridge} headed (4).
\end{proof}
